% SPDX-License-Identifier: Apache-2.0
%
% Datasheets: one per component the project has built. Built by
% //docs:datasheets. The tables of ports, state and children come from
% each component's own lowering, through //tools/datasheet; the prose is
% in docs/datasheets/, one file per component, and
% //tools/datasheet:coverage_test fails when a component has none.
\documentclass[journal,onecolumn,9pt]{IEEEtran}

\newcommand{\listingsize}{\normalsize}
% SPDX-License-Identifier: Apache-2.0
%
% The house preamble, shared by every document in this directory. Loaded
% after \documentclass, before \title. Keeping it in one file is what
% stops two articles drifting apart in font, listing style or figure
% style.
% Computer Modern. IEEEtran selects Times (ptm) by default, so the face has
% to be chosen explicitly.
%
% Latin Modern is the Computer Modern variety used here, and the choice is
% forced rather than aesthetic. Under T1 encoding, plain `cmr` has no Type 1
% outlines unless cm-super is installed, and it is not in the pinned
% distribution, so pdflatex falls back to EC bitmaps and every glyph in the
% PDF becomes a Type 3 font. That was measured: the first build produced a
% document with 10 Type 3 fonts and no Type 1. Latin Modern is Computer
% Modern redrawn with a full T1 repertoire and real .pfb outlines, so the
% same design comes out scalable.
\usepackage[T1]{fontenc}
\usepackage[utf8]{inputenc}
\usepackage{lmodern}
% microtype is not in the pinned TeX distribution. emergencystretch alone
% keeps the two-column measure from producing overfull lines.
\setlength{\emergencystretch}{3em}

\usepackage{amsmath}
\usepackage{amssymb}
\usepackage{amsthm}
\usepackage{booktabs}
\usepackage{array}
% enumitem is not in the pinned TeX distribution; the list spacing below
% does what \setlist would have done.
\usepackage{float}
\usepackage{xcolor}
\usepackage{listings}
\usepackage{tikz}
\usetikzlibrary{arrows.meta,positioning,fit,backgrounds,calc}
\usepackage[hidelinks,bookmarks=true]{hyperref}

% Numbered, definition-style box for a rule the reader is meant to apply.
\theoremstyle{definition}
\newtheorem{rulebox}{Rule}
\newtheorem{example}{Example}

% Unnumbered asides.
\theoremstyle{remark}
\newtheorem*{tip}{Tip}
\newtheorem*{pitfall}{Pitfall}

% Tighter lists, in place of enumitem's \setlist.
\setlength{\topsep}{2pt}
\setlength{\itemsep}{1pt}
\setlength{\parsep}{0pt}

% Code in the running text. The typewriter face under T1 ligates `<<`
% and `>>` into guillemets, so a shift printed as \code{<<} came out as
% a French quotation mark (issue 571). microtype's \DisableLigatures is
% not in the pinned distribution, but the primitive it rests on is
% pdfTeX's own: \pdfnoligatures strips every ligature from the font it
% is given, and the current font inside \texttt is the typewriter face
% at the size in use, so each size is stripped the first time \code
% sets it. Code wants no ligature at all: two quotes are two quotes.
\newcommand{\code}[1]{\texttt{\ifdefined\pdfnoligatures\pdfnoligatures\font\fi #1}}
\newcommand{\term}[1]{\emph{#1}}

% Syntax highlighting. Four roles, distinguishable in colour and still
% distinguishable in greyscale, because an IEEE article gets printed: the
% keyword is the only bold role, the comment is the only italic one, and the
% two remaining roles differ in lightness as well as in hue.
\definecolor{kwblue}{RGB}{20,60,150}
\definecolor{cmtgreen}{RGB}{90,120,90}
\definecolor{strred}{RGB}{150,50,40}
\definecolor{numgray}{gray}{0.45}
\definecolor{framegray}{gray}{0.70}
\definecolor{bgshade}{gray}{0.97}
% A darker fill for figure cells. The listing background has to stay very
% light to keep code readable; a figure cell has to be clearly shaded or the
% distinction it draws is invisible in print.
\definecolor{cellshade}{gray}{0.90}

% The listing size is the document's choice, because it depends on the
% column: a two-column article sets `\listingsize` to 6pt and keeps its
% listings within 70 columns, or to 5.6pt when it includes source that
% rustfmt sets at 80, which is what fits a column's frame (issue 1061);
% a one-column document keeps the body size and includes source files
% at 80. `//docs:frame_test` checks every listing line against its frame. Lines are never broken; a line that does not fit
% overflows the frame, which is visible, rather than wrapping, which is
% not.
\providecommand{\listingsize}{\scriptsize}

% `'` is deliberately not a string delimiter: the language uses it for loop
% labels, as Rust does, so treating it as a quote would swallow the rest of
% the line.
\lstdefinestyle{txhdl}{
  basicstyle=\ttfamily\listingsize,
  keywordstyle=\bfseries\color{kwblue},
  commentstyle=\itshape\color{cmtgreen},
  stringstyle=\color{strred},
  numberstyle=\tiny\color{numgray},
  morekeywords={mod,use,pub,const,type,struct,enum,trait,impl,fn,async,let,mut,
    match,if,else,loop,while,for,in,break,continue,return,as,await,
    module,bus,channel,transaction,protocol,pipeline,flow,seq,config,
    port,stage,pipe,instance,connect,bind,tag,domain,blackbox,foreign,
    property,role,signal,handler,true,false,
    sequence,interface,modport,implements,package,import,var,wire,func,proc,
    case,default,next,fork,branch,join,state,fsm},
  morecomment=[l]{//},
  morestring=[b]",
  showstringspaces=false,
  columns=fullflexible,
  keepspaces=true,
  breaklines=false,
  % Framed, so a listing is bounded rather than floating in the column.
  frame=single,
  rulecolor=\color{framegray},
  framesep=3pt,
  backgroundcolor=\color{bgshade},
  xleftmargin=3pt,
  xrightmargin=3pt,
  aboveskip=6pt,
  belowskip=6pt,
  % Every listing is numbered and captioned, so it can be referred to by
  % number. `caption` without `float` numbers the listing and leaves it
  % where it was written, which is what a listing in running prose needs.
  captionpos=b,
  abovecaptionskip=2pt,
  belowcaptionskip=6pt,
}

% Figures drawn as text. Framed the same way, and with the highlighting
% switched off, because a box-and-arrow diagram is not code and half its
% words turning blue reads as an accident. Setting a style adds keys rather
% than clearing the ones already set, so the three roles are neutralised by
% hand rather than by leaving them out. Program output is printed in this
% style too, and a netlist's assignment can run past the frame, so a line
% that does not fit is broken and the rest indented; source never is.
\lstdefinestyle{ascii}{
  basicstyle=\ttfamily\listingsize,
  keywordstyle=\ttfamily,
  commentstyle=\ttfamily,
  stringstyle=\ttfamily,
  columns=fullflexible,
  keepspaces=true,
  breaklines=true,
  breakindent=2em,
  frame=single,
  rulecolor=\color{framegray},
  framesep=3pt,
  backgroundcolor=\color{bgshade},
  xleftmargin=3pt,
  xrightmargin=3pt,
  aboveskip=6pt,
  belowskip=6pt,
}
\lstset{style=txhdl}

% House TikZ styles. `shorten` lives on the arrow styles rather than on each
% arrow, so every arrow in the document gets the gap, including ones added
% later. TikZ clips a path at a node's border, so without it an arrowhead
% butts against the box with no gap at all. Keep `node distance` well above
% twice the shorten, or the arrow shrinks to nothing.
\tikzset{
  box/.style={draw,rounded corners=2pt,align=center,inner sep=4pt,
              font=\footnotesize},
  boxf/.style={box,fill=bgshade},
  lbl/.style={font=\scriptsize\itshape,align=center},
  fl/.style={-{Stealth[length=4pt]},shorten >=2.5pt,shorten <=2.5pt},
  fld/.style={fl,dashed},
}


% The contents list: the class gives a section's number 2.75em, which
% a Roman numeral past thirty overruns onto the title, so the box is
% wider here, and a subsection's entry starts where a title does.
\makeatletter
\def\l@section#1#2{\addpenalty{\@secpenalty}\addvspace{1.0em plus 1pt}%
    \@tempdima 5.75em \begingroup \parindent \z@ \rightskip \@pnumwidth%
    \parfillskip-\@pnumwidth {\bfseries\leavevmode #1}\hfil\hbox to\@pnumwidth{\hss #2}\par%
    \endgroup}
\def\l@subsection{\@dottedtocline{2}{5.75em}{5em}}
\def\l@subsubsection{\@dottedtocline{3}{10.75em}{5.5em}}
\makeatother


\InputIfFileExists{buildstamp.tex}{}{}
\providecommand{\buildstamp}{Draft build (outside the Bazel flow).}

% The generated tables. A sheet that asks for a table the generator
% does not write stops the build, rather than printing nothing.
\input{tools/datasheet/ds_tables}
\makeatletter
\newcommand{\dsget}[2]{%
  \@ifundefined{ds@#1@#2}%
    {\PackageError{datasheets}{//tools/datasheet writes no #1 for #2}%
      {Add the component to tools/datasheet/main.rs.}}%
    {\csname ds@#1@#2\endcsname}}
\makeatother

% A datasheet's heading: the component and what it is, on a page of
% its own.
\newcommand{\datasheet}[2]{\clearpage\section{#1}\label{ds:#1}%
\noindent\textit{#2}\par\medskip}

% Where a component lives: its source, its Rust path, and the document
% that explains it.
\newcommand{\dsfacts}[3]{%
\begin{center}\footnotesize
\begin{tabular}{@{}ll@{}}
Source & #1 \\
Rust path & #2 \\
Explained in & #3 \\
\end{tabular}
\end{center}}

% The generated part of a sheet: the interfaces as a diagram, how the
% component is made and joined, and the tables. All of it is written
% by //tools/datasheet from the component's own lowering, so a sheet
% cannot say a port the netlist does not have (issue 173).
% A component's register table, written by //tools/datasheet from the
% map its peripheral declares with `regmap!` (issue 499), for a sheet
% whose component has one. The sheet gives the heading and says how the
% word is selected; this is the table, a row a register and a row a
% field under it.
\newcommand{\dsregs}[1]{\dsget{regs}{#1}}

\newcommand{\dstables}[1]{%
\subsection*{Interfaces}
\dsget{figure}{#1}
\subsection*{How it is instantiated}
As lowered for \dsget{type}{#1}. The ports are named in the order the
unit's \code{run} takes them.
\dsget{instance}{#1}
\subsection*{Ports, state and children}
As lowered for \dsget{type}{#1}: \dsget{counts}{#1}.
\dsget{ports}{#1}
\dsget{state}{#1}
\dsget{children}{#1}}

\title{TxHDL Datasheets}

\author{Filip Filmar and Dragiša Janković%
\thanks{Both authors are independent. Filip Filmar is the
corresponding author, at f@hdlfactory.com; Dragiša Janković is
reached at linkedin.com/in/dragisa-jankovic.}%
\thanks{This volume provides comprehensive datasheets for hardware components
implemented in the \code{HDL/txhdl} project repository. Port, state, and
submodule tables are extracted automatically during compilation directly from
lowered netlist definitions. The project overview~\cite{cover} indexes companion
documentation.}%
\thanks{The authors used Claude, an AI assistant, while developing
architectural concepts, documentation, and software. Verbatim prompts are
preserved in repository commit logs.}%
\thanks{\buildstamp}}

\markboth{TxHDL Datasheets}{Filmar and Janković: TxHDL Datasheets}

\begin{document}
\maketitle

\section*{How to Read a Datasheet}

Each component is implemented within the synthesizable lowered subset,
functioning simultaneously as cycle-accurate simulation logic and structural
RTL. Individual datasheets define component functionality, generic parameters,
port interface protocols, verification strategies, system instantiations, and
architectural limitations.

Interface tables are extracted systematically rather than transcribed
manually. The build system compiles each component using canonical repository
parameters, generating structural tables directly from the resulting netlists.
Channel ports occupy single table rows displaying payload bit widths; within
synthesized netlists, each channel expands to standard \code{\_data},
\code{\_valid}, and \code{\_ready} signals. Discrete wire ports map directly to
individual nets. Memory entries document data word widths alongside total depth.

Parameterized component families instantiated across varying port configurations
(such as routers multiplexing between two and eight peripherals) share a single
datasheet displaying representative topology from production designs.

% Channels and queues.
% SPDX-License-Identifier: Apache-2.0
% covers: Buffer
\datasheet{Buffer}{The channel as hardware: an elastic buffer of two.}

\dsfacts{\code{lib/parts/src/buffer.rs}}{\code{txhdl\_parts::buffer::Buffer}}%
{\code{//docs:runtime}, \code{//docs:examples}}

\subsection*{Function}

A channel in a simulation is a buffer of two between a sender and a
receiver. \code{Buffer} is that buffer as a netlist, for the places
where two lowered units meet on a channel: a board top, or a design
wired by hand. Its sender side is \code{tx\_data}, \code{tx\_valid} and
\code{tx\_ready}; its receiver side is \code{rx\_data}, \code{rx\_valid}
and \code{rx\_ready}.

\subsection*{Parameters}

\begin{center}\footnotesize
\begin{tabular}{@{}lp{0.7\textwidth}@{}}
\toprule
Parameter & Meaning \\
\midrule
\code{W} & The width of the words the channel holds, in bits. \\
\bottomrule
\end{tabular}
\end{center}

\dstables{Buffer}

\subsection*{Behaviour}

\begin{itemize}
\item \code{tx\_ready} is high when the tail is empty, as the last edge
left it. \code{rx\_valid} is high when the head is full, and
\code{rx\_data} is the head.
\item All three outputs are registers, so there is no combinational
path from one side to the other.
\item At an edge, a take moves the tail into the head, and a push goes
into the head if the head is then empty, else into the tail.
\item With both sides running every cycle, one word passes per cycle.
\end{itemize}

\subsection*{Verification}

\begin{itemize}
\item \code{buffer\_is\_the\_runtime\_channel} in \code{txhdl\_parts}
runs the buffer beside a runtime channel for 200 cycles of pseudorandom
offers and takes, and checks \code{ready}, \code{valid} and the head
after every cycle.
\item \code{ex\_buffer} does the same on a fixed pattern, and the build
simulates \code{Buffer<8>}'s netlist against that run under nvc and
Verilator: \code{//docs:buffer\_sim\_buffer8\_tb\_test} and
\code{//docs:buffer\_vsim\_test}.
\end{itemize}

\subsection*{Used by}

The example \code{ex\_buffer}. A unit of units does not need it, since
its netlist puts a \code{txhdl\_chan} of the same rules between its
children.

\subsection*{Limits}

It holds two words. A consumer that takes in bursts needs \code{Fifo}.

% SPDX-License-Identifier: Apache-2.0
% covers: ChanCdc
\datasheet{ChanCdc}{A channel from one clock to another, with Gray
coded pointers and two flops each way.}

\dsfacts{\code{lib/parts/src/cdc.rs}}{\code{txhdl\_parts::cdc::ChanCdc}}%
{\code{//docs:examples}}

\subsection*{Function}

Two clocks that are not the same clock have no fixed relation, so a
register written by one and read by the other is read at an arbitrary
moment in its life, including while it is changing. \code{ChanCdc}
transfers words across that clock boundary: a channel in on the writing clock, a
channel out on the reading clock, and a memory between them written
by one side and read by the other.

It is one unit with two processes, one per clock. What crosses
between them is only the two pointers, each through two flip-flops of
the clock that reads it.

\subsection*{Parameters}

\begin{center}\footnotesize
\begin{tabular}{@{}lp{0.68\textwidth}@{}}
\toprule
Parameter & Meaning \\
\midrule
\code{T} & The payload data type: a \code{Transaction} and \code{Value} type. \\
\code{AW} & The width of the address, in bits. \\
\code{D} & How many words it holds, which must be \code{1 << AW}. \\
\code{PW} & The width of the pointers, which must be \code{AW + 1}. \\
\code{W} & The writing clock. \\
\code{R} & The reading clock. \\
\bottomrule
\end{tabular}
\end{center}

\code{D} and \code{PW} are stated rather than computed, because an
expression in a const parameter's position needs nightly Rust, which
is the same reason \code{Fifo} takes both \code{AW} and \code{D}. The
tables below use \code{ChanCdc<U<32>, 4, 16, 5, DefaultClock,
ClkPix>}, the crossing of \code{ex\_cdc}: sixteen words of thirty-two
bits.

\dstables{ChanCdc}

\subsection*{Behaviour}

\begin{itemize}
\item \code{wbin} and \code{rbin} are the two pointers in binary, each
on its own clock. \code{wgray} and \code{rgray} are the same pointers
Gray coded, and they are what the other side samples.
\item \code{rgray\_w1} and \code{rgray\_w2} are the reading pointer
caught on the writing clock, and \code{wgray\_r1} and \code{wgray\_r2}
the writing pointer caught on the reading clock. The second of each
pair is the one its side reads.
\item Empty is \code{rgray} equal to \code{wgray\_r2}. Full is
\code{wgray} equal to \code{rgray\_w2} in every bit but the top two.
\item A word is taken when one is offered and it is not full, and
lands in the memory at the writing pointer. A word is offered out
when it is not empty, read from the memory at the reading pointer.
\end{itemize}

The pointers are Gray coded because a binary pointer changes several
bits at once: \code{0111} to \code{1000} is four bits, and a sample
caught among them can read any of sixteen values, most of which the
pointer never held. One bit per step means a sample caught mid-change
reads the value before or the value after and nothing else. Both are
safe, and safe in a particular direction: the reader concludes the
queue is emptier than it is and the writer that it is fuller, so each
errs towards refusing work rather than corrupting it.

The pointers are one bit wider than the address so that full and empty
are distinguishable, which equal pointers of the same width are not.

\subsection*{Verification}

\begin{itemize}
\item \code{ex\_cdc} runs the writing side on a clock of period four
and the reading side on one of period six with a phase of two, so the
edges drift past each other and every relative position happens
somewhere in the run. The build checks the netlist against that run
under nvc and under Verilator:
\code{//docs:cdc\_sim\_cdc\_tb\_test} and
\code{//docs:cdc\_vsim\_test}.
\item Those checks reach inside. Every register is compared by
hierarchical name as well as every port, so both Gray pointers and all
four synchroniser flops are checked tick by tick, and a pointer that
advanced at the wrong moment or a full test off by the wrap bit is a
failure rather than a difference nobody looks at.
\end{itemize}

\subsection*{Limits}

\textbf{Metastability is not modelled and cannot be.} Both sides are
deterministic and the sampling clock reads a register that has
settled, so the two flops are never seen to do anything. What the
co-simulation demonstrates is agreement about the pointer protocol.
The second flop is an argument about silicon, and no simulation in
either language says anything about it. A passing test beside a
synchroniser invites the opposite conclusion, which is why this
paragraph is here rather than in a comment.

\code{D} must be \code{1 << AW} and \code{PW} must be \code{AW + 1},
and nothing checks either.

It does not replace \code{lib/board/chan\_cdc.v}, of which the
flagship holds five instances. That is a separate decision and a
separate change: the hand-written file is proven on hardware and this
one is proven in simulation, and those are not the same claim.

% SPDX-License-Identifier: Apache-2.0
% covers: Fifo
\datasheet{Fifo}{A FIFO of a power of two words, with a channel at each end.}

\dsfacts{\code{lib/parts/src/fifo.rs}}{\code{txhdl\_parts::fifo::Fifo}}%
{\code{//docs:runtime}, \code{//docs:examples}, \code{//docs:station}}

\subsection*{Function}

A channel is a buffer of two. \code{Fifo} is the depth a channel does
not have, for a unit whose consumer takes in bursts or holds off. It
has a channel in, \code{inp}, a channel out, \code{out}, and \code{D}
words between. \code{//docs:station} puts one behind a reservation
station for that reason.

\subsection*{Parameters}

\begin{center}\footnotesize
\begin{tabular}{@{}lp{0.7\textwidth}@{}}
\toprule
Parameter & Meaning \\
\midrule
\code{T} & What it holds: a \code{Transaction} and \code{Value} type. \\
\code{AW} & The width of its pointers, in bits. \\
\code{D} & How many words it holds, which must be \code{1 << AW}. \\
\bottomrule
\end{tabular}
\end{center}

\code{D} is stated rather than computed, because an expression in a
const parameter's position needs nightly Rust. The tables below use
\code{Fifo<U<8>, 2, 4>}, the FIFO of \code{ex\_queue}: four bytes.

\dstables{Fifo}

\subsection*{Behaviour}

\begin{itemize}
\item \code{mem} holds the words. \code{head} is where the oldest word
is, and \code{tail} is where the next word lands.
\item The pointers are equal both when it is empty and when it is
full. \code{full} says which.
\item A word is taken from \code{inp} whenever \code{full} is clear,
and written at \code{tail}, which moves on by one.
\item When it is not empty and \code{out} has room, the word at
\code{head} is sent on \code{out}, and \code{head} moves on by one.
\item A word taken lands in the memory at the end of its step and is
offered the step after. The source states the latency as one cycle.
\item \code{full} clears in a cycle that sends a word. It sets in a
cycle that takes a word and sends none, when the moved \code{tail}
meets \code{head}.
\end{itemize}

\subsection*{Verification}

\begin{itemize}
\item \code{fifo\_keeps\_every\_word\_in\_order} in
\code{txhdl\_parts}, in \code{//lib/parts:parts\_test}, runs
\code{Fifo<U<8>, 2, 4>} against a queue. It offers words at random
for 200 cycles and takes none for the first 40, so the FIFO fills. It
checks every word comes out in order, that \code{full} was seen, and
that it is empty at the end.
\item \code{ex\_queue} bursts nine words into the same FIFO while the
sink takes nothing for eight cycles, as \code{//docs:examples}
describes. The build simulates the netlist, named \code{queue4},
against that run under nvc and Verilator:
\code{//docs:queue\_sim\_queue4\_tb\_test} and
\code{//docs:queue\_vsim\_test}.
\end{itemize}

\subsection*{Used by}

The example \code{ex\_queue}. \code{ex\_fifo} is a different unit: an
asynchronous FIFO of its own, defined in the example.

\subsection*{Limits}

It runs on one clock, \code{DefaultClock}. \code{D} must be
\code{1 << AW}, and nothing checks it. When \code{full} is set it
takes no word, even in a cycle in which it sends one.

% SPDX-License-Identifier: Apache-2.0
% covers: LineBuf
\datasheet{LineBuf}{A line of pixels on the clock that shows them,
filled from a channel and read by column.}

\dsfacts{\code{lib/parts/src/dma.rs}}{\code{txhdl\_parts::dma::LineBuf}}%
{\code{//docs:examples}}

\subsection*{Function}

The far end of a scanout. \code{LineFetch} reads a line from memory
on the bus clock and \code{ChanCdc} transfers the words to the pixel
clock; this is what they arrive in. A word taken lands at the next
place in the line, and the raster asks for a column and gets what is
there.

\subsection*{Parameters}

\begin{center}\footnotesize
\begin{tabular}{@{}lp{0.68\textwidth}@{}}
\toprule
Parameter & Meaning \\
\midrule
\code{LEN} & How many pixels a line holds. \\
\code{AW} & The width of a column index, in bits. \\
\code{C} & The clock it runs on, which is the pixel clock. \\
\bottomrule
\end{tabular}
\end{center}

The tables below use \code{LineBuf<32, 5, ClkPix>}, the buffer of
\code{ex\_scanout}.

\dstables{LineBuf}

\subsection*{Behaviour}

\begin{itemize}
\item A word is taken whenever one is offered. The buffer is the
consumer that the fetch is held up by, so never refusing is what
keeps the line filling.
\item \code{shown} holds the pixel the column named at the last edge,
and \code{pix} is that register.
\item \code{sol} puts the write back to the start of the line, so a
fetch that delivered too few words does not walk the buffer out of
step for ever.
\end{itemize}

\subsection*{Why the output is registered}

A block RAM has an output register, and a design that reads the
memory straight out to a port asks it to settle combinationally
within a pixel. A raster also wants a pixel that is stable for the
whole cycle it is shown rather than one that moves when the column
does. Both say the same thing, so the port is a register and not the
memory.

It costs one cycle of lead, which a raster has, since it knows the
column it will want next.

\subsection*{Verification}

\begin{itemize}
\item \code{ex\_scanout} joins \code{LineFetch}, \code{ChanCdc} and
this over two clocks of period four and period six, and reads
thirty-two pixels back a column at a time. The netlist is checked
against that run under nvc and Verilator:
\code{//docs:scanout\_sim\_linebuf\_tb\_test} and
\code{//docs:scanout\_vsim\_test}.
\item The memory it reads from holds its own address in each word, so
a pixel that came from the wrong burst, crossed badly, or landed out
of step says which place it came from.
\end{itemize}

\subsection*{Limits}

Writing and reading are independent, and nothing here says the fill
is far enough ahead of the beam. What "far enough" means belongs to
the video timing rather than to a buffer. A line that is not filled
in time shows the previous line's pixel at that column, which is a
visible artefact rather than a hang.

One word a pixel, which is what a framebuffer holds and what the
CDC channel transports. It does not pack several pixels to a word.

\code{LEN} and \code{AW} are not checked against each other.

% SPDX-License-Identifier: Apache-2.0
% covers: LineFetch
\datasheet{LineFetch}{A host on the bus that reads a run of memory in
bursts, into a channel.}

\dsfacts{\code{lib/parts/src/dma.rs}}{\code{txhdl\_parts::dma::LineFetch}}%
{\code{//docs:examples}}

\subsection*{Function}

A peripheral that wants a lot of memory should not have the core copy
it. The video peripheral's framebuffer is 19\,200 register writes for
a frame of the board's 160 by 120, and would be 307\,200 for a full
640 by 480. The Ethernet peripheral is the same shape a byte at a
time, and in both the processor's whole job during them is copying.

\code{LineFetch} is the alternative, and it is the half the video and
the MAC both want. Told an address and a count it issues read bursts
and puts the words that come back into a channel. What consumes that
channel is the peripheral's business: for the video it crosses to the
pixel clock through \code{ChanCdc}, and for a MAC it feeds the
transmitter.

\subsection*{Parameters}

\begin{center}\footnotesize
\begin{tabular}{@{}lp{0.68\textwidth}@{}}
\toprule
Parameter & Meaning \\
\midrule
\code{A} & The width of an address, in bits. \\
\code{I} & The width of a burst identifier, in bits. \\
\code{BEATS} & How many beats a burst asks for. \\
\code{WC} & The width of the word count, which says how long a run is. \\
\bottomrule
\end{tabular}
\end{center}

The identifier is a parameter and never a literal, so a link that
widens its identifiers costs this component nothing. The tables below
use \code{LineFetch<32, 2, 16, 16>}, the fetcher of \code{ex\_dma}:
sixteen-beat bursts on a thirty-two bit address.

\dstables{LineFetch}

\subsection*{Behaviour}

\begin{itemize}
\item A run starts when \code{go} is high and nothing is in flight.
\code{base} and \code{words} are read once, there, so a register
written under the engine does not move the run out from under it.
\item A burst goes out when none is in flight and words are still
wanted, asking for \code{BEATS} beats, the rest of the run, or the
words to the next 4\,KB boundary, whichever is fewest. AXI4 forbids a
burst that crosses one, and the board's DDR3 controller refuses it
(issue 1209).
\item A beat is taken when one is offered and the channel out has
room. That room is the backpressure: a consumer that stops taking
stops the bursts.
\item \code{running} is high while a run is unfinished.
\end{itemize}

Two things about the bus are the opposite of what they look like, and
both were written the other way first, and both deadlock.

\emph{The identifier arrives with the burst, not before it.} The host
takes an \code{Issue}, chooses a free identifier and says which on
\code{grant}, so a client sends the burst and catches the identifier
as it comes back. Asking for one and waiting is a deadlock: nothing
is granted until something is issued.

\emph{A read burst is finished by its last beat, not by \code{done}.}
The host puts a write's response on \code{done} and answers a read on
the beat channel. So the identifier goes back as the last beat is
taken, and the last beat is only taken when there is room to give it
back.

\subsection*{Verification}

\begin{itemize}
\item \code{ex\_dma} fetches sixty-four words through sixteen-beat
bursts across a real link, with \code{AxiHost}, \code{AxiPer} and a
memory on the far side. The build checks the netlist against that run
under nvc and Verilator:
\code{//docs:dma\_sim\_linefetch\_tb\_test} and
\code{//docs:dma\_vsim\_test}.
\item \code{a\_burst\_never\_crosses\_a\_4k\_boundary}, in
\code{//lib/parts:parts\_test}, starts a run 32 words below a 4\,KB
boundary with bursts of up to 128 and checks that the first is cut
there, that none crosses one, and that every word still comes, in
order. It fails without the cut.
\item The example asserts rather than prints. Every word arrived;
each word is the one at its own address, since the memory holds its
address in each word, so a word from the wrong burst or out of order
says so; and the run came in under a cycle bound a path spending an
address phase per word could not meet. Sixty-four words in 93 cycles
against a bound of 128.
\item The memory model there answers bursts rather than words. A
model that answered one beat per request would make a burst look like
a word: the words would still all arrive, the cycle count would still
pass, and nothing about bursting would be demonstrated.
\end{itemize}

\subsection*{Choosing \code{BEATS}}

\code{BEATS} is the lever that decides whether real memory latency
matters, and the number that looks like a measurement above does not
say so on its own.

The 93 cycles are measured against a memory model that answers a beat
the cycle after the request: no activation, no CAS, no refresh. So
that figure is a floor and not an estimate. What a real memory adds is
one first-beat latency per burst, and the burst count is
\code{ceil(words / BEATS)}, so the exposed latency is the burst count
multiplied by whatever the controller's latency is. One burst is in
flight at a time, so none of it is hidden.

That makes the parameter, rather than the engine, the answer to a
consumer that is faster than a raster. A frame of 1500 bytes is 375
words, and at gigabit the wire takes about 1200 cycles at 100\,MHz:

\begin{center}\footnotesize
\begin{tabular}{@{}lll@{}}
\toprule
\code{BEATS} & Bursts for a frame & Fits 1200 cycles while \\
\midrule
16 & 24 & latency is under about 27 cycles \\
256 & 2 & latency is under about 327 cycles \\
\bottomrule
\end{tabular}
\end{center}

A realistic DDR3 first-beat latency is tens of cycles, so sixteen-beat
bursts can lose that margin and 256-beat bursts cannot. The board's
scanout learnt the same: in sixteen-beat bursts a 640-word line under
the core's copying and the Ethernet port's sending took 4682 cycles
of the 3175 a line lasts, and every line was late. It fetches in bursts
of 64, \code{LineFetch<32, 2, 64, 16>}, and takes 1206; bursts of 128
took 942, but a load of the core's waiting behind one waited 153 cycles
rather than 94 (issue 1209). A MAC should use a longer one too. Both
are the same component with a different parameter, and neither needs a
second burst in flight, which is what keeps a second identifier, a
tracker and out-of-order answers out of the design.

\subsection*{Limits}

\textbf{It counts the beats of a burst itself and never reads the
returned beat's \code{last}.} So it cannot notice a peripheral that
answers fewer beats than the burst asked for: it takes what arrives
and waits for the rest, holding \code{running} high, with nothing to
say why. Against the Wishbone bridge before issue 471, which answered
any read with one beat, sixty four words became one and the run did
not end.

That is worth knowing before joining it to a peripheral written
elsewhere. The engine assumes AXI4 is obeyed, and a peripheral that
under-delivers is outside what it can detect rather than something it
reports.


\textbf{A burst never crosses a 4\,KB address boundary}, which is
AXI4's rule: the engine cuts a burst at the boundary and starts the
next from it (issue 1209). A run from a base aligned to \code{BEATS}
words never meets the cut, so it costs such a run nothing; it is there
so that a run from anywhere else is still legal. Before it, a caller
that started a run at an arbitrary address with a long \code{BEATS}
could send a burst the interconnect was entitled to answer however it
liked.

One burst is in flight at a time, deliberately. A second would want a
second identifier, a tracker to tell the two apart and an answer for
what happens when they come back out of order, and none of that buys
anything while the consumer is a raster taking one pixel per visible
cycle.

It reads and never writes. \code{done} is drained so that a write
answer arriving from elsewhere cannot block the host, and is
otherwise unused.

It runs on one clock, \code{DefaultClock}, which is the bus clock.
Crossing to a consumer on another clock is \code{ChanCdc}'s job and
not this component's.

Beats are words: \code{size} is fixed at two, four bytes a beat.

% SPDX-License-Identifier: Apache-2.0
% covers: LinePair
\datasheet{LinePair}{The pixel-clock end of a scanout from memory: two
lines, one shown while the other fills, each asked for a line ahead.}

\dsfacts{\code{lib/parts/src/scanout.rs}}%
{\code{txhdl\_parts::scanout::LinePair<LEN, AW, ROWS, TOTAL, STRIDE, C>}}%
{\code{//docs:examples}}

\subsection*{Function}

A video peripheral that shows a frame from memory cannot wait for
memory pixel by pixel: the beam does not stop. \code{LinePair} sits on
the pixel clock beside the raster and holds two lines. The beam reads
one while the other fills, and they change places at the start of each
line. At that start it asks for the line after the one about to be
shown, so each line has the entire duration of the preceding line to
arrive. It asks for nothing until \code{show} is high, and then it asks for a
frame's first line, from the base, before any other. Out of a reset
the next line's address is zero, and the line there is the boot
memory's, whose one-beat answer to a burst hung the fetch for good on
the board (issue 1178). So the picture starts with the first frame
whose last row begins after \code{show} rises.
The bus side of the same scanout is \code{ScanFetch}, and between the
two the request and the words cross clocks through \code{ChanCdc}
(issue 151).

The frame's base address is a register taken at the vertical sync, so
a host that draws into one buffer and shows another flips them with
one write. A column the beam shows before its word has arrived sets
\code{starved}, a sticky bit that \code{clear} resets, so a run can
show that no line starved. The pair keeps the addresses of the lines
it has asked for and not had whole, at most four. A line that gets no
word for two line times in a row sets \code{stuck}, sticky too, and
puts its address on \code{stuck\_at}, so a fetch that hangs says so
rather than leaving a black screen and a clear \code{starved}
(issue 1197).

\subsection*{Parameters}

\begin{center}\footnotesize
\begin{tabular}{@{}lp{0.68\textwidth}@{}}
\toprule
Parameter & Meaning \\
\midrule
\code{LEN} & The words in a visible line, at most \code{1 << AW}. \\
\code{AW} & The width of a column. \\
\code{ROWS} & The visible rows of a frame. \\
\code{TOTAL} & The rows of a frame, blanking included. \\
\code{STRIDE} & The bytes from one line's start in memory to the next's. \\
\code{C} & The pixel clock. \\
\bottomrule
\end{tabular}
\end{center}

The tables below use \code{LinePair<8, 3, 4, 6, 32, ClkPix>}, the pair
of \code{ex\_scanpair}: lines of eight words, four visible rows of
six.

\dstables{LinePair}

\subsection*{Behaviour}

\begin{itemize}
\item A word is taken whenever one is offered, into the line filling
at the next column. The pair is what the fetch is backpressured by,
and a line longer than \code{LEN} loses its tail rather than stalling.
\item At a line's start (\code{line} high) the two lines change
places, and the line filling's count starts again. A word taken in
that same cycle counts in the line just filled.
\item At a line's start a request goes out on \code{req}, if it has
room: the address of the line after the one about to be shown. On the
frame's last row it is the frame's first line, at the base taken at
the last vertical sync.
\item While \code{show} is low nothing is asked for. The frame's first
line is asked for at the start of the frame's last row with
\code{show} high, and the lines after it only once it has been, so the
first request after a reset or a hide is always the base's line.
With \code{show} low the pair disarms: it asks for nothing more, and
does not set \code{starved}. Shown again, it starts from the base's
line as out of a reset.
\item \code{pix} is the word at \code{col} in the line not filling, a
cycle after \code{col} names it, or \code{LATE}, magenta
(\code{0x00ff00ff}), if that word has not arrived (issue 1209). No
picture is drawn in magenta, so a late column shows as itself rather
than as the word the line before left there.
\item A word belongs to the line it was asked for, which the lines owed
say (issue 1233). A line that starts before it has come whole is
\code{late}, and the rest of its words go into the half the beam
reads, so its columns fill in as they come. Lines owed whose time has
passed are counted in \code{behind}, and the rest of their words are
dropped. Neither lands in the half filling, where the next row would
show it as its own.
\item Once the frame's first line has been asked for, a visible column
at or past the words that line had when it was swapped in sets
\code{starved}, until \code{clear}.
\item A line asked for goes after the lines still owed, and no line is
asked for while four are. A word counts against the oldest owed, and
that line's last word takes it off.
\item Two line starts in a row with lines owed and no word come set
\code{stuck}, and \code{stuck\_at} is the oldest owed line's address
then; while \code{stuck} is set it does not change. \code{clear}
clears \code{stuck} and leaves \code{stuck\_at}. The pair goes on
asking.
\end{itemize}

\subsection*{Verification}

\begin{itemize}
\item \code{ex\_scanpair} runs \code{LinePair} on a pixel clock
unrelated to the bus clock, with \code{ScanFetch}, \code{LineFetch},
\code{AxiHost}, \code{AxiPer}, a memory and \code{ChanCdc} both ways.
The memory stalls its beats, standing in for a refresh. Three frames
with a short stall in every line show every pixel as the word at its
place and never set \code{starved}; a fourth with one stall longer than
a line sets it, and it stays set until it is cleared.
\item \code{//docs:scanpair\_sim\_linepair\_linepair\_tb\_test} and
\code{//docs:scanpair\_vsim\_linepair\_test} simulate the netlist
against that run under nvc and Verilator.
\end{itemize}

\subsection*{Used by}

The example \code{ex\_scanpair}, and \code{ScanVideo}, which the
flagship holds in its video slot (issue 151).

\subsection*{Limits}

One line a request, and a line is \code{LEN} words; there is no
horizontal offset or scaling. A column of a line that arrives late
shows \code{LATE} until its word comes, and \code{starved} says that
it happened; a line that never comes sets \code{stuck}. It runs on one clock, \code{C}.

% SPDX-License-Identifier: Apache-2.0
% covers: LineStore
\datasheet{LineStore}{A host on the bus that writes a run of memory
from a channel, in bursts, to a count of bytes.}

\dsfacts{\code{lib/parts/src/dma.rs}}{\code{txhdl\_parts::dma::LineStore}}%
{\code{//docs:examples}}

\subsection*{Function}

The mirror of \code{LineFetch}, and the half an Ethernet receiver
wants. Words arrive on a channel and are written to memory as bursts,
so a peripheral receiving bulk data puts it away itself rather than
interrupting the core per word.

\subsection*{Parameters}

\begin{center}\footnotesize
\begin{tabular}{@{}lp{0.68\textwidth}@{}}
\toprule
Parameter & Meaning \\
\midrule
\code{A} & The width of an address, in bits. \\
\code{I} & The width of a burst identifier, in bits. \\
\code{BEATS} & How many beats a burst asks for. \\
\code{WC} & The width of the byte count. \\
\bottomrule
\end{tabular}
\end{center}

The tables below use \code{LineStore<32, 2, 16, 16>}, the engine of
\code{ex\_dmaw}.

\dstables{LineStore}

\subsection*{Behaviour}

\begin{itemize}
\item A run starts when \code{go} is high and nothing is in flight.
\code{base} and \code{bytes} are read once, there.
\item The beats wanted are \code{(bytes + 3) / 4}: a partial last
word is still a beat, it is just not a whole one.
\item A beat goes out when a word is offered and the beat channel has
room. A producer that stops offering stops the burst, which is the
backpressure in this direction.
\item The last beat of the run is strobed down to the bytes that are
real: \code{0x1}, \code{0x3} or \code{0x7}, and \code{0xf} when the
count is a whole number of words.
\end{itemize}

\subsection*{Two differences from \code{LineFetch}, both the bus}

\emph{A write burst is finished by its response, not by its last
beat.} The host answers a read on the beat channel and a write on
\code{done}, so the identifier goes back when the response arrives,
which may be well after the last beat went out.

\emph{AXI4 puts no identifier on the beats}, so a burst's beats
belong to the oldest address phase outstanding. A second burst issued
before the first one's beats were out would take them. \code{LineFetch}
keeps one burst in flight because a second buys nothing; here it is a
rule rather than a choice.

\subsection*{Verification}

\begin{itemize}
\item \code{ex\_dmaw} writes sixty-four words across a real link with
\code{AxiHost}, \code{AxiPer} and a memory on the far side, then
reads the memory back. The netlist is checked against that run under
nvc and Verilator: \code{//docs:dmaw\_sim\_linestore\_tb\_test} and
\code{//docs:dmaw\_vsim\_test}.
\item The assertions are on what the memory holds afterwards, not on
what the beats looked like: every word at its own address, nothing
past the run touched, and a last word holding one byte of the run
over three bytes that were there before. The memory is poisoned with
\code{0xaaaaaa00} first, because without it a correct strobe and an
ignored one both leave zeros there.
\item Sixty-four words in 94 cycles, under the same bound as the read
direction.
\end{itemize}

\subsection*{Limits}

\textbf{It counts the beats of a burst itself, and a write burst is
finished by its response rather than by its last beat.} So a
peripheral that answers before it has taken every beat leaves the
rest of the burst in the channel, and what follows is neither refused
nor written where it was meant to go. Against the Wishbone bridge
before issue 471 nine words of sixty four landed and the rest of
memory was left as it was.

That failure is silent: a read stall is immediately detectable because
the transaction never completes, whereas a premature write response
appears outwardly identical to successful completion.


The byte count's tail is the low two bits of \code{bytes}, so a run
of zero bytes asks for zero beats and does nothing, and nothing
checks that a caller meant that.

The same 4\,KB boundary rule as \code{LineFetch}: a burst must not
cross one, this does not check it, and \code{BEATS} of 256 is
1\,KB, so a 1\,KB aligned base satisfies it without a special case.

It runs on one clock, \code{DefaultClock}. It writes and never reads.

% SPDX-License-Identifier: Apache-2.0
% covers: NoBeats
\datasheet{NoBeats}{The write-beat channel of a host whose client only
reads, held and never offered on.}

\dsfacts{\code{lib/parts/src/dma.rs}}{\code{txhdl\_parts::dma::NoBeats}}%
{\code{//docs:datasheets}}

\subsection*{Function}

A host takes write beats from its client whether the client writes or
not, and a board joins every channel it makes to a unit: an end that
nothing holds is refused rather than left floating. A fetch engine
never writes, so something has to hold its host's write-beat channel.
\code{NoBeats} holds it and offers nothing on it.

\subsection*{Parameters}

None. The beat is the tree's write beat, \code{W<32, 4>}.

\dstables{NoBeats}

\subsection*{Behaviour}

\begin{itemize}
\item It never offers a beat. In the netlist its valid is low for
ever, which is what a client that never writes would drive.
\item \code{idle} is its one register, and nothing sets it. It is
there only because a unit is a struct of registers; it keeps nothing.
\end{itemize}

\subsection*{Verification}

It has no example and no test of its own; what it does is nothing. It
is in the board's netlist, and the board tests in
\code{cpu/vreteno/tests/board.rs} run with it in place.

\subsection*{Used by}

\code{Board}, as \code{fnobeats}, on the write-beat channel of the
fetch engine's host \code{fhost}.

\subsection*{Limits}

A host whose client does write must not be given one: its writes would
wait for beats that never come.

The mirror for a client that only writes is \code{NoReads}.

% SPDX-License-Identifier: Apache-2.0
% covers: NoReads
\datasheet{NoReads}{The read-data channel terminator for a host client that only performs writes.}

\dsfacts{\code{lib/parts/src/dma.rs}}{\code{txhdl\_parts::dma::NoReads}}%
{\code{//docs:datasheets}}

\subsection*{Function}

This component forms the complementary sink to \code{NoBeats}. A DMA store engine
never reads, so its host read-data channel has no receiver. Because board-level
interconnects wire every instantiated channel to a destination endpoint,
\code{NoReads} terminates that channel and consumes no data.

\subsection*{Parameters}

\begin{center}\footnotesize
\begin{tabular}{@{}lp{0.7\textwidth}@{}}
\toprule
Parameter & Meaning \\
\midrule
\code{I} & The host identifier width: the response type is \code{R<32, I>}.
It specifies no default because the derive macro rejects one (issue 1044). \\
\bottomrule
\end{tabular}
\end{center}

The board instantiates two instances: \code{NoReads<2>} for the Ethernet DMA store
engine, and \code{NoReads<1>} for the SD host engine whose shared host
allocates one bit of transaction identifier. The tables list \code{NoReads<2>}.

\dstables{NoReads}

\subsection*{Behaviour}

\begin{itemize}
\item It never accepts a response beat. In the synthesized netlist, its ready signal
remains permanently deasserted.
\item Deasserting ready introduces no interconnect stall because the client never transmits
read requests on the address channel. A client that issues no reads receives no responses.
\item The \code{idle} signal represents its single register, which remains unset as in
\code{NoBeats}.
\end{itemize}

\subsection*{Verification}

The block has no standalone testbench. It executes as part of the board
synthesis netlist in \code{cpu/vreteno/tests/board.rs}.

\subsection*{Used by}

\code{Board}, twice: as \code{snoreads} (\code{NoReads<2>}) on the read-data channel
of the Ethernet store engine host \code{shost}, and as \code{sdnoreads} (\code{NoReads<1>})
on the SD host engine interconnect \code{sdshost}.

\subsection*{Limits}

A host whose client issues read transactions must not connect to \code{NoReads}.
The first returned response beat would stall indefinitely in the channel, blocking
all subsequent bus traffic.

% SPDX-License-Identifier: Apache-2.0
% covers: ScanCtl, ScanTap
\datasheet{ScanCtl}{The scanout's registers: the frame's base, the bit
that shows the scanout, and the underflow and stuck bits, read and
cleared, with the address of a line that did not come.}

\dsfacts{\code{lib/parts/src/scanout.rs}}%
{\code{txhdl\_parts::scanout::ScanCtl}}%
{\code{//docs:examples}}

\subsection*{Function}

\code{ScanCtl} holds what a program sets and reads of a scanout from
memory (issue 151): where the next frame starts, whether the screen
shows the scanout or the video peripheral's own framebuffer, and
whether a column was ever shown before its word arrived, and whether
a line asked for never came, with that line's address (issue 1197).
Its three outputs are registers, for \code{LinePair} and
\code{VidMux} to read.

The underflow bit, the stuck bit and the address come back from
\code{LinePair} over a channel, as one \code{ScanState}, rather than
on wires. \code{ScanTap}, which this sheet covers too, puts them on it:
it reads \code{LinePair}'s three wires after \code{LinePair} has
driven them and sends them whenever the channel has room. \code{ScanCtl} drives
\code{clear}, which \code{LinePair} reads, so it runs before
\code{LinePair}, and a wire back would be read a step stale in the run
while the netlist reads it fresh. A channel commits at the end of the
step whichever unit ran first.

\subsection*{Register map}

The unit decodes the word from address bits 4 to 2. In
\code{ScanVideo} it sits above \code{1 << SCAN\_BIT}, \code{0x80}, of
the video slot, so on the board its words are at \code{0x3280} to
\code{0x3290}, the HAL's \code{map::SCAN}.

\dsregs{ScanCtl}

A write to \code{clear} of any value clears the underflow and stuck
bits; the word itself is not kept. \code{stuck\_at} is the
first line that did not come since the stuck bit was last clear, and a
clear leaves it until another line does not come.

\dstables{ScanCtl}

\subsection*{Behaviour}

\begin{itemize}
\item \code{base}, \code{mode} and \code{clear} are driven from the
registers \code{fbase}, \code{show} and \code{clr}; \code{clr} is high
for the one cycle after a write to \code{clear}.
\item \code{mode} also starts the fetch: in \code{ScanVideo} it is
\code{LinePair}'s \code{show}, so no line is read from memory until
the show bit is set (issue 1178).
\item The underflow bit is \code{seen}, the stuck bit
\code{seen\_stuck} and the address \code{seen\_at}, the last values
the channel brought. It can read high for a few cycles after a clear, the values
already in the channel, the same in the run and the netlist.
\item Reads and writes are answered the cycle they are taken.
\end{itemize}

\subsection*{Verification}

\code{ex\_scanvideo} sets the base and the scanout bit, reads the
underflow bit clear after a clean frame, set after a frame whose
memory stalls for a line, and clear again after a write to
\code{clear}. \code{//docs:scanvideo\_sim\_scanvideo\_tb\_test} and
\code{//docs:scanvideo\_vsim\_test} simulate it inside
\code{ScanVideo}'s netlist against that run.
\code{//cpu/vreteno:board\_test}'s
\code{a\_line\_that\_never\_comes\_is\_stuck} points the scanout at
memory that never answers and reads the stuck bit set with that
line's address.

\subsection*{Used by}

\code{ScanVideo}. On the board, \code{cpu/vreteno/rust/scanprobe.rs}
uses it through the HAL's \code{Scan}, and says \code{scan stuck} with
the address when the bit is set.

\subsection*{Limits}

One base, taken at each vertical sync; there is no interrupt, so a
program reads the underflow and stuck bits.

% SPDX-License-Identifier: Apache-2.0
% covers: ScanFetch
\datasheet{ScanFetch}{The bus-clock end of a scanout from memory: a line
request taken and \code{LineFetch} started on it.}

\dsfacts{\code{lib/parts/src/scanout.rs}}%
{\code{txhdl\_parts::scanout::ScanFetch<A, WC, LEN>}}%
{\code{//docs:examples}}

\subsection*{Function}

\code{LinePair}, on the pixel clock, asks for each line a line ahead;
the request crosses to the bus clock through \code{ChanCdc}, and
\code{ScanFetch} is what takes it there. It hands the line's address
and length to \code{LineFetch}'s \code{base}, \code{words} and
\code{go}, and takes the next request only once that fetch has started
and finished (issue 151).

\subsection*{Parameters}

\begin{center}\footnotesize
\begin{tabular}{@{}lp{0.68\textwidth}@{}}
\toprule
Parameter & Meaning \\
\midrule
\code{A} & The width of an address, in bits. \\
\code{WC} & The width of \code{LineFetch}'s word count. \\
\code{LEN} & The words in a line, the count every fetch asks for. \\
\bottomrule
\end{tabular}
\end{center}

The tables below use \code{ScanFetch<32, 16, 8>}, the fetcher's end of
\code{ex\_scanpair}.

\dstables{ScanFetch}

\subsection*{Behaviour}

\begin{itemize}
\item \code{at}, \code{count} and \code{start} go to \code{LineFetch}'s
\code{base}, \code{words} and \code{go}: the line's address, \code{LEN},
and a pulse of one cycle.
\item A request is taken from \code{lines} only when no fetch is
running, none is about to start, and none has been started that has
not yet said it is running.
\item The request is cut to \code{A} bits.
\end{itemize}

\subsection*{Verification}

\begin{itemize}
\item \code{ex\_scanpair}, as \code{LinePair}'s sheet says, with this
on the bus clock between the crossing and \code{LineFetch}.
\item \code{//docs:scanpair\_sim\_scanfetch\_scanfetch\_tb\_test} and
\code{//docs:scanpair\_vsim\_scanfetch\_test} simulate the netlist
against that run under nvc and Verilator.
\end{itemize}

\subsection*{Used by}

The example \code{ex\_scanpair}. Nothing on a board holds it yet: the
join with the video peripheral is the rest of issue 151.

\subsection*{Limits}

One line at a time: the next request waits for the fetch before it to
finish, so the pixel side must ask no faster than a line's fetch takes.
It runs on one clock, \code{DefaultClock}.

% SPDX-License-Identifier: Apache-2.0
% covers: ScanVideo, ScanTap, VidMux, txhdl_parts::hdmi::Raster
\datasheet{ScanVideo}{The video peripheral with a scanout from memory
beside it, programmed over one AXI-Lite slot.}

\dsfacts{\code{lib/parts/src/scanout.rs}}%
{\code{txhdl\_parts::scanout::ScanVideo<HV, HFP, HSW, HBP, VV, VFP,
VSW, VBP, SHIFT, AW, TOTAL, STRIDE>}}%
{\code{//docs:examples}}

\subsection*{Function}

\code{ScanVideo} is what the flagship's video slot holds (issue 151): the
video peripheral \code{Hdmi}, unchanged, and beside it a scanout that
shows a frame from memory, with a multiplexer that puts one picture or
the other on the pins. It is a unit of seven, on the pixel clock:

\begin{itemize}
\item \code{LiteSplit} divides the slot by address bit 7: \code{Hdmi}'s
words below \code{0x80}, \code{ScanCtl}'s from there.
\item \code{Raster}, in \code{lib/parts/src/hdmi.rs}, counts the beam
again with \code{Hdmi}'s own functions and parameters, for
\code{LinePair}. It drives the column, the visible part, the line and
frame starts and the row from registers. It is the scanout's raster,
not Razboj's rasteriser, which has a sheet of its own under the same
name (issue 1055).
\item \code{LinePair} holds two lines and asks for each a line ahead,
on \code{req}; their words arrive on \code{words}. Both channels
leave the unit for the bus clock. \code{ScanCtl}'s show bit, which
picks the scanout's picture at \code{VidMux}, is also
\code{LinePair}'s \code{show}, so the fetch starts only once the
scanout is shown (issue 1178).
\item \code{ScanTap} puts \code{LinePair}'s underflow bit on a channel
back to \code{ScanCtl}.
\item \code{VidMux} puts \code{Hdmi}'s pixel or the scanout's, a
\code{0x00RRGGBB} word, on the pins, with \code{Hdmi}'s syncs either
way.
\end{itemize}

The children run in an order that puts every wire's driver before its
reader: \code{Raster} and \code{ScanCtl}, then \code{LinePair}, then
\code{ScanTap} and \code{Hdmi}, then \code{VidMux}. The one path the
other way, the underflow bit, is a channel. Counting the beam twice is
what keeps that order possible, and \code{ex\_scanvideo} checks the two
counts agree on every cycle.

\subsection*{Parameters}

\begin{center}\footnotesize
\begin{tabular}{@{}lp{0.68\textwidth}@{}}
\toprule
Parameter & Meaning \\
\midrule
\code{HV} to \code{VBP}, \code{SHIFT} & \code{Hdmi}'s: the video mode, and
a framebuffer pixel's size on the screen. \\
\code{AW} & The width of a column; \code{HV} is at most \code{1 << AW}. \\
\code{TOTAL} & The rows of a frame, \code{VV + VFP + VSW + VBP}; the
\code{Default} refuses one that disagrees. \\
\code{STRIDE} & The bytes from one line to the next in memory. \\
\bottomrule
\end{tabular}
\end{center}

The tables below use the flagship's: 640 by 480 at 60~Hz, a
framebuffer pixel of four by four, ten bits of column, 525 rows and
4096 bytes a line, as \code{hdmi/src/netlist.rs} writes \code{scan\_video}.
A line in memory is Razboj's row of 1024 words, of which the first 640
are shown (issue 985).

\dstables{ScanVideo}

\subsection*{Verification}

\begin{itemize}
\item \code{ex\_scanvideo} runs the whole unit with \code{LineFetch},
\code{ScanFetch} and a memory that stalls on command. It paints a
pixel through \code{Hdmi} and sees it, shows two frames of scanout and
checks every visible pixel, reads the underflow bit over the slot as a
stall sets it and a write clears it, and asserts that \code{Raster}'s
count is \code{Hdmi}'s on every cycle, across a reset.
\item \code{//docs:scanvideo\_sim\_scanvideo\_tb\_test} and
\code{//docs:scanvideo\_vsim\_test} simulate its netlist against that
run under nvc and Verilator.
\item On the board, \code{scanprobe} shows a frame from the DDR3 and
reads the underflow bit idle and under load; both read clear
(\code{docs/board-checks.md}, PR 1041's session).
\end{itemize}

\subsection*{Used by}

\code{flagship/board/flagship.v}, in the video slot at \code{0x3200},
with \code{chan\_cdc} transferring \code{req} and \code{words} across
clock domains to the board's \code{ScanFetch} and \code{LineFetch} engines.

\subsection*{Limits}

The module operates in a single video mode determined by its synthesis
parameters. Scanout timing and framebuffer addressing share common sync
generators; scanout resolution matches the native video mode without
scaling. The underflow status bit propagates to the register interface
several cycles after \code{LinePair} asserts it due to domain crossing.

% SPDX-License-Identifier: Apache-2.0
% covers: Station2, Station3, Station4, Station5, Station6, Station7, Station8, Station9, Station10
\datasheet{Station}{Reservation stations of two to ten inputs: operands gathered by tag, sent as one line.}

\dsfacts{\code{lib/parts/src/station.rs}; generator \code{station\_text} in \code{lib/macros/lib.rs}}%
{\code{txhdl\_parts::station::Station2} to \code{Station10}}%
{\code{//docs:station}, \code{//docs:runtime}, \code{//docs:examples}}

\subsection*{Function}

A reservation station gathers the operands of an operation that arrive
on separate channels, out of order, and fires when all of them have.
Each input is a channel of \code{Tagged<TB, T>}: a tag of \code{TB}
bits above a value. The station has a line per tag value, with a cell
per input. An input lands in its cell of the line its tag names.

A line whose every cell holds is sent on the one output as a
\code{Line}\textit{N}: the tag, and a value per input, \code{v0}
upwards. \code{station!(StationN,
N)} writes each station with its line struct beside it, since the
lowering reads a body and not a loop over inputs. \code{station.rs}
writes \code{Station2} to \code{Station10}; the macro itself accepts 2
to 16.

\subsection*{Parameters}

\begin{center}\footnotesize
\begin{tabular}{@{}lp{0.7\textwidth}@{}}
\toprule
Parameter & Meaning \\
\midrule
\code{TB} & The tag width, in bits. It is two at least. \\
\code{L} & The count of lines, which is \code{1 << TB}. \\
\code{T}\textit{i} & The value type of input \textit{i}, for \textit{i}
from 0 to \textit{N} less one. The types may differ. \\
\bottomrule
\end{tabular}
\end{center}

\code{L} is stated rather than computed, because a computed width in a
type needs nightly Rust. The tag has two bits at least, since a one-bit
value is a single wire in VHDL and cannot index a line. The tables
below use \code{Station3<4, 16, U<2>, U<16>, U<16>>}, the station of
\code{ex\_compose}: sixteen lines, an operation and two operands.

\dstables{Station}

\subsection*{Behaviour}

Each input \textit{i} has \code{mem}\textit{i}, its cells, a word per
line, and \code{occ}\textit{i}, a bit per line that says which cells
hold. Each cycle:

\begin{itemize}
\item An input's cell is free when its tag's bit of the input's mask
is clear.
\item Taking an input completes its line when, for every other input,
that line's cell holds, or that input offers the same tag now and its
own cell is free.
\item The line sent is the completing input's of lowest index, if any
input completes and the output has room. At most one line leaves per
cycle.
\item An input is taken when its cell is free and taking it completes
nothing, or completes the line being sent. An input's \code{ready} is
that take.
\item So an input never waits on another input's cell. It waits on its
own cell, on the output's room, or while an input of lower index
completes a different line in the same cycle.
\item A taken input writes its cell and sets its bit. The line sent
clears its bit in every mask.
\item The values sent come from the cells. The value of an input that
completes the line in this cycle goes straight through instead.
\end{itemize}

\subsection*{Verification}

\begin{itemize}
\item Two tests in \code{txhdl\_parts}, in
\code{//lib/parts:parts\_test}. \code{station2\_follows\_the\_rule} runs
\code{Station2<2, 4, U<8>, U<16>>} against the rule written with loops
for 400 cycles of random offers and takes. It checks both masks after
every cycle and every line the sink takes.
\code{station3\_lowest\_input\_wins} sets up two lines that complete in
one cycle, and checks that the lower input's line goes first.
\item \code{ex\_station} runs a \code{Station2<2, 4, U<8>, U<16>>}
beside \code{station.v}, the same station written by hand in Verilog,
and compares every line the two send. The build simulates the
station's netlist against that run under nvc and Verilator:
\code{//docs:station\_sim\_station\_tb\_test} and
\code{//docs:station\_vsim\_test}.
\item \code{ex\_compose} lowers \code{Top}, which holds the
\code{Station3} of the tables and the stream ALU as children. The build
simulates the netlist of the top against that run:
\code{//docs:compose\_sim\_top\_tb\_test} and
\code{//docs:compose\_vsim\_test}.
\end{itemize}

\subsection*{Used by}

\code{Station2} in the example \code{ex\_station}. \code{Station3} in
the example \code{ex\_compose}, in front of the stream ALU, whose
requests are \code{Line3} values.

\subsection*{Limits}

A station is not an arbiter. If an input's next offer names a cell it
already holds, it waits until that line completes. If the other inputs
never offer that tag, it waits for ever, as \code{//docs:station}
states. Each input must, in time, see every tag its lines wait on. One
line leaves per cycle at most, and a line waits for room on the output;
\code{//docs:station} puts a \code{Fifo} behind the station where the
consumer takes in bursts.

% The AXI link.
% SPDX-License-Identifier: Apache-2.0
% covers: AxiHost
\datasheet{AxiHost}{The host side of an AXI4 link: bursts out, identifiers handed out, answers back.}

\dsfacts{\code{lib/parts/src/bus/axi.rs}}{\code{txhdl\_parts::bus::axi::AxiHost}}%
{\code{//docs:axi}, \code{//docs:noc}}

\subsection*{Function}

An AXI4 link in \code{txhdl\_parts} has a tracker at each end, between
a client and the five AXI channels. \code{AxiHost} is the tracker on
the host side. A host client gives it an \code{Issue} on \code{issue}:
an address phase with no identifier, and a bit that says whether it is
a read. \code{AxiHost} gives the burst an identifier, sends it on
\code{ar} or \code{aw}, and tells the client the identifier on
\code{grant}.

It also passes the client's write beats from \code{wbeat} to \code{w},
turns each \code{B} beat into a \code{Done} on \code{done}, and passes
each read beat from \code{r} to \code{rdata}. The client hands an
identifier back on \code{release} when it has taken that burst's
answer. \code{axi}, \code{axi\_units} and \code{axi\_to\_unit} make the
channels it runs on. Its ports are named types, \code{HostIn} and
\code{HostOut}, so a design may supply a unit of its own instead.

\subsection*{Parameters}

\begin{center}\footnotesize
\begin{tabular}{@{}lp{0.7\textwidth}@{}}
\toprule
Parameter & Meaning \\
\midrule
\code{A} & The address width, in bits. \\
\code{D} & The data width, in bits. \\
\code{S} & The write strobe width, one bit per byte lane, so \code{D / 8}. \\
\code{I} & The identifier width, in bits. \\
\code{NIDS} & How many identifiers it hands out, so \code{1 << I}. It is
the width of \code{busy}, and the most bursts in flight at once. \\
\bottomrule
\end{tabular}
\end{center}

\code{S} and \code{NIDS} are stated rather than computed, because an
expression in a const parameter's position needs nightly Rust. The
source states that a link whose \code{S} and \code{NIDS} disagree with
its \code{D} and \code{I} will not behave. The tables below use
\code{AxiHost<32, 32, 4, 2, 4>}, the tracker of Vreteno's link and of
its board.

\dstables{AxiHost}

\subsection*{Behaviour}

\begin{itemize}
\item At each rising edge it looks at the issue at the head of
\code{issue}. The identifier that burst would take is \code{next}.
\item The burst goes out when that identifier's bit of \code{busy} is
clear, \code{grant} has room, and the address channel it needs has
room: \code{ar} for a read, \code{aw} for a write.
\item A burst that goes out is sent with \code{next} as its \code{id}.
Its \code{addr}, \code{len}, \code{size}, \code{burst}, \code{lock},
\code{cache}, \code{prot}, \code{qos} and \code{region} are the
issue's.
\item In that cycle \code{grant} sends the identifier, \code{next}
moves on by one, and the identifier's bit of \code{busy} is set. At
most one burst goes out per cycle.
\item In a cycle in which the identifier \code{next} names is
outstanding, \code{next} moves on by one as well. So the identifiers
go round, and a burst waits only while every identifier is
outstanding; it finds a free one within \code{NIDS} cycles.
\item \code{release} is taken whenever it offers, and it clears the
bit of the identifier it names. An identifier is therefore outstanding
from its address phase until the client has taken its answer.
\code{//docs:axi} explains why freeing it on the response would be
wrong.
\item A write beat passes from \code{wbeat} to \code{w} unchanged when
\code{w} has room. AXI4 puts no identifier on \code{w}, so the source
says a client sends a burst's beats before it issues the next burst.
\item A \code{B} beat becomes a \code{Done} of the same \code{id} and
\code{resp} when \code{done} has room. A read beat passes from
\code{r} to \code{rdata} unchanged when \code{rdata} has room.
\end{itemize}

\subsection*{Verification}

\begin{itemize}
\item Four tests in \code{bus::axi} run \code{AxiHost<16, 32, 4, 2, 4>}
and \code{AxiPer<16, 32, 4, 2>} with a host client and a peripheral
client on the two ends:
\begin{itemize}
\item \code{a\_read\_reads\_what\_the\_write\_wrote};
\item \code{answers\_come\_back\_out\_of\_order};
\item \code{a\_transaction\_is\_answered\_by\_another\_process};
\item \code{every\_burst\_is\_answered\_and\_reads\_agree\_with\_writes},
which issues 24 bursts of one to four words to a peripheral that
\code{serve} runs with four slots.
\end{itemize}
\item \code{bus::axi::fuzz} runs the two trackers with a relay that
stalls at random in every channel, a seeded client that holds up to
four bursts of one to sixteen beats and awaits them in its own order,
and a seeded peripheral that answers out of order and with errors. It
checks that every burst is answered once, that reads read the writes,
that \code{last} falls on the final beat, that no identifier is handed
out while its answer stands uncollected, and that no run deadlocks, on
24 seeds; \code{//lib/parts:axi\_fuzz\_test}, manual, runs more.
\item The tests of \code{bus::axi::sim} put the same two trackers in
front of the simulated RAM:
\begin{itemize}
\item \code{a\_burst\_reads\_what\_a\_burst\_wrote};
\item \code{a\_strobe\_covers\_the\_lanes\_it\_says};
\item \code{a\_fixed\_burst\_stays\_on\_one\_word};
\item \code{an\_address\_past\_the\_end\_is\_answered};
\item \code{an\_image\_is\_loaded\_and\_read\_back}.
\end{itemize}
All of these run in \code{//lib/parts:parts\_test}.
\item \code{ex\_axi} issues three bursts and has them answered out of
order. The build simulates the netlist of \code{AxiHost<16, 32, 4, 2,
4>} against that run under nvc and Verilator:
\begin{itemize}
\item \code{//docs:axi\_sim\_axi\_host\_axi\_host\_tb\_test};
\item \code{//docs:axi\_vsim\_axi\_host\_test}.
\end{itemize}
\item The tests of the router, the AXI-Lite bridge, the Wishbone bridge
and the network on chip in \code{txhdl\_parts} run it in simulation.
So do \code{//ddr3:ddr3\_test}, \code{//gpu/razboj:razboj\_test},
\code{//soc:demo\_test} and \code{//cpu/vreteno:lockstep\_test}.
\end{itemize}

\subsection*{Used by}

The examples \code{ex\_axi}, \code{ex\_axi\_serve}, \code{ex\_axi\_router},
\code{ex\_axi\_lite}, \code{ex\_axi\_reg}, \code{ex\_axi\_wb},
\code{ex\_eth} and \code{ex\_hdmi}. Vreteno's run in
\code{cpu/vreteno/src/run.rs} and \code{cpu/vreteno/src/main.rs}, and
the board's design, \code{Board}, as its field \code{host}. The
four-corner system in \code{soc/src/lib.rs}, once for the core and
once for Razboj. Razboj's simulation in \code{gpu/razboj/src/sim.rs}.

\subsection*{Limits}

At most \code{NIDS} bursts are in flight at once. A client that never
releases an identifier stops the unit when the turn reaches that
identifier. It sends \code{lock}, \code{cache}, \code{qos} and
\code{region} as the client gave them and acts on none of them; nothing
in \code{bus::axi} makes an exclusive access. It does not check
\code{S} and \code{NIDS} against \code{D} and \code{I}.

% SPDX-License-Identifier: Apache-2.0
% covers: AxiPer
\datasheet{AxiPer}{The peripheral side of an AXI4 link: two address channels merged into one stream of requests.}

\dsfacts{\code{lib/parts/src/bus/axi.rs}}{\code{txhdl\_parts::bus::axi::AxiPer}}%
{\code{//docs:axi}, \code{//docs:vreteno}, \code{//docs:noc}}

\subsection*{Function}

\code{AxiPer} is the tracker at the peripheral end of an AXI4 link. It
merges \code{aw} and \code{ar} into one stream of \code{PerReq} on
\code{req}, each with a bit that says whether it is a read. It passes
the write beats from \code{w} to \code{wd}. It turns the client's
\code{Answer} on \code{ans} into \code{B} beats on \code{b}, and
passes the client's read beats from \code{rb} to \code{r}.

It allocates no identifiers. A response leaves with the identifier the
client gave it. Its ports are named types, \code{PerIn} and
\code{PerOut}, so a design may supply a unit of its own instead.
\code{//docs:axi} notes that a peripheral may also read the five AXI
channels directly and leave \code{AxiPer} out.

\subsection*{Parameters}

\begin{center}\footnotesize
\begin{tabular}{@{}lp{0.7\textwidth}@{}}
\toprule
Parameter & Meaning \\
\midrule
\code{A} & The address width, in bits. \\
\code{D} & The data width, in bits. \\
\code{S} & The write strobe width, one bit per byte lane, so \code{D / 8}. \\
\code{I} & The identifier width, in bits. \\
\bottomrule
\end{tabular}
\end{center}

There is no \code{NIDS}, since this end hands out no identifiers. The
tables below use \code{AxiPer<32, 32, 4, 2>}, the tracker in front of
Vreteno's data memory and timer.

\dstables{AxiPer}

\subsection*{Behaviour}

\begin{itemize}
\item At most one request leaves per cycle, and only when \code{req}
has room.
\item \code{rr} says which address channel goes first: zero for
\code{ar}. When both offer, the one that goes first is taken. When
only one offers, that one is taken.
\item Taking a read sets \code{rr} to one, and taking a write sets it
to zero. So two channels that both stream take turns, and neither
starves the other.
\item The request has \code{read} high when it came from \code{ar}.
Its \code{id}, \code{addr}, \code{len}, \code{size}, \code{burst},
\code{lock}, \code{cache}, \code{prot}, \code{qos} and \code{region}
are the address phase's.
\item A write beat passes from \code{w} to \code{wd} unchanged when
\code{wd} has room, independently of the requests.
\item An answer on \code{ans} becomes a \code{B} of the same \code{id}
and \code{resp} when \code{b} has room. A read beat passes from
\code{rb} to \code{r} unchanged when \code{r} has room.
\end{itemize}

\subsection*{Verification}

\begin{itemize}
\item Four tests in \code{bus::axi} run \code{AxiHost<16, 32, 4, 2, 4>}
and \code{AxiPer<16, 32, 4, 2>} with a client on each end:
\begin{itemize}
\item \code{a\_read\_reads\_what\_the\_write\_wrote};
\item \code{answers\_come\_back\_out\_of\_order};
\item \code{a\_transaction\_is\_answered\_by\_another\_process};
\item \code{every\_burst\_is\_answered\_and\_reads\_agree\_with\_writes}.
\end{itemize}
\item \code{bus::axi::fuzz} runs the two trackers with a relay that
stalls at random in every channel and a seeded client and peripheral
on the ends, and checks every burst answered once, reads against
writes, \code{last} on the final beat, and no deadlock, on 24 seeds;
\code{//lib/parts:axi\_fuzz\_test}, manual, runs more.
\item The tests of \code{bus::axi::sim} put the two trackers in front
of the simulated RAM:
\begin{itemize}
\item \code{a\_burst\_reads\_what\_a\_burst\_wrote};
\item \code{a\_strobe\_covers\_the\_lanes\_it\_says};
\item \code{a\_fixed\_burst\_stays\_on\_one\_word};
\item \code{an\_address\_past\_the\_end\_is\_answered};
\item \code{an\_image\_is\_loaded\_and\_read\_back}.
\end{itemize}
All of these run in \code{//lib/parts:parts\_test}.
\item \code{ex\_axi} has three bursts answered out of order. The build
simulates the netlist of \code{AxiPer<16, 32, 4, 2>} against that run
under nvc and Verilator:
\begin{itemize}
\item \code{//docs:axi\_sim\_axi\_per\_axi\_per\_tb\_test};
\item \code{//docs:axi\_vsim\_axi\_per\_test}.
\end{itemize}
\item The tests of the router, the Wishbone bridge and the network on
chip in \code{txhdl\_parts} run it in simulation. So do
\code{//ddr3:ddr3\_test}, \code{//gpu/razboj:razboj\_test},
\code{//soc:demo\_test} and \code{//cpu/vreteno:lockstep\_test}.
\end{itemize}

\subsection*{Used by}

The examples \code{ex\_axi}, \code{ex\_axi\_serve},
\code{ex\_axi\_router}, \code{ex\_axi\_reg} and \code{ex\_axi\_wb}.
Vreteno's run in \code{cpu/vreteno/src/run.rs} and
\code{cpu/vreteno/src/main.rs}, in front of the data memory and the
timer. The board's design, \code{Board}, as its fields \code{pdmem},
\code{ptimer}, \code{prom} and \code{pflash}. The four-corner system in
\code{soc/src/lib.rs}, and Razboj's simulation in
\code{gpu/razboj/src/sim.rs}.

\subsection*{Limits}

It does not count beats, look at \code{last}, or match write beats to
requests. A client does that: \code{Per::accept} gathers a write's
beats after its request. It does not check that an answer's identifier
belongs to an open request. It passes \code{lock}, \code{cache},
\code{qos} and \code{region} to the client and acts on none of them.

% SPDX-License-Identifier: Apache-2.0
% covers: Router
\datasheet{Router}{The AXI4 router: one host and N peripherals, decoded by an address map.}

\dsfacts{\code{lib/parts/src/bus/router.rs}}%
{\code{txhdl\_parts::bus::router::Router}}%
{\code{//docs:axi}, \code{//docs:vreteno}, \code{//docs:noc}}

\subsection*{Function}

A router puts several peripherals behind one host. It sits in the five
channels of an AXI4 link, so nothing on either side changes when it is
put there. It takes the host's \code{aw}, \code{ar} and \code{w}, and
gives each peripheral \textit{i} its own, element \textit{i} of the
arrays \code{aws}, \code{ars} and \code{ws}. It takes each peripheral's
answers from \code{bs} and \code{rs}, and gives the host one \code{b}
and one \code{r}. In the netlist element \textit{i} of an array port is
a port of its own, \code{aws\_}\textit{i} and so on.

One unit serves every count of peripherals: the peripheral sides are
arrays of \code{N}, and the lowering unrolls the loops over them when
\code{lowered} runs (issue 500). Each peripheral's range is its entry
of the map, a type \code{M} implementing \code{txhdl::map::AddrMap<N>}
whose constant \code{RANGES} is \code{N} pairs of a base and a mask,
read the same way (issue 593). A burst whose address is in no range is
answered \code{DecErr} by the router itself. The router was a family of
units, \code{Router2} to \code{Router8}, written by \code{router!},
until \#648.

\subsection*{Parameters}

\begin{center}\footnotesize
\begin{tabular}{@{}lp{0.7\textwidth}@{}}
\toprule
Parameter & Meaning \\
\midrule
\code{N} & The number of peripherals. \\
\code{M} & The address map: a type implementing \code{AddrMap<N>},
whose \code{RANGES[i]} is peripheral \textit{i}'s base and mask. It
holds no signal; the router keeps it as \code{PhantomData}. \\
\code{A} & The address width, in bits. \\
\code{D} & The data width, in bits. \\
\code{S} & The write strobe width, one bit per byte lane, so \code{D / 8}. \\
\code{I} & The identifier width, in bits. \\
\bottomrule
\end{tabular}
\end{center}

\subsection*{Address map}

An address \code{a} is peripheral \textit{i}'s when
\code{(a \& mask) == base} for \code{RANGES[i]} and for no range before
it: the first range that matches decides. The router decodes the
address of each address phase. Every address in no range is a hole.
The tables below are generated for the board's router,
\code{BoardRouter} in \code{cpu/vreteno/src/board.rs}, which is
\code{Router<8, BoardMap, 32, 32, 4, 4>} with \code{BoardMap} beside
it:

\begin{center}\footnotesize
\begin{tabular}{@{}llll@{}}
\toprule
Peripheral & Base & Mask & On the board \\
\midrule
0 & \code{0x1000} & \code{0xffff\_f000} & The data memory \\
1 & \code{0x0200\_0000} & \code{0xffff\_0000} & The timer \\
2 & \code{0x3000} & \code{0xffff\_f000} & The AXI-Lite bridge to the
serial port and its neighbours \\
3 & \code{0x4000\_0000} & \code{0xc000\_0000} & The DDR3 memory \\
4 & \code{0x0c00\_0000} & \code{0xfc00\_0000} & The interrupt controller \\
5 & \code{0x0000\_0000} & \code{0xffff\_f000} & The boot memory \\
6 & \code{0x1000\_0000} & \code{0xffff\_0000} & The debug module \\
7 & \code{0x2000\_0000} & \code{0xff00\_0000} & The configuration flash,
read through a window \\
\bottomrule
\end{tabular}
\end{center}

\dstables{Router}

\subsection*{Behaviour}

\begin{itemize}
\item \textbf{Write address.} The head of \code{aw} is taken when no
write burst is running (\code{wbusy} clear) and the first peripheral
whose range matches has room on its \code{aws} port, or, for a hole,
when no write \code{DecErr} is owed. It is sent to that peripheral
alone.
\item Taking it sets \code{wbusy}, and records the peripheral in
\code{wsel}, a bit each, and the identifier in \code{wid}.
\code{whole} records that no range matched.
\item \textbf{Write beats.} While \code{wbusy} is set, beats pass from
\code{w} to the peripheral in \code{wsel} when it has room. A burst to
a hole has no bit in \code{wsel}, and its beats are taken and dropped.
The beat marked \code{last} clears \code{wbusy}.
\item So a second write address phase, to any peripheral, waits for
the first burst's last beat. AXI4 puts no identifier on \code{w}, and
this keeps two bursts' beats apart.
\item \textbf{Write to a hole.} When its last beat has gone, the router
sets \code{berr} and \code{berr\_id}. It sends a \code{B} of
\code{DecErr} with that identifier when \code{b} has room. That answer
goes before any peripheral's.
\item \textbf{Read address.} The head of \code{ar} is taken when the
first peripheral whose range matches has room on its \code{ars} port,
or, for a hole, when \code{rerr} is clear. Several reads may be
outstanding at once.
\item \textbf{Read to a hole.} Taking it sets \code{rerr}, keeps its
identifier in \code{rerr\_id}, and sets \code{rerr\_left} to its
\code{len} plus one. The router then sends that many beats of
\code{DecErr} and zero data on \code{r}, the last marked \code{last}.
\item \textbf{Read beats.} When no read burst is going up
(\code{rbusy} clear), the router's own error burst goes first. Else the
lowest peripheral that offers a beat is chosen.
\item A peripheral's beat that goes up without \code{last} sets
\code{rbusy} and records the peripheral in \code{rsel}. Its beats then
go up alone until its \code{last}, so \code{last} still marks the end
of one burst.
\item \textbf{Write responses.} When the router owes no \code{DecErr},
the lowest peripheral that offers a \code{B} sends it to the host. Each
response keeps the identifier the peripheral gave it.
\end{itemize}

\subsection*{Verification}

\begin{itemize}
\item Two tests in \code{bus::router} run \code{Router<2, TwoMap, 16,
32, 4, 3>} between a host tracker and two peripheral trackers, in
\code{//lib/parts:parts\_test}. \code{each\_burst\_reaches\_its\_range}
reads from both ranges and reads three beats from a hole, and checks
the hole's answer is \code{DecErr} in three beats.
\code{write\_beats\_do\_not\_interleave} writes three beats to each
peripheral, one burst after the other, and checks each peripheral got
its own burst whole and in order.
\item \code{ex\_axi\_router} issues five bursts before collecting any:
a read to each of two peripherals, a write to the third, and a read and
a write to a hole. The build simulates the netlist of \code{Router<3,
ThreeMap, 16, 32, 4, 3>} against that run under nvc and Verilator:
\code{//docs:axi\_router\_sim\_axi\_router\_tb\_test} and
\code{//docs:axi\_router\_vsim\_test}.
\item Vreteno's run lowers its router, \code{Router<3, RunMap, 32, 32,
4, 2>} in \code{cpu/vreteno/src/main.rs}, and the build simulates that
netlist against the run: \code{//docs:vreteno\_sim\_router\_router\_tb\_test}
and \code{//docs:vreteno\_vsim\_router\_test}.
\code{//cpu/vreteno:lockstep\_test} runs the same router in simulation.
\item \code{four\_peripherals\_share\_one\_node} in \code{bus::noc} puts
a \code{Router<4, NocMap, ..>} behind a node of the network, with a
memory in each range. Two hosts write and read every peripheral and
read a hole.
\end{itemize}

\subsection*{Used by}

\code{Router<3, RunMap, ..>} in Vreteno's run and its lockstep and
pair tests; \code{Router<4, RunMap, ..>}, with the boot memory, in
\code{cpu/vreteno/src/run.rs}; \code{Router<8, BoardMap, ..>} on the
board; \code{Router<4, NocMap, ..>} in the network's test;
\code{Router<3, ThreeMap, ..>} in \code{ex\_axi\_router}; and
\code{Router<2, TwoMap, ..>} in the router's own tests.

\subsection*{Limits}

It serves one host. Several hosts contending for one peripheral are not
handled here, as \code{//docs:axi} states. Write bursts pass one at a
time, whichever peripherals they go to. The merge of responses is by
fixed priority, lowest peripheral first. A mask compares only its own
bits: on an address wider than 16 bits, a mask of \code{0xf000} also
matches its base plus every multiple of \code{0x10000}. Nothing checks
that ranges do not overlap; an address in two ranges goes to the first.
The router decodes on the address alone, and ignores \code{region} and
\code{qos}.

% SPDX-License-Identifier: Apache-2.0
% covers: Arbiter, Arbiter2, Arbiter3, Arbiter4, Arbiter5, Arbiter6, Arbiter7, Arbiter8
\datasheet{Arbiter}{The AXI4 arbiter: up to eight hosts merged onto one peripheral link, answered by identifier.}

\dsfacts{\code{lib/parts/src/bus/arbiter.rs}}%
{\code{txhdl\_parts::bus::arbiter::Arbiter<N, ..>}, and \code{Arbiter2} to \code{Arbiter8}, which name the counts}%
{\code{//docs:axi}}

\subsection*{Function}

An arbiter puts several hosts on one peripheral link. It is the
mirror of the \code{Router}: the router decides on the address and
fans out, the arbiter decides whose turn it is and merges. It takes
\code{aws\_}\textit{i}, \code{ars\_}\textit{i} and \code{ws\_}\textit{i} from
each host \textit{i} and gives the peripheral side one \code{aw}, one
\code{ar} and one \code{w}; it takes that side's \code{b} and
\code{r} and gives each host its own, \code{bs\_}\textit{i} and \code{rs\_}\textit{i}.

Two things follow from AXI4 rather than from any choice made here.

An answer must reach the host that asked, and the peripheral does not
know there is more than one. The peripheral returns only the
identifier from the address phase, so the arbiter writes the host's
port number above the host's own identifier on the way out and takes
it off again on the way back. The peripheral side's identifier is
therefore wider than a host's.

The write data channel contains no identifier at all: a beat belongs
to the oldest write address phase that has not finished. So the
arbiter holds that channel for the host whose write address phase it
granted until that burst's last beat, and a second host's address
phase waits. Reads need no such lock.

It is one unit for any count of hosts, \code{N} a const parameter:
the hosts' channels are arrays of ports, and the grants are loops
over them that the lowering unrolls for each \code{N} (issue 500).
\code{Arbiter2} to \code{Arbiter8} are type aliases that name the
counts.

\subsection*{Parameters}

\begin{center}\footnotesize
\begin{tabular}{@{}lp{0.68\textwidth}@{}}
\toprule
Parameter & Meaning \\
\midrule
\code{N} & The number of hosts, two to eight. \\
\code{A} & The address width, in bits. \\
\code{D} & The data width, in bits. \\
\code{S} & The write strobe width, one bit per byte lane, so \code{D / 8}. \\
\code{I} & The identifier width on each host, in bits. \\
\code{J} & The identifier width on the peripheral side. It must be at
least \code{I} plus the bits the port number needs, which is one bit
for two hosts and three for five to eight. Nothing checks it; too
small a \code{J} loses the top of the tag. \\
\code{FIXED} & Zero for round robin, anything else for fixed
priority. \\
\bottomrule
\end{tabular}
\end{center}

\subsection*{The identifier}

A host's address phase goes out encoding
\code{(port << I) | id}, where \code{id} is the identifier that host's
own tracker allocated. An answer's identifier is read the other way:
\code{id >> I} names the host, and its low \code{I} bits are given
back to that host unchanged.

So a host's identifiers stay its own, and what the peripheral sees is
a wider space with the hosts side by side in it. The tables below are
generated for \code{Arbiter2<16, 32, 4, 2, 5, 0>}, the example's, whose
hosts provide two bits and whose peripheral side provides five.

\dstables{Arbiter}

\subsection*{Behaviour}

\begin{itemize}
\item \textbf{Read address.} Each cycle the hosts offering on
\code{ar} are considered, and one is granted if the peripheral side
has room. Several reads may be outstanding at once, from different
hosts, since each contains its identifier.
\item \textbf{Write address.} The same, except that no host is granted
while \code{wbusy} is set. Taking a phase sets \code{wbusy} and
records the host in \code{wsel}, a bit each.
\item \textbf{Write beats.} While \code{wbusy} is set, beats pass from
the host in \code{wsel} and from no other. The beat marked \code{last}
clears \code{wbusy}, and the next write address phase may then be
granted, to any host.
\item \textbf{Answers.} A beat on \code{b} or \code{r} is given to the
host its identifier's top bits name, when that host has room, with the
tag taken off. A host that is not ready holds that beat, and the
peripheral side is not taken from until it is.
\item \textbf{The turn.} With \code{FIXED} zero, \code{rturn} and
\code{wturn} hold where to start looking, and each moves past the host
that won, so a host that keeps asking cannot shut another out. With
\code{FIXED} not zero every host is eligible always, so the lowest
numbered one that is offering wins.
\end{itemize}

\subsection*{Verification}

\begin{itemize}
\item \code{//lib/parts:parts\_test} drives two hosts and four:
\code{each\_answer\_goes\_home} has two hosts read at once from a
memory that answers with the address it was asked for, so a host given
another's beats would see an address it never asked for;
\code{write\_beats\_do\_not\_interleave} has both write four beats at
the same moment and checks the memory was given two whole bursts;
\code{four\_hosts\_share\_the\_link} has four read as fast as they can
for twelve hundred cycles and checks the counts come out within one of
each other.
\item \code{//lib/examples:ex\_axi\_arbiter} is the same shape as a
run, and the build simulates the netlist against it under nvc and
Verilator: \code{//docs:axi\_arbiter\_sim\_axi\_arbiter\_tb\_test} and
\code{//docs:axi\_arbiter\_vsim\_test}.
\end{itemize}

\subsection*{Used by}

\code{Board}, three times. As \code{arb}, an
\code{Arbiter7<32, 32, 4, 2, 5, 0>} in front of its router, it takes
the core's tracker on port 0; on port 1 \code{jarb}, which shares it
between the JTAG master's pins (issue 241) and the debug transport's
bridge (issue 154); the Ethernet port's fetch and store engines on
ports 2 and 3 (issue 151); the video scanout's fetch on port 4 (issue
151); on port 5 \code{sdarb}, the SD card host's two engines
(issue 912); and on port 6 Razboj's rasteriser (issue 985). Its peripheral side has five bits of identifier, two for
a host's own and three for the port, so up to eight hosts fit without
widening anything past it (issue 1022). \code{jarb} and \code{sdarb}
are each an \code{Arbiter2<32, 32, 4, 1, 2, 0>}, giving their two hosts
one bit of identifier each. The SD engines never offer at the same
time, so neither takes a turn from the other.

\subsection*{Limits}

The turn is three bits, which is what eight hosts need and no more,
and nothing checks that \code{N} is at most eight.
\code{J} is not checked against \code{I} and the count. There is no
arbitration on quality of service, though the address phase provides
the field. A host that stops taking its answers stops the peripheral
side for every host, since the arbiter holds the beat rather than
dropping it.

% SPDX-License-Identifier: Apache-2.0
% covers: LiteBridge
\datasheet{LiteBridge}{The AXI4 to AXI-Lite bridge: one AXI4 link to any count of AXI-Lite peripherals, by an address map.}

\dsfacts{\code{lib/parts/src/bus/axi\_lite.rs}}%
{\code{txhdl\_parts::bus::axi\_lite::LiteBridge<N, M, ..>}, with \code{M} an address map, \code{txhdl::map::AddrMap<N>}}%
{\code{//docs:axi}, \code{//docs:vreteno}, \code{//docs:eth}, \code{//docs:hdmi}}

\subsection*{Function}

AXI-Lite is AXI4 with no identifier, no burst and no \code{last}, so
one transaction is one beat on each channel it uses. A bridge stands
between an AXI4 link and that many AXI-Lite peripherals. On its AXI4
side it takes the \code{aw}, \code{ar} and \code{w} an \code{AxiPer}
would take, and gives the \code{b} and \code{r} it would give. So it
goes where an \code{AxiPer} would: behind a host's tracker, or behind
one of a router's peripheral ports.

Its peripheral side is arrays of \code{N} ports: it drives
\code{aws}, \code{ars} and \code{ws}, and reads \code{bs} and
\code{rs}, entry \textit{i} of each being peripheral \textit{i}'s, so
the netlist's ports are \code{aws\_}\textit{i} and the like. An
address phase there is a \code{LiteAddr}: \code{addr} and
\code{prot}. A write word is a \code{LiteW}: \code{data} and
\code{strb}. \code{axi\_lite} makes the five channels of such a link.
One unit serves every count; the lowering unrolls its loops over the
arrays when it runs (issue 500). Until issue 647 a macro,
\code{lite\_bridge!}, wrote \code{LiteBridge1} to \code{LiteBridge8}
instead, one unit per count, with each range a pair of const
parameters.

\subsection*{Parameters}

\begin{center}\footnotesize
\begin{tabular}{@{}lp{0.7\textwidth}@{}}
\toprule
Parameter & Meaning \\
\midrule
\code{A} & The address width, in bits, on both sides. \\
\code{D} & The data width, in bits, on both sides. \\
\code{S} & The write strobe width, one bit per byte lane, so \code{D / 8}. \\
\code{I} & The identifier width of the AXI4 side, in bits. \\
\code{N} & The count of peripherals, one or more. \\
\code{M} & The address map: a type implementing \code{AddrMap<N>},
whose \code{RANGES} holds each peripheral's base and mask, in the
order of the ports. \\
\bottomrule
\end{tabular}
\end{center}

\subsection*{Address map}

An address \code{a} is peripheral \textit{i}'s when
\code{(a \& mask) == base} for \code{(base, mask) = M::RANGES[}\textit{i}\code{]}.
The bridge decodes a burst's first address, and the lowest range that
matches wins. An address in no range is a hole. The tables below are
generated for \code{LiteBridge<1, SerialMap, 32, 32, 4, 2>}, the bridge
in front of Vreteno's serial port, whose map is this one:

\begin{center}\footnotesize
\begin{tabular}{@{}llll@{}}
\toprule
Peripheral & \code{BASE} & \code{MASK} & In Vreteno \\
\midrule
0 & \code{0x3000} & \code{0xf000} & The serial port \\
\bottomrule
\end{tabular}
\end{center}

The board's design uses a bridge of one peripheral per slot of its
peripheral page, \code{SlotMap}, 256 bytes each from \code{0x3000}
with mask \code{0xffff\_ff00}; the page's table in \code{//docs:vreteno}
is written from \code{SlotMap} at build time and lists them. It also has
bridges of one at \code{0x0c00\_0000} (mask \code{0xfc00\_0000}, in
front of its interrupt controller) and at \code{0x1000\_0000} (mask
\code{0xffff\_0000}, the debug module). \code{ex\_axi\_lite} uses a
bridge of two with ranges at \code{0x1000} and \code{0x2000}, mask
\code{0xf000} each.

\dstables{LiteBridge}

\subsection*{Behaviour}

\begin{itemize}
\item \textbf{Taking a burst.} It takes a burst only when none is in
progress (\code{busy} clear). When both address channels offer,
\code{rfirst} says which goes first. After a read the write channel
goes first, and after a write the read channel.
\item Taking a burst records its identifier in \code{xid}, its address
in \code{addr}, its \code{len} in \code{left}, its \code{prot}, and
its peripheral in \code{sel}, a bit each. A hole has no bit set.
\code{sel} is \code{N} bits wide; a bridge of one has a one-bit
\code{sel}, which the netlist writes as a single bit.
\item \code{stride} is how far the address moves per beat:
\code{1 << size} bytes, or zero for a \code{Fixed} burst. Every other
burst kind moves as \code{Incr} does.
\item \textbf{A write beat.} The bridge takes the burst's next beat
from \code{w} when the chosen peripheral has room on both its
\code{aw} and its \code{w}. It sends that peripheral \code{addr} and
\code{prot} on \code{aw}, and the beat's \code{data} and \code{strb}
on \code{w}. To a hole, the beat is taken and dropped.
\item \textbf{A read beat.} The bridge sends \code{addr} and
\code{prot} on the chosen peripheral's \code{ar} when it has room.
\item Once a beat's request is out (\code{sent}), the bridge waits for
that peripheral's answer. A hole's answer is there at once:
\code{DecErr}, with zero data for a read.
\item \textbf{A read's answer} goes up on \code{r} as a beat of the
burst, when \code{r} has room. It has the burst's identifier, the word,
the peripheral's response, and \code{last} on the final beat.
\item \textbf{A write's answer} is taken when \code{b} has room, or when
more beats are to come. \code{wresp} keeps the first response that is
not \code{Okay}. After the last beat one \code{B} goes up, with the
burst's identifier and that response.
\item After each answer, \code{addr} moves on by \code{stride} and
\code{left} counts down. The answer to the last beat clears
\code{busy}.
\end{itemize}

\subsection*{Verification}

\begin{itemize}
\item \code{bursts\_cross\_the\_bridge\_a\_beat\_at\_a\_time} in
\code{bus::axi\_lite}, in \code{//lib/parts:parts\_test}, puts
\code{LiteBridge<2, TwoMap, 16, 32, 4, 2>} behind
a host tracker, with two simulated memories. It writes three words and
reads four, reads one word three times in a fixed burst, uses the
second peripheral, and sends a write and a read of two beats to a hole.
It checks every word and response, and counts ten transactions on the
first memory and two on the second.
\item \code{ex\_axi\_lite} runs the same kind of bursts against two
banks of registers written as hardware. The build simulates the
bridge's netlist against that run under nvc and Verilator:
\begin{itemize}
\item \code{//docs:axi\_lite\_sim\_axi\_lite\_bridge\_axi\_lite\_bridge\_tb\_test};
\item \code{//docs:axi\_lite\_vsim\_axi\_lite\_bridge\_test}.
\end{itemize}
\item Vreteno's run lowers its serial port's bridge, the
bridge of one of the tables, and the build simulates that netlist
against the run: \code{//docs:vreteno\_sim\_ubridge\_ubridge\_tb\_test}
and \code{//docs:vreteno\_vsim\_ubridge\_test}.
\item \code{//cpu/vreteno:lockstep\_test} and \code{//soc:demo\_test} run
a bridge of one in simulation.
\end{itemize}

\subsection*{Used by}

A bridge of two in the example \code{ex\_axi\_lite}. A bridge of one
in front of the serial port in Vreteno's run
(\code{cpu/vreteno/src/run.rs} and \code{cpu/vreteno/src/main.rs}) and
in \code{soc/src/lib.rs}; in the board's design, a bridge of six as its
field \code{puart} and bridges of one as \code{pplic} and \code{pdm}.
A bridge of one also puts each peripheral of the examples behind a
host's tracker: \code{ex\_eth}, \code{ex\_hdmi}, \code{ex\_sd} and the
rest.

\subsection*{Limits}

It serves one burst at a time, for all its peripherals together, since
AXI-Lite has no identifier to tell two apart. A burst's peripheral is
chosen on its first address, and every beat of the burst goes there.
It treats a \code{Wrap} burst as \code{Incr}. It counts a write's beats
from \code{len} and does not read \code{last} on \code{w}. Of the
address phase only \code{addr} and \code{prot} reach a peripheral;
\code{size} sets the stride, and \code{lock}, \code{cache},
\code{qos} and \code{region} are dropped.

% SPDX-License-Identifier: Apache-2.0
% covers: LiteSplit
\datasheet{LiteSplit}{One AXI-Lite slot split in two by an address bit,
so that two peripherals share a slot.}

\dsfacts{\code{lib/parts/src/bus/lite\_split.rs}}%
{\code{txhdl\_parts::bus::lite\_split::LiteSplit<A, D, S, BIT>}}%
{\code{//docs:examples}}

\subsection*{Function}

The board's peripheral page gives each peripheral one slot of 256 bytes,
and a slot is one AXI-Lite link. \code{LiteSplit} puts two peripherals
in one slot: a transaction whose address has bit \code{BIT} low goes to
the first, \code{lo}, and one whose bit is high to the second,
\code{hi}. Neither peripheral sees the bit, since each decodes only the
low bits of its own words (issue 151).

\subsection*{Parameters}

\begin{center}\footnotesize
\begin{tabular}{@{}lp{0.68\textwidth}@{}}
\toprule
Parameter & Meaning \\
\midrule
\code{A} & The address width. \\
\code{D} & The data width. \\
\code{S} & The strobe width, \code{D / 8}. \\
\code{BIT} & The address bit that picks the side. \\
\bottomrule
\end{tabular}
\end{center}

The tables below use \code{LiteSplit<32, 32, 4, 7>}, the split inside
\code{ScanVideo}: the video peripheral below \code{0x80} of the slot,
the scanout's registers from there.

\dstables{LiteSplit}

\subsection*{Behaviour}

\begin{itemize}
\item A read is taken when none is out and its side's address channel
has room, and goes to that side; \code{rpend} and \code{rside} remember
it until its word comes back, from that side only, and goes up when
the host has room for it.
\item A write's address and word are taken together, when none is out
and both of its side's channels have room; \code{wpend} and
\code{wside} remember it until its response comes back from that side.
\item One read and one write may be out at once, but never two of
either, so two answers on one channel never pass each other.
\end{itemize}

\subsection*{Verification}

\begin{itemize}
\item \code{ex\_scanvideo} reaches both sides through it: a pixel
painted through \code{Hdmi}'s cursor below \code{0x80}, and the
scanout's base, control, status and clear above it.
\item \code{//docs:scanvideo\_sim\_scanvideo\_tb\_test} and
\code{//docs:scanvideo\_vsim\_test} simulate \code{ScanVideo}'s
netlist, this unit inside it, against that run.
\end{itemize}

\subsection*{Used by}

\code{ScanVideo}, which the flagship holds in the video slot at
\code{0x3200}.

\subsection*{Limits}

One outstanding read and one outstanding write, which is what the
AXI4-to-AXI-Lite bridge in front of it issues; a host that pipelines
several waits for each answer. Two sides only.

% SPDX-License-Identifier: Apache-2.0
% covers: AxiWb
\datasheet{AxiWb}{A bridge from the peripheral end of an AXI4 link to the host lines of a pipelined Wishbone bus.}

\dsfacts{\code{lib/parts/src/bus/wb.rs}}{\code{txhdl\_parts::bus::wb::AxiWb}}%
{\code{//docs:axi}, \code{//docs:vreteno}}

\subsection*{Function}

A controller written elsewhere often speaks Wishbone rather than AXI.
\code{AxiWb} puts such a controller on a link. It is a peripheral
client written as hardware. On one side it holds the channel ends an
\code{AxiPer} gives a client, as one port, \code{bus}, a
\code{PerPort} of the four channels: requests on \code{bus\_req} and write beats on
\code{bus\_w} in, answers on \code{bus\_ans} and read beats on \code{bus\_r} out.
On the other side it drives the host lines of a pipelined Wishbone
bus, \code{cyc}, \code{stb}, \code{we}, \code{adr}, \code{dat} and
\code{sel}, and reads \code{stall}, \code{ack} and the peripheral's
word, \code{rdat}.

It takes one burst at a time and puts each of its words on the lines
as its own request, at consecutive word addresses. The DDR3
controller of Vreteno's board is behind one.

\subsection*{Parameters}

\begin{center}\footnotesize
\begin{tabular}{@{}lp{0.7\textwidth}@{}}
\toprule
Parameter & Meaning \\
\midrule
\code{A} & The link's address width, in bits. \\
\code{I} & The link's identifier width, in bits. \\
\code{AW} & The width of the Wishbone word address, in bits. \\
\bottomrule
\end{tabular}
\end{center}

Words are 32 bits with four lanes, fixed in the source. The tables
below use \code{AxiWb<32, 2, 28>}. There \code{AW} is \code{ddr3::AW},
28 bits, a gigabyte of words, which was the board's DDR3 memory's word
address while this bridge was in front of it.

\dstables{AxiWb}

\subsection*{Behaviour}

\code{stage} says where the request is. The source names five values.

\begin{itemize}
\item \textbf{0, waiting.} A request on \code{bus\_req} is taken when one
offers. The bridge keeps its identifier and its word address, sets
\code{wsel} to all four lanes, and latches \code{left}, the count of
words the burst still has to put on the lines, as \code{len + 1}. A
read goes to stage 2 at once, and a write to stage 1.
\item \textbf{1, a write waiting for its beat.} The beat on \code{bus\_w}
is taken when it offers. Its \code{data} goes to \code{wdat} and its
\code{strb} to \code{wsel}. Then stage 2.
\item \textbf{2, offered.} \code{cyc} and \code{stb} are high. A cycle
with \code{stall} low is the one the peripheral takes the request in,
and the stage moves to 3.
\item \textbf{3, taken.} \code{cyc} alone is high until \code{ack}.
An \code{ack} in the cycle of the take counts too. On \code{ack} the
stage moves to 4, and a read keeps the word on \code{rdat} in
\code{wdat}.
\item \textbf{4, answering.} A read sends a beat on \code{bus\_r}: the
identifier, \code{wdat}, \code{Okay}, and \code{last} high only on
the burst's last word. If words remain the stage returns to 2 at the
next word address with \code{left} one lower; otherwise to 0.
\item A write is answered once, after its last word. Each
acknowledgement before that sends the bridge back to stage 1 for the
next beat, at the next word address. The \code{Answer} of \code{Okay}
on \code{bus\_ans} is sent from stage 4, and each answer waits for room on
its channel.
\item \code{we} is high unless the request is a read. \code{adr},
\code{dat} and \code{sel} are \code{wadr}, \code{wdat} and \code{wsel}.
Every line depends only on the bridge's registers, so nothing
combinational passes from the link to the lines or back.
\item The Wishbone address is a word address: the burst's byte address
shifted down by two and cut to \code{AW} bits. A peripheral decoded at
a region of the link's map sees offsets from the region's base, as
long as the base lies above those bits. A write at
\code{0x4000\_0010} lands in word 4.
\end{itemize}

\subsection*{Verification}

\begin{itemize}
\item Two tests in \code{bus::wb}, in \code{//lib/parts:parts\_test},
run the bridge on a link, behind \code{AxiHost} and \code{AxiPer},
against \code{WbMem}, a simulated Wishbone memory. The first,
\code{a\_read\_gets\_what\_a\_write\_put}, writes and reads one word.
The second, \code{slow\_answers\_and\_a\_calibration}, writes and
reads back 24 words at pseudorandom addresses, against a memory that
stalls for its first 40 cycles and answers after 5.
\item \code{ex\_axi\_wb} writes three words and reads them back, against
a memory that stalls for 10 cycles and answers after 2. The build
simulates the netlist of \code{AxiWb<32, 2, 28>} against that run under
nvc and Verilator: \code{//docs:axi\_wb\_sim\_axi\_wb\_tb\_test} and
\code{//docs:axi\_wb\_vsim\_test}.
\end{itemize}

\subsection*{Used by}

The example \code{ex\_axi\_wb}, and
\code{a\_burst\_is\_served\_word\_by\_word} in the module's own
tests, which writes sixteen words and reads them back as one burst
each way. Sixteen is what \code{LineFetch} issues by default.
Until issue~\#1023 put the board's DDR3 memory behind the controller's
own AXI4 port, \code{Ddr3Per} held this bridge in front of a Wishbone
wrapper of the controller.
That test asserts the COUNT of beats and not only their values: a
bridge that answered one beat would return the right word in it, so
checking the words alone passes on the fault.

\subsection*{Impact of the single-beat bridge (prior to issue 471)}

This behavior is documented here because joining a burst-issuing client
to a peripheral that returns fewer beats than requested risks protocol
failure that neither engine can detect.

\code{LineFetch} decides a burst is finished by its own beat counter
and never reads the returned beat's \code{last}. Against a bridge
that answered one beat it took that word and waited for the rest for
ever, holding \code{running} high: sixty four words became one, and
the run did not end.

\code{LineStore} fared worse and quietly. A write burst is finished
by its response, so the bridge answered after one beat while the rest
of the burst sat in the channel: nine words of sixty four landed and
the rest of memory was left as it was, with nothing to say why.

\subsection*{Limits}

It serves a burst one request at a time rather than as a pipelined
run of several. That removes a link round trip per word and not the
word itself: a burst of sixteen is sixteen Wishbone request and
acknowledge handshakes, one after another. Issuing several requests
before their acknowledgements would remove more, and wants somewhere
to hold answers the link is not ready for, which the bridge does not
have. It always answers \code{Okay}, and it reads no
error line from the peripheral. Words are 32 bits with four lanes.
Address bits above the \code{AW} bits it keeps are dropped.

% SPDX-License-Identifier: Apache-2.0
% covers: AxiPins, AxiHostPins
\datasheet{AxiPins}{A host's AXI4 pins joined to the AXI link's channels.}

\dsfacts{\code{lib/parts/src/bus/axi\_pins.rs}}%
{\code{txhdl\_parts::bus::axi\_pins::AxiPins}}%
{\code{//docs:pcie}}

\subsection*{Function}

A core from elsewhere that is an AXI4 host has pins: a \code{valid} and
a \code{ready} per channel and a wire per field, named as AXI4 names
them. \code{AxiPins} joins those pins to the link's five channels, so it
goes where an \code{AxiHost} would, with an \code{AxiPer} or a router
behind it. Its inputs are the side \code{inp: AxiPinsIn}: the host's
24 pins as one field, \code{pins}, an \code{AxiHostPins}, and the
\code{b} and \code{r} channels from the link. One rule names the ports
of a struct side, \code{name\_field} for a side written
\code{name: S} (issue 580), so its input ports are
\code{inp\_pins\_awid} and the rest beside \code{inp\_b} and
\code{inp\_r}.
\code{AxiHostPins} derives \code{Ports}, so a design that has a host's
pins among its own ports holds them as one field too and passes them
whole (issue 579). Its outputs
are the side \code{outp: AxiPinsOut}: the \code{aw}, \code{ar} and
\code{w} channels to the link, and the eleven pins the host reads, so
its output ports are \code{outp\_aw}, \code{outp\_ar}, \code{outp\_w},
\code{outp\_awready} and the rest.

\code{axi\_pins::sim} has \code{PinHost}, a host on the pins written as
a simulation, and \code{pin\_host} and \code{pins}, which make the
wires.

\subsection*{Parameters}

\begin{center}\footnotesize
\begin{tabular}{@{}lp{0.7\textwidth}@{}}
\toprule
Parameter & Meaning \\
\midrule
\code{A} & The address width, in bits. The tables use 32. \\
\code{D} & The data width, in bits. The tables use 64. \\
\code{S} & The strobe width, \code{D / 8}. The tables use 8. \\
\code{I} & The identifier width, in bits. The tables use 4. \\
\bottomrule
\end{tabular}
\end{center}

The tables are generated for \code{AxiPins<32, 64, 8, 4>}, the widths of
the PCIe endpoint's master as \code{PcieBar} uses it.

\dstables{AxiPins}

\subsection*{Behaviour}

\begin{itemize}
\item It holds no state; every output is a function of the step's
inputs and the channels.
\item \code{awready}, \code{arready} and \code{wready} are the room on
\code{aw}, \code{ar} and \code{w}.
\item When \code{awvalid} is high and \code{aw} has room, an address
phase goes onto \code{aw} with the pins' \code{id}, \code{addr},
\code{len}, \code{size}, \code{lock}, \code{cache} and \code{prot}, the
burst type decoded from its two bits, and \code{qos} and \code{region}
zero. The same holds for \code{arvalid} and \code{ar}, and for
\code{wvalid} and \code{w}, with \code{data}, \code{strb} and
\code{last}.
\item \code{bvalid} is high while \code{b} holds a response, with
\code{bid} and \code{bresp} its fields, the response encoded in two
bits. It is taken in a step \code{bready} is high.
\item \code{rvalid}, \code{rid}, \code{rdata}, \code{rresp} and
\code{rlast} are the head of \code{r} the same way, taken in a step
\code{rready} is high.
\end{itemize}

\subsection*{Verification}

\begin{itemize}
\item \code{a\_host\_on\_pins\_reaches\_a\_peripheral\_through\_the\_link}
in \code{bus::axi\_pins}, in \code{//lib/parts:parts\_test}, writes three
words through the pins, an \code{AxiPer} and a simulated RAM, reads four
back with their identifier, and reads past the RAM's end, answered
\code{SlvErr}.
\item \code{ex\_axi\_pins} runs \code{AxiPins<32, 64, 8, 4>} the same
way, and the build simulates its netlist against that run under nvc and
Verilator: \code{//docs:axi\_pins\_sim\_axi\_pins\_tb\_test} and
\code{//docs:axi\_pins\_vsim\_test}.
\item \code{//pcie:bar\_test} drives it inside \code{PcieBar}.
\end{itemize}

\subsection*{Used by}

\code{PcieBar}, as its field \code{pins}, between the PCIe endpoint's
master and a peripheral tracker. \code{Board}, as its field
\code{jtag}, between Vivado's JTAG-to-AXI master on the board's top
and the arbiter in front of the router (issue 241). The example \code{ex\_axi\_pins}.

\subsection*{Limits}

Its \code{ready} outputs follow the channels' room in the same step, so
a path runs from a channel's state through the unit to the host's pins
with no register. It has no \code{AxQOS}, \code{AxREGION} or user pins,
and drives \code{qos} and \code{region} as zero.

% SPDX-License-Identifier: Apache-2.0
% covers: AxiPerPins, AxiPerDriven
\datasheet{AxiPerPins}{The AXI link's peripheral end joined to an AXI4 peripheral's pins.}

\dsfacts{\code{lib/parts/src/bus/axi\_per\_pins.rs}}%
{\code{txhdl\_parts::bus::axi\_per\_pins::AxiPerPins}}%
{\code{//docs:pcie}}

\subsection*{Function}

A core from elsewhere that is an AXI4 peripheral, such as AMD's DDR3
controller generated with its AXI4 port, has pins: a \code{valid} and a
\code{ready} per channel and a wire per field. \code{AxiPerPins} is
\code{AxiPins} the other way round. It stands on the link's peripheral
end, where an \code{AxiPer} would, and puts the link's channels onto
those pins, holding nothing.

Its inputs are the side \code{inp: AxiPerPinsIn}: the eleven pins the
peripheral drives as one field, \code{pins}, an \code{AxiPerDriven},
and the \code{aw}, \code{ar} and \code{w} channels from the link, so
its input ports are \code{inp\_pins\_awready} and the rest beside
\code{inp\_aw}, \code{inp\_ar} and \code{inp\_w}. \code{AxiPerDriven}
derives \code{Ports}, so a design holding a peripheral's pins among its
own ports passes them on whole. Its outputs are the side
\code{outp: AxiPerPinsOut}: the \code{b} and \code{r} channels back to
the link, and the twenty-six pins the peripheral reads.

\code{axi\_per\_pins::sim} has \code{PinRam}, a memory on the pins
written as a simulation, and \code{pins}, which makes the wires.
\code{PinRam} can take a warm-up, before which it is not ready, and a
read and a write latency, and it can take an address past its end
modulo its size; that is how \code{ddr3::Ddr3} models AMD's controller.

\subsection*{Parameters}

\begin{center}\footnotesize
\begin{tabular}{@{}lp{0.7\textwidth}@{}}
\toprule
Parameter & Meaning \\
\midrule
\code{A} & The address width, in bits. The tables use 32. \\
\code{D} & The data width, in bits. The tables use 32. \\
\code{S} & The strobe width, \code{D / 8}. The tables use 4. \\
\code{I} & The identifier width, in bits. The tables use 5. \\
\bottomrule
\end{tabular}
\end{center}

The tables are generated for \code{AxiPerPins<32, 32, 4, 5>}, the
widths \code{Ddr3Per} uses it with on the board. The example
\code{ex\_axi\_per\_pins} uses \code{AxiPerPins<30, 32, 4, 5>}, the
controller's own 30-bit address.

\dstables{AxiPerPins}

\subsection*{Behaviour}

\begin{itemize}
\item It holds no state; every output is a function of the step's
inputs and the channels.
\item An address phase at the head of \code{aw} is \code{awvalid} and
the \code{aw} pins: \code{awid}, \code{awaddr}, \code{awlen},
\code{awsize}, the burst type encoded in two bits, \code{awlock},
\code{awcache}, \code{awprot} and \code{awqos}. It is taken off the
channel in a step \code{awready} is high. The same holds for
\code{ar}, and for \code{w} with \code{wdata}, \code{wstrb} and
\code{wlast}.
\item A response the peripheral offers, \code{bvalid} with
\code{bid} and \code{bresp}, the response decoded from its two bits,
goes onto \code{b} in a step \code{b} has room, and \code{bready} is
that room. The same holds for \code{rvalid} and \code{r}, with
\code{rid}, \code{rdata}, \code{rresp} and \code{rlast}, and
\code{rready}.
\item Every field of the link's address phase reaches a pin except
\code{region}: AMD's controller has no region pins, and a peripheral
joined to one link needs none, since the region is how an interconnect
tells decoded ranges of one peripheral apart. \code{AxQOS} is forwarded.
\end{itemize}

\subsection*{Verification}

\begin{itemize}
\item \code{a\_host\_on\_the\_link\_reaches\_a\_peripheral\_on\_pins} in
\code{bus::axi\_per\_pins}, in \code{//lib/parts:parts\_test}, writes
three words through an \code{AxiHost}, the unit and a \code{PinRam},
reads four back across the burst's ends, and reads past the memory's
end, answered \code{SlvErr}.
\item \code{a\_timed\_memory\_waits\_and\_then\_answers\_late} there
holds \code{PinRam}'s timing: a warm-up holds a write off, and a latency
delays an answer by itself and no more.
\item \code{ex\_axi\_per\_pins} runs \code{AxiPerPins<30, 32, 4, 5>} with
a host on the link and a \code{PinRam} on the pins, and the build
simulates its netlist against that run under nvc and Verilator:
\code{//docs:axi\_per\_pins\_sim\_axi\_per\_pins\_tb\_test} and
\code{//docs:axi\_per\_pins\_vsim\_test}.
\item \code{//ddr3:ddr3\_test} drives it inside \code{Ddr3Per}, and
\code{//ddr3/sim:wrapper\_test}, which is manual, checks the
controller's side of the same pins under Vivado's simulator.
\end{itemize}

\subsection*{Used by}

\code{Ddr3Per}, as its field \code{pins}, between the board's router
and AMD's DDR3 controller. The example \code{ex\_axi\_per\_pins}.

\subsection*{Limits}

Its \code{ready} outputs are the channels' room in the same step, so
a path runs from a channel's state through the unit to the
peripheral's pins with no register. It has no \code{AxREGION} or user
pins, so a peripheral that wants a region gets it as a constant where
it is instantiated.

% The network on chip.
% SPDX-License-Identifier: Apache-2.0
% covers: Switch
\datasheet{Switch}{A five-port switch of the network on chip, one per virtual channel.}

\dsfacts{\code{lib/parts/src/bus/noc/switch.rs}}{\code{txhdl\_parts::bus::noc::switch::Switch}}%
{\code{//docs:noc}}

\subsection*{Function}

A switch has a two-way link on each of its four sides and a fifth port, the exit,
where something that is not a node attaches. Its inputs are
\code{n\_in}, \code{s\_in}, \code{w\_in}, \code{e\_in} and
\code{x\_in}, and its outputs are \code{n\_out}, \code{s\_out},
\code{w\_out}, \code{e\_out} and \code{x\_out}. Each port passes one
\code{Pkt}, a beat of one AXI channel addressed to a node. The switch
routes in dimension order, X before Y. Its place in the lattice is two more
inputs, \code{col} and \code{row}, which whoever places it holds at
constants with \code{tie}, so every switch of a lattice is one type
(issue 635). A \code{Node} is two switches,
one per virtual channel.

\subsection*{Parameters}

\begin{center}\footnotesize
\begin{tabular}{@{}lp{0.7\textwidth}@{}}
\toprule
Parameter & Meaning \\
\midrule
\code{XB} & The width of a column coordinate, and of \code{col}, in
bits. A lattice is at
most \code{1 << XB} columns wide. \\
\code{YB} & The width of a row coordinate, and of \code{row}, in bits. A lattice is at
most \code{1 << YB} rows high. \\
\code{A} & The address width of the AXI link. \\
\code{D} & The data width of the AXI link. \\
\code{S} & The strobe width, which is \code{D / 8}. \\
\code{I} & The identifier width of the AXI link. \\
\bottomrule
\end{tabular}
\end{center}

The tables below are for \code{Switch<2, 2, 32, 32, 4, 2>}: two-bit
coordinates, on the link that \code{//soc} uses.

\dstables{Switch}

\subsection*{Behaviour}

\begin{itemize}
\item The switch holds one thing, \code{turn}, three bits: which
input the choice starts from. The port a packet leaves by holds
nothing, being a function of the packet's \code{dx} and \code{dy} and
of \code{col} and \code{row}; what needed a register is which input goes
when several could (issue 315).
\item The switch compares coordinates at \code{CW}, 8 bits. It sends a
packet east when \code{dx} is above \code{col}, and west when \code{dx}
is below it. In the right column, it sends the packet south when
\code{dy} is above \code{row}, and north when \code{dy} is below it. At its own node it sends the packet out of the exit.
\item Each cycle the switch works out, for every input with a packet,
the port that packet wants and whether that port is ready. An input
whose port is full does not hold up the others.
\item Of the inputs that can move, the switch takes one per cycle: the
first at or after \code{turn}, going round the five, and \code{turn}
becomes the one after whichever moved. So an input waits on at most
four others, where the fixed priority it replaced could starve one
outright: a busy link kept the exit from moving at all, since the exit
was last and the links were never all quiet.
\item The switch sends the packet it takes unchanged.
\end{itemize}

\subsection*{Verification}

\begin{itemize}
\item Two tests in \code{lib/parts/src/bus/noc.rs}, in
\code{//lib/parts:parts\_test}, drive one switch at 0,0 alone.
\code{a\_busy\_link}\allowbreak\code{\_and\_the\_exit\_take\_turns} has a link and the exit
offer every cycle for the east port, and asserts that each moves more
than fifty of two hundred and that the two differ by at most two.
\code{two\_inputs\_of\_a\_switch}\allowbreak\code{\_contend\_for\_their\_shared\_output}
measures two inputs leaving by one port, both offering every cycle and
one offering every third.
\item Eight more tests of the network in the same file run it inside
four nodes, and \code{a\_core\_draws\_a\_solid\_and\_a\_gpu\_fills\_it},
in \code{//soc:demo\_test}, runs a core, a GPU, a memory and a serial
port on a two by two lattice of nodes.
\item The switch's netlist is checked inside \code{Mesh}, two to a
node: \code{ex\_mesh} lowers a mesh of six nodes, and
\code{//docs:mesh\_sim\_mesh\_tb\_test} and \code{//docs:mesh\_vsim\_test}
run it under nvc and Verilator against the Rust run, at the mesh's
ports.
\end{itemize}

\subsection*{Used by}

\code{Node}, which holds two: \code{req} for requests and \code{rsp}
for responses.

\subsection*{Limits}

The switch moves one packet per cycle, not one per port.
\code{//docs:noc} measures the cost on the system of \code{//soc}: the
rasteriser draws its list in 36\,516 cycles on a direct link and in
78\,391 cycles across the lattice.

% SPDX-License-Identifier: Apache-2.0
% covers: Node
\datasheet{Node}{A node of the network on chip: a switch per virtual channel.}

\dsfacts{\code{lib/parts/src/bus/noc/node.rs}}{\code{txhdl\_parts::bus::noc::node::Node}}%
{\code{//docs:noc}}

\subsection*{Function}

A node is two \code{Switch} units: \code{req}, for what a host sends
and a peripheral receives, and \code{rsp}, for what a peripheral sends
and a host receives. The two share nothing. So a request held up behind
a full buffer cannot hold up the response that would empty it.

The node's place in the lattice is two inputs, \code{col} and
\code{row}, which it passes to both switches. Whoever places the node
holds them at constants with \code{tie}, so every node of a lattice is
one type (issue 635), and \code{//docs:noc} states that synthesis
folds the constants into the comparisons that route a packet.

Its other twenty ports are written out one by one. A port name starts
with \code{q} for the request channel or \code{p} for the response
channel. The side follows: \code{n}, \code{s}, \code{w}, \code{e}, or
\code{x} for the exit. Last is \code{i} or \code{o}. So \code{qei} is
the request input from the east, and \code{pxo} is the response output
to the exit.

\subsection*{Parameters}

\begin{center}\footnotesize
\begin{tabular}{@{}lp{0.7\textwidth}@{}}
\toprule
Parameter & Meaning \\
\midrule
\code{XB} & The width of a column coordinate, and of \code{col}, in
bits. A lattice is at most \code{1 << XB} columns wide. \\
\code{YB} & The width of a row coordinate, and of \code{row}, in bits.
A lattice is at most \code{1 << YB} rows high. \\
\code{A} & The address width of the AXI link. \\
\code{D} & The data width of the AXI link. \\
\code{S} & The strobe width, which is \code{D / 8}. \\
\code{I} & The identifier width of the AXI link. \\
\bottomrule
\end{tabular}
\end{center}

The tables below are for \code{Node<2, 2, 32, 32, 4, 2>}, the node on
the link that \code{//soc} uses.

\dstables{Node}

\subsection*{Behaviour}

\begin{itemize}
\item The node has no logic of its own. It runs \code{req} on the ten
\code{q} ports and \code{rsp} on the ten \code{p} ports, side by side,
both at the place \code{col} and \code{row} say. What it holds is its
two switches' state, which is six bits: a \code{turn} of three in
each.
\item Each switch routes in dimension order and moves one packet per
cycle, choosing among the inputs that can move by a round robin of its
own, as the \code{Switch} sheet states. The two switches take turns
independently: a busy request channel does not decide which response
moves.
\item A host's bridge sends requests into \code{qxi} and takes
responses from \code{pxo}. A peripheral's bridge takes requests from
\code{qxo} and sends responses into \code{pxi}.
\end{itemize}

\subsection*{Verification}

\begin{itemize}
\item Eight tests in \code{lib/parts/src/bus/noc.rs}, in
\code{//lib/parts:parts\_test}, run four nodes on a two by two lattice
with both bridges, both AXI trackers and memory. Among them,
\code{a\_burst\_crosses\_the\_lattice\_and\_its\_answer\_comes\_back}
writes a word from the host at 0,0 to a memory at 1,1 and reads it
back; \code{two\_hosts\_share\_a\_memory\_and\_each\_answer\_goes\_home}
puts hosts at 0,0 and 1,0 on that memory; and
\code{four\_peripherals\_share\_one\_node} puts a router of four
memories behind the node at 1,1. The others cross the lattice with
reads of four beats, writes of two and of sixteen, a fixed write and a
wrapping one.
\item \code{a\_core\_draws\_a\_solid\_and\_a\_gpu\_fills\_it}, in
\code{//soc:demo\_test}, runs the core, the GPU, a memory and a serial
port on the four corners of one lattice.
\item The node's netlist is checked inside \code{Mesh}: \code{ex\_mesh}
lowers a mesh of six nodes, and \code{//docs:mesh\_sim\_mesh\_tb\_test}
and \code{//docs:mesh\_vsim\_test} run it under nvc and Verilator
against the Rust run, at the mesh's ports.
\end{itemize}

\subsection*{Used by}

The tests of the network in \code{lib/parts/src/bus/noc.rs}, and
\code{soc::run} in \code{soc/src/lib.rs}, each with four nodes wired by
\code{mesh::lattice}; and \code{Mesh}, which holds \code{N} of them.

\subsection*{Limits}

\code{mesh::lattice} makes the channels in simulation and does not
lower; \code{Mesh} is the same wiring as a unit that does. On a lattice
that \code{mesh::lattice} makes, a link at the edge drives a channel
nobody reads. Each of the two switches moves one packet per cycle, not
one per port.

% SPDX-License-Identifier: Apache-2.0
% covers: Mesh
\datasheet{Mesh}{A lattice of network nodes as one unit, which lowers.}

\dsfacts{\code{lib/parts/src/bus/noc/mesh.rs}}{\code{txhdl\_parts::bus::noc::mesh::Mesh}}%
{\code{//docs:noc}}

\subsection*{Function}

\code{Mesh<W, H, N, ..>} is \code{N} nodes, \code{W} wide and \code{H}
high, in row-major order, held as one field, \code{nodes:
Units<Node, N>} (issue 635). Node \code{i} is at column \code{i \% W}
and row \code{i / W}, which the mesh gives it by tying its \code{col}
and \code{row} inputs to those constants.

The mesh's ports are the nodes' exits, four arrays of \code{N}:
\code{qx\_in} the requests into each node, \code{px\_in} the responses
into each, and \code{qx\_out} and \code{px\_out} what leaves by each.
Everything between the nodes is inside. The links are one array of
channels per direction of travel and per virtual channel. Node
\code{i}'s output to a direction is channel \code{i} of that
direction's array, and its neighbour reads that channel on the input
from the opposite side.

\subsection*{Parameters}

\begin{center}\footnotesize
\begin{tabular}{@{}lp{0.7\textwidth}@{}}
\toprule
Parameter & Meaning \\
\midrule
\code{W} & The width of the lattice, in nodes. \\
\code{H} & The height of the lattice, in nodes. \\
\code{N} & The number of nodes, which must be \code{W * H}. \\
\code{XB} & The width of a column coordinate, in bits: at least enough
for \code{W - 1}. \\
\code{YB} & The width of a row coordinate, in bits: at least enough
for \code{H - 1}. \\
\code{A} & The address width of the AXI link. \\
\code{D} & The data width of the AXI link. \\
\code{S} & The strobe width, which is \code{D / 8}. \\
\code{I} & The identifier width of the AXI link. \\
\bottomrule
\end{tabular}
\end{center}

The tables below are for \code{Mesh<3, 2, 6, 2, 1, 8, 8, 1, 2>}, the
mesh \code{ex\_mesh} makes: three nodes wide and two high, on a link
of eight-bit addresses and data.

\dstables{Mesh}

\subsection*{Behaviour}

\begin{itemize}
\item The mesh has no logic of its own and no register. What it holds
is its nodes' state, six bits each.
\item Node \code{i} reads from the north what the node above sent
down, from the south what the node below sent up, from the west what
the node to its west sent east, and from the east what the node to its
east sent west.
\item A link at an edge wraps to the other end of its row or column,
so that every channel has one driver and one reader. Routing is
dimension order, so no packet for a place on the lattice ever takes
one: nothing leaves the east edge eastward, since no column is east of
it.
\item \code{Default} refuses a mesh whose \code{N} is not
\code{W * H}, with a message naming both, since Rust cannot yet
compute \code{N} from \code{W} and \code{H}.
\end{itemize}

\subsection*{Verification}

\begin{itemize}
\item \code{ex\_mesh} runs \code{Mesh<3, 2, 6, ..>}. Every node sends a
request to each of the other five, and a response to the node at the
opposite corner. The run asserts that each packet leaves by the exit
of the node it was addressed to, and that all thirty requests and six
responses arrive.
\item The netlist, six node instances and the channels between them,
is checked against that run under nvc and Verilator, by
\code{//docs:mesh\_sim\_mesh\_tb\_test} and \code{//docs:mesh\_vsim\_test},
at every port of the mesh.
\end{itemize}

\subsection*{Used by}

\code{ex\_mesh}. The tests of the network and \code{//soc} wire their
lattices with \code{mesh::lattice}, which does not lower.

\subsection*{Limits}

The lattice does not reject packets addressed beyond its boundaries,
such as columns of \code{W} or greater. Wrapped links remain idle only
while every destination lies within the grid. The nodes' registers are
checked at the mesh ports only, not inside the core logic. Each node
switch moves one packet per cycle rather than one packet per port.

% SPDX-License-Identifier: Apache-2.0
% covers: HostBridge
\datasheet{HostBridge}{The host side of an exit: AXI from a host in, packets into the network out.}

\dsfacts{\code{lib/parts/src/bus/noc/bridge.rs}}{\code{txhdl\_parts::bus::noc::bridge::HostBridge}}%
{\code{//docs:noc}}

\subsection*{Function}

\code{HostBridge} stands where an AXI host's five channels would go to
a peripheral, and sends them across the network instead. It reads the
host tracker's \code{aw}, \code{ar} and \code{w}, and drives its
\code{b} and \code{r}. On the network side it sends request packets on
\code{req} into its node's exit, and takes response packets from
\code{rsp}. It addresses each request by a map of three ranges.

\subsection*{Parameters}

\begin{center}\footnotesize
\begin{tabular}{@{}lp{0.7\textwidth}@{}}
\toprule
Parameter & Meaning \\
\midrule
\code{X}, \code{Y} & The column and row of the bridge's node. The
bridge stamps them on every request as its source. \\
\code{XB}, \code{YB} & The widths of a column and a row coordinate, in
bits. \\
\code{A} & The address width of the AXI link. \\
\code{D} & The data width of the AXI link. \\
\code{S} & The strobe width, which is \code{D / 8}. \\
\code{I} & The identifier width of the AXI link. \\
\code{B0}, \code{M0}, \code{X0}, \code{Y0} & The first range. An
address whose bits under the mask \code{M0} equal \code{B0} goes to
the node at \code{X0}, \code{Y0}. \\
\code{B1}, \code{M1}, \code{X1}, \code{Y1} & The second range, in the
same way. The first range wins where both match. \\
\code{B2}, \code{M2}, \code{X2}, \code{Y2} & The third range, the
default route. The doc comment says to give it a mask of zero. The
bridge's \code{run} does not read \code{B2} or \code{M2}: it sends
every address the first two ranges miss to \code{X2}, \code{Y2}. \\
\bottomrule
\end{tabular}
\end{center}

\subsection*{Address map}

The tables below are for \code{soc::Bridge<0, 0>}, the bridge of the
core at 0,0 in \code{soc/src/lib.rs}. That type is
\code{HostBridge<0, 0, 2, 2, 32, 32, 4, 2, 0x3000, 0xffff\_f000, 1, 1,
0x1000, 0xffff\_f000, 0, 1, 0, 0, 0, 1>}. Both hosts of \code{//soc}
use this map.

\begin{center}\footnotesize
\begin{tabular}{@{}llllp{0.4\textwidth}@{}}
\toprule
Range & Base & Mask & Node & What is there \\
\midrule
0 & \code{0x3000} & \code{0xffff\_f000} & 1,1 & The serial port. \\
1 & \code{0x1000} & \code{0xffff\_f000} & 0,1 & The program's data, in
the memory. \\
2 & \code{0} & \code{0} & 0,1 & Everything else, the display list and
the framebuffer included, in the memory. \\
\bottomrule
\end{tabular}
\end{center}

The masks reach the whole address. \code{//docs:noc} states why: with
a mask of \code{0xf000}, the framebuffer's rows at \code{0x13000} and
above went to the serial port.

\dstables{HostBridge}

\subsection*{Behaviour}

\begin{itemize}
\item A write goes when \code{aw} offers an address phase, \code{w}
offers a beat, and \code{req} has room. The bridge takes both and sends
one packet with \code{chan} \code{W}: the address phase from
\code{aw} with \code{len} zero, the data and strobe from \code{w}, and
\code{last} high. A write of one beat is done with that.
\item A write whose \code{len} is not zero is split (issue 125). Its
first beat goes as above, and the phase is held: \code{left}, the
beats still to go; \code{saddr}, the address the next one is for;
\code{sid}, its identifier; \code{sfix}, whether the burst is fixed;
and \code{ssize}, \code{slock}, \code{scache}, \code{sprot},
\code{sqos} and \code{sregion}, sent again with every beat. Each beat
after the first goes when \code{w} offers it and \code{req} has room,
as a single-beat write at \code{saddr}, which then moves by a beat's
width, \code{S} bytes, or not at all for a fixed burst. The beats are
routed by their own addresses.
\item The answers to a split burst are merged: \code{pend} counts
those still to come and \code{err} whether any was not \code{OKAY}. A
\code{B} packet under the burst's identifier is counted rather than
passed on, and the one that completes the count is answered to
\code{b} in its place, \code{SLVERR} if \code{err} was set and
\code{OKAY} otherwise, when \code{b} is ready. A long write whose
first beat is offered while another is still being merged waits;
writes of one beat and reads do not.
\item A wrapping burst is refused rather than split, since the wrap is
not computed here: \code{eat}, whether its beats are still arriving;
\code{owe}, whether it is waiting to be answered; and \code{bad}, the
identifier to answer it under. The bridge takes its address phase,
sends no packet for it, drops every beat until the one marked
\code{last}, and answers \code{b} with \code{SLVERR} in the first
cycle \code{b} is ready and nothing else is being answered. Before
the split every long write was refused so, which was issue 280, and
before the refusal it hung.
\item A read goes when no write goes in that cycle, \code{ar} offers an
address phase, and \code{req} has room. The bridge takes it and sends
one packet with \code{chan} \code{Ar}.
\item The bridge sends at most one request packet per cycle, and a
write before a read.
\item Every request packet has the destination the map gives for its
address, the source \code{X}, \code{Y}, and the host's own identifier.
\item A response packet with \code{chan} \code{B} goes to \code{b} as
its identifier and response, when \code{b} is ready. Any other response
packet goes to \code{r} as its identifier, data, response and
\code{last}, when \code{r} is ready.
\end{itemize}

\subsection*{Verification}

\begin{itemize}
\item Eight tests in \code{lib/parts/src/bus/noc.rs}, in
\code{//lib/parts:parts\_test}, put host bridges on a two by two
lattice with 16-bit addresses. Their map sends \code{0x1000} and
\code{0x2000}, each under the mask \code{0xf000}, and the default, all
to the node at 1,1.
\item \code{a\_burst\_crosses\_the\_lattice\_and\_its\_answer\_comes\_back}
has one host, at 0,0.
\item \code{two\_hosts\_share\_a\_memory\_and\_each\_answer\_goes\_home}
has hosts at 0,0 and 1,0 on one memory.
\item \code{four\_peripherals\_share\_one\_node} has hosts at 0,0 and
1,0, and four memories behind a router at 1,1.
\item \code{a\_write\_of\_two\_beats\_crosses\_the\_lattice\_as\_two\_writes}
asks for two words at once and finds both in the memory, each at its
word, and \code{OKAY}.
\item \code{a\_fixed\_write\_lands\_every\_beat\_on\_one\_word} sends
three beats fixed and finds the last of them at the one word and the
next word untouched.
\item \code{two\_hosts\_write\_sixteen\_beats\_at\_once\_and\_every\_word\_lands}
has hosts at 0,0 and 1,0 each writing sixteen words to one memory
at once and reading them back: every word in its place and in order,
with the two bursts cut into each other on the way.
\item \code{a\_wrapping\_write\_is\_refused\_and\_the\_one\_after\_it\_is\_served}
asks for a wrapping burst, finds \code{SLVERR} and a memory nobody
wrote, and then a write of one beat served, so the dropped beats left
the channel clean.
\item \code{a\_core\_draws\_a\_solid\_and\_a\_gpu\_fills\_it}, in
\code{//soc:demo\_test}, runs \code{soc::Bridge} at 0,0 for the core
and at 1,0 for the rasteriser.
\item No waveform in \code{docs/BUILD.bazel} lowers the bridge, so its
netlist has no co-simulation under nvc or Verilator.
\code{//docs:showcase} lists the network's netlists as not yet
simulated.
\end{itemize}

\subsection*{Used by}

\code{soc::Bridge} in \code{soc/src/lib.rs}, for the core at 0,0 and
the rasteriser at 1,0. The tests of the network in
\code{lib/parts/src/bus/noc.rs}.

\subsection*{Limits}

A write of one beat is one packet, and a burst of more crosses as
that many single-beat writes: a peripheral that means something by a
burst, a controller that opens a row for one, sees the beats and not
the burst. One burst is split at a time, so a second long write waits
for the first's answer. The address moves by the link's width and
not by \code{size}, which nothing in this tree acts on. A wrapping
burst is refused with \code{SLVERR} rather than executed. The map has
three ranges. The last one takes every address the others miss. So no
address is unknown to the network, and neither bridge has to answer for
one.

% SPDX-License-Identifier: Apache-2.0
% covers: PerBridge
\datasheet{PerBridge}{The peripheral side of an exit: request packets in, AXI to a peripheral out.}

\dsfacts{\code{lib/parts/src/bus/noc/bridge.rs}}{\code{txhdl\_parts::bus::noc::bridge::PerBridge}}%
{\code{//docs:noc}}

\subsection*{Function}

\code{PerBridge} is the other end of \code{HostBridge}. It takes
request packets on \code{req} from its node's exit and unpacks them
into \code{aw}, \code{ar} and \code{w} for a peripheral's tracker. It
takes the tracker's answers on \code{b} and \code{r}, and packs them on
\code{rsp} back to the node each request came from.

It also renames. An identifier belongs to the host that made it, so
two hosts may use the same one. The bridge gives every request a local
identifier of its own, and remembers against it the source node and
the identifier that node used. The answer comes back under the local
identifier, and the bridge sends the packet home under the original.

\subsection*{Parameters}

\begin{center}\footnotesize
\begin{tabular}{@{}lp{0.7\textwidth}@{}}
\toprule
Parameter & Meaning \\
\midrule
\code{X}, \code{Y} & The column and row of the bridge's node. The
bridge stamps them on every answer as its source. \\
\code{XB}, \code{YB} & The widths of a column and a row coordinate, in
bits. \\
\code{A} & The address width of the AXI link. \\
\code{D} & The data width of the AXI link. \\
\code{S} & The strobe width, which is \code{D / 8}. \\
\code{I} & The identifier width of the AXI link. \\
\code{NIDS} & How many local identifiers the bridge gives out, and so
how many bursts the peripheral has in flight at most. \\
\bottomrule
\end{tabular}
\end{center}

The tables below are for \code{PerBridge<0, 1, 2, 2, 32, 32, 4, 2,
4>}, the bridge of the memory at 0,1 in \code{soc/src/lib.rs}.

\dstables{PerBridge}

\subsection*{Behaviour}

\begin{itemize}
\item The local identifier on offer is \code{turn}. It is free when its
bit in \code{busy} is low.
\item A request packet with \code{chan} \code{W} is a write. The bridge
takes it when \code{turn} is free and both \code{aw} and \code{w} are
ready. It sends the address phase on \code{aw} under the local
identifier, and the data and strobe on \code{w} with \code{last} high.
\item The bridge treats any other request packet as a read. It takes it
when \code{turn} is free and \code{ar} is ready, and sends the address
phase on \code{ar} under the local identifier.
\item On taking a request, the bridge stores the source column in
\code{sx}, the source row in \code{sy} and the original identifier in
\code{oid}, each at the local identifier. It sets that identifier's
\code{busy} bit and adds one to \code{turn}.
\item A request waits while the identifier whose turn it is has not
been answered. The bridge does not look for another free identifier.
\item The bridge sends at most one answer packet per cycle, when
\code{rsp} has room, and a write response before a read beat. The
packet goes to the \code{sx} and \code{sy} stored under the answer's
identifier, under the \code{oid} stored there, with the source
\code{X}, \code{Y}.
\item A write response goes with \code{chan} \code{B}, its response,
and \code{last} high. A read beat goes with \code{chan} \code{R}, its
data, its response and its \code{last}.
\item The bridge clears an identifier's \code{busy} bit when it sends a
write response, or the last beat of a read, under that identifier.
\end{itemize}

\subsection*{Verification}

\begin{itemize}
\item Three tests in \code{lib/parts/src/bus/noc.rs}, in
\code{//lib/parts:parts\_test}, put a
\code{PerBridge<1, 1, 2, 2, 16, 32, 4, 2, 4>} in front of the memory
at 1,1.
\code{a\_burst\_crosses\_the\_lattice\_and\_its\_answer\_comes\_back}
writes and reads one word from one host.
\code{two\_hosts\_share\_a\_memory\_and\_each\_answer\_goes\_home} checks
that each answer reaches the host that asked. \code{//docs:noc} states
that this test failed before the renaming was written.
\code{four\_peripherals\_share\_one\_node} puts a router of four
memories behind the bridge, and checks each word and the router's
\code{DecErr} for an address none of them has.
\item \code{a\_core\_draws\_a\_solid\_and\_a\_gpu\_fills\_it}, in
\code{//soc:demo\_test}, runs a bridge at 0,1 for the memory and one at
1,1 for the serial port.
\item No waveform in \code{docs/BUILD.bazel} lowers the bridge, so its
netlist has no co-simulation under nvc or Verilator.
\code{//docs:showcase} lists the network's netlists as not yet
simulated.
\end{itemize}

\subsection*{Used by}

\code{soc::run} in \code{soc/src/lib.rs}: at 0,1 in front of the
memory's \code{AxiPer}, and at 1,1 in front of the serial port, behind
\code{LiteBridge<1, ..>}. The tests of the network in
\code{lib/parts/src/bus/noc.rs}.

\subsection*{Limits}

A write request is one packet with one beat, so a burst of more than
one beat does not cross the network. At most \code{NIDS} bursts are in
flight at the peripheral. \code{turn} is a register of \code{I} bits
and \code{busy} has \code{NIDS} bits. Every use in the tree sets
\code{NIDS} to 4 with \code{I} of 2. The bridge sends one answer per
cycle.

% Peripherals for the board.
% SPDX-License-Identifier: Apache-2.0
% covers: EthTx
\datasheet{EthTx}{The transmit half of the Ethernet MAC: whole frames from a channel onto GMII.}

\dsfacts{\code{lib/parts/src/eth.rs}}{\code{txhdl\_parts::eth::EthTx}}%
{\code{//docs:eth}}

\subsection*{Function}

\code{EthTx} takes a frame's bytes from the channel \code{tx}, each an
\code{EthByte}: the byte and whether it is the frame's last. A frame
there is the destination address, the source address, the type or
length and the payload. The unit stores the whole frame, then sends it
on GMII: \code{txd}, a byte per cycle, and \code{tx\_en}. It adds the
preamble, the start of frame delimiter, padding up to sixty bytes and
the four bytes of the CRC-32 frame check sequence. The unit runs on one
clock, so on a board it can run on the transmit clock.

\subsection*{Parameters}

\code{EthTx} has no type parameters. Its frame memory holds
\code{FRAME\_MAX} bytes, 2048, a constant of the module.

\dstables{EthTx}

\subsection*{Behaviour}

The unit is two processes. The store takes the frame in; the sender is
the wire's protocol written as the sequence it is: wait for a whole
frame, the preamble and the delimiter, the frame's bytes, zeros up to
sixty, the check sequence and the gap, a byte a cycle, each a turn of
a loop. The lowering numbers the sender's waits into a state register
the unit does not declare, \code{at\_wait}, keeps its loops' counts in
\code{for0} to \code{for3}, and holds the two outputs between states in
\code{txd\_held} and \code{tx\_en\_held}; all are in the tables above.

\begin{itemize}
\item While no frame is whole, the store takes a byte from \code{tx}
whenever one is offered, and stores it in \code{frame} at \code{fill}.
The byte marked last sets \code{whole}.
\item A whole frame starts the sender: seven bytes of \code{0x55}, then
\code{0xd5}, with the CRC register set to \code{CRC\_INIT}, every bit
set. The frame follows, the register running over it.
\item A frame of fewer than sixty bytes is followed by zeros up to
sixty, with the register still running.
\item The check sequence is the register complemented, least
significant byte first.
\item \code{tx\_en} is high from the preamble to the end of the check
sequence. Twelve bytes of gap follow with \code{tx\_en} low. Then
\code{sent} is up for one cycle, which adds one to \code{frames} and
lets the store forget the frame and take the next.
\item The store takes no byte from \code{tx} while a frame is whole,
from its last byte stored to the end of its gap, so it holds the client
off. The client may offer bytes as slowly as it likes.
\item \code{crc\_byte} steps the register by a whole byte: the byte
added to it, then eight shifts with the polynomial \code{0xedb88320}
going in whenever the bit that leaves is set.
\end{itemize}

\subsection*{Verification}

\begin{itemize}
\item The tests below are in \code{lib/parts/src/eth.rs}, and run in
\code{//lib/parts:parts\_test}.
\item \code{the\_byte\_wide\_crc\_is\_the\_bitwise\_one} checks
\code{crc\_byte} against the register stepped a bit at a time, on every
byte from four registers.
\item \code{the\_crc\_of\_a\_known\_string} checks \code{crc32} against
the catalogue value and the residue.
\item \code{a\_frame\_goes\_out\_as\_the\_model\_says\_and\_comes\_back}
loops the transmitter into \code{EthRx}, and holds the bytes on the
wire to \code{wire\_bytes}, for frames of 20 and 100 bytes.
\item \code{frames\_at\_the\_minimum\_length} does the same for frames
of 59, 60 and 61 bytes.
\item \code{ex\_eth} runs the unit behind \code{EthLite} with the wire
looped back, and asserts that it sent three frames. The build simulates
the netlist of \code{EthTx} against that run under nvc and Verilator:
\code{//docs:eth\_sim\_eth\_tx\_eth\_tx\_tb\_test} and
\code{//docs:eth\_vsim\_eth\_tx\_test}.
\end{itemize}

\subsection*{Used by}

The example \code{ex\_eth}. The echo design for the AX7A200B:
\code{//eth:netlist} writes the Verilog of \code{eth\_tx}, and
\code{eth/board/eth\_echo.v} instantiates it on the 125~MHz transmit
clock. \code{//eth:echo\_synth} and \code{//eth:echo\_pnr} put that
design through Vivado to a bitstream. \code{//docs:eth} states that it
has not yet run on the board.

\subsection*{Limits}

The unit holds one frame. A frame longer than \code{FRAME\_MAX} bytes
keeps overwriting its last byte. Only gigabit works: at 100 or
10~Mbit a PHY sends a nibble per cycle, and the unit does not take a
byte over two cycles.

% SPDX-License-Identifier: Apache-2.0
% covers: EthRx
\datasheet{EthRx}{The receive half of the Ethernet MAC: frames off GMII, checked, onto a channel.}

\dsfacts{\code{lib/parts/src/eth.rs}}{\code{txhdl\_parts::eth::EthRx}}%
{\code{//docs:eth}}

\subsection*{Function}

\code{EthRx} takes a frame off GMII: \code{rxd}, a byte per cycle,
\code{rx\_dv} and \code{rx\_er}. It stores the whole frame and checks
its CRC-32 frame check sequence. Only then does it offer the frame's
bytes on the channel \code{rx}, each an \code{EthByte}, without the
preamble and the check sequence. It drops a frame whose check fails.
While it offers a frame, \code{rx\_len} says how many bytes the frame
has, so a consumer that must know the length before it has taken the
frame can read it; \code{EthSlots} does.
The unit runs on one clock, so on a board it can run on the receive
clock the PHY recovers from the line.

\subsection*{Parameters}

\code{EthRx} has no type parameters. Its frame memory holds
\code{FRAME\_MAX} bytes, 2048, a constant of the module.

\dstables{EthRx}

\subsection*{Behaviour}

The unit is two processes. The receiver watches the wire every cycle;
the offerer is the sequence of giving a frame away: wait for a full
frame, then its bytes without the check sequence, one per turn, each
held until the consumer takes it. The lowering numbers the offerer's
waits into a state register the unit does not declare, \code{at\_wait},
keeps its loop's count in \code{for0}, and holds \code{rx\_len} between
states in \code{rx\_len\_held}; all are in the tables above.

\begin{itemize}
\item While hunting, that is while it is neither receiving, holding a
full frame nor skipping, the receiver starts a frame on \code{rx\_dv}
high with \code{0xd5} on \code{rxd}: it sets \code{receiving}, clears
\code{len} and \code{bad}, and sets the CRC register to
\code{CRC\_INIT}. It passes over \code{0x55}.
\item Any other byte with \code{rx\_dv} high while hunting is the middle
of a frame whose start the unit missed. The receiver sets \code{skip}
and ignores the line until \code{rx\_dv} falls.
\item While receiving, the receiver stores a byte per cycle for as long
as \code{rx\_dv} stays high. The CRC register runs over every byte, the
check sequence included. \code{bad} records any cycle with
\code{rx\_er} high.
\item When \code{rx\_dv} falls, the frame is good if \code{bad} is low,
the register holds \code{CRC\_RESIDUE}, and \code{len} is more than
four. A good frame sets \code{full}. The receiver drops any other frame
and adds one to \code{dropped}.
\item A full frame starts the offerer. \code{rx\_len} reads the frame's
length without the check sequence while it is offered, and zero at
every other time. The offerer sends \code{len} less four bytes on
\code{rx}, each held until \code{rx} takes it, and marks the last. Then
\code{given} is up for one cycle, which adds one to \code{frames} and
clears \code{full}, and the receiver hunts again.
\item A frame that arrives while a frame is full is dropped: the
receiver sets \code{skip} and adds one to \code{dropped}, once for that
frame.
\end{itemize}

\subsection*{Verification}

\begin{itemize}
\item The tests below are in \code{lib/parts/src/eth.rs}, and run in
\code{//lib/parts:parts\_test}. Each loops \code{EthTx} into the
receiver a cycle late.
\item \code{a\_frame\_goes\_out\_as\_the\_model\_says\_and\_comes\_back}
checks that a frame of 20 bytes comes back padded to sixty, and a frame
of 100 bytes comes back as it went.
\item \code{frames\_at\_the\_minimum\_length} does the same at 59, 60
and 61 bytes.
\item \code{a\_corrupted\_frame\_is\_dropped} flips a byte of the second
of three frames. It checks that the other two come back and that one
frame is counted dropped.
\item \code{a\_frame\_on\_a\_busy\_receiver\_is\_dropped} holds the
client off. It checks that the second frame is dropped and counted, and
that the first is intact.
\item \code{ex\_eth} runs the unit with one byte of the second frame
flipped on the wire, and asserts that the receiver dropped one frame.
The build simulates the netlist of \code{EthRx} against that run under
nvc and Verilator: \code{//docs:eth\_sim\_eth\_rx\_eth\_rx\_tb\_test}
and \code{//docs:eth\_vsim\_eth\_rx\_test}.
\end{itemize}

\subsection*{Used by}

The example \code{ex\_eth}. The echo design for the AX7A200B:
\code{//eth:netlist} writes the Verilog of \code{eth\_rx}, and
\code{eth/board/eth\_echo.v} instantiates it on the PHY's receive
clock. \code{eth/board/eth\_cdc.v} takes its bytes across to the
transmit clock. \code{//eth:echo\_synth} and \code{//eth:echo\_pnr} put
that design through Vivado to a bitstream. \code{//docs:eth} states
that it has not yet run on the board.

\subsection*{Limits}

The unit holds one frame. A client that has not yet taken a frame
loses the next one. \code{len} stops at 2047, so the unit stores any
further bytes of a longer frame in the last slot. Only gigabit works:
at 100 or 10~Mbit a PHY sends a nibble per cycle, and the unit does not
take a byte over two cycles. The unit reads \code{rx\_er} only while
receiving.

% SPDX-License-Identifier: Apache-2.0
% covers: EthLite
\datasheet{EthLite}{The Ethernet peripheral on AXI-Lite: three words, a byte per transaction.}

\dsfacts{\code{lib/parts/src/eth.rs}}{\code{txhdl\_parts::eth::EthLite}}%
{\code{//docs:eth}}

\subsection*{Function}

\code{EthLite} stands between an AXI-Lite host and the two halves of
the MAC. It takes bytes from AXI-Lite writes and sends them to
\code{EthTx} on \code{tx}. It answers AXI-Lite reads with the bytes
\code{EthRx} offers on \code{rx}. Its AXI-Lite side is one port,
\code{bus}, a \code{LitePort} with 32-bit addresses and words:
\code{bus\_aw}, \code{bus\_ar} and \code{bus\_w} in, and \code{bus\_b}
and \code{bus\_r} out. \code{irq} is high while a received byte waits.

\subsection*{Parameters}

\code{EthLite} has no type parameters.

\subsection*{Register map}

The unit decodes the word from address bits 3 to 2, so the words are at
offsets \code{0x0}, \code{0x4} and \code{0x8} of its range.
\code{ex\_eth} puts them at \code{0x1000}, \code{0x1004} and
\code{0x1008}.

\dsregs{EthLite}

A write to \code{txbyte} is answered only when the transmitter takes
the byte, so a program that writes a frame a byte at a time is held
back by the wire and not by polling. A read of \code{rxbyte} takes the
oldest byte, which is what its access says; when nothing waits it
reads with \code{valid} low.

\dstables{EthLite}

\subsection*{Behaviour}

\begin{itemize}
\item A write goes when \code{bus\_aw} and \code{bus\_w} both offer
and \code{bus\_b} is ready. A write to the transmit word also waits
until \code{tx} is ready, so the unit answers it only when the
transmitter has taken the
byte. A client that writes faster than the transmitter takes is held
off rather than losing bytes.
\item A write to the transmit word sends bits 7 to 0 and bit 8 of the
written data on \code{tx}. The unit answers every write
\code{Okay}; a write to any other word does nothing else.
\item A read goes when \code{bus\_ar} offers and \code{bus\_r} is ready.
The unit
answers every read \code{Okay}. A read of word 1 or word 3 answers
zero.
\item A read of the receive word takes the waiting byte, if one waits,
and answers with bit 9 high. With none waiting it answers with bit 9
low and takes nothing.
\item \code{irq} is high while \code{rx} offers a byte.
\item The unit holds no state. It does not read the write strobe.
\end{itemize}

\subsection*{Verification}

\begin{itemize}
\item \code{EthLite} has no \code{\#[test]} of its own.
\item \code{ex\_eth} puts it behind \code{LiteBridge<1, ..>} and an
\code{AxiHost}, in front of \code{EthTx} and \code{EthRx} on a looped
wire. The host sends three frames, each as one fixed burst to the
transmit word, and reads them back in fixed bursts of sixteen from the
receive word. The run asserts that the first comes back padded to
sixty bytes, the second, corrupted on the wire, does not come back, and
the third comes back as it went.
\item The build simulates the netlist of \code{EthLite} against that
run under nvc and Verilator:
\code{//docs:eth\_sim\_eth\_lite\_eth\_lite\_tb\_test} and
\code{//docs:eth\_vsim\_eth\_lite\_test}.
\end{itemize}

\subsection*{Used by}

The example \code{ex\_eth}. \code{//docs:eth} states that the
peripheral is not yet on Vreteno's board, and the echo design for the
AX7A200B does not use it.

\subsection*{Limits}

The unit moves one byte per transaction. The AXI-Lite bridge makes a
fixed burst one transaction per beat at the same address, so a host
sends a whole frame as one burst. The unit checks no more of the
address than the word: the bridge in front sends it only its own range.

% SPDX-License-Identifier: Apache-2.0
% covers: EthSlots
\datasheet{EthSlots}{The Ethernet slot registers on AXI-Lite: LiteEth's
map, with the frame buffers in main memory.}

\dsfacts{\code{lib/parts/src/ethslots.rs}}%
{\code{txhdl\_parts::ethslots::EthSlots}}{\code{//docs:examples}}

\subsection*{Function}

\code{EthSlots} is the register map of LiteX's LiteEth in hardware,
at LiteX's offsets, so that Linux's \code{litex\_liteeth}, whose
offsets are fixed, drives it unchanged (issue 1203), and Zephyr's
\code{drivers/ethernet/eth\_litex\_liteeth.c} takes them from the
header \code{//tools/regmap} writes. Bringing up Ethernet is a port of
a driver that already works rather than a design nobody has driven.

Its AXI-Lite side is one port, \code{bus}, a \code{LitePort} with
32-bit addresses and words: \code{bus\_aw}, \code{bus\_ar} and
\code{bus\_w} in, and \code{bus\_b} and \code{bus\_r} out. Ten
further ports face the engines that move the bytes and the unit that
receives them: \code{tx\_busy}, \code{rx\_busy}, \code{rx\_len},
\code{rx\_which} and \code{rx\_drops} in, and \code{tx\_base},
\code{tx\_bytes}, \code{tx\_start}, \code{rx\_base} and
\code{rx\_full} out. \code{irq} is high while a received frame is
unacknowledged and the receive interrupt is enabled.

It holds no frame's bytes. There are two slots each way, used in turn,
and each is a flat buffer of 2048 bytes in main memory rather than
inside the peripheral, so that \code{LineFetch} and \code{LineStore}
move the bytes instead of the processor. The unit's work is to say
which buffer and how many bytes, to tell the driver when a frame has
landed, and to keep the frames that have landed and not been
acknowledged in order, a slot and a length each, at most one a slot
(issue 1313).

Two slots a direction rather than a descriptor ring, because Zephyr's
driver model never hands a chain to hardware:
\code{net\_pkt\_read(pkt, buf, len)} and \code{net\_pkt\_write(pkt,
buf, len)} take a flat pointer and a byte count, so the hardware only
ever sees one contiguous region per frame.

\subsection*{Parameters}

\code{BASE}, where the four buffers begin. A slot is \code{BASE + n *
2048}, receive slots first and transmit slots 4096 bytes above them.
The board uses \code{0x4100\_0000}, which \code{isa::ETH\_BUF\_BASE}
states.

The two directions are separate regions and have to be. A slot number
alone would make transmit slot zero and receive slot zero the same
address, and a frame arriving would land on one waiting to go out.

\textbf{\code{BASE} must be 1 KB aligned, and nothing checks it.} A
slot is then 1 KB aligned too, since the slots are 2 KB apart, and a
256 beat burst of words is exactly 1 KB, so no burst can cross AXI4's
4 KB boundary. An unaligned \code{BASE} puts that obligation on the
caller, and a burst that straddles a boundary is a protocol violation
rather than a wrong address, so it will not look like a fault in this
unit.

\subsection*{Register map}

The unit decodes the word from address bits 5 to 2, so the words are
at offsets \code{0x0} to \code{0x3c} of its range. The board puts
them at \code{0x3400}, which \code{isa::ETH\_BASE} states.

\dsregs{EthSlots}

The fourteen words are LiteX's, in LiteX's order, and words 14 and
15 are named by no register. Three words exist for the drivers and
hold nothing here: \code{tx\_level} reads zero, since one transmit
runs at a time, and \code{tx\_ev\_status} reads zero, since the
transmit raises no event. \code{rx\_errors} reads \code{rx\_drops},
the frames the receiving side dropped because both slots held a frame
the driver had not acknowledged; it read zero, and said none were
dropped, before issue 1313. \code{rx\_ev\_status} reads the same bit as
\code{rx\_ev\_pending}. \code{tx\_ev\_pending} is never set, as
the behaviour below says, and \code{tx\_ev\_enable} is kept so the
driver can write it.

The write-only words, \code{tx\_start}, \code{tx\_slot} and
\code{tx\_length}, and the unnamed 14 and 15 read zero. They once
read as \code{tx\_ev\_enable}, the read path's fall-through, so a
driver that wrote \code{tx\_length} and read it back got an unrelated
bit (issue 454); every word the read names now, and the rest are
zero.

\emph{The transmit side has no interrupt.} LiteEth's driver disables
it and polls \code{tx\_ready}, because Zephyr's send may block, so
\code{tx\_ev\_pending} is never set by anything and
\code{tx\_ev\_enable} gates nothing. Both words exist because the
driver acknowledges them. A driver that waited on
\code{tx\_ev\_pending} would wait for ever.

\dstables{EthSlots}

\subsection*{Behaviour}

\begin{itemize}
\item A read is answered in the cycle the request is taken, and a
write is answered with \code{bus\_b}. The unit never stalls: it needs
nothing the other AXI-Lite peripherals do not.
\item \textbf{An arrival is the store engine falling idle, not the
frame being received.} \code{LineStore} holds \code{running} high
until the write response comes back from memory, so its fall is the
first moment the frame is certainly there. The unit takes that level
rather than a completion pulse, deliberately: a pulse could be joined
to the receiver by somebody reading the map and not the comment, and
the interrupt would then tell a driver to read a frame whose last
beats are still in flight.
\item \textbf{Frames wait in order (issue 1313).} An arrival puts
\code{rx\_which} and \code{rx\_len} behind any frame already waiting.
\code{rx\_slot} and \code{rx\_length} read the oldest, and
\code{rx\_ev\_pending} reads one while any waits. Writing one to it
takes the oldest off, so a driver that acknowledges one frame an
interrupt, as Zephyr's and Linux's do, is interrupted again for the
next. Before issue 1313 an arrival replaced the slot and length the
driver had been told of, and a driver busy with one frame was told of
the last of a burst only.
\item \textbf{An acknowledgement in the same cycle as an arrival does
not lose the frame.} The oldest frame is taken off and the new one put
behind what is left, so the second frame comes to the front.
\item \textbf{A frame with no slot is dropped and counted.}
\code{rx\_full} is high while both slots hold a frame, or while one
does and the frame just stored is about to be counted. The receiving
side, \code{FrameIn}, then takes the next frame's bytes, stores none
of them, and counts it, and \code{rx\_errors} reads the count. A
frame dropped at the input is one the sender retransmits; a frame
landed on one being read was two frames lost, silently.
\item \code{tx\_start} is a request rather than a pulse the engine
must catch. Writing one sets a flag; the unit drives \code{tx\_start}
high while the flag is set and the engine is not busy, and clears the
flag when the engine goes busy. So the request lasts until it is
taken.
\item \code{tx\_base} is \code{BASE + 4096 + tx\_slot * 2048} and
\code{rx\_base} is \code{BASE + rx\_which * 2048}. The receiving side
is told where the slot it is filling begins, and which slot that is
comes from the engine, so that a frame lands somewhere the driver is
not reading.
\item \code{irq} is high while any frame waits and
\code{rx\_ev\_enable} is set. It is a level, and the board's
interrupt controller asks again for a level still high on the
complete.
\end{itemize}

\subsection*{What a driver must do}

\textbf{Read the last word of a frame back before writing
\code{tx\_start}.} Stores on this machine are posted, so without
an order the fetch engine may read the slot before the frame has
landed and send what was there before. A \code{fence} now waits
until every posted store has been answered (issue 432), and either
the fence or the read orders them; the read is what was proven on
the board, and \code{ethtx.rs} keeps it.

What makes the obligation dischargeable is measured rather than
assumed: this link answers a load behind an unanswered store to the
same address with the stored value, which
\code{a\_load\_behind\_an\_unanswered\_store\_to\_one\_address} in
\code{lib/parts/src/bus/axi/sim.rs} asserts, and \code{//docs:axi}
states.

\textbf{That is this interconnect's behaviour and not AXI's
guarantee.} AXI does not order a read behind a write, and a driver
written for another fabric on the strength of this paragraph would be
wrong. The reason the read works here is the link, not the memory
being a serialisation point.

\subsection*{Verification}

\begin{itemize}
\item \code{only\_word\_13\_reads\_tx\_ev\_enable}, in
\code{//lib/parts:parts\_test}, sets \code{tx\_ev\_enable} and reads
every word, so a word that mirrored it again would fail (issue 454).
\item \code{ex\_ethslots} drives the map as
\code{eth\_litex\_liteeth.c} will, in the order that driver does it:
\code{tx\_ready} one while idle and zero while a frame is going, no
interrupt while the store is in flight and one after it, and the
pending bit tested both ways, since a write of zero must not
acknowledge.
\item It also tests the cycle where an acknowledgement meets an
arrival. The cycle to issue at was measured against both orders of
the drives rather than reasoned about, because two earlier attempts
at that test passed against the faulty order.
\item It lands a third frame while the second waits, and checks that
it waits behind it, that \code{rx\_full} is up, that
\code{rx\_errors} reads the receiving side's count, and that two
acknowledgements take the two frames off in order (issue 1313).
\item The build simulates the netlist against that run under nvc and
Verilator: \code{//docs:ethslots\_sim\_ethslots\_tb\_test} and
\code{//docs:ethslots\_vsim\_test}.
\item On the board, with the engines behind it, the two Ethernet tests
in \code{cpu/vreteno/tests/board.rs} run programs on the core that
use its registers as the driver does: one waits for an arrival and
reads the frame from the slot and length it names, and the other
sends two frames back to back, waiting on \code{tx\_ready} between
them.
\item \code{frames\_beyond\_the\_slots\_are\_dropped\_and\_counted\_not\_lost},
in the same file, sends eight full frames back to back to a driver
that spends five thousand loop turns on each: it is told of three,
intact, and the other five are counted. Before issue 1313 it was told
of two, both overwritten under it, and \code{rx\_errors} read zero.
\end{itemize}

\subsection*{Used by}

The board, at \code{cpu/vreteno/src/board.rs}, on the fifth AXI-Lite
slot at \code{0x3400}, with its arrival as the third source of
\code{Plic3}. Its engine-facing ports go to the fetch engine and
\code{FrameOut} on the way out, and to \code{FrameIn} and the store
engine on the way in (issue 151). The example \code{ex\_ethslots}.

\subsection*{Limits}

It moves no frames itself: the engines behind it do. Without them,
as in \code{ex\_ethslots}, a driver binds and reads and writes every
register, and no frame arrives or leaves.

Two slots a direction is queue depth and not throughput: a burst
longer than two frames, while the driver is busy, loses the rest,
counted. A ring, if
it is ever wanted, is an addition rather than a rewrite: a descriptor
states the same three values \code{tx\_slot}, \code{tx\_length} and
\code{tx\_start} already state.

The unit checks no more of the address than the word, because the
bridge in front of it sends it only its own range.

% SPDX-License-Identifier: Apache-2.0
% covers: FrameOut
\datasheet{FrameOut}{Words from a fetch engine, bytes to the Ethernet
transmitter, the last of them the frame's last.}

\dsfacts{\code{lib/parts/src/ethdma.rs}}{\code{txhdl\_parts::ethdma::FrameOut}}%
{\code{//docs:examples}}

\subsection*{Function}

\code{LineFetch} moves 32-bit words, because that is what the bus
moves. A frame is a count of bytes, from 64 to 1518 of them, so its
last word holds one, two or three real bytes as often as four.
\code{FrameOut} is that difference on the way out, and nothing else:
it does not touch the bus and does not know where the frame is in
memory.

It takes a word at a time on its input channel and emits its bytes,
lowest first, as \code{EthByte}s. It marks the last byte of the
\emph{frame}, not the last byte of the last word, so the bytes that
sit above the frame's end in a partial last word are never sent.

\code{bytes} is the frame's length and \code{go} asks for a frame.
\code{running} is high while one is going. \code{nwords} is the word
count a fetch engine needs for a frame of \code{bytes}, that is
\code{(bytes + 3) >> 2}.

\subsection*{Parameters}

None. The byte count is 16 bits wide and the words are 32.

\dstables{FrameOut}

\subsection*{Behaviour}

\begin{itemize}
\item \code{go} is a level, not a pulse. A request is taken when the
unit is idle, and \code{bytes} is latched then, so the register block
may hold \code{go} until \code{running} says it was taken.
\item A run ends when the bytes emitted reach the count latched at
the start. A count of zero ends the run as it starts and sends
nothing.
\item A word is taken when the last one is spent, which leaves one
cycle in five with no byte. That costs nothing on the wire: the
transmitter stores a whole frame before it sends any of it, so this
side is never what the line waits for.
\item \code{nwords} is driven from \code{bytes} as it is, not from the
count latched at the start, because the fetch engine reads it in the
cycle it is told to go. The division is here because a board is only
wires between its units, and this is the one unit that stands between
a count of bytes and a count of words.
\end{itemize}

\subsection*{Verification}

\begin{itemize}
\item \code{ex\_ethdma} sends two frames, of 101 and 67 bytes, from
\code{FrameOut} through the real \code{EthTx} and \code{EthRx} to
\code{FrameIn}, and checks every word that comes back. Neither length
divides by four, and the two leave one and three real bytes in their
last words.
\item Two frames, because one cannot catch a byte side that reads past
the count. Such a unit still marks the right byte as the last, so the
first frame is the right length and passes every check. The bytes it
read past that frame wait in the transmitter, which saw no last byte
for them, and go out in front of the next frame.
\item The netlist is checked against that run under nvc and
Verilator: \code{//docs:ethdma\_sim\_frame\_out\_frame\_out\_tb\_test}
and \code{//docs:ethdma\_vsim\_frame\_out\_test}.
\item \code{ex\_framelen} runs it again in front of \code{FrameLen},
with the same two frames.
\item On the board,
\code{two\_frames\_written\_to\_memory\_leave\_the\_port\_as\_written\_in\_order}
in \code{cpu/vreteno/tests/board.rs} has the core write two frames of
23 and 18 bytes and ask for both back to back, and it checks the bytes
that leave the port.
\end{itemize}

\subsection*{Used by}

\code{Board}, as \code{fout}, between the fetch engine \code{fetch}
and \code{EthShare}. \code{EthSlots} gives it \code{bytes} and
\code{go}, and reads \code{running} as \code{tx\_busy}. The examples
\code{ex\_ethdma} and \code{ex\_framelen}.

\subsection*{Limits}

It trusts the count. A count larger than the frame in memory sends
whatever the fetch engine read past it, and nothing here can tell.

It runs on one clock, \code{DefaultClock}.

% SPDX-License-Identifier: Apache-2.0
% covers: FrameIn
\datasheet{FrameIn}{The Ethernet receiver's bytes packed into words for
a store engine, with the frame's length and the slot it lands in.}

\dsfacts{\code{lib/parts/src/ethdma.rs}}{\code{txhdl\_parts::ethdma::FrameIn}}%
{\code{//docs:examples}}

\subsection*{Function}

The mirror of \code{FrameOut}. It takes a frame's bytes on
\code{rx} and offers them to \code{LineStore} as 32-bit words on
\code{words}, lowest lane first. It does not touch the bus and does
not know where the frame lands in memory.

It needs the frame's length on \code{len} before the first byte,
because the store engine is told how many bytes to write when it
starts. \code{EthRx} offers exactly that, since it takes and checks a
whole frame before offering any of it, and so does \code{FrameLen}.

It also says four things to the register block: \code{count}, the
length of the frame taken; \code{store}, high while a frame is being
taken and stored, which starts the store engine; \code{which}, the
slot the frame lands in, which alternates; and \code{drops}, the
frames it dropped because no slot was free (issue 1313). The register
block tells it so on \code{no\_slot}.

\subsection*{Parameters}

None. The length is 16 bits wide, the words 32, and there are two
slots.

\dstables{FrameIn}

\subsection*{Behaviour}

\begin{itemize}
\item A frame starts when a byte is offered, \code{len} is not zero,
and \code{hold} is low. The length is latched there, and the slot
moves to the other one.
\item \textbf{A frame that starts while \code{no\_slot} is high is
dropped and counted (issue 1313).} Its bytes are taken, so the stream
behind it moves on, and none is stored: \code{store} stays low, no
word is offered, the slot does not move, and \code{drops} counts it.
The register block raises \code{no\_slot} while both slots hold a
frame the driver has not acknowledged, so a frame no longer lands on
one being read.
\item Bytes are packed from the top of the word down, so after four
the first is in the low lane.
\item The last word of a frame is usually partial, and it is sent with
its real bytes shifted into the low lanes, where the store engine's
strobe expects them. The engine strobes the rest away, so nothing past
the frame is written.
\item \textbf{\code{hold} keeps it from starting the next frame while
the store engine is busy with the one before.} This unit finishes
handing over words well before the engine's write is answered, and the
register block reads \code{count} and \code{which} when the engine
goes idle. A new frame started early would change both under the
frame still being stored, and the driver would be told the new frame's
length and slot for the old frame's bytes. Only a run with a real
store engine behind this unit can see that; it was found by reading
the board join, not by a test.
\end{itemize}

\subsection*{Verification}

\begin{itemize}
\item \code{ex\_ethdma} is \code{FrameOut}, the real \code{EthTx} and
\code{EthRx}, and this unit in a line. Two frames of 101 and 67 bytes
go through it, and every word that lands is checked, the partial last
words included. The run also checks that four words past the second
frame are untouched, that the receiver's length for the last frame
arrived, and that the second frame took the other slot.
\item A third frame, of 77 bytes, goes through with \code{no\_slot}
high: none of its words lands, \code{drops} counts one, and the slot
stays where the second frame left it.
\item The netlist is checked against that run under nvc and
Verilator: \code{//docs:ethdma\_sim\_frame\_in\_frame\_in\_tb\_test}
and \code{//docs:ethdma\_vsim\_frame\_in\_test}.
\item \code{ex\_framelen} runs it behind \code{FrameLen} instead of
\code{EthRx}, with the same two frames.
\item On the board,
\code{a\_frame\_received\_lands\_in\_memory\_and\_the\_core\_reads\_it\_back}
in \code{cpu/vreteno/tests/board.rs} injects a frame on the wire and
has the core read it back from its slot.
\end{itemize}

\subsection*{Used by}

\code{Board}, as \code{fin}, behind \code{FrameLen} and in front of the
store engine \code{store}. \code{hold} is the store engine's
\code{running}, \code{count}, \code{which} and \code{drops} go to
\code{EthSlots} as \code{rx\_len}, \code{rx\_which} and
\code{rx\_drops}, and \code{no\_slot} is its \code{rx\_full}. The examples
\code{ex\_ethdma} and \code{ex\_framelen}.

\subsection*{Limits}

Its ports are named apart from \code{FrameOut}'s on purpose. A
testbench reads the trace by port name, so two lowered units in one
run that both name a port \code{out} bind to the same recorded signal
(issue 462).

A frame whose length is zero is never started, and its bytes stay in
the channel.

It runs on one clock, \code{DefaultClock}.

% SPDX-License-Identifier: Apache-2.0
% covers: FrameLen
\datasheet{FrameLen}{An Ethernet frame held whole and handed on with its
length, so that a receiver's length can be found again after a clock
crossing.}

\dsfacts{\code{lib/parts/src/ethdma.rs}}{\code{txhdl\_parts::ethdma::FrameLen}}%
{\code{//docs:examples}}

\subsection*{Function}

\code{FrameIn} must know a frame's length before the store engine
starts. \code{EthRx} knows it and says so, but on the board the
receiver is in the top, on the clock the PHY recovers, and only the
bytes cross to the core's clock.

The length cannot follow the bytes on a wire of its own, because it
would arrive on the wrong clock and out of step with the frame it
describes. It cannot ride inside the stream either, because the remote
peripheral reads the same stream and would read the length as a byte
of its frame.

So the length is found again after the crossing. \code{FrameLen}
takes a frame a byte at a time, holds all of it, and then offers it
back with its count on \code{len}, which is exactly what \code{EthRx}
offers. It is the receiver's output side without the receiver:
\code{FrameIn} joined to it cannot tell the difference.

\subsection*{Parameters}

None. It holds \code{FRAME\_MAX} bytes, 2048, the same as the
receiver.

\dstables{FrameLen}

\subsection*{Behaviour}

\begin{itemize}
\item It takes bytes until one is marked last, then hands the frame
on, the held bytes in order and the last one marked. It takes no byte
while it is handing one on.
\item \code{len} is the count held while a frame is being handed on,
and zero otherwise, as the receiver's own length is, so that a stale
count cannot be read as a current one.
\end{itemize}

\subsection*{Verification}

\begin{itemize}
\item \code{ex\_framelen} puts it between \code{FrameOut} and
\code{FrameIn}, in the place \code{EthRx} has in \code{ex\_ethdma},
and sends the same two frames, of 101 and 67 bytes. It checks every
word stored, that the length \code{FrameIn} was told for the last frame
is 67, and that the slots alternated.
\item The netlist is checked against that run under nvc and
Verilator: \code{//docs:framelen\_sim\_framelen\_tb\_test} and
\code{//docs:framelen\_vsim\_test}.
\item On the board, it is in the path of
\code{a\_frame\_received\_lands\_in\_memory\_and\_the\_core\_reads\_it\_back}
in \code{cpu/vreteno/tests/board.rs}.
\end{itemize}

\subsection*{Used by}

\code{Board}, as \code{flen}, between \code{EthShare}'s second user
and \code{FrameIn}. The example \code{ex\_framelen}.

\subsection*{Limits}

It costs a frame of memory and a frame of latency, since the frame was
already held whole once, before the crossing. That is the price of the
length not crossing with it.

A frame longer than 2047 bytes is handed on as its first 2047, the
last of them marked. The bytes after those are written over the one
slot of memory that is not handed on, rather than wrapping round onto
the start of the frame.

It runs on one clock, \code{DefaultClock}.

% SPDX-License-Identifier: Apache-2.0
% covers: EthShare
\datasheet{EthShare}{One Ethernet wire and two users: frames split by
EtherType on the way in, merged a frame at a time on the way out.}

\dsfacts{\code{lib/parts/src/ethshare.rs}}{\code{txhdl\_parts::ethshare::EthShare}}%
{\code{//docs:examples}}

\subsection*{Function}

The board has one Ethernet port and two things that want it: the
remote peripheral of issue 297, whose transport is a frame per
transaction, and the port a Zephyr driver talks to through
\code{EthSlots}. Neither can have the wire to itself.

Nothing new has to be invented to share it, because Ethernet
multiplexes several protocols on one wire and identifies each by type:
bytes twelve and thirteen of every frame are its EtherType. The remote
peripheral sends \code{0x88b5}, which IEEE leaves for local use, and
Zephyr's stack sends only the standard types, so the two cannot be
confused.

Its three inputs are the wire and the two users' frames to send, and
its three outputs the wire and the two users' frames received. The
first user is the one whose frames are of type \code{TYPE}; the
second takes every other frame.

\subsection*{Parameters}

\code{TYPE}, the first user's EtherType. The tree uses
\code{0x88b5}, the remote peripheral's, which
\code{remote::eth::ETHERTYPE} states; the tables are generated for
that.

\dstables{EthShare}

\subsection*{Behaviour}

\begin{itemize}
\item \textbf{Receiving holds a frame's first fourteen bytes before
either user sees any of it.} A frame's user is not known until its
fourteenth byte, and the thirteen before it belong to that user too.
So the unit holds them, reads the type from the last two, then hands
the held bytes on and passes the rest of the frame through to the
same user. Every frame received waits for its fourteenth byte.
\item Holding does not wait for either user. Handing on does: a user
that is not ready stops the frame, and the wire behind it.
\item \textbf{A frame that ends before its type is known goes to the
second user}, not to nobody. The receiver checks every frame and pads
a short one, so this should not arrive; but waiting for a fourteenth
byte that never comes would hold the wire for ever, which is worse
than delivering a frame nobody wants.
\item A frame of exactly fourteen bytes learns its user and its end
in the same cycle, and goes to the user its type names.
\item \textbf{Sending takes turns only when both users offer.} A frame
cannot be interleaved with another, so the two are merged a frame at
a time. When both offer, the one that did not send last goes first,
so neither can starve the other whichever sends more. A user with the
wire to itself is not slowed.
\end{itemize}

\subsection*{Verification}

\begin{itemize}
\item \code{ex\_ethshare} sends five frames in on the wire: two of the
first user's type and one of another, which is the split; one that
ends after ten bytes, before its type is known; and one of exactly
fourteen bytes. It checks that each user got its own frames whole and
in order. Both users also offer two frames each in the same cycle,
and the run checks that the four go out whole, taking turns.
\item The short frame is adversarial on purpose. Its last byte is
\code{0xb5}, the low half of the first user's type, and the frame
before it left \code{0x88} as the high half, so a unit that read the
type without checking that the frame had reached it would see
\code{0x88b5} and send the frame to the wrong user. The first version
of the run used bytes that did not form that type, and it passed with
the check removed.
\item The netlist is checked against that run under nvc and
Verilator: \code{//docs:ethshare\_sim\_ethshare\_tb\_test} and
\code{//docs:ethshare\_vsim\_test}.
\item On the board, both board tests of the Ethernet port in
\code{cpu/vreteno/tests/board.rs} go through it, one in each
direction.
\end{itemize}

\subsection*{Used by}

\code{Board}, as \code{share}, an \code{EthShare<0x88b5>} between the
wire and its two users: the remote peripheral's link first and the
Ethernet port's \code{FrameLen} and \code{FrameOut} second. The
example \code{ex\_ethshare}.

\subsection*{Limits}

It checks the type and nothing else. It does not look at addresses, so
both users see every frame of their type, whoever it was for.

It holds sixteen bytes of header memory and uses fourteen.

It runs on one clock, \code{DefaultClock}.

% SPDX-License-Identifier: Apache-2.0
% covers: Hdmi
\datasheet{Hdmi}{The video peripheral for an HDMI encoder chip: a raster, a framebuffer, and AXI-Lite to paint it.}

\dsfacts{\code{lib/parts/src/hdmi.rs}}{\code{txhdl\_parts::hdmi::Hdmi}}%
{\code{//docs:hdmi}}

\subsection*{Function}

\code{Hdmi} drives the parallel inputs of an encoder chip such as the
SiI9134 on the Alinx AX7A200B. It is one unit on the pixel clock. It
counts the raster, reads its framebuffer under the beam, and drives
\code{rgb}, 24 bits with red on top, \code{hsync}, \code{vsync} and
\code{de} from registers. A host paints the framebuffer through the
unit's AXI-Lite side, one port, \code{bus}, a \code{LitePort} with
32-bit addresses and words: \code{bus\_aw}, \code{bus\_ar} and
\code{bus\_w} in, \code{bus\_b} and \code{bus\_r} out.

\subsection*{Parameters}

\begin{center}\footnotesize
\begin{tabular}{@{}lp{0.7\textwidth}@{}}
\toprule
Parameter & Meaning \\
\midrule
\code{HV} & Visible columns. \\
\code{HFP} & Columns of horizontal front porch. \\
\code{HSW} & Columns of horizontal sync. \\
\code{HBP} & Columns of horizontal back porch. \\
\code{VV} & Visible rows. \\
\code{VFP} & Rows of vertical front porch. \\
\code{VSW} & Rows of vertical sync. \\
\code{VBP} & Rows of vertical back porch. \\
\code{SHIFT} & A framebuffer pixel covers \code{1 << SHIFT} by
\code{1 << SHIFT} pixels of the screen. \\
\bottomrule
\end{tabular}
\end{center}

The tables below are for \code{Hdmi<640, 16, 96, 48, 480, 10, 2, 33,
2>}. Those are the constants of \code{hdmi::vga}, 640 by 480 at 60~Hz
as VESA states it, with \code{SHIFT} of 2. That is the mode of
\code{hdmi/src/netlist.rs}, and its picture is 160 by 120.

\subsection*{Register map}

The unit decodes the word from address bits 3 to 2, so the words are at
offsets \code{0x0}, \code{0x4} and \code{0x8} of its range.
\code{ex\_hdmi} puts them at \code{0x1000}, \code{0x1004} and
\code{0x1008}.

\dsregs{Hdmi}

\code{colour} is four bits each of red, green and blue, red at the top.
A write to \code{status} is ignored, and a read of \code{pixel} or of
\code{0xc}, which no register names, reads zero.

\dstables{Hdmi}

\subsection*{Behaviour}

\begin{itemize}
\item \code{hc} counts the column and \code{vc} the row, each over the
visible part and the blanking. \code{frames} goes up by one at the end
of each frame.
\item At each edge the unit reads the framebuffer at the pixel under
the beam into \code{px}. \code{//docs:hdmi} calls this the synchronous
read a block RAM has. The framebuffer address is the row,
\code{vc >> SHIFT}, in the top seven bits and the column,
\code{hc >> SHIFT}, in the low eight.
\item At the same edge the unit registers the syncs in \code{hs\_q} and
\code{vs\_q}, low inside the sync pulse, and the enable in
\code{de\_q}, high over the visible part. So they meet the pixel they
belong to.
\item \code{rgb} is \code{px} with each 4-bit colour widened to 8 bits
by repeating its bits, so 15 is full scale.
\item A write to the cursor word sets \code{cx} and \code{cy}.
\item A write to the pixel word puts the colour into the framebuffer at
the cursor, and moves the cursor to the next column. From the last
column it moves to the first column of the next row, and from the last
row to the first. A fixed burst to this word paints a run of pixels.
\item A write goes when \code{bus\_aw} and \code{bus\_w} both offer
and \code{bus\_b} is ready. A read goes when \code{bus\_ar} offers and
\code{bus\_r} is ready. The
unit answers every transaction \code{Okay}. A read of word 2 or word 3
answers zero, and a write to word 0 or word 3 does nothing else.
\end{itemize}

\subsection*{Verification}

\begin{itemize}
\item \code{the\_raster\_is\_where\_the\_mode\_puts\_it}, in
\code{//lib/parts:parts\_test}, runs \code{Hdmi<8, 2, 3, 3, 6, 1, 2, 2,
1>} for two frames. It checks that exactly one lag of the outputs
behind the counters fits \code{de}, \code{hsync} and \code{vsync}
together.
\item \code{ex\_hdmi} runs \code{Hdmi<16, 2, 3, 3, 12, 1, 2, 2, 1>}
behind \code{LiteBridge<1, ..>} and an \code{AxiHost}. The host sets the
cursor to the corner and paints all 48 pixels of the 8 by 6
framebuffer in one fixed burst. The run rebuilds the last whole frame
from the pixels \code{de} marks, and asserts every one of its 192
screen pixels against the colour painted.
\item The build simulates that unit's netlist against the run under
nvc and Verilator: \code{//docs:hdmi\_sim\_hdmi\_video\_hdmi\_video\_tb\_test}
and \code{//docs:hdmi\_vsim\_hdmi\_video\_test}. The co-simulated mode
is the example's, not the one in the tables above.
\end{itemize}

\subsection*{Used by}

The example \code{ex\_hdmi}. The demonstration for the AX7A200B:
\code{//hdmi:netlist} writes the Verilog of \code{hdmi\_video} in the
640 by 480 mode, with a test picture in the framebuffer's initial
values. \code{hdmi/board/hdmi\_demo.v} instantiates it on a 25.2~MHz
pixel clock, with its AXI-Lite inputs held idle.
\code{//hdmi:demo\_synth} and \code{//hdmi:demo\_pnr} put the design
through Vivado to a bitstream. \code{//docs:hdmi} states that it has
not yet run on the board.

\subsection*{Limits}

A framebuffer pixel is 12 bits, and the framebuffer holds
\code{FB\_WORDS}, $2^{15}$ words, so the picture is at most 256
columns by 128 rows. \code{hc} and \code{vc} are 12 bits wide. The mode
is fixed when the type is named. The whole unit runs on the pixel
clock, bus included. A host on another clock needs a crossing of the
five AXI-Lite channels, and nothing in the tree provides one. The
framebuffer cannot be read over the bus. On the board, \code{//docs:hdmi}
marks two things as unconfirmed: which of the chip's data inputs are
red, green and blue, and the I/O standard of the chip's pins.

% SPDX-License-Identifier: Apache-2.0
% covers: I2cInit
\datasheet{I2cInit}{The I2C master that resets and configures an HDMI encoder chip.}

\dsfacts{\code{lib/parts/src/hdmi.rs}}{\code{txhdl\_parts::hdmi::I2cInit}}%
{\code{//docs:hdmi}}

\subsection*{Function}

\code{I2cInit} holds the chip's reset low on \code{nreset}, releases
it, waits, and then writes a table of register values as I2C
transactions. The table is the SiI9134's configuration,
\code{SII9134\_WRITES}, with two entries. The unit drives the two
open-drain lines through \code{scl\_low} and \code{sda\_low}, and reads
the data line back on \code{sda\_in}. \code{done} rises when the table
is written, and \code{failed} is high when the chip did not acknowledge
a byte.

The unit has no bus. It writes this table, one transaction per entry.
\code{//docs:hdmi} and the source both state that no document the part
was written from gives the chip's registers. The values are
unconfirmed until checked against the SiI9134's documentation or the
board vendor's example.

\begin{center}\footnotesize
\begin{tabular}{@{}llllp{0.4\textwidth}@{}}
\toprule
Entry & Address & Register & Value & Meaning, unconfirmed \\
\midrule
0 & \code{0x72} & \code{0x08} & \code{0x35} & System control: out of
power down, with the input bus as the board wires it. \\
1 & \code{0x7a} & \code{0x2f} & \code{0x00} & The output as DVI: no
HDMI packets, only the picture. \\
\bottomrule
\end{tabular}
\end{center}

The address is the chip's I2C address as eight bits, write bit
included.

\subsection*{Parameters}

\begin{center}\footnotesize
\begin{tabular}{@{}lp{0.7\textwidth}@{}}
\toprule
Parameter & Meaning \\
\midrule
\code{DIV} & Cycles in a quarter of an I2C bit. A bit is four
quarters. \\
\code{HOLD} & Cycles the reset is held low, and cycles the unit waits
after releasing it. \\
\bottomrule
\end{tabular}
\end{center}

The tables below are for \code{I2cInit<63, 2520000>}, the master of
\code{hdmi/src/netlist.rs}. Its doc comment states that a quarter of 63
cycles at 25.2~MHz makes a bit at 100~kHz, and that 2\,520\,000 cycles
are a tenth of a second.

\dstables{I2cInit}

\subsection*{Behaviour}

The unit is two processes. The quarter clock counts \code{tick} from 0
to \code{DIV} less one and round again, every cycle, and drives
\code{done} from \code{written} and \code{failed} from \code{nak}. The
other is the whole job written as the sequence it is: the reset held,
the reset released, then a transaction per entry, a quarter of a bit
per wait. The lowering numbers its waits into a state register the unit
does not declare, \code{at\_wait}, keeps its loops' counts in
\code{for0} to \code{for5}, and holds the three lines it drives
between states in \code{nreset\_held}, \code{scl\_low\_held} and
\code{sda\_low\_held}; all are in the tables above. \code{byte} is
\code{byte\_rw} in the netlist, since SystemVerilog reserves the word.

\begin{itemize}
\item \code{nreset} is low for \code{HOLD} cycles, then high, and the
unit waits \code{HOLD} cycles more before the first transaction.
\item A transaction is the start, the address, the register and the
value, each a byte followed by its acknowledge, and the stop, every
step a whole number of quarters of \code{DIV} cycles.
\item At the start the data line falls while the clock is high, then
the clock falls. At the stop the clock rises while the data line is
low, then the data line rises.
\item In a data bit the data line is set while the clock is low, and
the clock is high through the two middle quarters, so the data changes
only while the clock is low. Bytes go out most significant bit first,
from \code{shift}.
\item In an acknowledge the unit releases the data line, and reads
\code{sda\_in} at the end of the clock's high half, as it pulls the
clock low again. A high line sets \code{nak}, which \code{failed}
follows. Nothing clears \code{nak}.
\item \code{scl\_low} or \code{sda\_low} high pulls its line low.
\item After the stop of the last entry, \code{written} is set,
\code{done} rises a cycle later, and the sequence waits for ever.
\end{itemize}

\subsection*{Verification}

\begin{itemize}
\item \code{the\_chip\_is\_configured\_over\_i2c}, in
\code{//lib/parts:parts\_test}, runs \code{I2cInit<3, 5>} against
\code{I2cDevice}, a model of a device that acknowledges every byte and
records what it was sent. It checks that \code{done} rises, that
\code{failed} stays low, that the reset is released at cycle 5, and
that the device received \code{SII9134\_WRITES}.
\item \code{ex\_hdmi} runs \code{I2cInit<4, 10>} against the same model
and asserts the same table, \code{done}, and \code{failed} low.
\item The build simulates that unit's netlist against the run under
nvc and Verilator: \code{//docs:hdmi\_sim\_hdmi\_i2c\_hdmi\_i2c\_tb\_test}
and \code{//docs:hdmi\_vsim\_hdmi\_i2c\_test}.
\end{itemize}

\subsection*{Used by}

The example \code{ex\_hdmi}. The demonstration for the AX7A200B:
\code{//hdmi:netlist} writes the Verilog of \code{hdmi\_i2c}, and
\code{hdmi/board/hdmi\_demo.v} instantiates it with open-drain pads on
the chip's I2C lines and \code{done} and \code{failed} on LEDs.
\code{//hdmi:demo\_synth} and \code{//hdmi:demo\_pnr} put the design
through Vivado to a bitstream. \code{//docs:hdmi} states that it has
not yet run on the board.

\subsection*{Limits}

There is one table, and its values are unconfirmed. The unit keeps
writing after a byte is not acknowledged, and it writes the table once.
It has no input for the clock line. \code{tick} is 16 bits wide, so
\code{DIV} is at most 65\,536; the counters of the two holds are as
wide as \code{HOLD} needs, 22 bits in the tables.

% SPDX-License-Identifier: Apache-2.0
% covers: I2c
\datasheet{I2c}{A general I2C master on AXI-Lite: one piece of a transaction per command, on two open drain lines.}

\dsfacts{\code{lib/parts/src/i2c.rs}}{\code{txhdl\_parts::i2c::I2c}}%
{\code{//docs:examples}}

\subsection*{Function}

\code{I2cInit} of the HDMI part replays a fixed list of writes and
cannot read. This master is the other kind: a program says what one
piece of a transaction is and the master does it, so a driver can
address a device, write to it, turn the bus around and read from it.

One command is at most a start, one byte either way, and a stop. A
driver that wants to read a register writes the address with
\code{start|write}, then asks for \code{start|read|nack}, whose start
is the repeated start. Between the commands of one transaction the
lines are parked with the clock low, since a data line that moves
while the clock is high is a start or a stop to every device on the
bus.

\subsection*{Register map}

Offsets from the block's base.

\dsregs{I2c}

A write to \code{cmd} starts the command it describes, and \code{cmd}
reads zero. A write of ones to \code{state} clears \code{fired},
\code{nack} and \code{lost}; \code{busy} is read only.

A quarter of a bit takes \code{div + 1} cycles, so the bus runs at the
system clock over \code{4 * (div + 1)}. From 100 MHz, 249 is
100 kbit/s and 62 is 400 kbit/s.

Bit 4 of \code{cmd} is set for a NACK, which is what a master sends
after the last byte it wants; a master that acknowledges the last byte
tells the device to send another.

\dstables{I2c}

\subsection*{The two lines}

Both are open drain, as I2C requires: a device pulls a line low and
nobody drives one high. \code{scl\_low} and \code{sda\_low} high pull
their line low, and a line nobody pulls is high through the board's
pull-up, so a board wrapper drives a pad from each output and reads
the pad back into \code{scl\_in} and \code{sda\_in}.

The master watches both inputs, and two things follow that a master
driving the lines itself could not do.

A device that holds the clock low is stretching it, and the master
waits for the line to rise rather than counting cycles. And a data
line that reads low while this master released it is another master
winning the bus, which sets arbitration lost and ends the command.

\subsection*{Behaviour}

The master is two processes. The bus process answers the registers
every cycle, counts the cycles of a quarter of a bit in \code{tick},
drives \code{scl\_low} and \code{sda\_low} from \code{scl\_pull} and
\code{sda\_pull}, and latches what a command found. The engine is the
transaction written as the sequence it is: a wait for a command, then
the start, the eight bits, the acknowledge and the stop, a quarter of
a bit per wait, with the lines set for each quarter before its wait.
The lowering numbers the waits into a state register the unit does not
declare, \code{at\_wait}, and the eight bits are a counted loop with a
counter of its own, \code{for0}; both are in the tables above, and
there is no \code{step} or \code{quarter} register to read.

\begin{itemize}
\item Writing \code{cmd} while no command runs records what was asked
and the byte to send, clears \code{nack\_hit} and \code{lost\_hit},
sets \code{busy}, and sets \code{held} on a start.
\item A quarter ends when \code{tick} reaches the divider and the
clock is not being stretched: a device holding the clock low where
the master released it holds the quarter with it.
\item On a write, the acknowledge reads the data line: low is the
device acknowledging, high sets \code{nack\_hit}.
\item On a read, the byte arrives in \code{shift} and is put in
\code{data}, which the word at \code{0x08} reads. The master then
drives the answer bit 4 asked for.
\item A master writing a one that reads back a zero has lost the bus
to another master. It sets \code{lost\_hit}, releases both lines for
the rest of the byte, whose bits the loop still counts through, and
skips the acknowledge and the stop.
\item At the end the lines are parked, the clock low while the
transaction stays open and both released when it is over or the bus
was lost. \code{done} is up for one cycle, when the bus process
latches \code{fired}, \code{nack} and \code{lost}; \code{busy} falls a
cycle after it. \code{fired} is the done bit, which raises the
interrupt while the enable is set.
\item Out of reset both lines are released: \code{scl\_pull} and
\code{sda\_pull} say whether the engine pulls a line, and reset to
not pulling (issue 606).
\end{itemize}

\subsection*{Verification}

\begin{itemize}
\item \code{//lib/parts:parts\_test} drives the master against the
device model in the same file: a write reaches the device, a read
takes a byte from it, an address nobody answers sets not acknowledged,
a device that stretches the clock is waited for, and another master
that wins the third bit of the address makes this one report
arbitration lost and let go of the data line.
\item \code{//lib/examples:ex\_i2c} addresses a device, writes to it,
turns the bus around and reads back, and the build simulates the
netlist against that run under nvc and Verilator.
\end{itemize}

\subsection*{Used by}

Nothing in the tree yet. The HDMI part still uses \code{I2cInit} for
its fixed setup; issue 148 notes that \code{I2cInit} could become a
program for this master instead, and that a board's EEPROM, which
often holds the Ethernet address, is what wants a master that can
read.

\subsection*{Limits}

Ten-bit addressing is not here: an address is a byte like any other,
so a driver writes what it likes, but nothing helps it. There is no
general call handling, no clock synchronisation with another master
beyond losing arbitration, and no timeout on a stretched clock, so a
device that holds the clock low for ever holds the master with it.

% SPDX-License-Identifier: Apache-2.0
% covers: Mdio
\datasheet{Mdio}{An MDIO master on AXI-Lite: one clause 22 frame per command, to read or write an Ethernet PHY's management registers.}

\dsfacts{\code{lib/parts/src/mdio.rs}}{\code{txhdl\_parts::mdio::Mdio}}%
{\code{//docs:examples}}

\subsection*{Function}

An Ethernet PHY keeps its configuration in 32 registers of 16 bits,
and MDIO is how a MAC reaches them: IEEE 802.3, clause 22. MDC is a
clock the master always drives. MDIO is one data line, which the
master drives except while the PHY answers a read.

A frame is 32 ones of preamble, then 32 bits, high bit first: the
start \code{01}, the operation (\code{10} a read, \code{01} a write),
the PHY's address and the register's, five bits each, two bits of
turnaround, and the sixteen bits of the word. On a read the master
lets go of the line at the turnaround.

\subsection*{Register map}

Offsets from the block's base.

\dsregs{Mdio}

A write to \code{cmd} starts the frame it describes, and \code{cmd}
reads zero. \code{data} is the word the last read took, once
\code{busy} has fallen. A write of one to \code{fired} clears it, and
a command clears it itself.

Half of an MDC cycle is \code{div + 1} cycles, so MDC is the system
clock over \code{2 * (div + 1)}. 802.3 allows MDC up to 2.5\,MHz.
\code{div} resets to zero, far too fast for any PHY, so a driver
writes \code{ctrl} before its first frame; \code{phyregs} uses 39,
1.25\,MHz from 100\,MHz.

\dstables{Mdio}

\subsection*{The line}

\code{mdio\_out} is the level and \code{mdio\_oe} high drives it, so
a board wrapper makes the pad three-state from the two and reads the
pad back into \code{mdio\_in}. 802.3 asks for a pull-up on the line,
so a line nobody drives reads one, and a read of an address with no
PHY reads \code{ffff}.

\subsection*{Behaviour}

The master is two processes, as the I2C master is. The bus process
answers the registers every cycle, counts the cycles of a half in
\code{tick}, and marks a frame done. The engine waits for a write to
\code{cmd} and loads two registers of 32 bits from it: \code{frame},
the bits after the preamble, and \code{drive}, a one for each bit the
master drives, all of a write and the first fourteen of a read. Then
come two counted loops of 32 bits, the preamble and the frame, each
bit a low half and a high half.

\begin{itemize}
\item A bit goes onto the line while MDC is low; the PHY takes it as
MDC rises.
\item The master takes each bit at the end of the low half, just
before MDC rises again, a whole MDC cycle after the PHY changed it.
802.3 gives the PHY up to 300\,ns for that, and at 1\,MHz there are
1000.
\item Every bit of the frame shifts into \code{rx}; after a read the
last sixteen are the word.
\item At the end the clock is parked low and the line let go.
\code{done} is up for one cycle, when the bus process sets
\code{fired}; \code{busy} falls a cycle after it.
\end{itemize}

\subsection*{Verification}

\begin{itemize}
\item \code{//lib/parts:parts\_test} drives the master against
\code{sim::MdioPhy}, a PHY at one address with 32 registers that
takes a bit as MDC rises and pages its vendor registers by register
31 when it has pages: the identifier read, all 32 registers read
back, an absent PHY reading all ones, a write read back, a paged
register read through the page select with page 0 left alone, the
master and the PHY never driving the line in the same cycle with
every half of MDC at least \code{div + 1} cycles, and the lowering.
\item \code{//lib/examples:ex\_mdio} scans for its PHY and reads it,
and the build simulates the netlist against that run under nvc and
Verilator (\code{//docs:mdio\_sim\_mdio\_tb\_test},
\code{//docs:mdio\_vsim\_test}).
\item \code{//cpu/vreteno:board\_test} runs \code{phyregs} and
\code{phydelay} on the board's design against the model, holding the
registers the JL2121 on the AX7A200B answered with.
\end{itemize}

\subsection*{Used by}

The Vreteno board, on the peripheral page's eighth slot at
\code{0x3700}, with its lines out to the flagship's JL2121 PHY.
\code{cpu/vreteno/rust/phyregs.rs} reads the PHY's 32 registers, and
\code{phydelay.rs} its RGMII delay bits on page 3336 (issues 864 and
869).

\subsection*{Limits}

Clause 22 only: there is no clause 45 frame, so a PHY's extended
register space is reached only through whatever paging the PHY
offers. One frame per command, and no interrupt; a program polls
\code{busy}. The master drives MDIO push-pull while it drives it,
which is what 802.3 describes, and leaves the pull-up to the board.

% SPDX-License-Identifier: Apache-2.0
% covers: Sd
\datasheet{Sd}{A native-mode SD card host on AXI-Lite: one command and
at most one block at a time, the block through a buffer of 128 words.}

\dsfacts{\code{lib/parts/src/sd.rs}}{\code{txhdl\_parts::sd::Sd}}%
{\code{//docs:examples}}

\subsection*{Function}

Operating an SD card in SPI mode limits throughput to several hundred
kilobytes per second over a single data line. In contrast, SD native
mode uses a bidirectional protocol over dedicated physical lines: a command
line, \code{CMD}, conveys 48-bit commands and returns 48-bit or 136-bit
responses. Data transfers use one or four data lines (\code{DAT}) to
transfer 512-byte blocks, protected by per-line CRC-16 checksums. Outgoing
commands include a CRC-7 checksum, as do all responses with the exception
of \code{ACMD41}, whose CRC field is defined as all ones.

\code{Sd} implements this native host controller behind an AXI-Lite port,
\code{bus}. Software configures the command argument and writes the command
control word. The command word selects the transfer parameters: response
type (none, short, or long), data transfer direction (read block or write
block), busy line monitoring, and CRC checking. The controller dispatches
the command, samples the response, transfers block data through its internal
buffer via the \code{data} register, and updates execution state in \code{status}.

Multi-block transfers issue an initial command (\code{CMD18} or \code{CMD25}),
stream data blocks across subsequent cycles, and terminate with \code{CMD12}.
Because streaming reads at 25~MHz overloads register polling loops, the host
integrates dedicated DMA paths to system memory (issue 912). Two DMA engines
from issue 151 interface with the controller: a \code{LineStore} instance
for reads and a \code{LineFetch} instance for writes. Setting a non-zero count
in \code{blocks} initiates automated multi-block DMA: the internal buffer
acts as a ring buffer between physical SD lines and the DMA engine. Read data
streams to the store engine as the card presents bytes; write data streams
from the fetch engine, dispatching each block once fully buffered. The transfer
completes only after the DMA engine retires. Commands perform either reads or
writes, never both simultaneously, ensuring mutual exclusion between the DMA
engines.

The card's side is eight wires: \code{sclk}, the command line in, out
and its output enable, the four data lines in, out and their output
enable, and \code{irq}. The engines' side is nine more:
\code{dma\_in} and \code{dma\_out}, the words from the fetch engine
and to the store engine; \code{dma\_at}, \code{dma\_bytes} and
\code{dma\_words}, where in memory and how much; \code{store\_go}
and \code{fetch\_go}, the starts; and \code{store\_busy} and
\code{fetch\_busy}, the engines' busy lines.

\subsection*{Parameters}

None.

\subsection*{Register map}

The word is address bits 5 to 2. A word the map does not name reads as
zero.

\dsregs{Sd}

\code{cmd}'s \code{resp} is 0 for none, 1 for short and 2 for long.
Writing \code{cmd} starts the command. Writing a one to
\code{status}'s \code{done} clears it and the four faults with it.
\code{dma} is where in memory the next command's blocks start, and
\code{blocks}' \code{count} how many it moves, zero for the buffer as
before; \code{left} counts the blocks still to move and \code{run}
says an engine is still moving words.

A half of a card clock takes \code{div + 1} cycles, so the card's
clock is the system's over \code{2 * (div + 1)}: at 100~MHz, 124 is
the 400~kHz a card starts at, 1 is 25~MHz, the fastest this host
runs, and 0 runs as 1 (issue 929).

\dstables{Sd}

\subsection*{Behaviour}

\begin{itemize}
\item \textbf{A write to \code{cmd} while a command runs is ignored.}
A program reads \code{status} first.
\item The host drives its lines on the falling edge of the clock it
makes and samples the card's on the rising edge. The clock runs only
while a command runs, and is low otherwise.
\item A response that has not started 64 card clocks after the command
ends is a timeout. A block, the CRC status token after a written
block, and the end of the card's busy are each given $2^{24} - 1$ card
clocks.
\item A CRC is checked the way a shift register checks one: the
received CRC goes into the same register after the data, and what is
left is zero when the two agree. For a long response the CRC covers
the 120 bits after the header. With four lines each line has its own
CRC-16, and any of the four being wrong sets bit 5.
\item A written block is accepted only when the card's CRC status
token is \code{010}; anything else sets bit 5. A command whose
response CRC is wrong goes no further: its block does not move.
\item Starting a read sets the write pointer to zero, and starting a
write sets the read pointer to zero, so a block fills or empties the
buffer from its start. A read of \code{data} moves the read pointer
and a write moves the write pointer, each modulo 128.
\item \code{irq} is done and the interrupt enable together, so it
stays high until the program clears done.
\item Every AXI-Lite access is answered \code{OKAY}.
\end{itemize}

\subsection*{Verification}

\begin{itemize}
\item \code{ex\_sd} does what a driver does, through an AXI4 host and
\code{LiteBridge<1, ..>}, against \code{SdCard}, the model of a card in \code{sd.rs},
stepped on the wires every cycle. It brings the card up with
\code{CMD0}, \code{CMD8}, \code{ACMD41}, \code{CMD2}, \code{CMD3},
\code{CMD7} and \code{ACMD6}, reads block 3 on one line, then writes
block 5 on four and reads it back, and checks both blocks.
\item The netlist is checked against that run under nvc and
Verilator: \code{//docs:sd\_sim\_sd\_tb\_test} and
\code{//docs:sd\_vsim\_test}.
\item \code{ex\_sddma} moves blocks through memory with real
\code{LineStore<32, 1, 16, 16>} and \code{LineFetch<32, 1, 16, 16>}
engines behind an \code{Arbiter2}, the card sending its blocks back to
back, and its netlist is checked against that run:
\code{//docs:sddma\_sim\_sd\_dma\_tb\_test} and
\code{//docs:sddma\_vsim\_test}.
\item \code{//lib/parts:parts\_test} holds eleven tests of it:
\begin{itemize}
\item the card comes up and gives its CID;
\item a block is read on one line;
\item a block is written and read back on four lines, the card's CRC
status accepting it;
\item a card slow to drive is read and written at 25~MHz;
\item a divider of zero runs as one;
\item several blocks each way, as a command, data-only transfers and
\code{CMD12};
\item a spoiled response CRC sets bit 3 and the next command is
clean; a command the card does not answer sets bit 2; a spoiled block
sets bit 5;
\item a card wants its clocks before its first command;
\item the CRC-7 of \code{CMD0} with a zero argument is \code{0x4a} and
the CRC-16 of 512 bytes of \code{0xff} is \code{0x7fa1}, the
specification's own examples;
\item blocks read back to back go to memory through the store engine;
\item blocks written from memory through the fetch engine reach the
card.
\end{itemize}
\end{itemize}

\subsection*{Used by}

The example \code{ex\_sd}, and the board: the Vreteno board holds it
in the peripheral page's ninth slot, at \code{0x3800}, with the card
slot's pins in bank 16 (E13 the clock, E14 the command line, D15,
D14, F14 and F13 the data lines) on every top that holds the board.
Its engines are the board's \code{sdstore} and \code{sdfetch}, sharing
one host, \code{sdshost}, through \code{sdarb}, an
\code{Arbiter2<32, 32, 4, 1, 2, 0>}, on the main arbiter's sixth port.
\code{cpu/vreteno/rust/sdprobe.rs} brings a card up and prints its CID
and block 0, then reads blocks 0 to 3 with one \code{CMD18} straight
into the DDR3 and compares them with the same blocks read one at a time
through the buffer; \code{//cpu/vreteno:board\_test} runs it against
\code{SdCard} and expects \code{dma same}, and
\code{docs/board-checks.md} says how to run it on the board.

\subsection*{Limits}

Only reading has met a real card: on 2026-10-03 the board read a
32~GB card's CID and block 0, its partition table, twice (issue 153),
and on 2026-10-04 read them again at 25~MHz, line for line the same
(issue 929).
No write has been sent to a card on the board, and \code{sdprobe}
sends none.

One command at a time. Without a count in \code{blocks} a command
moves at most one block, through a register; with one, its blocks go
through memory by the engines, which only a read or a write uses at a
time.

\code{SdCard}, the model, is not a component: nothing lowers it, and
it is only as right as the tests that check the host against it.

It runs on one clock, \code{DefaultClock}.

% Vreteno and its devices.
% SPDX-License-Identifier: Apache-2.0
% covers: Vreteno
\datasheet{Vreteno}{A three-stage RV32IMAC core with machine, supervisor and user modes and an AXI host port.}

\dsfacts{\code{cpu/vreteno/src/core.rs}}{\code{vreteno32::core::Vreteno}}%
{\code{//docs:vreteno}, \code{//docs:noc}, \code{//docs:tapeout}}

\subsection*{Function}

\code{Vreteno} is a RISC-V core for the RV32I base set, the M
extension and the C extension's compressed instructions. It is one
process with three stages. The fetch stage reads the instruction at
the program counter into the instruction register, expanding a
compressed one into the thirty-two bit instruction it stands for. The
execute stage decodes that word, computes, and issues
loads and stores on the bus. The writeback stage extends a loaded
word, writes the register file and retires the instruction.

The core holds three memories: the register file, 32 words; the
instruction memory, \code{IMEM\_WORDS} or 1024 words; and a data RAM
of 64~KiB at \code{DRAM\_BASE}, \code{0x1\_0000}, the window a program
keeps its stack in (issue 1275). \code{Vreteno::with} loads a program
into the instruction memory. The data memory and every device are on
the bus, outside the core. Above
the instruction memory the fetch reads from the bus, into a buffer of
two words, and the words of the data memory and of DDR3 through an
instruction cache (issue 1021): sixteen kilobytes, 1024 lines of
four words, one way, looked up after translation and so indexed and
tagged by the physical address (issue 1290), in \code{ic\_data} and
\code{ic\_tag}. A hit fills both words of the buffer when the line
holds the next one, and the buffer shifts its second word into the
first at the clock edge, as the program counter moves into it
(issues 1187 and 1279).

The data RAM is on the core's own port and not on the bus: four lanes
of a byte, \code{dl0} to \code{dl3}, each written at one address and
read at another, as a block RAM is. A load in its window reads the
lanes into \code{dl\_word} at the edge and takes the word from there in
writeback, the cycle after it ran; a store writes the lanes under its
strobes in the cycle it runs. Whether an access is in the window comes
from the upper half of \code{rs1} and the carry out of the lower
halves' sum, without the full add.

The bus side is an AXI host written as hardware. The core issues a
burst on \code{issue}, sends a store's beat on \code{wbeat}, and hands
identifiers back on \code{release}. It receives \code{grant},
\code{done} and \code{rdata} from its tracker, an \code{AxiHost}. The
other inputs are \code{rst}, the external interrupt \code{irq}, the
timer's line \code{tirq} and the software interrupt \code{sirq}, which
the controller raises from \code{msip}; and the debug module's, the
halt and resume requests \code{haltreq} and \code{resumereq} and its
register access, \code{dbg\_regno}, \code{dbg\_wdata} and
\code{dbg\_we}, honoured in debug mode (issue 154). The outputs are
\code{halt}, \code{instr}, the word retiring, \code{wb}, a
\code{Writeback} of \code{done}, \code{rd} and \code{val}, \code{dbg},
debug mode, and \code{dbg\_rdata}, the register the module asked for.

Two additional inputs supply the machine timer count, \code{time}
(accessed by \code{rdtime}), and \code{seirq}, the supervisor external
interrupt line registered alongside MEIP as \code{mip.SEIP} (issue 1094).

Eleven ports join the core to its memory management unit, an
\code{Mmu8}, inside a \code{Hart} (issue 1014). The core tells the
unit \code{mmu\_satp}; \code{mmu\_prv}, the mode a data access is
checked as, which is \code{MPP} under \code{MPRV} (issue 1105);
\code{mmu\_sum} and \code{mmu\_mxr} from \code{mstatus};
\code{mmu\_flush}, raised for a write of
\code{satp} or an \code{sfence.vma}; and its two requests,
\code{ireq} for a fetch and \code{dreq} for a load or a store. The
unit answers on \code{ires} and \code{dres}, each a \code{Res} of a
physical address, a page fault or an access fault, a cycle after the
request. Its walker's reads come in on \code{ptw} as addresses, go out
on the core's own \code{issue}, and their words go back to the unit on
\code{pte}, so the hart needs no bus port of its own for the page
tables. The core believes an answer only at the second cycle of a
request, which is right both in simulation, where the unit sees a
request a step late, and in hardware.

\subsection*{Parameters}

\begin{center}\footnotesize
\begin{tabular}{@{}lp{0.7\textwidth}@{}}
\toprule
Parameter & Meaning \\
\midrule
\code{IW} & The width of the AXI identifier on \code{rdata}, \code{done},
\code{grant} and \code{release}, in bits. The tables use 2, as the
board and the runs do. \\
\code{DW} & The data RAM's words a lane, a power of two, 16384 by
default, which is 64~KiB. The netlist the documents simulate and the
layout maps uses 4, where the window's addresses wrap, since the cell
library has no RAM. The tables use the default. \\
\bottomrule
\end{tabular}
\end{center}

\subsection*{Register map}

The control registers, by CSR number. Software reaches them through
the six CSR instructions. They are not on the bus.

\begin{center}\footnotesize
\begin{tabular}{@{}llp{0.6\textwidth}@{}}
\toprule
Number & Register & What the core keeps \\
\midrule
\code{0x300} & \code{mstatus} & Bits 3 and 7 only; a write is masked
with \code{0x88}. \\
\code{0x304} & \code{mie} & Bits 11 and 7 only, the external and the
timer enable. \\
\code{0x305} & \code{mtvec} & The handler's address; a write clears
bits 1 and 0. \\
\code{0x340} & \code{mscratch} & Entire 32-bit register. \\
\code{0x341} & \code{mepc} & The trap's address; a write clears
bit 0. \\
\code{0x342} & \code{mcause} & Entire 32-bit register. \\
\code{0x343} & \code{mtval} & Entire 32-bit register. \\
\code{0x344} & \code{mip} & Bit 11 is stored: \code{irq} sets it and
a write keeps only that bit. Bit 7 reads as \code{tirq}. \\
\code{0x7c0} & \code{mhalt} & Nothing: it reads as zero. A write of an
odd value stops the core when that instruction retires. \\
\code{0x301} & \code{misa} & Nothing: it reads
\code{0x4000\_1104}, which is \code{MXL} of 1 and the letters
\code{C}, \code{I} and \code{M}. Writable and ignored, as the
specification allows. \\
\code{0xf11} & \code{mvendorid} & Nothing: zero, which means
unimplemented. Read only. \\
\code{0xf12} & \code{marchid} & Nothing: zero. Read only. \\
\code{0xf13} & \code{mimpid} & Nothing: zero. Read only. \\
\code{0xf14} & \code{mhartid} & Nothing: zero, this machine's one
hart. Read only. \\
\bottomrule
\end{tabular}
\end{center}

A CSR instruction that names any other number is an illegal word and
traps.

\dstables{Vreteno}

\subsection*{Behaviour}

\begin{itemize}
\item With \code{rst} high, the fetch counter goes to zero, the
instruction register is emptied, and the load wait and the sequencer
are cleared.
\item The fetch reads the instruction memory at bits 11 to 2 of the
program counter and at the word after. The halfword at the counter is
the upper half of its word when bit 1 is set. If its low two bits are
not both set it is a compressed instruction, expanded in the fetch;
otherwise the halfword after it is the instruction's upper half. The
counter moves on by two or four.
\item A fetch above the instruction memory goes to the bus when its
physical page is a device's, and to the cache when it is the data
memory's or DDR3's. The cache reads the line's tag and the word the
cycle after, at the address the fetch registered in \code{ic\_pa}. A
hit puts the word in the fetch buffer. A miss reads the line as one
burst of four beats, \code{len} 3, once nothing else of the core's is
on the bus, writes each beat into the line, makes the line valid on the
last, and looks the word up again. A refused beat ends the fetch with
the word's access fault, and leaves the line invalid.
\item \code{fence.i} and \code{rst} clear the cache's tags, one line a
cycle, and a lookup waits until they are clear. A store does not reach
the cache, so code written by stores runs after a \code{fence.i}.
\item \code{ir\_c} says the instruction in execute was compressed.
\code{jal} and \code{jalr} link the address two bytes on for a
compressed one and four for the rest.
\item An instruction in execute that reads the register the writeback
stage is about to write gets that value directly. When that value is a
load's, the instruction waits one cycle and reads the register file
instead.
\item A taken branch, \code{jal}, \code{jalr}, a trap and \code{mret}
redirect the fetch. The word fetched in that cycle is squashed, which
costs one bubble.
\item A load or store in the data RAM's window does not go to the bus.
A load's word comes from \code{dl\_word} the cycle after it ran, and
never sets \code{dev\_wait}; a store writes the lanes in the cycle it
runs and is not counted in \code{stores\_out}; \code{lr.w} and
\code{sc.w} act there as on the bus, and an AMO takes its word into
\code{wb\_dev} in its first cycle in writeback and stores to the lanes.
A load there clears \code{wb\_err}, since nothing refuses it.
\item Any other load or store whose address has any bit set above bit
11 goes on the bus as a burst of one beat, size 2, \code{INCR}. A store's beat
has the lanes it covers as its strobe.
\item A load or store waits for room on both \code{issue} and
\code{wbeat}, whatever its address.
\item A store is posted and the core runs on. A load sets
\code{dev\_wait} and sits in writeback, holding the pipeline, until
its answer lands in \code{wb\_dev}.
\item When a write response and a read beat arrive together, the
write's identifier goes back first and the read's waits a cycle. A
read's identifier goes back on its last beat, so once for a cache
fill's four. The core drops what it receives on \code{grant}.
\item \code{ecall} traps with cause 11 and \code{ebreak} with cause 3,
which takes its own address as \code{mtval}. A word the core does not
know traps with cause 2, and \code{mtval} takes the word. A halfword
that is no RV32C instruction expands to itself, which no thirty-two bit
instruction is, so it traps with cause 2 and \code{mtval} takes the
halfword. An unaligned load traps with cause 4 and an unaligned store
with cause 6, and \code{mtval} takes the address that was asked for.
Every other trap writes zero to \code{mtval}.
\item On a trap, \code{mepc} takes the instruction's address, bit 7 of
\code{mstatus} takes bit 3, bit 3 is cleared, and the fetch goes to
\code{mtvec}. \code{mret} sets bit 3 from bit 7, sets bit 7, and goes
to \code{mepc}.
\item The external interrupt is ready when bit 11 is set in both
\code{mie} and \code{mip}. The software interrupt is ready when bit 3
of \code{mie} is set and \code{sirq} is high, and the timer's when bit
7 of \code{mie} is set and \code{tirq} is high. Any of the three is
taken when bit 3 of \code{mstatus} is set, in place of the instruction
in execute. The external interrupt, cause \code{0x8000\_000b}, wins
over the software interrupt, cause \code{0x8000\_0003}, which wins over
the timer's, cause \code{0x8000\_0007}.
\item A multiply or divide starts the sequencer and stalls in execute.
A multiply is one step and a division thirty-two. \code{//docs:vreteno}
states that a multiply takes three cycles and a division thirty-three.
An interrupt cancels the sequencer, and the instruction starts it
again after the handler returns.
\item The sequencer does not start while a device load in writeback is
still waiting for its answer.
\item A write of an odd value to \code{mhalt} stops the fetch and
parks the program counter one word past that instruction, which
retires first. \code{halt} rises as it retires, and stays high.
\item \code{stall}, \code{int\_take} and \code{redirect} are wires kept
as fields, so the trace and the netlist show them.
\end{itemize}

\subsection*{Verification}

\begin{itemize}
\item \code{//cpu/vreteno:lockstep\_test} runs the core with the router, the
data memory, the timer and the serial port against the reference model
in \code{model.rs}. After every cycle the program counter, registers
1 to 31, the eight control registers that hold something and the halt
must agree. At the
halt the data memory, the timer's compare and the serial port's bytes
must agree. Its tests are \code{demo\_program},
\code{a\_multiply\_right\_after\_a\_load},
\code{compressed\_instructions}, \code{an\_unaligned\_access\_traps},
\code{the\_data\_ram\_on\_the\_cores\_own\_port}, a store and a load of
one word at once, a load through a pointer just loaded, bytes and
halves, both arms of the window's compare, \code{lr}/\code{sc} and an
AMO there, and \code{random\_programs}, which
runs 64 seeds of 200 instructions each with both lengths mixed and
counts that the mix is there.
\item \code{//cpu/vreteno:isa\_test} compares the core's expander with
the decoder's on every compressed halfword, and
\code{//cpu/vreteno:compressed\_test} checks the decoder against the
GNU disassembler on all 49\,152 of them.
\item \code{//cpu/vreteno:hello\_test} (\code{rust\_runs\_on\_vreteno})
and \code{//cpu/vreteno:hello\_cc\_test}
(\code{cpp\_runs\_on\_vreteno}) run compiled programs on the machine
of \code{run.rs} and check what the serial line says and that the core
halted.
\item \code{//cpu/vreteno:board\_test} and \code{//soc:demo\_test} run
the core inside \code{Board} and on the lattice.
\item \code{//cpu/vreteno:icache\_test} runs issue 1009's loop of three
instructions from the instruction memory and from the data memory,
through the cache, at 1.334 and 3.349 cycles an instruction counted
from the reset. The second includes the 1024 cycles in which the cache
clears its tags, which the loop waits for, since the cache is sixteen
kilobytes (issue 1290); it was 3.093 at four kilobytes, 4.093 before
one cache hit filled both words of the buffer and the shift moved to
the edge (issues 1187 and 1279), and 14.655 from the data memory
before the cache. It also runs the same loop across two lines, and a
routine rewritten by a store and called again after
\code{fence.i}, which runs the word stored.
\item The build simulates the netlist of \code{Hart<2>}, the core and its
unit, under the entity name \code{vreteno}, with the demonstration
program in the core's instruction memory, against the trace of
\code{//cpu/vreteno:vreteno} under nvc and Verilator:
\code{//docs:vreteno\_sim\_vreteno}\allowbreak\code{\_vreteno\_tb\_test} and
\code{//docs:vreteno\_vsim\_vreteno\_test}.
\item \code{//cpu/vreteno:vreteno\_synth} and
\code{//cpu/vreteno:vreteno\_pnr} synthesise and place the core on an
\code{xc7a200tfbg484-2}; both are manual. \code{//docs:vreteno}
reports 2151 LUTs as logic, 48 LUTs as distributed RAM, 489
flip-flops, two block RAM tiles and four DSP blocks after synthesis. It
reports that the routed core meets a 10~ns clock with 0.296~ns of
setup slack.
\item \code{//cpu/vreteno/asic:vreteno\_cells} maps the core onto the
Nangate45 cells and builds with everything.
\code{//cpu/vreteno/asic:}\allowbreak\code{vreteno\_pnr} places and routes it, and is
manual.
\end{itemize}

\subsection*{Used by}

\code{Board}, as \code{cpu}. The demonstration
\code{//cpu/vreteno:vreteno}, the lockstep test, and \code{run.rs},
which the \code{hello} and \code{hello\_cc} runs call. The system in
\code{soc}, where the core is the host at corner 0,0 of the lattice.

\subsection*{Limits}

The boot memory is 1024 words, 4096 bytes, and nothing loads it at
run time: a netlist gets its contents through \code{init}. A program
counter at or above \code{IMEM\_BYTES}, the boot memory's end, is
fetched from the bus instead (issue 134), so a program larger than the
boot memory runs from wherever it was loaded, the DDR3 memory on the
board.

Every load and store goes to the bus, the boot memory included, which
answers reads and refuses writes (issue 268), except those in the data
RAM's window. Nothing but the core reaches the data RAM: not the debug
module's system bus access, Razboj or a DMA engine, so a debugger
reads it through the core. A fetch from the window is refused.

A trap handler must start at a whole word: \code{mtvec} keeps no
bits 1 and 0. The core has machine mode only and the fourteen control
registers above, in three kinds: eight it keeps, \code{mhalt}, which
is its own, and five that say what the machine is and hold nothing. A
write to one of the four whose number begins \code{0xf} is an illegal
instruction, which is what read only means in the specification;
\code{misa} is writable and ignores what is written. The two counters answer as well:
\code{mcycle} at \code{0xb00} with its high half at \code{0xb80},
and \code{minstret} at \code{0xb02} with \code{0xb82}, each 64 bits
and each writable. \code{mcycle} counts while the core runs and stops
when the halt stops it; \code{minstret} counts what retires. They are
the one part of the core the lockstep test leaves alone, for the
reason \code{//docs:vreteno} gives, and
\code{//cpu/vreteno:counters\_test} checks them instead. It has one load in
flight at a time. On the board, \code{irq} is the interrupt
controller's line.

% SPDX-License-Identifier: Apache-2.0
% covers: Hart
\datasheet{Hart}{The Vreteno core and its Sv32 unit, joined, with the core's ports and nothing more.}

\dsfacts{\code{cpu/vreteno/src/hart.rs}}{\code{vreteno32::hart::Hart}}%
{\code{//docs:vreteno}}

\subsection*{Function}

\code{Hart<IW>} is a unit of two units: a \code{Vreteno<IW>} core and
the \code{Mmu8} that translates its addresses for Sv32 (issue 1014).
Its ports are the core's, all of them and no others, so everything
outside sees one core: the bus side, the interrupts, the debug
module's lines and the retire outputs.

Inside, the core tells the unit \code{satp}, the mode a data access is
checked as, \code{sum}, \code{mxr} and the flush, and asks on its two
request wires; the unit answers on two answer wires. The walker's page
table reads go to the core on a channel, which the core puts on its
own bus port as loads, and the words read come back to
the unit on a second channel. The hart therefore needs no bus port of
its own for the page tables.

\code{Hart::with} loads a program into the core's instruction memory.

\subsection*{Parameters}

\begin{center}\footnotesize
\begin{tabular}{@{}lp{0.7\textwidth}@{}}
\toprule
Parameter & Meaning \\
\midrule
\code{IW} & The width of the AXI identifier, passed to the core. The
tables use 2, as the board and the runs do. \\
\bottomrule
\end{tabular}
\end{center}

\dstables{Hart}

\subsection*{Behaviour}

\begin{itemize}
\item The unit goes first in the join. Its answers are registers, and
the core reads them in the same step.
\item The core's requests are registers too, and the unit reads them a
step late in simulation and on time in hardware. So the core waits two
cycles on every request before it believes an answer, which is right
in both, since the unit holds its answer while the request is held.
\item The hart adds no register and no logic of its own: every
register in it is the core's or the unit's.
\end{itemize}

\subsection*{Verification}

\begin{itemize}
\item \code{run.rs}'s machine, which \code{//cpu/vreteno:lockstep\_test},
\code{//cpu/vreteno:hello\_test}, \code{//cpu/vreteno:icache\_test} and
the core's other tests use, runs a \code{Hart}. Lockstep holds it
against the reference model after every cycle, with paging directed
and over random tables.
\item The build simulates \code{Hart<2>}'s netlist, under the entity
name \code{vreteno} and with the demonstration program in the core's
instruction memory, against the trace of \code{//cpu/vreteno:vreteno}
under nvc and Verilator: \code{//docs:vreteno\_sim\_vreteno}\allowbreak\code{\_vreteno\_tb\_test}
and \code{//docs:vreteno\_vsim\_vreteno\_test}.
\item \code{//cpu/vreteno:openocd\_test} runs a stock OpenOCD against a
hart on the simulated board, and \code{//cpu/vreteno:pair\_test} runs
two harts in step.
\item \code{//cpu/vreteno:vreteno\_synth} and
\code{//cpu/vreteno:vreteno\_pnr}, which are manual, synthesise and
place the hart's netlist; on the board it is in every flagship.
\end{itemize}

\subsection*{Used by}

\code{Hart<2>} as \code{cpu} in \code{Board}, and so in the flagship;
two of them in \code{Pair}; and the machine of \code{run.rs}.

\subsection*{Limits}

One hart: there is no second core sharing the unit or the bus port.
The core's fetch and data requests each take the unit's cycle of
latency, which the core's own rules account for (\code{Vreteno}'s
sheet).

% SPDX-License-Identifier: Apache-2.0
% covers: Mmu, Mmu8
\datasheet{Mmu}{Sv32 address translation: a fetch TLB, a data TLB and a walker that never writes, answering in a register a cycle after each request.}

\dsfacts{\code{lib/parts/src/mmu.rs}}{\code{txhdl\_parts::mmu::Mmu}}%
{\code{//docs:vreteno}}

\subsection*{Function}

\code{Mmu<E>} translates a core's virtual addresses for Sv32, the
RISC-V page-based virtual memory of thirty-two bits (issue 1014). It
has two translation lookaside buffers of \code{E} entries each, one
for fetches and one for data, and one walker that fills them from the
page tables.

The core asks on \code{ireq} for a fetch and on \code{dreq} for a load
or a store, with a virtual address it already holds in a register. The
answer is a register too: on \code{ires} or \code{dres}, a cycle after
the request, a \code{Res} that says the access is allowed at a
physical address, or is a page fault, or an access fault. A miss
answers none of the three. The walker then reads the page table, a
word at a time, by sending the entry's address on \code{ptw} and
taking the entry, or the failure of its read, on \code{pte}. It fills
an entry, and the answer comes a cycle after the fill. The core holds
its request steady until one of the three comes. An answer belongs to
the request of the cycle before it.

The other inputs are \code{satp}, the mode the access is checked as,
\code{prv}, the two \code{mstatus} bits \code{sum} and \code{mxr}, and
\code{flush}, which empties both buffers.

\subsection*{Parameters}

\begin{center}\footnotesize
\begin{tabular}{@{}lp{0.7\textwidth}@{}}
\toprule
Parameter & Meaning \\
\midrule
\code{E} & The entries in each buffer, from 1 to 16. The tree uses 8, as
\code{Mmu8}, which is what the tables lower. \\
\bottomrule
\end{tabular}
\end{center}

\dstables{Mmu}

\subsection*{Behaviour}

\begin{itemize}
\item Translation is on when \code{satp.MODE} is set and \code{prv} is
not machine mode. Otherwise the address passes through unchecked, with
the same latency of one cycle, so the core has one rule.
\item Pages are four kilobytes. A leaf at the first level is a four
megabyte megapage, which compares the upper ten bits of the page
number only. Each buffer is fully associative and replaced in turn.
\item The checks are made on a hit: read, write and execute against the
access, the user bit against the mode and \code{sum}, and \code{mxr}
for a load from an execute-only page.
\item The walker never writes. An entry whose accessed bit is clear, or
whose dirty bit is clear on a store, is a page fault, and the kernel
sets the bit (Svade).
\item A flush empties both buffers. Address space identifiers are not
kept, which the specification allows, so \code{satp.ASID} is ignored.
A request in the cycle of a flush is not answered; the core holds it,
and it is answered from the emptied buffers, by a walk. A flush during
a walk drops what the walk finds.
\item When both ports miss at once, the data port is walked first.
\item A walk that ends in a fault answers for that page while its port
goes on asking for it, which the core does until it takes the trap,
and is not kept after.
\item Physical addresses are thirty-two bits. An entry or a
\code{satp} whose page number reaches past them is an access fault, as
is a table read that nothing answers.
\end{itemize}

\subsection*{Verification}

\begin{itemize}
\item \code{translate} in \code{mmu.rs} is the reference: the
specification's walk, with no buffer. In \code{mmu::tests}, in
\code{//lib/parts:parts\_test}, the unit is held to it against tables
in a memory that answers the walker after a few cycles:
\code{bare\_and\_machine\_mode}\allowbreak\code{\_pass\_the\_address}\allowbreak\code{\_through\_a\_cycle\_later};
\code{a\_page\_is\_walked\_once}\allowbreak\code{\_and\_then\_hits\_a\_cycle\_later};
\code{a\_megapage\_maps}\allowbreak\code{\_four\_megabytes\_from\_one\_entry};
\code{every\_fault}\allowbreak\code{\_is\_the\_specifications};
\code{an\_address\_past}\allowbreak\code{\_thirty\_two\_bits}\allowbreak\code{\_or\_nowhere}\allowbreak\code{\_is\_an\_access\_fault};
\code{a\_fault\_is\_not\_kept}\allowbreak\code{\_after\_its\_request};
\code{a\_flush\_empties\_both\_tlbs};
\code{a\_flush\_during\_a\_walk}\allowbreak\code{\_drops\_what\_it\_finds};
\code{a\_request\_in\_the\_cycle}\allowbreak\code{\_of\_a\_flush\_is\_answered\_after\_it};
\code{entries\_are\_replaced\_in\_turn};
\code{the\_data\_port\_walks\_first}; and
\code{random\_requests}\allowbreak\code{\_agree\_with\_the\_reference}, six hundred random
requests over random tables, on both ports and in every mode.
\code{the\_eight\_entry\_unit\_lowers} writes \code{Mmu8}'s Verilog.
\item \code{ex\_mmu} runs \code{Mmu8} the way a kernel meets it, and
the build simulates its netlist against that run under nvc and
Verilator: \code{//docs:mmu\_sim\_mmu8\_tb\_test} and
\code{//docs:mmu\_vsim\_test}.
\item In the core, \code{//cpu/vreteno:lockstep\_test} runs Sv32
directed and over random page tables against the reference model
(\code{sv32\_translates\_and\_faults}\allowbreak\code{\_where\_the\_tables\_say},
\code{random\_programs\_under\_paging}), and Linux runs under it on the
board (issue 279).
\item \code{//lib/parts:mmu\_pnr}, which is manual, places and routes
\code{Mmu8} alone, out of context, against the core's ten nanosecond
clock: it meets it with 2.440\,ns to spare, in 493 LUTs and 919
flip-flops (\code{//docs:vreteno}).
\end{itemize}

\subsection*{Used by}

\code{Mmu8} in \code{Hart}, the core and its unit, in
\code{cpu/vreteno/src/hart.rs}, and so in every design that holds a
hart: the board, the flagship and \code{Pair}. \code{Mmu8} in the
example \code{ex\_mmu}.

\subsection*{Limits}

No address space identifiers. No hardware update of the accessed and
dirty bits. The buffers are held in flip-flops, not block RAM, so
\code{E} stays small. Physical addresses are thirty-two bits, with no
Sv32 thirty-four bit physical space.

% SPDX-License-Identifier: Apache-2.0
% covers: Dmem
\datasheet{Dmem}{Vreteno's data memory: a block RAM of words behind an AXI peripheral end.}

\dsfacts{\code{cpu/vreteno/src/dmem.rs}}{\code{vreteno32::dmem::Dmem}}%
{\code{//docs:vreteno}}

\subsection*{Function}

\code{Dmem} is the data memory of the Vreteno machine, a device on its
bus. It holds \code{W} words, \code{DMEM\_WORDS}, 1024, by default,
in four memories of a byte, \code{lane0} to \code{lane3}, one per
lane of the word. A store of a byte or a half writes only its lanes and
reads nothing. The board had a second, \code{Dmem<5, 14, 16384>}, the
stack memory at \code{0x1\_0000} (issue 1278), until the stack window
became a RAM in the core, on its own port (issue 1275).

It is the peripheral client of an \code{AxiPer}, and its link is one
port, \code{bus}, a \code{PerPort} of the four channels: it takes
requests on \code{bus\_req} and write beats on \code{bus\_w}, and
answers on \code{bus\_ans} for a write and \code{bus\_r} for a read.
It answers a read burst of incrementing words a beat a cycle, which
is what the core's instruction cache fills its lines with (issue 1021),
and write bursts, whose beats go to consecutive words and which it
answers once, after the last beat (issue 1208).

\subsection*{Parameters}

\begin{center}\footnotesize
\begin{tabular}{@{}lp{0.7\textwidth}@{}}
\toprule
Parameter & Meaning \\
\midrule
\code{I} & The width of the AXI identifier on \code{bus\_req},
\code{bus\_ans} and \code{bus\_r}, in bits. The tables use 4, as the board
does since its bus has two hosts behind an arbiter (issue 241). \\
\code{AW} & The bits that number the words, 10 by default. \\
\code{W} & The words, \code{1 << AW}, 1024 by default. \\
\bottomrule
\end{tabular}
\end{center}

\subsection*{Address map}

The word a burst names is bits \code{AW + 1} to 2 of its address,
11 to 2 by default. Lane $k$ is
byte $k$ of the word: \code{lane0} is bits 7 to 0, \code{lane3} bits 31
to 24. The memory looks at no other address bits.

On the board, \code{BoardRouter} sends it the page at \code{0x1000}
with mask \code{0xffff\_f000}, so addresses \code{0x1000} to
\code{0x1ffc} are words 0 to 1023. The model's \code{DATA\_BASE} is
the same \code{0x1000}.

\dstables{Dmem}

\subsection*{Behaviour}

\begin{itemize}
\item A write is taken when no other write is held. Its first word
and identifier go into \code{paddr} and \code{pid}, and \code{pend}
is set.
\item The held write's beats are taken as they arrive, the last only
when \code{bus\_ans} has room. Each lane whose strobe bit is set takes
its byte of the beat, and \code{paddr} moves to the next word. The
last beat clears \code{pend}, and \code{bus\_ans} sends \code{OKAY}
with \code{pid}.
\item A read is taken when no write is waiting and the answer register
is free or is being emptied in this cycle. The four lanes are read
into \code{word} at the edge, the synchronous read a block RAM has.
\item The read is answered on \code{bus\_r} the cycle after, with
\code{rid} and \code{OKAY}. A burst's \code{len} goes into
\code{rleft} and the next word's index into \code{raddr}; while
\code{rleft} is not zero, each beat sent reads the next word into
\code{word} and counts \code{rleft} down, and the beat with
\code{rleft} at zero has \code{last} set. A new read is taken only
when the answer register is free or is sending a burst's last beat.
\item While a write is held, no request of either kind is taken.
\item \code{Dmem::with} fills the lanes from bytes before the first
cycle. \code{data\_word} reads a word back for a run or a test.
\end{itemize}

\subsection*{Verification}

\begin{itemize}
\item \code{//cpu/vreteno:lockstep\_test} compares every word of the model's
data memory with the four lanes at the halt, in \code{demo\_program},
\code{a\_multiply\_right\_after\_a\_load} and \code{random\_programs}.
\item \code{//cpu/vreteno:hello\_test}, \code{//cpu/vreteno:hello\_cc\_test}
and \code{//cpu/vreteno:board\_test} put the compiled programs'
constants in it with \code{Dmem::with}, and check what the programs
print.
\item The build simulates \code{Dmem<2>}'s netlist against the trace
of \code{//cpu/vreteno:vreteno} under nvc and Verilator:
\code{//docs:vreteno\_sim\_dmem}\allowbreak\code{\_dmem\_tb\_test} and
\code{//docs:vreteno\_vsim\_dmem\_test}.
\end{itemize}

\subsection*{Used by}

\code{Board}, as \code{dmem} behind the tracker \code{pdmem}. The
demonstration \code{//cpu/vreteno:vreteno}, the lockstep test and
\code{run.rs}. \code{board\_netlist} writes a program's data into
\code{lane0} to \code{lane3} of the board's netlist.

\subsection*{Limits}

It does not read a write burst's length: the last beat ends it. It
checks no address, since the router sends it only its own range, so
an address beyond word \code{W} $-$ 1 wraps. One write is held at a time. Nothing
on the machine fills it at run time except stores, so a program's
constants have to be in it before the first cycle.
\code{cpu/vreteno/board/ax7a200.xdc} constrains the four lanes to block
RAM.

% SPDX-License-Identifier: Apache-2.0
% covers: Rom
\datasheet{Rom}{Vreteno's boot memory on the bus: the 1024 words the core fetches, readable and not writable, behind an AXI peripheral end.}

\dsfacts{\code{cpu/vreteno/src/rom.rs}}{\code{vreteno32::rom::Rom}}%
{\code{//docs:vreteno}}

\subsection*{Function}

\code{Rom} is the boot memory of the Vreteno machine as a device on
its bus: the same \code{IMEM\_WORDS}, 1024 words, that the core
fetches from inside itself, at address zero, so that a load can read
a constant that sits beside the code, and a debugger or a host on the
bus can read the program that is there. The core keeps its own copy
for the fetch, which reads it in the same cycle; this copy answers the
bus a cycle after a read is taken, as the data memory does. Both are
filled with the same image before the first cycle, so they cannot
disagree, and neither is written at run time.

It is the peripheral client of an \code{AxiPer}, and its link is one
port, \code{bus}, a \code{PerPort} of the four channels: it takes
requests on \code{bus\_req} and write beats on \code{bus\_w}, and
answers on \code{bus\_ans} for a write and \code{bus\_r} for a read.
It answers single-beat bursts,
which is all the core makes.

\subsection*{Parameters}

\begin{center}\footnotesize
\begin{tabular}{@{}lp{0.7\textwidth}@{}}
\toprule
Parameter & Meaning \\
\midrule
\code{I} & The width of the AXI identifier on \code{bus\_req},
\code{bus\_ans} and \code{bus\_r}, in bits. The tables use 4, as the board
does since its bus has two hosts behind an arbiter (issue 241). \\
\bottomrule
\end{tabular}
\end{center}

\subsection*{Address map}

The word a burst names is bits 11 to 2 of its address. The memory
looks at no other address bits.

On the board, \code{BoardRouter} sends it the page at \code{0} with
mask \code{0xffff\_f000}, so addresses \code{0x0} to \code{0xffc} are
words 0 to 1023, which are the words the core fetches from
\code{0x0} to \code{0xffc}. \code{run.rs} puts it at the same place.

\dstables{Rom}

\subsection*{Behaviour}

\begin{itemize}
\item A read is taken when no write is waiting for its beat and the
answer register is free or is being emptied in this cycle. The word
is read into \code{word} at the edge, with its identifier in
\code{rid}, and \code{answer} is set.
\item The read is answered on \code{bus\_r} the cycle after, with
\code{rid}, \code{OKAY} and \code{last} set.
\item A write is taken when no other write is waiting for its beat.
Its identifier goes into \code{pid} and \code{pend} is set; nothing
is written.
\item The held write's beat is taken when it has arrived and
\code{bus\_ans} has room, and dropped, so that the write data channel is
not left holding it; \code{bus\_ans} sends \code{SLVERR} with \code{pid},
which is what a peripheral that was reached and refused says.
\item While a write waits for its beat, no request of either kind is
taken.
\item \code{Rom::with} fills the words from a program before the
first cycle: the same words \code{Vreteno::with} puts in the core.
\end{itemize}

\subsection*{Verification}

\begin{itemize}
\item \code{the\_boot\_memory\_is\_readable\_and\_not\_writable}, in
\code{//cpu/vreteno:board\_test}, runs the board under a program
that loads a word of itself from address zero, stores over it, loads
it again and says whether it kept; the first load is checked to
have returned the program rather than the zero a hole answers.
\item \code{//cpu/vreteno:hello\_test},
\code{//cpu/vreteno:hello\_cc\_test} and
\code{//cpu/vreteno:traps\_test} run compiled programs whose
constants the linker put beside the code, in \code{.rodata}, which
their loads read through it.
\item Its netlist is part of the board's, which
\code{//cpu/vreteno:vreteno\_board\_synth} synthesises and the manual
\code{//cpu/vreteno/board/sim:}\allowbreak\code{board\_test} runs under Vivado's
simulator; it is not simulated on its own.
\end{itemize}

\subsection*{Used by}

\code{Board}, as \code{rom} behind the tracker \code{prom}, at the
router's first slot. \code{run.rs}, which every compiled program's
test runs on, as \code{Rom<IW>}.

\subsection*{Limits}

It serves bursts of one beat and does not look at a burst's length. It
checks no address, since the router sends it only its own range, so
an address beyond word 1023 wraps. One write is held at a time. It is
a second copy of the program, in block RAM beside the core's, and the
two are the same only because the netlist fills both from one image:
nothing at run time can change either.

% SPDX-License-Identifier: Apache-2.0
% covers: Timer
\datasheet{Timer}{Vreteno's core-local interruptor: a count, a compare, a software interrupt and two lines, behind an AXI peripheral end.}

\dsfacts{\code{cpu/vreteno/src/timer.rs}}{\code{vreteno32::timer::Timer}}%
{\code{//docs:vreteno}, \code{//docs:axi}}

\subsection*{Function}

\code{Timer} is the core-local interruptor of the Vreteno machine, the
CLINT. It holds a 64-bit count, \code{mtime}, a 64-bit compare,
\code{mtimecmp}, and one bit, \code{msip}, and it drives two lines. It
raises \code{tirq} while the count has reached the compare, which the
core reads as bit 7 of \code{mip}, and it raises \code{sirq} while
\code{msip} is set, which the core reads as bit 3. A program raises its
own software interrupt by writing \code{msip} and clears it by writing
it back to zero; nothing else touches that bit.

The three registers are at the offsets every RISC-V platform puts them
at, so that a stock port of an operating system finds them where it
looks. That is what makes the window 64\,KiB: \code{mtime} sits at
\code{0xbff8}, far from the compare, and the two cannot share a page.

It is the peripheral client of an \code{AxiPer}, and its link is one
port, \code{bus}, a \code{PerPort} of the four channels: requests
come on \code{bus\_req} and write beats on \code{bus\_w}, and answers
go out on \code{bus\_ans} and \code{bus\_r}. The input \code{rst}
clears the count and
nothing else.

\subsection*{Parameters}

\begin{center}\footnotesize
\begin{tabular}{@{}lp{0.7\textwidth}@{}}
\toprule
Parameter & Meaning \\
\midrule
\code{I} & The width of the AXI identifier on \code{bus\_req},
\code{bus\_ans} and \code{bus\_r}, in bits. The tables use 4, as the board
does since its bus has two hosts behind an arbiter (issue 241). \\
\bottomrule
\end{tabular}
\end{center}

\subsection*{Register map}

Offsets from the controller's base, which is
\code{isa::CLINT\_BASE}, \code{0x0200\_0000}, under the mask
\code{isa::CLINT\_MASK}, \code{0xffff\_0000}: a window of 64\,KiB.

\dsregs{Timer}

The map is declared as \code{clint}, with fourteen select bits: an
address's bits 15 to 2 select the word. A word is addressed only at its
first byte. An address with either low bit set selects no word, so it
reads zero and takes no write, as the reference model has it, and a
store of a byte or a half to a register's second lane is dropped. A
word the map does not name reads zero, and a write to it is dropped
(issue 681). The other 31 bits of a write to \code{msip} are dropped
and read back as zero.

The controller selects the word by the whole 16-bit offset, bits 15 to
0, and looks at no bit above them: the router sends it only the bursts
in its window, so nothing higher has to be decoded. The offsets are far
apart, which is why the decode is on the offset rather than on two bits
of it. Every offset in the window that is not one of the four above
reads and writes as \code{msip} does; there are five words in 64\,KiB
and nothing answers for the rest.

\dstables{Timer}

\subsection*{Behaviour}

\begin{itemize}
\item The count adds one every cycle. With \code{rst} high it goes to
zero.
\item A read is taken when \code{bus\_r} has room and no write is waiting
for its beat. It is answered on \code{bus\_r} in the cycle it is taken,
with the word addressed, \code{OKAY} and \code{last} set.
\item A write is taken when no write is waiting. The map's write
enables for the word it names, one bit a register, go into
\code{pwe}, and its identifier into \code{pid}.
\item The held write's beat is taken when it has arrived and
\code{bus\_ans} has room. The lanes the strobe covers replace those bytes of
the word, the other bytes stay, and \code{bus\_ans} sends \code{OKAY}.
\item Every cycle \code{pending} takes whether \code{mtime} is at or
above \code{mtimecmp}, compared as unsigned numbers. \code{tirq} is
\code{pending}, a register's output, so the core sees in one cycle
what the timer decided in the cycle before.
\item \code{mtimecmp} starts at zero and \code{rst} does not clear it,
so the line is high until software writes a compare above the count.
\item \code{sirq} is \code{msip} itself, with no register between, so
the core sees a program's own software interrupt in the cycle the write
lands rather than the cycle after. Nothing but a write to \code{0x0000}
changes that bit: the controller neither sets it nor clears it.
\item The count is driven every cycle, and a write to one of its halves
is driven after that, so a write in the same cycle as a tick keeps the
written word and the tick is lost.
\end{itemize}

\subsection*{Verification}

\begin{itemize}
\item \code{//cpu/vreteno:lockstep\_test} hands \code{pending} and
\code{msip} to the model for each instruction and checks
\code{mtimecmp} against the model at the halt, after draining the bus,
since a store is posted and the compare lands some cycles after the
core issued it. \code{demo\_program} sets the compare to 150 and counts
the timer's interrupt. \code{random\_programs} reads the count and
writes the compare 64 cycles beyond it, in sixty-four seeded programs.
\item \code{a\_program\_that\_interrupts\_itself} writes
\code{msip}, takes the trap, clears the bit in the handler and counts;
it is the software interrupt's own test.
\item \code{//cpu/vreteno:pair\_test} and \code{//cpu/vreteno:board\_test}
run the controller on a bus beside a core, though neither is about it.
\item The build simulates \code{Timer<2>}'s netlist against the trace
of \code{//cpu/vreteno:vreteno} under nvc and Verilator:
\code{//docs:vreteno\_sim\_timer\_timer\_tb\_test} and
\code{//docs:vreteno\_vsim\_timer\_test}.
\end{itemize}

\subsection*{Used by}

\code{Board}, as \code{timer} behind the tracker \code{ptimer}, with
\code{tirq} and \code{sirq} wired to the core's two inputs of those
names, in the 64\,KiB from \code{0x0200\_0000}. The demonstration
\code{//cpu/vreteno:vreteno}, the lockstep test and \code{run.rs}.

\subsection*{Limits}

The counter increments once per clock cycle. The controller supports
single-beat bursts only: transaction fields such as \code{len} and \code{size}
are uninspected, every read response asserts \code{last}, and multi-beat
bursts are not supported. Address decoding ignores bits above bit 15 because
the interconnect router filters traffic outside the peripheral's 64\,KiB window.
One write request is held at a time, during which incoming reads stall.
Writes update 32 bits of a 64-bit register per access, requiring two store
instructions to update a full compare threshold. The counter and compare registers
reside beyond the 12-bit immediate displacement range of a single base register,
requiring separate base pointers in software routines (\code{program.rs}).
All transactions return \code{OKAY} without generating decode errors.

% SPDX-License-Identifier: Apache-2.0
% covers: Plic, Plic1, Plic2, Plic3, Plic4, Plic5, Plic6, Plic7, Plic8
\datasheet{Plic}{The platform-level interrupt controller: one to thirty-one sources and two targets, machine and supervisor, at the RISC-V PLIC's offsets.}

\dsfacts{\code{lib/parts/src/plic.rs}}%
{\code{txhdl\_parts::plic::Plic<N, EDGE>}, and \code{Plic1} to \code{Plic8}, which name the counts}%
{\code{//docs:vreteno}}

\subsection*{Function}

A controller gathers the interrupt lines of several sources into a
line for each of two targets, a hart's machine mode and its supervisor
mode (issue 1013), and lets software ask which source to serve. Each
source \textit{i}, numbered from 1, is an input line,
\code{srcs\_}\textit{i}\code{-1}, of the array of ports \code{srcs}. The
output \code{irq} is the machine target's line and \code{sirq} the
supervisor target's, \code{mip.SEIP}'s. The two targets share the
sources, their gateways and the pending bits; each has its own enable
bits, threshold and line. Software reaches the controller as an AXI-Lite
peripheral whose link is one port, \code{bus}, a \code{LitePort} of
the five channels: reads come on \code{bus\_ar} and are answered on
\code{bus\_r}, and writes come on \code{bus\_aw} and \code{bus\_w}
and are answered on \code{bus\_b}. It sits behind a
\code{LiteBridge<1, ..>}, which sends it only its own range.
The input \code{rst} clears the requests, both lines, the priorities,
and both targets' enable bits and thresholds.

It is one unit for any count of sources, \code{N} a const parameter
(issue 500): the sources are an array of ports and the priorities an
array of registers, \code{prio\_0} for source 1 onward, and the
lowering unrolls the loops over them for each \code{N}. The pending,
enable, active, held and previous-line words are 32 bits whatever
\code{N} is, bit \textit{i} for source \textit{i}, and the bits above
\code{N} stay zero; a source's number is 5 bits. \code{Plic1} to
\code{Plic8} are type aliases that name counts. \code{plic.rs} also names the
offsets, \code{PENDING}, \code{ENABLE}, \code{THRESHOLD} and
\code{CLAIM} for the machine target, \code{S\_ENABLE},
\code{S\_THRESHOLD} and \code{S\_CLAIM} for the supervisor target, and
\code{priority(i)}.

\subsection*{Parameters}

\begin{center}\footnotesize
\begin{tabular}{@{}lp{0.7\textwidth}@{}}
\toprule
Parameter & Meaning \\
\midrule
\code{N} & The count of sources, from 1 to 31, the most one word of
pending bits holds. Nothing checks it. \\
\code{EDGE} & A bit per source: bit \textit{i} set makes source
\textit{i} ask on a rising edge rather than while its line is high.
Bit 0 and bits above \textit{N} are ignored. The tables use 0, as the
board does. \\
\bottomrule
\end{tabular}
\end{center}

\subsection*{Register map}

Offsets from the controller's base, which are the RISC-V PLIC's for
two targets, contexts 0 and 1. The controller decodes the low 22 bits of an address. On
the board the base is \code{0x0c00\_0000}.

\begin{center}\footnotesize
\begin{tabular}{@{}llp{0.55\textwidth}@{}}
\toprule
Offset & Word & Access \\
\midrule
\code{0x4} $\times$ \textit{i} & Source \textit{i}'s priority, 3 bits &
Read and write, for \textit{i} from 1 to \textit{N}. \\
\code{0x1000} & The pending bits, bit \textit{i} for source \textit{i} &
Read only. \\
\code{0x2000} & The machine target's enable bits, the same way & Read
and write; bit 0 does not stick. \\
\code{0x2080} & The supervisor target's enable bits & The same. \\
\code{0x20\_0000} & The machine target's threshold, 3 bits & Read and
write. \\
\code{0x20\_0004} & The machine target's claim and complete & A read
claims; a write completes. \\
\code{0x20\_1000} & The supervisor target's threshold, 3 bits & Read
and write. \\
\code{0x20\_1004} & The supervisor target's claim and complete & The
same. \\
\bottomrule
\end{tabular}
\end{center}

Every other offset, offset 0 among them, reads as zero and ignores a
write.

\dstables{Plic}

\subsection*{Behaviour}

\begin{itemize}
\item \textbf{The gateway.} A level source asks while its line is high.
An edge source asks when its line is high and was low a cycle before
(\code{prev}), or when it holds an edge (\code{held}).
\item A request goes forward when its source is not \code{active}.
Going forward sets the source's bits in \code{pending} and in
\code{active}.
\item An edge that comes while its source is \code{active} sets its bit
in \code{held}, one deep, and it asks as soon as the source is no
longer active.
\item \textbf{The choice}, made for each target. A source is a
candidate for a target when it is pending, enabled in that target's
bits, \code{enable} or \code{senable}, and its priority is above that
target's threshold, \code{threshold} or \code{sthreshold}. The best
candidate has the highest priority, and the lowest number wins a tie.
A priority of 0 is never above a threshold, so it never interrupts.
\item \textbf{A claim.} A read of a target's claim word,
\code{0x20\_0004} or \code{0x20\_1004}, answers with that target's best
candidate's number, or 0 for none, and clears that source's
\code{pending} bit. The source stays \code{active}. So a source
enabled for both targets is served by whichever claims first, and the
other then finds it no longer pending. One read is taken a cycle, so
the two never claim together.
\item \textbf{A complete.} A write of a number from 0 to \textit{N} to
either claim word clears that source's \code{active} bit, whichever
target claimed it. A larger number changes nothing. A level source
whose line is still high then asks again at once.
\item \textbf{The lines.} \code{asserted} is set each cycle when the
machine target has a best candidate, and \code{sasserted} when the
supervisor target has one. \code{irq} follows \code{asserted} and
\code{sirq} follows \code{sasserted}, so each line is a register's
output and lags its choice by a cycle.
\item \textbf{The bus.} A read is answered in the cycle it is taken, and
a write is taken when its address and its word are both there, and
answered \code{Okay} at once. Every answer is \code{Okay}.
\end{itemize}

\subsection*{Verification}

\begin{itemize}
\item In \code{plic}, in \code{//lib/parts:parts\_test}, against a
\code{Plic3}:
\begin{itemize}
\item \code{the\_highest\_priority\_is\_claimed\_first\_and\_the\_lowest\_number\_wins\_a\_tie};
\item \code{a\_priority\_at\_or\_below\_the\_threshold\_does\_not\_interrupt};
\item \code{a\_claimed\_source\_asks\_again\_only\_after\_its\_complete};
\item \code{an\_edge\_source\_asks\_once\_per\_edge\_and\_a\_level\_source\_while\_high};
\item \code{the\_supervisor\_target\_is\_a\_target\_of\_its\_own}: its
enables and threshold read back, its line asks, and its claims and
completes follow the same rules;
\item \code{a\_source\_enabled\_for\_both\_targets\_is\_claimed\_once}:
both lines ask, the first claim takes it and the other target finds
none, and a complete by either ends the service.
\end{itemize}
\code{the\_eight\_source\_controller\_lowers} writes \code{Plic8}'s
Verilog.
\item \code{ex\_plic} runs \code{Plic3<0b100>} against a host that
claims and completes, last on the supervisor target, where source 3
enabled for it alone raises \code{sirq} and not \code{irq}. The build
simulates its netlist against that run under nvc and Verilator: \code{//docs:plic\_sim\_plic3\_tb\_test}
and \code{//docs:plic\_vsim\_test}.
\item \code{input\_comes\_by\_interrupt\_one\_byte\_each}, in
\code{//cpu/vreteno:board\_test}, runs the board's \code{Plic3} under a
compiled program that takes the serial port's input by interrupt.
\end{itemize}

\subsection*{Used by}

\code{Plic3<0>} in the board's design, \code{Board}, as its field
\code{plic} behind the bridge \code{pplic}: source 1 is the serial
port's receive interrupt, source 2 the board's \code{irq} input,
source 3 the Ethernet peripheral's, and
\code{irq} is the core's machine external interrupt, \code{mip.MEIP},
and \code{sirq} its supervisor external interrupt, \code{mip.SEIP}:
the core takes it as \code{seirq} and registers it beside MEIP
(issue 1094), which is how Linux in supervisor mode takes the serial
port's and the Ethernet port's interrupts. \code{Plic3<0b100>} in the example \code{ex\_plic}.

\subsection*{Limits}

It has two targets, contexts 0 and 1, and no others. Priorities are 3
bits, so from 0 to 7. A write ignores its strobe and writes the whole
field. The controller checks no address above bit 21, and relies on its
bridge for the rest. An edge source holds one edge during a service; a
second is lost.

% SPDX-License-Identifier: Apache-2.0
% covers: Dm
\datasheet{Dm}{The RISC-V debug module on the bus: halt, resume and status, and the general registers, \code{dpc} and \code{dcsr} read and written by the abstract command.}

\dsfacts{\code{cpu/vreteno/src/debug.rs}}{\code{vreteno32::debug::Dm}}%
{\code{//docs:vreteno}}

\subsection*{Function}

The RISC-V debug specification puts a debugger's reach into a core
behind a debug module, a small address space of 32-bit registers
bridged from the cable transport. \code{Dm} is that module as far as
a debugger on the bus needs it, with the address space on the bus:
an AXI-Lite peripheral whose registers are the specification's, at
their numbers, each a word at four times its number. Through it the
JTAG host on the board's bus halts the core from the hardware manager,
reads its registers and resumes it (issue 154), and a transport of its
own reaches the same registers the same way when there is one.

The core's side is three lines and a word: a halt request and a
resume request in, debug mode out, and the register access, a number
in the specification's space, the word to write and whether to,
answered by the word read. The core lends the module its register
file's first read port and its write port while it is halted, and its
CSR read takes the module's number then.

\subsection*{Parameters}

None. One hart, so \code{hartsel} is not decoded; 32-bit registers,
so the command's \code{aarsize} must be 2.

\subsection*{Register map}

Offsets from the module's base, which on the board is
\code{0x1000\_0000}, a router port of its own; the part selects a
register by bits 8 to 2 of the address, the specification's number,
and looks at no other bits. \code{debug::at(number)} gives a
register's address on the board.

\begin{center}\footnotesize
\begin{tabular}{@{}lllp{0.45\textwidth}@{}}
\toprule
Offset & Number & Register & Access \\
\midrule
\code{0x10} & \code{0x04} & \code{data0} & Read and write, the
abstract command's one data word: what a read brought back, or what
a write sends. Written while a command is in flight it is the busy
error and unchanged. \\
\code{0x40} & \code{0x10} & \code{dmcontrol} & Read and write.
\code{haltreq} bit 31, a level; \code{resumereq} bit 30, an action,
which asks if the core is halted as the write lands and clears the
acknowledgement, and does nothing as a zero; \code{dmactive} bit 0,
which clear holds the module at nothing. \\
\code{0x44} & \code{0x11} & \code{dmstatus} & Read only.
\code{allhalted} and \code{anyhalted} bits 9 and 8, \code{allrunning}
and \code{anyrunning} bits 11 and 10, \code{allresumeack} and
\code{anyresumeack} bits 17 and 16, \code{authenticated} bit 7, and
the version, 2 for the 0.13 specification, in bits 3 to 0. \\
\code{0x48} & \code{0x12} & \code{hartinfo} & Reads as zero: no
scratch registers and no data window. \\
\code{0x58} & \code{0x16} & \code{abstractcs} & \code{busy} bit 12,
\code{cmderr} bits 10 to 8, cleared by writing ones, and
\code{datacount} 1 in bits 3 to 0. \\
\code{0x5c} & \code{0x17} & \code{command} & Write only, the
access-register command: \code{cmdtype} 0, \code{aarsize} 2,
\code{transfer} bit 17, \code{write} bit 16, \code{regno} bits 15 to
0, \code{0x1000} to \code{0x101f} a general register and below that
a CSR. \code{debug::access(regno, write)} builds one. \\
\code{0x100} & \code{0x40} & \code{haltsum0} & Read only, bit 0 the
one hart, halted. \\
\bottomrule
\end{tabular}
\end{center}

\dstables{Dm}

\subsection*{Behaviour}

\begin{itemize}
\item \code{haltreq} and \code{resumereq} go to the core as written
and held, and \code{halted} from the core is what \code{dmstatus}
reports. A resume request drops and is acknowledged when the core has
left debug mode; a halt request stays as written, and a debugger
clears it before it resumes.
\item A command is accepted when the module is active, nothing is in
flight and no error stands, and refused otherwise as the
specification says: \code{cmderr} 1 while one is in flight, 2 for a
command the module does not do, 4 while the core runs; under a
standing error a command is ignored. A command takes two cycles, one
for the number to reach the core and one for the word to come back
into \code{data0} or go in from it.
\item The module reads any general register and any CSR the core's
CSR read knows, and writes the general registers, \code{dpc} and
\code{dcsr}; a write of another CSR is \code{cmderr} 2.
\item The arms of its one \code{with!} are ordered for what may fire
together: the core's completion of a resume before the debugger's
write, so a request written in the cycle the core resumes clears the
acknowledgement rather than losing to it (issue 439).
\item A read is answered on \code{bus\_r} in the cycle its address beat is
taken, with \code{OKAY}; a write needs its address and data beats and
room on \code{bus\_b} in one cycle, and is answered \code{OKAY} there.
Inactive, the module holds nothing and answers zeros.
\end{itemize}

\subsection*{Verification}

\begin{itemize}
\item Thirteen tests in \code{cpu/vreteno/src/debug.rs}, in
\code{//cpu/vreteno:isa\_test}, play the core: they set halted and
the word read, drive the module over AXI-Lite, and read its lines
back. Among them:
\begin{itemize}
\item \code{a\_register\_is\_read\_and\_written\_by\_the\_abstract\_command};
\item \code{an\_unsupported\_command\_or\_a\_running\_core\_is}\-\code{\_an\_error\_that\_sticks};
\item \code{a\_resume\_request\_written\_as\_the\_core\_resumes}\-\code{\_clears\_the\_ack};
\item \code{data0\_written\_while\_a\_command\_is\_in\_flight}\-\code{\_is\_a\_busy\_error}.
\end{itemize}
The last two were constructed to fail against the arms as they were
before issue 439.
\item \code{the\_debug\_module\_halts\_reads\_writes\_and\_resumes\_the\_core},
in \code{//cpu/vreteno:board\_test}, puts a debugger on the board's
JTAG pins, a master model following a plan of writes, reads and
waits, against a program counting to three thousand: it halts the
core, reads the count and \code{dpc}, writes the count to its end,
resumes, and the program finishes at once.
\item \code{//cpu/vreteno:vreteno\_board\_dm\_probe} does the same from
Vivado's hardware manager through the JTAG host, against whatever
program runs; its transcript is owed to the document once the board
is reached again.
\end{itemize}

\subsection*{Used by}

\code{Dm} in the board's design, \code{Board}, as its field
\code{dmod} behind the bridge \code{pdm}, a \code{LiteBridge<1, ..>} on the
router's seventh port at \code{0x1000\_0000}; its lines go to the
core's \code{haltreq}, \code{resumereq}, \code{dbg\_regno},
\code{dbg\_wdata} and \code{dbg\_we}, and take the core's \code{dbg}
and \code{dbg\_rdata}.

\subsection*{Limits}

One hart. The access-register command only: no program buffer, no
system bus access and no abstract access to memory, which the bus the
module sits on reaches directly. A CSR other than \code{dpc} and
\code{dcsr} is read but not written. No authentication, no reset of
the hart through \code{ndmreset}, no \code{havereset}. Not the transport:
the module is reached over the bus, and the debug transport module on
the cable is issue 154's third step. The module sits at
\code{0x1000\_0000} and not nearer the interrupt controller, whose
64~MiB window from \code{0x0c00\_0000} the router would match first.

% SPDX-License-Identifier: Apache-2.0
% covers: Dtm
\datasheet{Dtm}{The RISC-V debug transport over JTAG, behind the FPGA's BSCANE2, so that a stock OpenOCD, and gdb through it, reaches the core's debug module from the board's own cable.}

\dsfacts{\code{lib/parts/src/dtm.rs}}{\code{txhdl\_parts::dtm::Dtm}}%
{\code{//docs:examples}, \code{//docs:vreteno}}

\subsection*{Function}

The FPGA's JTAG port belongs to the device, so a design reaches it
through \code{BSCANE2}, one of the user scan chains of the device's
TAP. OpenOCD speaks to a transport there with its BSCAN tunnel,
\code{riscv use\_bscan\_tunnel 5}: it sets the device's instruction
register to \code{USER4}, and encapsulates each scan of the transport's own
TAP inside one data scan of that chain, framed as
OpenOCD 0.12.0's \code{riscv.c} frames it:

\begin{center}\footnotesize
\begin{tabular}{@{}lp{0.62\columnwidth}@{}}
\toprule
Scan & First bit shifted to last \\
\midrule
instruction & \code{0}, the width 5 in 7 bits, the register in 5,
\code{000} \\
data & \code{1}, the width $W$ in 7 bits, the data in $W + 1$,
\code{000} \\
\bottomrule
\end{tabular}
\end{center}

The payload is one bit longer than the register because what comes
back is a clock late: the device's TAP presents the bit that left on
the edge after. The unit shifts the first $W$ bits into the register
and lets the last go by.

The unit runs on \code{Tck}, the cable's clock, which \code{BSCANE2}
hands the fabric. Its request and answer cross to the system's clock
beside it, so it has one clock.

\subsection*{Registers}

The specification's, at its numbers.

\begin{center}\footnotesize
\begin{tabular}{@{}lllp{0.4\columnwidth}@{}}
\toprule
IR & Register & Width & What it is \\
\midrule
\code{0x01} & \code{idcode} & 32 & \code{0x00154001}, the transport's
own; selected after reset \\
\code{0x10} & \code{dtmcs} & 32 & version 1, \code{abits} 7, an idle
hint of one, \code{dmistat}; \code{dmireset} at bit 16 and
\code{dmihardreset} at bit 17 \\
\code{0x11} & \code{dmi} & 41 & address 7, data 32, op 2 \\
other & bypass & 1 & zero \\
\bottomrule
\end{tabular}
\end{center}

\dstables{Dtm}

\subsection*{Behaviour}

\begin{itemize}
\item A \code{dmi} scan that asks for a read (op 1) or a write (op 2)
sends its 41 bits on \code{req}, if nothing is in flight and no error
stands. The answer comes back on \code{ans}: the word above a two-bit
status, 0 for done and 2 for failed.
\item A scan that asks while an access is in flight, or that captures
\code{dmi} meanwhile, sets \code{dmistat} to 3, busy. OpenOCD answers
busy by idling longer before its next scan.
\item An answer of failed sets \code{dmistat} to 2.
\item Either error sticks, and no access goes while it stands, until
a write of \code{dtmcs} sets \code{dmireset}, which clears it, or
\code{dmihardreset}, which clears it and drops the access in flight.
\item The \code{reset} pin, the device's test-logic reset, selects
\code{idcode} again.
\item A \code{dmi} capture shows the address and the word of the last
access, with busy as its op while an access is in flight and
\code{dmistat} otherwise.
\end{itemize}

\subsection*{Verification}

\begin{itemize}
\item \code{//lib/parts:parts\_test} drives the device's side of
\code{BSCANE2} bit by bit: the \code{idcode} after reset,
\code{dtmcs}, bypass for an unknown instruction, a \code{dmi} write
then read, busy while an access is in flight until \code{dmireset},
and a failed access sticking.
\item \code{//lib/examples:ex\_dtm}, and the build simulates the
netlist against its run under nvc and Verilator
(\code{//docs:dtm\_sim\_dtm\_tb\_test}, \code{//docs:dtm\_vsim\_test}).
\item \code{//cpu/vreteno:openocd\_test} puts stock OpenOCD 0.12.0 on
the whole simulated board behind a model of the Artix-7's TAP: it
examines the hart, halts it, reads \code{pc} and memory, steps,
writes and resumes; the same \code{dtmcs} scans bit by bit as
OpenOCD's driver makes them; and gdb loading a program and stopping
at a breakpoint.
\end{itemize}

\subsection*{Used by}

The Vreteno board, on \code{USER4}, in every top. Its requests reach
the debug module through \code{DtmBridge}.

\subsection*{Limits}

One access in flight at a time, which is all the specification's
\code{dmi} has. Only the nested-TAP framing OpenOCD 0.12.0 uses by default is understood.

% SPDX-License-Identifier: Apache-2.0
% covers: DtmBridge
\datasheet{DtmBridge}{The debug transport's other half, on the system's clock: each access the transport sends becomes a read or a write on the board's link, and the specification's system bus access is served here.}

\dsfacts{\code{lib/parts/src/dtm.rs}}{\code{txhdl\_parts::dtm::DtmBridge}}%
{\code{//docs:vreteno}}

\subsection*{Function}

\code{Dtm} turns OpenOCD's scans into \code{dmi} accesses: a 7-bit
address, a word, and an op. The bridge takes each one, after it has
crossed to the system's clock, and does one of two things with it.

\begin{itemize}
\item An access to \code{sbcs} (\code{0x38}), \code{sbaddress0}
(\code{0x39}) or \code{sbdata0} (\code{0x3c}) is the debug module's
system bus access (0.13, section 3.12), served here, where the bus
master is. It is answered at once, and the access it starts runs on
the link behind it, as the specification has it.
\item Any other address is one of the debug module's registers. The
module sits on the link at \code{DM}, each register a word at four
times its number; the access goes on the link and is answered when
its word comes back.
\end{itemize}

One access is on the link at a time, which is all a transport asks
for, so the bridge's host needs identifiers of one bit.

\subsection*{Parameters}

\begin{center}\footnotesize
\begin{tabular}{@{}lp{0.62\columnwidth}@{}}
\toprule
Parameter & What it is \\
\midrule
\code{DM} & the debug module's base on the link; the board gives
\code{0x1000\_0000} \\
\code{I} & the width of the host's identifiers; the board gives 1 \\
\bottomrule
\end{tabular}
\end{center}

\subsection*{System bus registers}

\code{sbcs} reads version 1, a 32-bit address, and accesses of 8, 16
and 32 bits. Its settings are \code{sbreadonaddr} (bit 20),
\code{sbaccess} (bits 17 to 19), \code{sbautoincrement} (bit 16) and
\code{sbreadondata} (bit 15); \code{sberror} (bits 12 to 14) and
\code{sbbusyerror} (bit 22) are cleared by writing ones to them.

\dstables{DtmBridge}

\subsection*{Behaviour}

\begin{itemize}
\item A system bus access goes on a write of \code{sbaddress0} under
\code{sbreadonaddr}, a write of \code{sbdata0}, or a read of
\code{sbdata0} under \code{sbreadondata}, and only with no error
standing.
\item A narrow access goes on the link as a word at the aligned
address. A write positions data into its active byte lanes with only
their strobes set; a read moves its lane down into \code{sbdata0} and
clears the rest. \code{sbautoincrement} steps by the width.
\item An access not aligned to its width sets \code{sberror} to 3,
one wider than 32 bits to 4, and one the bus refuses to 7, other,
since the bus does not say which kind of refusal it was. The first two
never reach the bus.
\item The width is kept with the access, so a write of \code{sbcs}
while it runs does not change it.
\end{itemize}

\subsection*{Verification}

\begin{itemize}
\item \code{//lib/parts:parts\_test} runs the bridge against a model
memory that honours the strobes and checks that the link transmits word
addresses: a module register as a word at four times its number;
\code{sbcs}; reads following the address and the data; a write
landing; a byte and a halfword read and written on every lane;
autoincrement by the width; the misaligned and oversized refusals;
OpenOCD's \code{c.ebreak} round trip, which leaves the word as it was;
and a refusal by the bus sticking in \code{sberror} until cleared.
\item \code{//cpu/vreteno:openocd\_test} puts stock OpenOCD 0.12.0 and
gdb 16.3 on the whole simulated board through the transport and this
bridge: OpenOCD reads and writes memory by system bus access, and
gdb loads a program into the DDR3 and stops it at a breakpoint. A
third test has the JTAG-to-AXI master read the data memory while the
bridge reads by system bus access, and both get their words.
\end{itemize}

\subsection*{Used by}

The Vreteno board, as \code{DtmBridge<0x1000\_0000, 1>}. Its host and
the JTAG-to-AXI master's pins share the arbiter's second port through
an \code{Arbiter2}, taking turns (issue 154).

\subsection*{Limits}

The debug module has no program buffer, so a debugger reaches memory
by the system bus access served here, which is the setting OpenOCD is
given, \code{riscv set\_mem\_access sysbus}. One access at a time.

% SPDX-License-Identifier: Apache-2.0
% covers: Gpio
\datasheet{Gpio}{General purpose pins on AXI-Lite: a value out, a value in, a direction, and an interrupt per pin.}

\dsfacts{\code{lib/parts/src/gpio.rs}}{\code{txhdl\_parts::gpio::Gpio}}%
{\code{//docs:examples}}

\subsection*{Function}

\code{Gpio<N>} puts \code{N} pins behind an AXI-Lite link. A program
chooses which way each pin faces, drives the ones that face out, reads
the ones that face in, and may ask for an interrupt from any of them.

The part drives no pin itself. It offers the three wires a pad wants,
\code{drive}, the value; \code{dirs}, whether to drive it; and
\code{pins}, the value read back, and a board's top level joins them
to a tristate buffer. A \code{Pad} models no signals in a simulation, so
a part written around one could not be simulated or tested.

\code{pins} is taken through two flip-flops, \code{sync0} then
\code{sync1}, before anything reads it. Nothing else in the part reads
the port.

\subsection*{Parameters}

\begin{center}\footnotesize
\begin{tabular}{@{}lp{0.7\textwidth}@{}}
\toprule
Parameter & Meaning \\
\midrule
\code{N} & How many pins there are. Every register is \code{N} bits
wide and reads back into the low \code{N} bits of a word; the bits
above are zero. The tables use 8, as \code{ex\_gpio} does. \\
\bottomrule
\end{tabular}
\end{center}

\subsection*{Register map}

Offsets from the peripheral's base. The part selects the word by bits
4 to 2 of the address and looks at no other bits, so the router or the
bridge in front of it decides the base.

\dsregs{Gpio}

Every register is \code{N} bits wide, and the word read back has zeros
above them. A write to the register \code{pins}, which is read only, is taken
and ignored. The register reads \code{sync1}, the second of the two
flip-flops, and not the port of the same name.

VHDL reserves \code{out}, so the netlist calls the register
\code{out\_rw} in both targets (issue 497).

\dstables{Gpio}

\subsection*{Behaviour}

\begin{itemize}
\item Every cycle \code{sync0} takes \code{pins}, \code{sync1} takes
\code{sync0}, and \code{seen} takes \code{sync1}. So a read of
\code{sync1} is what the pin held two cycles before, and \code{seen}
is what it held three cycles before, which is what an edge is measured
against.
\item A pin fires while it is enabled in \code{ie} and either its
\code{kind} bit is zero and \code{sync1} equals \code{pol}, or its
\code{kind} bit is one and \code{sync1} crossed \code{seen} in the
direction \code{pol} names.
\item \code{status} takes what it held, less any bit a write to
\code{0x18} set to one in the same cycle, plus every pin firing that
cycle. A level therefore sets its bit again in the cycle after it is
cleared, while the level lasts. Clearing is by a written one and not a
written zero, so two programs clearing different pins cannot lose each
other's work.
\item \code{irq} is high while any bit of \code{status} is set. It is
combinational from \code{status}, which is a register, so it follows
the register by no cycles.
\item A read is answered on \code{bus\_r} in the cycle it is taken, with
\code{OKAY}. A write is taken when its address phase, its beat and
room on \code{bus\_b} are all there, and is answered \code{OKAY} in that
cycle.
\item An address that selects no word reads as zero and a write to it
is answered and does nothing.
\end{itemize}

\subsection*{Verification}

\begin{itemize}
\item \code{//lib/examples:ex\_gpio} sets the direction, drives a
pattern, reads a pin through the synchroniser, takes an interrupt on a
rising edge and not on the fall, clears the status, and asks for a
level and finds the bit back. Every step is asserted against what the
program read.
\item The build simulates \code{Gpio<8>}'s netlist against that run
under nvc and Verilator: \code{//docs:gpio\_sim\_gpio\_tb\_test} and
\code{//docs:gpio\_vsim\_test}.
\end{itemize}

\subsection*{Used by}

Nothing in the tree yet. It is meant for a board's LEDs and keys, and
for the header pins, behind an AXI-Lite bridge; issue 145 says so.

\subsection*{Limits}

The interrupt is one line for all \code{N} pins; which pin raised it
is in \code{status}. There is no debounce, so a key wired to a pin
bounces into an edge interrupt many times. The part does not drive a
pad, so a board that wants one writes the tristate buffer itself.
\code{N} is at most the width of a bus word, since every register
reads back in one.

% SPDX-License-Identifier: Apache-2.0
% covers: CfgFlash, FlashPins, Startup
\datasheet{CfgFlash}{The configuration flash reached from user logic after the bitstream has loaded, through \code{STARTUPE2}.}

\dsfacts{\code{lib/parts/src/cfgflash.rs}}{\code{txhdl\_parts::cfgflash::CfgFlash}}%
{\code{//docs:examples}}

\subsection*{Function}

On a 7-series part the configuration flash's clock is the dedicated
CCLK pin, which user logic reaches only through the \code{STARTUPE2}
primitive's \code{USRCCLKO} (issue 312). The primitive brings two
rules, from UG470: its \code{EOS} says when configuration is over, and
until then the flash is not the user's; and the first three clocks on
\code{USRCCLKO} after \code{EOS} never reach the pin, so a command whose
clocks start there reaches the chip three bits short.

\code{CfgFlash} answers both, in two children joined as a unit of
units. \code{prim} is \code{Startup}, the primitive as the foreign
module \code{startup} in \code{lib/parts/hdl/startup.v}: synthesised,
it instantiates \code{STARTUPE2} and registers \code{EOS} twice;
simulated, in Rust, Verilog and VHDL alike, it is a model that raises
\code{eos} after \code{EOS\_AFTER} cycles, swallows the three clocks, and
brings out the pin as \code{cclk} for a model of the chip. \code{pins}
is \code{FlashPins<GUARD>}, between the chip and the two parts that
reach it, the SPI master \code{Spi} and the window \code{FlashWin}.

\subsection*{Parameters}

\begin{center}\footnotesize
\begin{tabular}{@{}lp{0.7\textwidth}@{}}
\toprule
Parameter & Meaning \\
\midrule
\code{GUARD} & One refuses write enable, \code{0x06}, from the master;
zero passes it. The tables use 1, as the board and the example do. \\
\bottomrule
\end{tabular}
\end{center}

\code{EOS\_AFTER}, eight, and \code{SWALLOWED}, three, are constants of
\code{cfgflash.rs} and belong to the model.

\subsection*{Register map}

None. The master and the window have their own; this part has wires.

\dstables{CfgFlash}

\subsection*{Behaviour}

\begin{itemize}
\item Until \code{eos} the pins hold the chip unselected and the clock
low, and \code{w\_rst} high, so the window takes nothing off its bus.
\item After \code{eos} they turn the clock eight times, four clocks, with
the chip not selected, so the three the part swallows are spent on
nothing. Then \code{w\_rst} falls.
\item From then on the wires go to whichever part selects the chip
first, the window before the master when both ask in one cycle, and
stay with it until it lets go of its select. The other part's select
is not passed on.
\item With \code{GUARD} set, the pins count the master's command a bit
a rise of its clock. On the seventh rise of a command whose seven bits
begin write enable they raise the select, so the chip has a command
shorter than eight bits, which it discards. \code{refused} is high
until the master lets go. Every command that programs or erases a
flash, or writes its registers, is ignored by the chip without write
enable, so that one refusal is the door to all of them.
\item Every output of the pins is a register, and the clock passes the
primitive as well, so the chip sees the master's lines a cycle late
and its clock two. A master must run at a half bit of at least four
cycles, \code{DIV} of three, for the chip's answer to be on \code{miso}
when it takes it. The chip's \code{miso} goes to both parts directly.
\item Only mode 0 is supported: the clock idles low.
\item The unit runs the pins before the primitive, and the primitive
drives \code{eos} a step ahead, since the two drive each other through
wires and a wire promises nothing about which process ran first.
\end{itemize}

\subsection*{Verification}

\begin{itemize}
\item \code{//lib/examples:ex\_cfgflash} puts the master behind its
bridge and the window on its own link, both on one \code{FlashDevice}
through \code{CfgFlash<1>}. A window read asked before \code{eos} waits
and is answered; the master reads the identity \code{20 ba 18} whole,
the first command after the swallowed clocks; its write enable is cut
short and \code{refused} rises once; a second window read is answered.
The chip took the identity and the two reads whole and one command
short.
\item \code{//docs:cfgflash\_sim\_cfgflash\_tb\_test} and
\code{//docs:cfgflash\_vsim\_test} simulate the netlist against that
run under nvc and Verilator, with the startup model in both languages.
\item \code{//cpu/vreteno:board\_test} runs
\code{the\_configuration\_flash\_answers\_on\_the\_board} through
\code{Board}, and \code{//docs:vreteno} reports the same program on the
board: the identity, the latch clear and the bitstream's sync word.
\end{itemize}

\subsection*{Used by}

\code{Board}, as \code{cfgflash}, with \code{spi} and \code{flashwin}
on its other side and the configuration pins on the top.
\code{ex\_cfgflash}.

\subsection*{Limits}

The guard refuses write enable only from the master; the window only
reads. A command that a chip takes without write enable is not
refused. The window and the master share the chip by select, so a
master command while the window holds the chip is not passed on and
reads back what the line holds, and a program that uses the master
must not do so while the window is being read. Synthesised,
\code{cclk} is held low and means nothing; it exists for simulation.

% SPDX-License-Identifier: Apache-2.0
% covers: FlashWin
\datasheet{FlashWin}{A read-only window onto an SPI flash: a read on the bus becomes a fast read on the wires, and the word comes back.}

\dsfacts{\code{lib/parts/src/flashwin.rs}}{\code{txhdl\_parts::flashwin::FlashWin}}%
{\code{//docs:examples}}

\subsection*{Function}

The SPI master in \code{spi} is what a program drives a byte at a
time. \code{FlashWin} is the other way of reaching the same chip: an
address range in which an ordinary read is answered with what the
flash holds there, with no registers to drive and nothing for software
to sequence.

That is what a bootloader needs, and what a core fetching from the bus
needs, because neither can run software to read the software it is
about to run.

A read at an address sends the fast read command \code{0x0b}, the
three bytes of that address with its low two bits cleared, one byte
the chip throws away, and then takes four, which are the word, lowest
address in the lowest byte. Nine bytes on the wires for four on the
bus.

\subsection*{Parameters}

\begin{center}\footnotesize
\begin{tabular}{@{}lp{0.7\textwidth}@{}}
\toprule
Parameter & Meaning \\
\midrule
\code{DIV} & The clock: a half of a bit takes \code{DIV + 1} cycles,
so the flash is clocked at the system's rate over
\code{2 * (DIV + 1)}. The tables use 1, as the example does; a board
divides further. \\
\code{I} & The width of the AXI identifier on \code{bus\_req},
\code{bus\_ans} and \code{bus\_r}, in bits. The tables use 2, as every
design here does. \\
\bottomrule
\end{tabular}
\end{center}

\subsection*{Register map}

None. The window has no registers: every address in its range is a
word of the flash, and the range's base is the router's business
rather than the part's, as it is for the data memory.

\dstables{FlashWin}

\subsection*{Behaviour}

\begin{itemize}
\item A read is taken when nothing is on the wires and no answer is
waiting. Its address is kept with the low two bits cleared, since a
word is what the bus asks for.
\item The nine bytes go out under one select. The select is low from
the cycle the read is taken and rises when the ninth byte is done, so
each read is a command the chip sees whole, which is what a real chip
requires.
\item The wires are mode 0, and only mode 0: the clock idles low, the
chip takes a bit as the clock rises and the master moves one as it
falls. Every flash answers a fast read in that mode, so the part needs
neither the master's two mode bits nor a register to hold them.
\item The four bytes of data are assembled as they arrive, the first
into the low eight bits. The answer goes out on \code{bus\_r} in the first
cycle the channel has room after the ninth byte.
\item A write is taken, its beat is dropped and it is answered
\code{SLVERR}. The window is read only, and saying so is what makes it
safe to fetch from: nothing on the bus can change what the instruction
stream is reading, so the question of what a fetch sees when a write
lands in the same cycle does not arise. That question is issue 268.
\item \code{rst} stops whatever is on the wires and clears the answer
and the held write, and while it is high nothing is taken off the bus:
a request waits on its channel and is served once the reset ends
(issue 813). \code{CfgFlash} holds the window in reset until the
flash is ready after configuration.
\end{itemize}

\subsection*{Verification}

\begin{itemize}
\item \code{//lib/examples:ex\_flashwin} puts the window behind an AXI
link, reads two words the chip model holds and checks each against
what was put there, then writes and finds \code{SLVERR}.
\item The chip is \code{FlashDevice} from \code{spi}, the model that
already checks the master, so the two ways of reaching a flash are
checked against one description of one chip rather than two that could
drift apart.
\item The build simulates the netlist against that run under nvc and
Verilator: \code{//docs:flashwin\_sim\_flashwin\_tb\_test} and
\code{//docs:flashwin\_vsim\_test}.
\item \code{//lib/parts:parts\_test} runs
\code{a\_read\_asked\_in\_reset\_is\_answered\_after\_it}: a read asked
while \code{rst} is high is answered with the word once it falls.
\end{itemize}

\subsection*{Used by}

\code{Board}, as \code{flashwin} behind its own \code{AxiPer} on the
router's eighth port, a 16~MiB window at \code{0x2000\_0000} onto the
configuration flash through \code{CfgFlash} (issue 312), and
\code{ex\_cfgflash} beside the SPI master. The bootloader of issue 135
is meant to read programs out of it.

\subsection*{Limits}

A word costs nine bytes on the wires, which is about three hundred
cycles at the divider the example uses: slow beside a memory and fast
beside a resynthesis. There is no cache in front of it and no burst:
every word is a command of its own, and a read of several beats is not
served. Only mode 0, and only the fast read command, so a chip that
wants another command for another purpose wants the master instead.
The address is 24 bits, matching the three address bytes the command specifies, so the
window reaches 16\,MiB. Nothing is refused by address: every address
in the range is asked of the chip, and a chip that has nothing there
answers what it answers.

% SPDX-License-Identifier: Apache-2.0
% covers: Remote, Ask, Answer
\datasheet{Remote}{A memory-mapped bridge forwarding bus transactions over channels to an external software endpoint.}

\dsfacts{\code{lib/parts/src/remote.rs}}{\code{txhdl\_parts::remote::Remote}}%
{\code{//docs:examples}}

\subsection*{Function}

\code{Remote} connects to an AXI-Lite interconnect as a memory-mapped target
without decoding register state internally. Every transaction accepted by the
core is forwarded on \code{out} as an \code{Ask} message, and the corresponding
response returns on \code{back} as an \code{Answer}. The remote endpoint is
decoupled from the peripheral logic: in simulation, it connects to a software
model within the test process; on physical hardware, it connects via packet
frames to software running on a remote host.

This architecture enables hardware-software co-development. A target peripheral
can be developed and verified as software against the actual system bus. Once
the software implementation is proven, a dedicated hardware module replaces it
without modifying host bus transactions. Issue 297 documents this workflow.

The bus interface presents a multi-cycle peripheral latency: exactly one
transaction remains outstanding at a time, and the transaction channel is held
until the response arrives. This adheres to AXI-Lite protocol semantics and
accommodates network transit times.

\subsection*{Parameters}

\begin{center}\footnotesize
\begin{tabular}{@{}lp{0.7\textwidth}@{}}
\toprule
Parameter & Meaning \\
\midrule
\code{T} & The timeout limit in clock cycles before the peripheral responds
with \code{SLVERR}. This prevents an unresponsive remote endpoint from
permanently hanging the system bus. The datasheet tables use 40.
\code{ex\_remote} uses 1000 cycles, covering round-trip latency through two MACs.
\code{Board} sets this to 100,000,000 cycles (one second at 100~MHz), providing
sufficient margin for network round trips to a host workstation. \\
\bottomrule
\end{tabular}
\end{center}

\subsection*{Register map}

None. Every address within the peripheral's allocated aperture is interpreted
by the target program, and the target address is forwarded verbatim in the
\code{Ask} payload.

\dstables{Remote}

\subsection*{Behaviour}

\begin{itemize}
\item Exactly one transaction is outstanding at a time; \code{busy} indicates
active transfer, while \code{writing} indicates transaction direction.
\item Write transactions take priority over read transactions. Because writes
and reads arrive on separate channels, prioritizing writes prevents read streams
from starving pending store operations.
\item Each transaction includes a \code{tag} that increments with each dispatched
request. Responses arriving on \code{back} are accepted only when their tag matches
the currently outstanding transaction.
\item Tag validation guarantees safe timeout handling. When a transaction times
out after \code{T} cycles, the peripheral returns \code{SLVERR} to the bus and
retires the transaction. If a delayed response arrives later, its tag no longer
matches the active transaction, and the core discards it rather than misattributing
it to subsequent bus traffic.
\item The transaction tag advances when a transaction completes rather than when
it is issued, cleanly separating late responses from the subsequent request.
\item Read transactions complete on \code{bus\_r} with the returned data word,
and write transactions complete on \code{bus\_b}; either response returns
\code{SLVERR} if the remote target indicates failure or if the cycle timer expires.
\end{itemize}

\subsection*{Verification}

\begin{itemize}
\item Four unit tests in \code{lib/parts/src/remote.rs} (\code{//lib/parts:parts\_test}):
\begin{itemize}
\item \code{a\_word\_written\_to\_the\_program\_is\_read\_back\_from\_it}
verifies standard loopback;
\item \code{the\_program\_says\_which\_transactions\_fail} validates
remote error signaling;
\item \code{a\_slow\_program\_is\_waited\_for} tests latency tolerance;
\item \code{a\_program\_that\_never\_answers\_does\_not\_stop\_the\_bus}
confirms timeout recovery.
\end{itemize}
\item \code{//lib/examples:ex\_remote} executes the complete pipeline comprising
\code{RemoteLink} and dual Ethernet MACs. Netlists synthesized from this design
are simulated under both nvc and Verilator via
\code{//docs:remote\_sim\_remote\_remote\_tb\_test} and
\code{//docs:remote\_vsim\_remote\_test}.
\item \code{//cpu/vreteno:board\_test} executes the RISC-V core against the board
design: software writes a data word to \code{0x3300} and reads it back, with the
testbench harness serving response frames.
\end{itemize}

\subsection*{Used by}

\code{ex\_remote}, and \code{Board} at address \code{0x3300}, where a
\code{RemoteLink<1>} connects it to the board's Ethernet MAC for hardware-software
co-development over the network.

\subsection*{Limits}

The block supports AXI-Lite only with a single outstanding transaction, causing
bus accesses to block for the full remote round-trip latency. It implements no
transport layer directly: transmitting \code{Ask} requests to an external machine
requires \code{RemoteLink} and an Ethernet MAC, or equivalent channel adapters.
The timeout counter measures clock cycles rather than absolute wall-clock time;
operating at lower clock frequencies increases the elapsed timeout duration proportionally.

% SPDX-License-Identifier: Apache-2.0
% covers: RemoteLink
\datasheet{RemoteLink}{The wire under a remote peripheral: each transaction one Ethernet frame, and each frame that answers it an answer.}

\dsfacts{\code{lib/parts/src/remote/eth.rs}}{\code{txhdl\_parts::remote::eth::RemoteLink}}%
{\code{//docs:flagship}}

\subsection*{Function}

\code{RemoteLink} interfaces between the two channels of \code{Remote}
and the Ethernet MAC. An \code{Ask} transaction accepted on \code{ask}
is converted into a 27-byte frame transmitted on \code{tx}. An incoming
response frame received on \code{rx} is decoded into an \code{Answer}
returned on \code{back}. The Ethernet MAC described in \code{//docs:eth}
transmits and receives these frames over the physical link, keeping the
underlying network transparent to both host and peripheral.

The module ignores all other traffic on the wire. A frame belongs to
this peripheral when it specifies this protocol's EtherType, the
response message code, and the matched device number. Any mismatch in
these fields terminates processing for that frame.

\subsection*{Parameters}

\begin{center}\footnotesize
\begin{tabular}{@{}lp{0.7\textwidth}@{}}
\toprule
Parameter & Meaning \\
\midrule
\code{DEV} & The device number. It goes in every frame this sends and
must match every frame it accepts, so several of these can share one
wire and one program, each answered by the number it was asked under.
It is a parameter rather than a register because a design knows how
many peripherals it has. The tables use 3, as \code{ex\_remote} does;
the board uses 1. \\
\bottomrule
\end{tabular}
\end{center}

\subsection*{The frame}

27 bytes, before whatever padding the MAC adds to reach the minimum
frame. An answer is read at the same offsets, so a program answers by
changing four bytes of what it was sent.

\begin{center}\footnotesize
\begin{tabular}{@{}llp{0.52\textwidth}@{}}
\toprule
Bytes & Field & What \\
\midrule
0 to 6 & destination & Broadcast. The peripheral does not know where
the program is, and the program may move. \\
6 to 12 & source & \code{02:00:00:00:00:\textit{dev}}, locally
administered. \\
12 to 14 & type & \code{0x88b5}, which IEEE leaves for local
experimental use, so no registered protocol can be mistaken for it. \\
14 & kind & \code{0x01} asks, \code{0x81} answers. \\
15 & device & \code{DEV}. \\
16 & tag & The transaction's, which the answer repeats. \\
17 & flags & Bit 0 is a write on an ask, an error on an answer. \\
18 to 22 & address & Most significant byte first. \\
22 to 26 & word & The word written, or the word answered. \\
26 & strobe & Which lanes a write means. \\
\bottomrule
\end{tabular}
\end{center}

\dstables{RemoteLink}

\subsection*{Behaviour}

\begin{itemize}
\item One frame goes out at a time. An \code{Ask} is taken when
nothing is going out, and its fields are held in registers while the
frame is spelled from them a byte per cycle that \code{tx} has room,
\code{at} saying which byte is next.
\item The bytes are a function of \code{at} and those registers
rather than a buffer, so the frame costs no memory.
\item A frame coming in is counted by \code{got}, and four bytes are
checked as they pass: the two of the type, the kind, and the device.
\item What the check sets is \code{alien} rather than a match, because
a register starts at zero and a frame is this peripheral's until a
byte of it says otherwise. It is cleared on the byte after the last,
so the next frame starts innocent.
\item The tag, the error bit and the four bytes of the word are
latched as they arrive, at bytes 16, 17 and 22 to 25.
\item An \code{Answer} is sent on \code{back} on the last byte of a
frame that nothing objected to. A frame that is not this
peripheral's ends with nothing sent, and its bytes are read and
dropped rather than left to block the MAC.
\item Sending and receiving are independent: a frame can arrive while
one is going out.
\end{itemize}

\subsection*{Verification}

\begin{itemize}
\item Five tests in \code{lib/parts/src/remote/eth.rs}, in
\code{//lib/parts:parts\_test}.
\code{a\_transaction\_leaves\_as\_a\_frame} asserts all 27 bytes of a
write, one by one.
\code{a\_read\_leaves\_with\_no\_word\_and\_no\_strobe} is the other
kind.
\code{an\_answer\_frame\_becomes\_an\_answer} is the way back.
\code{a\_frame\_that\_is\_not\_this\_devices\_is\_ignored} is another
device's frame and another protocol's.
\code{the\_bus\_is\_served\_by\_a\_program\_that\_speaks\_frames} puts
a \code{Remote} in front of it and serves the bus from frames alone.
\item \code{//lib/examples:ex\_remote} runs the whole chain, two MACs
and a wire between them, and the build simulates this netlist against
that run under nvc and Verilator:
\code{//docs:remote\_sim\_remote\_link\_remote\_link\_tb\_test}
and \code{//docs:remote\_vsim\_remote\_link\_test}.
\item \code{//tools/remote:remote\_test} asserts the same 27 bytes
from the Go end, against the same constants rather than against this
implementation, so the two spellings of the protocol are held apart.
\end{itemize}

\subsection*{Used by}

\code{ex\_remote}, and \code{Board}, where a \code{RemoteLink<1>}
connects the remote peripheral at \code{0x3300} to the board's
Ethernet interface.

\subsection*{Limits}

The link transmits one frame at a time with a single transaction in
flight, matching the issue rate of \code{Remote}. The frame length is
fixed at 27 bytes and holds one data word; bursts therefore generate
as many frames as beats. The protocol provides no retransmission and
no sequence numbers beyond the transaction tag. If the physical link
drops a frame, the transaction remains unanswered until terminated by
the peripheral timeout counter. Frames lack cryptographic authentication,
allowing any station on the local network that transmits the protocol
with the matching device identifier to access the peripheral address space.

% SPDX-License-Identifier: Apache-2.0
% covers: Spi
\datasheet{Spi}{An SPI master on AXI-Lite: one byte each way at a time, under a chip select a program holds.}

\dsfacts{\code{lib/parts/src/spi.rs}}{\code{txhdl\_parts::spi::Spi}}%
{\code{//docs:examples}}

\subsection*{Function}

\code{Spi} drives the four wires of an SPI link: \code{sclk}, the
clock; \code{mosi}, the byte going out; \code{miso}, the byte coming
in; and \code{cs\_n}, the chip select, low while a chip is held.

A program writes a byte to \code{data}, waits for the transfer to
finish, and reads \code{data} for the byte that arrived the other way.
SPI has no direction: the eight clocks that send a byte are the eight
that receive one, so every transfer is a swap, and a driver that wants
to read sends a byte it does not care about.

The chip select is a bit in \code{ctrl} that a program sets and
clears, not something a transfer does. Every SPI command is several
bytes under one unbroken select, so a master that dropped it between
bytes would start a new command each time.

\subsection*{Parameters}

None. The divider and the mode are registers rather than parameters,
so one instance can reach chips that want different clocks.

\subsection*{Register map}

Offsets from the peripheral's base. The part selects the word by bits
3 and 2 of the address and looks at no other bits.

\dsregs{Spi}

A write to \code{data} starts a transfer, and a read returns the byte
that came in. \code{0xc}, which no register names, reads zero.

A half of a bit takes \code{div + 1} cycles, so the clock is the
system's over \code{2 * (div + 1)}.

\dstables{Spi}

\subsection*{Behaviour}

\begin{itemize}
\item A write to \code{data} while no transfer is running loads
\code{txb}, sets \code{busy}, and clears \code{tick}, \code{half} and
\code{count}. A write while one is running is answered and does
nothing, so a driver reads \code{state} first.
\item While \code{busy}, \code{tick} counts to \code{div}. The cycle
it reaches it, \code{half} turns over and \code{count} adds one.
\code{sclk} is \code{cpol} exclusive-or \code{half}, so it idles at
\code{cpol}.
\item The bit is taken on the leading turn when \code{cpha} is clear
and on the trailing turn when it is set. \code{rxb} shifts left and
takes \code{miso}.
\item \code{txb} shifts on the other turn, with one exception: when
\code{cpha} is set the first leading turn does not shift, because the
first bit is already on the wire and moving it would lose the last bit
of the byte. That is the turn at \code{count} zero.
\item After sixteen turns, eight bits, \code{busy} clears,
\code{half} returns to zero and \code{fired} is set. \code{irq} is
\code{fired} and the interrupt enable.
\item \code{mosi} is the top bit of \code{txb} and \code{cs\_n} is the
inverse of the select bit, both combinational.
\item A read is answered on \code{bus\_r} in the cycle it is taken, with
\code{OKAY}. A write is taken when its address phase, its beat and
room on \code{bus\_b} are all there.
\end{itemize}

\subsection*{Verification}

\begin{itemize}
\item \code{//lib/examples:ex\_spi} reads the chip's identity with
\code{0x9f} and four bytes from an address with the \code{0x0b} fast
read, against \code{FlashDevice}, and asserts both against what was
put in the model.
\item The build simulates \code{Spi}'s netlist against that run under
nvc and Verilator: \code{//docs:spi\_sim\_spi\_tb\_test} and
\code{//docs:spi\_vsim\_test}.
\end{itemize}

\subsection*{Used by}

\code{Board}, as \code{spi} on the seventh slot of the peripheral page at
\code{0x3600}, where it reaches the configuration flash through
\code{CfgFlash} (issue 312); \code{ex\_cfgflash}, beside \code{FlashWin};
and \code{ex\_spi}. Issue 146 means it for an SD card in SPI mode and
for sensors on the header as well.

\subsection*{Limits}

One chip select, so a second chip wants a second instance or a pin
from elsewhere. One byte at a time, so a driver spends a bus
transaction per byte and the interrupt is per byte, not per command.
No quad or dual mode, and no direct memory access: reading a large
image through this costs the core a read and a write for every byte.
The board's configuration flash is reached through the
\code{STARTUPE2} primitive rather than an ordinary pin, which this
part does not do.

% SPDX-License-Identifier: Apache-2.0
% covers: Syscon
\datasheet{Syscon}{The system control block: identity, build stamp, reset cause, soft reset, and a word that survives a reset.}

\dsfacts{\code{lib/parts/src/syscon.rs}}{\code{txhdl\_parts::syscon::Syscon}}%
{\code{//docs:examples}}

\subsection*{Function}

\code{Syscon} answers three questions only the hardware can answer:
which design is running, why it last restarted, and where to leave a
word for whatever runs next. It also drives the system's reset line.

Four of its words are constants from the type's parameters, so a build
can put a commit and a date into the hardware and a program can read
back which bitstream it is on. Three are registers: the cause, the
reset key and the scratch.

\subsection*{Parameters}

\begin{center}\footnotesize
\begin{tabular}{@{}lp{0.68\textwidth}@{}}
\toprule
Parameter & Meaning \\
\midrule
\code{ID} & The design's identifier, read at \code{0x00}. \\
\code{VERSION} & Its version, read at \code{0x04}. \\
\code{STAMP0} & The build stamp's low word, read at \code{0x08}. A
build would put a commit here. \\
\code{STAMP1} & The build stamp's high word, read at \code{0x0c}. A
build would put a date here. \\
\code{KEY} & The word that must be written to \code{0x14} for a reset
to happen. Any other value is ignored. \\
\bottomrule
\end{tabular}
\end{center}

\subsection*{Register map}

Offsets from the block's base. The block selects the word by bits 4 to
2 of the address and looks at no other bits.

\dsregs{Syscon}

The four constants are the parameters \code{ID}, \code{VERSION},
\code{STAMP0} and \code{STAMP1}. A write of \code{KEY} to \code{reset}
asks for a reset, a write of anything else does nothing, and
\code{reset} reads zero.

\dstables{Syscon}

\subsection*{Behaviour}

\begin{itemize}
\item Each cycle \code{cause} takes what it held, less any bit a write
to \code{0x10} set to one in that cycle, plus a bit for each reason
asserting now: \code{por}, \code{button}, \code{wdog}, and a write of
\code{KEY} to \code{0x14}.
\item The reasons accumulate. A system reset by the button and then by
software reads both bits until a program clears them, so the case
where two things went wrong is not lost.
\item A write of \code{KEY} to \code{0x14} sets \code{hold} to eight.
While \code{hold} is not zero it counts down by one a cycle. The reset
is held for eight cycles rather than one so that every part of a
design sees it.
\item \code{srst} is high while any of \code{por}, \code{button},
\code{wdog} or the key write is asserting, or while \code{hold} is not
zero. It is combinational from those and from a register.
\item Nothing in this block is cleared by \code{srst}. That is what
lets \code{cause} and \code{scratch} survive a reset, and it is why
the block sits outside whatever \code{srst} drives.
\item A read is answered on \code{bus\_r} in the cycle it is taken, with
\code{OKAY}. A write is taken when its address phase, its beat and
room on \code{bus\_b} are all there.
\end{itemize}

\subsection*{Verification}

\begin{itemize}
\item \code{//lib/examples:ex\_syscon} reads the identity and the
build stamp, reads the power-on cause, leaves a word in
\code{scratch}, clears the cause, asks for a reset with the key, and
checks that the cause comes back as software while the scratch word is
unchanged. It then writes a wrong key and checks that nothing moved,
and presses the button to see the two causes side by side.
\item The build simulates the netlist against that run under nvc and
Verilator: \code{//docs:syscon\_sim\_syscon\_tb\_test} and
\code{//docs:syscon\_vsim\_test}.
\end{itemize}

\subsection*{Used by}

Nothing in the tree yet. Issue 147 asks for it on the board, where the
watchdog of issue 149 would drive \code{wdog} and the bootloader of
issue 135 would use \code{scratch}.

\subsection*{Limits}

The four constants are type parameters, so putting a real commit in
them wants a build that generates the instantiation; this part offers
the words and does not fill them. The reset is eight cycles, not a
parameter. There is no lock, so a program can clear a cause another
program wanted to read. \code{reset} reads as zero rather than as what
was last written.

% SPDX-License-Identifier: Apache-2.0
% covers: Wdog
\datasheet{Wdog}{A watchdog timer: a countdown software has to keep refreshing, and a reset when it stops.}

\dsfacts{\code{lib/parts/src/wdog.rs}}{\code{txhdl\_parts::wdog::Wdog}}%
{\code{//docs:examples}}

\subsection*{Function}

A program that hangs stays hung until somebody presses the button, and
on a board nobody is standing next to, nobody does. So the hardware
counts down and software has to keep saying that it is still there.
When it stops saying so, \code{Wdog} asks the system control block of
\code{Syscon} for a reset, which records the watchdog as the cause, so
a program that restarts can read why it did.

Three architectural features distinguish this watchdog from a general-purpose timer:

First, a watchdog refresh must supply the parameter \code{KEY}. Errant
software frequently performs arbitrary memory writes; requiring an explicit key
prevents accidental watchdog resets from masking program corruption. Writing an
incorrect key is treated as an immediate fault rather than being ignored.

Second, windowed mode enforces a lower bound on refresh timing: writing to
\code{feed} while the counter remains above \code{sill} triggers a fault. A
software loop stuck in an infinite polling loop is dysfunctional even though it
periodically executes instructions; windowed timing detects early refreshes that
standard countdown timers miss.

Third, the configuration lock bit is write-once. Once asserted, \code{ctrl},
\code{load}, and \code{sill} become read-only until hardware reset, protecting
critical timing thresholds from corruption.

\subsection*{Parameters}

\begin{center}\footnotesize
\begin{tabular}{@{}lp{0.68\textwidth}@{}}
\toprule
Parameter & Meaning \\
\midrule
\code{KEY} & The validation word required for a refresh. Writing any other
value to \code{feed} triggers an immediate reset fault. \\
\bottomrule
\end{tabular}
\end{center}

\subsection*{Register map}

Offsets from the block's base. The block selects the word by bits 4 to
2 of the address and looks at no other bits.

\dsregs{Wdog}

Once \code{lock} is set, writes to \code{ctrl}, \code{load} and
\code{sill} are ignored, and the lock bit itself only ever sets. A write
of \code{KEY} to \code{feed} refreshes the count, and a write of
anything else is a failure; \code{feed} reads zero. In window mode, a
refresh while \code{count} is above \code{sill} is a failure too.

The field is \code{sill} and not \code{open}, because \code{open} is a
reserved word of VHDL and the lowering refused such a name when the
part was written (issue 77). It now escapes one instead, to
\code{open\_rw} (issue 497), and the field keeps the name it has.

\dstables{Wdog}

\subsection*{Behaviour}

\begin{itemize}
\item While \code{ctrl} bit 0 is set, \code{count} falls by one a
cycle. A refresh loads it from \code{load}.
\item A write of \code{KEY} to \code{feed} is a refresh, unless window
mode is on and \code{count} is above \code{sill}, in which case it is
a failure. A write of anything else is a failure in either mode.
\item When \code{count} reaches zero: if warning is on and has not
happened, the warning bit sets, \code{irq} rises and the count
reloads, which gives a program that is merely slow one more timeout to
notice. Otherwise it is a failure.
\item A failure sets the failed bit, drives \code{rst\_req} and holds it
for eight cycles, and reloads the count. The reload matters: the count
is at zero when a timeout fails, so a watchdog that did not reload
would ask for a reset on every cycle from then on and hold the machine
down rather than restart it.
\item \code{irq} is a level, standing while the warning bit is set and
warning is on. \code{rst\_req} is a request to \code{Syscon}, not a reset
line.
\item Nothing in this block is reset by its own request, which is what
lets a program read \code{status} after the restart.
\end{itemize}

\subsection*{Verification}

\begin{itemize}
\item \code{//lib/examples:ex\_wdog} sets a timeout, watches the count
fall between two reads and come back after a refresh, stops refreshing
and takes the warning and then the reset, writes a wrong key, refreshes
while the window is shut, and locks the watchdog on and fails to turn
it off.
\item The build simulates the netlist against that run under nvc and
Verilator: \code{//docs:wdog\_sim\_wdog\_tb\_test} and
\code{//docs:wdog\_vsim\_test}.
\end{itemize}

\subsection*{Used by}

Nothing in the tree yet. \code{Syscon} has the \code{wdog} line this
drives, and issue 149 asks for a board program that stops refreshing
and reads \code{watchdog} as the cause after it restarts.

\subsection*{Limits}

\code{count} and \code{load} are sixteen bits, so the timeout is at
most 65535 cycles. The reset request is eight cycles and not a
parameter. \code{feed} reads as zero rather than as what was last
written. There is no second stage: the warning and the reset are the
only two things a timeout can do.

% SPDX-License-Identifier: Apache-2.0
% covers: Pwm
\datasheet{Pwm}{Four pulse width modulated outputs on AXI-Lite: one counter, a duty a channel, and a shadow so a width never changes under a pulse.}

\dsfacts{\code{lib/parts/src/pwm.rs}}{\code{txhdl\_parts::pwm::Pwm}}%
{\code{//docs:examples}}

\subsection*{Function}

A pin is on or off, and much of what a board drives wants a value
between: an LED at a third of its brightness, a fan at half its speed, a
servo told where to stand. \code{Pwm} switches four pins fast enough
that the load cannot follow, and what the load sees is the fraction of
each period a pin was high.

One counter runs the period for all four channels, and a channel is
high while the counter is below its duty. That comparison is what makes
the ends flat: a duty of zero never wins it and a duty of the period or
more always does, so neither end leaves a pulse one cycle wide.

\subsection*{Parameters}

None. The period, the duties and the four polarity bits are registers,
so one instance drives four channels of different widths; the channel
count and the 16-bit period are fixed in the unit.

\subsection*{Register map}

Offsets from the peripheral's base. The part selects the word by bits 4
to 2 of the address and looks at no other bits, so the bridge in front
of it decides the base and the map repeats every 32 bytes.
\code{pwm::duty(i)} gives channel \code{i}'s offset, which is
\code{8 + 4i}.

\dsregs{Pwm}

\code{ctrl} is eight bits wide. \code{CTRL\_ENABLE} and
\code{CTRL\_CENTRE} are its first two fields' masks, and
\code{pwm::polarity(i)} names channel \code{i}'s bit of \code{pol}. A
read of \code{period} or of a duty word returns the shadow, not the
value the counter is running, while \code{count} is the live counter,
and a write to it is taken and does nothing. \code{0x1c} is named by no
register and reads zero.

A write takes the low 16 bits of the word for a period or a duty and
the low 8 for \code{ctrl}, and the rest are dropped.

\dstables{Pwm}

\subsection*{Behaviour}

\begin{itemize}
\item Edge aligned, which is \code{CTRL\_CENTRE} clear, the counter
runs from zero to the period less one and returns to zero, so a period
is that many cycles and every pulse starts at the same moment.
\item Centre aligned it counts up and then down, holding one cycle at
each end, so a period is twice the count and a pulse is centred in it,
as one run rather than two at the ends. A bank of outputs driven that
way does not switch all at once, which is what a supply prefers.
\item A write of a period or a duty goes to a shadow, and every shadow
is loaded together when the counter wraps. A width that changed under
the counter would make one pulse of neither the old width nor the new,
which a load that averages sees as a step nobody asked for and a servo
reads as a command.
\item While the counter is stopped the shadows are loaded every cycle
and the counter is held at zero, so a period written then takes effect
at once: no pulse is under way to spoil.
\item A channel's pin is the comparison exclusive-or'd with its
polarity bit. While the counter is stopped the pins are the four
polarity bits themselves, so a design chooses what a stopped output
rests at.
\item A period of zero stops the counter, since a period shorter than a
cycle is no period at all.
\item A read is answered on \code{bus\_r} in the cycle its address beat is
taken, with \code{OKAY}. A write needs its address beat, its data beat
and room on \code{bus\_b} in the same cycle, and is answered \code{OKAY} in
that cycle.
\end{itemize}

\subsection*{Verification}

\begin{itemize}
\item Six tests in \code{lib/parts/src/pwm.rs}, in
\code{//lib/parts:parts\_test}, drive it over AXI-Lite and count the
cycles each pin was high.
\code{a\_duty\_is\_the\_cycles\_high\_in\_a\_period} finds three and seven
of ten, four periods over.
\code{the\_ends\_are\_flat} finds a duty of zero never high and a duty of
the period, or past it, never low.
\code{polarity\_turns\_a\_channel\_over} finds the same duty high for
nine cycles and for twenty-one.
\code{a\_duty\_written\_mid\_period\_makes\_no\_runt} widens a duty in the
middle of a period and finds every whole pulse two wide or eight, and
the last of them eight.
\code{a\_centred\_pulse\_is\_one\_run\_a\_period} finds four of eight,
centred, as one run of eight in sixteen.
\code{a\_stopped\_counter\_sits\_at\_the\_polarity} finds the pins at
their polarity bits with the enable clear.
\item \code{//lib/examples:ex\_pwm} drives it through a
\code{LiteBridge<1, ..>} on an AXI4 link and prints what the program
measured: the duties as cycles high, the pulses around a change of
width, a channel turned over, and the centred pulses.
\item The build simulates the netlist against that run under nvc and
Verilator: \code{//docs:pwm\_sim\_pwm\_tb\_test} and
\code{//docs:pwm\_vsim\_test}.
\end{itemize}

\subsection*{Used by}

Nothing in the tree yet. Issue 248 asks for a program on the board that
fades an LED with it.

\subsection*{Limits}

Four channels of 16 bits, neither a parameter. The period is in system
clock cycles, since there is no prescaler, so the slowest period is
65\,535 cycles. Byte strobes are ignored. Nothing is refused: an
address that names no word reads as zero or as the counter, and every
read and write is answered \code{OKAY}. A read of a period or a duty
returns what was written rather than what the counter is using, so
there is no way to read the value in use. There is no interrupt, no
dead time between complementary outputs, no fault input, no phase
between channels and no period of its own for any of them. A period of
zero written while the counter is enabled cannot take effect, because
the load of the shadows is itself gated on the counter running; the
enable has to come off first.

% SPDX-License-Identifier: Apache-2.0
% covers: Tee
\datasheet{Tee}{One stream to two receivers in step: a word is taken when both have room, and neither can run a word ahead of the other.}

\dsfacts{\code{lib/parts/src/redundant.rs}}{\code{txhdl\_parts::redundant::Tee}}%
{\code{//docs:vreteno}}

\subsection*{Function}

Fault-tolerant lockstep systems execute redundant processor instances
and compare outputs to detect silent computational divergence. This
comparison remains valid only if both execution paths receive identical
inputs on identical clock cycles. \code{Tee<T>} enforces this lockstep
invariance on a streaming channel: it accepts incoming data words from its
single input port, duplicates them across both output ports, and accepts the
next word only after both receivers acknowledge receipt.

The block performs replication and synchronization without altering payload
order or performing comparisons; error verification occurs downstream in
\code{Check}.

\subsection*{Parameters}

\code{T}, the payload, which must be a \code{Transaction} and a
\code{Value}. The tables use the tree's read answer,
\code{R<32, 2>}. \code{Pair} also tees \code{Done<2>} and
\code{Grant<2>}, and the example and the tests tee a byte.

\dstables{Tee}

\subsection*{Behaviour}

\begin{itemize}
\item The word is registered rather than passed through. A unit that
read a channel and drove one in the same cycle would be a
combinational path from one side's handshake to the other's, which is
the shape this repository keeps out of its channels, and its netlist
could not be replayed against the run cycle by cycle. So a word taken
in one cycle is offered from the next.
\item \code{full} says a word is held. Each side takes it in the first
cycle in which that side has room and has not taken it already, which
\code{took\_a} and \code{took\_b} record.
\item A new word is taken when the held one has been taken by both, or
when none is held. So the tee refuses its input for as long as the
slower side has not caught up, and the faster side cannot be offered
the next word early.
\item With both sides ready in every cycle the tee passes one word a
cycle, and the two sides are always on the same word.
\end{itemize}

\subsection*{Verification}

\begin{itemize}
\item \code{both\_sides\_of\_a\_tee\_see\_every\_word} and
\code{a\_slow\_side\_holds\_the\_other\_back}, in
\code{//lib/parts:parts\_test}. The second gives one side a receiver
that takes a word every other cycle and finds the two sides' lists
equal, which is the property the part exists for.
\item \code{//lib/examples:ex\_redundant} runs two adders behind one
tee and gives one of them a different addend halfway through. The build
simulates the netlist against that run under nvc and Verilator:
\code{//docs:redundant\_sim\_tee\_tee\_tb\_test} and
\code{//docs:redundant\_vsim\_tee\_test}.
\end{itemize}

\subsection*{Used by}

\code{Pair}, three times: the read answers, the write completions and
the identifiers the tracker hands out, each to both cores.

\subsection*{Limits}

Two receivers, not \code{N}. One word of storage, so one cycle between
taking a word and offering it, and no queue behind that. A side that
never becomes ready stops the tee, which is the intended behaviour and
also means the part has no timeout. There is no reset port: the
registers start at their defaults.

% SPDX-License-Identifier: Apache-2.0
% covers: Check
\datasheet{Check}{Two streams to one, with a line that says they differed: a word moves when both offer one, the first copy's goes on.}

\dsfacts{\code{lib/parts/src/redundant.rs}}{\code{txhdl\_parts::redundant::Check}}%
{\code{//docs:vreteno}}

\subsection*{Function}

\code{Check<T>} is the other half of running one thing twice. It takes
a word from each copy in the same cycle, passes the first copy's on,
and raises \code{differs} when the two are not equal. The line is
raised in the cycle the difference is seen and stays raised until
\code{rst}.

It compares the whole payload, since \code{T} is a \code{PartialEq},
rather than a checksum of it.

\subsection*{Parameters}

\code{T}, the payload, a \code{Transaction}, a \code{Value} and a
\code{PartialEq}. The tables use \code{Issue<32>}, the burst a core
issues. \code{Pair} also checks \code{W<32, 4>} and \code{Grant<2>},
and the example and the tests check a byte.

\dstables{Check}

\subsection*{Behaviour}

\begin{itemize}
\item A word moves only when both inputs offer one and the output has
room. Both are taken in that cycle and the first's is sent.
\item The comparison is made only when a word moves. Two copies that
are merely out of step, one offering and the other not, stall the part
rather than raise the line: a copy that has stopped shows up on the
wires a design compares beside this, not here.
\item \code{differed} is set in the cycle a difference moves and is
cleared only by \code{rst}. The \code{differs} port is that register
or this cycle's difference, so the line is up in the same cycle rather
than the one after.
\item The differing word is not dropped and the stream is not stopped:
the first copy's word goes on whatever the second said.
\end{itemize}

\subsection*{Verification}

\begin{itemize}
\item \code{two\_that\_agree\_pass\_through\_with\_the\_line\_down} and
\code{a\_difference\_raises\_the\_line\_and\_keeps\_it\_up}, in
\code{//lib/parts:parts\_test}. The second feeds two lists that differ
in their middle word and finds the first list on the output and the
line still up at the end.
\item \code{//lib/examples:ex\_redundant}, and the netlist simulated
against its run under nvc and Verilator:
\code{//docs:redundant\_sim\_check}\allowbreak\code{\_check\_tb\_test} and
\code{//docs:redundant\_vsim\_check\_test}.
\end{itemize}

\subsection*{Used by}

\code{Pair}, three times: the bursts the two cores issue, the write
beats they send and the identifiers they release.

\subsection*{Limits}

It says that there was a difference, not which copy is wrong and not
which word it was in; a design that wants either keeps its own record.
It does not stop the stream on a difference. It keeps no word: the two
heads are compared as they pass, so a difference is seen and not
recorded. There is no timeout: two copies that stop offering stall the
output for ever.

% SPDX-License-Identifier: Apache-2.0
% covers: Tracer
\datasheet{Tracer}{A ring of the last N things that happened, read out over AXI-Lite after the machine has stopped.}

\dsfacts{\code{lib/parts/src/tracer.rs}}{\code{txhdl\_parts::tracer::Tracer}}%
{\code{//docs:examples}}

\subsection*{Function}

\code{Tracer<N>} takes one 128-bit entry a cycle and keeps the last
\code{N} of them. What it holds is always the most recent window: when
the ring is full the oldest goes. A host reads the window out over
AXI-Lite, oldest first, by writing a cursor and reading the four words
of the entry the cursor names.

It is meant for what a machine was doing when it stopped, which is why
the readout is separate from the capture: a design raises \code{halt},
the buffer freezes, and the window is read at leisure afterwards over
the same bus the rest of the system uses.

An entry is 128 bits because a retirement is. A program counter, an
instruction, the register written and the word written come to a
hundred and one bits, and a width that is a whole number of bus words
is one the readout walks without shifting.

\subsection*{Parameters}

\begin{center}\footnotesize
\begin{tabular}{@{}lp{0.7\textwidth}@{}}
\toprule
Parameter & Meaning \\
\midrule
\code{N} & How many entries the ring holds. It is a power of two, since
the ring wraps by masking. The tables use 8, as the example and the
tests do. \\
\bottomrule
\end{tabular}
\end{center}

\subsection*{Register map}

Offsets from the peripheral's base. The buffer selects the word by bits
4 to 2 of the address and looks at no other bits, so the bridge in
front of it decides the base, and the seven words repeat every 32
bytes.

\dsregs{Tracer}

\code{CTRL\_RUN} and \code{CTRL\_FREEZE} are the masks of
\code{ctrl}'s two fields. \code{0x0c} is named by no register and reads
zero. A write to \code{count} or to a data word is taken and does
nothing.

\code{tracer::word(k)} gives the offset of word \code{k}, which is
\code{0x10 + 4k}.

\dstables{Tracer}

\subsection*{Behaviour}

\begin{itemize}
\item An entry is stored when the run bit is set, \code{take} is high,
and the buffer is not freezing. It goes into the slot \code{head}
names, \code{head} moves on by one and wraps by masking, and
\code{count} rises until it reaches \code{N} and then stays there.
\item What is held runs from the oldest to the slot before the head, so
a read at cursor \code{k} is the entry \code{k} after the oldest:
the slot is \code{(head - count + cursor)} masked to the ring.
\item The freeze is a level, not an edge. While \code{CTRL\_FREEZE} is
set and \code{halt} is high, the run bit is cleared, and it is cleared
again in every such cycle, so a host cannot start the buffer while
\code{halt} is still up. A host that reads \code{ctrl} back therefore
learns why the buffer stopped.
\item A write of the run bit while the buffer is stopped also zeroes
\code{count} and \code{head}: a host that starts the buffer starts a
new window rather than adding to the old one.
\item A read is answered on \code{bus\_r} in the cycle its address beat is
taken, with \code{OKAY}. It returns the registers as the last edge left
them, so a read in the same cycle as a write of the same word returns
the old value.
\item A write needs its address beat and its data beat in the same
cycle, and room on \code{bus\_b}; it is answered \code{OKAY} in that cycle.
The buffer holds no write address between cycles.
\item The cursor is a register, so a host writes it in one transaction
and reads the entry in the next.
\end{itemize}

\subsection*{Verification}

\begin{itemize}
\item Five tests in \code{lib/parts/src/tracer.rs}, in
\code{//lib/parts:parts\_test}, drive a \code{Tracer<8>} through a
\code{LiteHost}. \code{what\_is\_held\_is\_the\_last\_n\_entries} offers
twelve into a ring of eight and finds the fifth to the twelfth.
\code{fewer\_than\_a\_ring\_are}\allowbreak\code{\_held\_oldest\_first} offers three and
finds three, in order. \code{a\_halt\_freezes\_the\_window} offers three,
raises \code{halt}, offers six more, and finds the three and a cleared
run bit. \code{a\_start\_opens\_a\_new\_window} finds the count and the
window reset by a start. \code{a\_stopped\_buffer\_takes\_nothing} finds
nothing held when the buffer was never started.
\item \code{//lib/examples:ex\_tracer} puts it behind a
\code{LiteBridge<1, ..>} on an AXI4 link at \code{0x1000}, offers twelve
entries, raises \code{halt} and offers four more that the buffer must
refuse.
\item The build simulates \code{Tracer<8>}'s netlist against that run
under nvc and Verilator:
\code{//docs:tracer\_sim\_tracer\_tb\_test} and
\code{//docs:tracer\_vsim\_test}.
\end{itemize}

\subsection*{Used by}

Nothing in the tree yet. It is meant for the core's retirement wires,
which the board computes and throws away; issue 240 says so.

\subsection*{Limits}

\code{N} must be a power of two. The entry width is 128 bits and the
bus 32, and neither is a parameter. The buffer takes one entry a cycle,
since the ring has one write port. Byte strobes are ignored: a write of
\code{ctrl} or \code{cursor} takes the word's low bits whatever the
strobe says. Nothing is refused: an offset that names no word reads as
zero or as \code{ctrl}, and every read and every write is answered
\code{OKAY}. The cursor is not masked when it is written, only when the
slot is formed, so a cursor past \code{count} reads a slot that is in
the ring but not in the window. There is no reset port, no interrupt,
no timestamp and no flag that says the ring wrapped.

% SPDX-License-Identifier: Apache-2.0
% covers: Uart
\datasheet{Uart}{Vreteno's serial port: an AXI-Lite peripheral with SiFive's register map, a line each way, a transmit queue of eight and a receive queue of sixty-four.}

\dsfacts{\code{cpu/vreteno/src/uart.rs}}{\code{vreteno32::uart::Uart}}%
{\code{//docs:vreteno}, \code{//docs:noc}}

\subsection*{Function}

\code{Uart} sends bytes on \code{tx} and receives them on \code{rx}.
A frame is one start bit, eight data bits least significant first and
one stop bit, or two with \code{nstop}. Each bit lasts \code{div} plus
one cycles. The lines rest high.

Its registers are SiFive's \code{sifive,uart0} (issue 1011), so
Linux's \code{serial/sifive.c}, OpenSBI's console and Zephyr's
\code{uart\_sifive} drive it unchanged. A byte written to
\code{txdata} joins a transmit queue of eight. A byte received lands
in a receive queue of sixty-four, and a read of \code{rxdata} takes
the oldest. SiFive's queue holds eight, but Linux's driver on this core
takes longer to answer the interrupt than eight characters take to
arrive at 115200 baud, so a line pasted at its shell lost all but its
first eight (issue 1153). \code{irq} is high while a watermark that \code{ie} enables is
passed.

It is an AXI-Lite peripheral. It sits behind a \code{LiteBridge<1, ..>},
which hands it one transaction at a time with no identifier and no
burst. Its link is one port, \code{bus}, a \code{LitePort} of the
five channels: reads come on \code{bus\_ar} and are answered on
\code{bus\_r}; writes come on \code{bus\_aw} and \code{bus\_w} and
are answered on \code{bus\_b}. The input \code{rst} stops both sides,
empties both queues and puts the controls back to their reset values.

\subsection*{Parameters}

\begin{center}\footnotesize
\begin{tabular}{@{}lp{0.7\textwidth}@{}}
\toprule
Parameter & Meaning \\
\midrule
\code{DIV} & The bit period \code{div} starts with, in cycles; the
register resets to \code{DIV} less one. The tables use 868, which
\code{uart.rs} states is 115200 baud at 100~MHz, as the board runs.
The checked runs use 4. \\
\bottomrule
\end{tabular}
\end{center}

\subsection*{Register map}

Offsets from the port's base. On the board the base is \code{0x3000},
which is also \code{isa::UART\_BASE}. The port selects the word by
bits 4 to 2 of the address.

\dsregs{Uart}

The map is declared as \code{serial}. Bit 31 of \code{txdata} reads
one while the transmit queue is full, and a byte written then is
dropped. A read of \code{rxdata} takes the oldest byte; bit 31 reads
one when there was none, and the low byte then reads zero. \code{ip}'s
\code{txwm} is set while fewer than \code{txcnt} bytes wait to go
out, and its \code{rxwm} while more than \code{rxcnt} wait to be read.
Writes to \code{rxdata} and \code{ip} are ignored, and \code{0x1c},
which no register names, reads zero.

Two resets differ from SiFive's, so that a program that never writes
the controls finds the port running: \code{txen} and \code{rxen}
start set, and \code{div} starts at \code{DIV} less one, 867 for
115200 baud at 100~MHz. \code{ie} starts clear, as SiFive's does, so
\code{irq} stays low until a program sets a watermark's enable. It
started with \code{rxwm} set until issue~\#1137: Linux's driver asks
for its interrupt before its port is registered, and a receive enable
that started set let the loader's bytes leave a request pending in the
interrupt controller, which Linux took before its port was ready
(issue~\#1136).

\dstables{Uart}

\subsection*{Behaviour}

\begin{itemize}
\item A read is taken when \code{bus\_r} has room and a request is on
\code{bus\_ar}, and it is answered with \code{OKAY} in the same cycle.
\item A write is taken when \code{bus\_b} has room and both
\code{bus\_aw} and \code{bus\_w} hold a request. It is answered with
\code{OKAY} in the same cycle, whichever word it names.
\item A byte written to \code{txdata} with the queue not full goes in
at the queue's tail. It also goes into \code{last}, and \code{sent}
counts it.
\item When the line is idle, the queue holds a byte and \code{txen}
is set, the oldest byte leaves the queue. Its frame loads into
\code{shift}, and \code{bits} into ten, or eleven with \code{nstop}.
\item While \code{bits} is not zero the port is busy and \code{tx} is
bit 0 of \code{shift}. Every \code{div} plus one cycles the frame
shifts down and \code{bits} counts down. When idle, \code{tx} is high.
\item The \code{rx} line passes through one register, \code{line},
before the logic.
\item When the port is not receiving, \code{rxen} is set and
\code{line} is low, a frame begins. The port samples \code{line} in
the middle cycle of each of the ten bits.
\item A high where the start bit should be is a false start, and the
frame is dropped. The eight data bits shift into \code{rx\_shift}.
\item A high stop bit lands the byte. If the receive queue holds
sixty-four, the byte is dropped and \code{dropped} counts it; else it goes
in at the tail and \code{received} counts it. A low stop bit lands
nothing.
\item A push and a removal in the same cycle leave either queue's
count unchanged.
\item \code{irq} is \code{ie}'s \code{txwm} and \code{ip}'s
\code{txwm}, or \code{ie}'s \code{rxwm} and \code{ip}'s \code{rxwm}.
\item With \code{rst} high, \code{bits}, \code{tick},
\code{rx\_bits} and \code{rx\_tick} go to zero, both queues empty, and
the controls and \code{div} return to their reset values.
\end{itemize}

\subsection*{Verification}

\begin{itemize}
\item \code{//cpu/vreteno:lockstep\_test} puts the terminal of
\code{term.rs} on the two lines. At the halt it checks \code{sent} and
\code{last} against the bytes the model was given, and that
\code{dropped} is zero. In \code{demo\_program} the port must say
\code{OK} and a newline, receive three bytes, and echo them.
\item \code{//cpu/vreteno:hello\_test},
\code{//cpu/vreteno:hello\_cc\_test} and \code{//cpu/vreteno:board\_test}
check what compiled programs print on \code{tx}. The harnesses run the
line on after the halt, while the transmit queue empties.
\item \code{//cpu/vreteno:board\_test}'s
\code{a\_line\_typed\_back\_to\_back\_waits\_whole\_in\_the\_receive\_queue}
types forty-eight bytes at the line's full speed while the program is
busy, and only then has it read and echo them. Every byte must come
back, which a queue of eight would not allow.
\item \code{//cpu/vreteno:zephyr\_port\_test} checks that the map's
offsets and masks are the ones Zephyr's stock \code{uart\_sifive}
driver uses, that Zephyr's device tree names the port
\code{sifive,uart0} with its clock, and that the console is that port.
\item The build simulates \code{Uart<4>}'s netlist, the entity
\code{uart4}, against the trace of \code{//cpu/vreteno:vreteno} under
nvc and Verilator: \code{//docs:vreteno\_sim\_uart4\_uart4\_tb\_test}
and \code{//docs:vreteno\_vsim\_uart4\_test}. The run also writes
\code{Uart<868>}'s netlist, which no test simulates.
\end{itemize}

\subsection*{Used by}

\code{Board}, as \code{uart} behind the bridge \code{puart}, with
\code{irq} as source 1 of the interrupt controller. The
demonstration \code{//cpu/vreteno:vreteno}, the lockstep test and
\code{run.rs}. The system in \code{soc}, where it is the serial port at
corner 1,1 of the lattice. Zephyr's console on Vreteno, through the
stock \code{uart\_sifive} driver.

\subsection*{Limits}

A byte written while the transmit queue is full is dropped, and a
byte received with the receive queue full is dropped and counted.
\code{sent}, \code{received} and \code{dropped} are eight-bit counters
that are not on the bus. The port checks no address beyond bits 4 to
2, since the bridge sends it only its own range. There is no parity
bit and no flow-control line.

% SPDX-License-Identifier: Apache-2.0
% covers: Ddr3Per
\datasheet{Ddr3Per}{The board's DDR3 memory as an AXI peripheral: AMD's MIG controller with its own AXI4 port, joined to the link by \code{AxiPerPins}, and making the design's clock.}

\dsfacts{\code{ddr3/src/lib.rs}}{\code{ddr3::Ddr3Per}}%
{\code{//docs:vreteno}, \code{//docs:pcie}}

\subsection*{Function}

\code{Ddr3Per} puts the DDR3 memory of the Alinx AX7A200B on an AXI
link. It is a unit of units with two children. \code{pins} is
\code{AxiPerPins<32, 32, 4, 5>} of \code{txhdl\_parts}, which puts the
link's peripheral end onto an AXI4 peripheral's pins and holds
nothing. \code{ctl} is \code{Ddr3}, AMD's MIG 7 Series controller,
which the build generates with its AXI4 port as
\code{//ddr3:ddr3\_mig\_axi}, behind the wrapper
\code{ddr3/hdl/ddr3\_axi32.v}. Thirty-seven AXI4 signals join the two.

There is no tracker in front of it. The controller keeps its own
transactions, cuts a burst into the memory's commands, and packs
thirty-two bit beats into its 256-bit native interface, so the
peripheral takes the link's channels as they are: \code{aw},
\code{ar} and \code{w} in, \code{b} and \code{r} out.

\code{Ddr3} is a foreign unit. Its netlist is an instance of the
module \code{ddr3\_axi32}, which the lowering names and does not write.
Its \code{run} is a model for simulation: \code{PinRam} of
\code{AxiPerPins} on the controller's pins, which holds the whole
gigabyte and is not ready until it has calibrated. The model does not
drive the memory's pins.

On the memory side its ports are the board's clock and the
controller's reset in, the design's clock and its reset out, since the
controller makes them, the \code{calib} output, and the memory's pins,
with \code{dq}, \code{dqs} and \code{dqs\_n} as pads both ways.

\subsection*{Parameters}

None. The controller calibrates in hardware, and its simulation
library provides the variant with the fast calibration, so neither
wants a parameter. The widths are the controller's: 32-bit data, five
identifier bits, and a 30-bit address, the memory's whole gigabyte.

\subsection*{Address map}

The link's address is 32 bits and the controller's 30, so the wrapper
drops the top two bits. On the board, \code{BoardRouter} sends the
peripheral every address that matches \code{0x4000\_0000} under the
mask \code{0xc000\_0000}, and dropping the two bits leaves the offset
from \code{0x4000\_0000}.

\dstables{Ddr3Per}

\subsection*{Behaviour}

\begin{itemize}
\item An address phase or a write beat at the head of its channel is
the controller's \code{valid} and fields, taken in the step its
\code{ready} is high. A response or a read beat the controller offers
goes onto its channel in the step the channel has room, and that room
is the controller's \code{ready}.
\item Every field of the link's address phase reaches a pin except the
region, which the controller has no pin for.
\item The wrapper gives the controller its AXI reset from the user
clock's reset, inverted and registered, never requests a refresh,
self-refresh or ZQ calibration, which the controller makes itself, and
gives it a constant temperature.
\item The controller makes the clocks. It takes the board's 200~MHz on
\code{sys\_clk}, which is also its delay reference, runs the memory
at 400~MHz, and hands out \code{ui\_clk}, the 100~MHz the design
runs on, with \code{ui\_rst} high until that clock is good. In the
netlist the controller's \code{i\_ui\_clk} is the design's clock,
which is that same clock fed back. \code{sys\_rst} is active high, a
pulse at power-on on the board.
\item \code{calib} is the controller's \code{o\_calib\_complete}.
\item In simulation the model is not ready for \code{MODEL\_WARMUP},
64 cycles. It offers a read's first beat \code{MODEL\_READ\_LATENCY},
24 steps, after it sees the address phase, and a write's response
\code{MODEL\_WRITE\_LATENCY}, 2 steps, after it sees the last beat. It
sees what the pins part drove a step after the pins part drove it.
Those are the controller's latencies as \code{//ddr3/sim:wrapper\_test}
measured them against the Micron models: a read's first beat 25
cycles after its address phase was taken, 27 for a burst of sixteen
and 63 right after a write, and a write's response 3 cycles after its
last beat.
\end{itemize}

\subsection*{Verification}

\begin{itemize}
\item \code{//ddr3:ddr3\_test} runs \code{words\_go\_in\_and\_come\_back}.
Behind an \code{AxiHost}, it writes words at \code{0x4000\_0000},
\code{0x4000\_0104} and \code{0x7fff\_fffc}, the first while the model
still calibrates. It reads them back, writes and reads back a burst of
sixteen, and checks that \code{calib} did not rise before the
warm-up.
\item \code{//ddr3:ddr3\_test} also runs
\code{the\_controller\_is\_instantiated\_and\_not\_written}. It checks
that the Verilog and the VHDL of \code{Ddr3Per} instantiate
\code{ddr3\_axi32}. The controller's clock must be on \code{clk}, the
board's clock and the clock it makes on the ports, and its pads on
the ports. The pins part's module must be written and the
controller's must not.
\item \code{//ddr3:ddr3\_test} runs the bandwidth tests of
\code{ddr3::bw}. With bursts in flight, reads take a word a cycle at
any latency up to 32 cycles, and one burst at a time pays the latency
once. Writes take a beat a cycle. The peripheral with its model
behaves as the pins part does at the model's latency.
\item \code{//cpu/vreteno:board\_test} runs
\code{the\_memory\_test\_runs\_on\_the\_board}, a compiled program that
writes and reads the memory through the model and must print
\code{ddr3 ok}. It also runs \code{the\_netlist\_holds\_the\_controller},
and \code{the\_ddr3\_path\_is\_timed\_by\_the\_core}, which finds a load
from the memory costing the model's read latency less two cycles more
than a load from the data memory.
\item \code{//ddr3/sim:wrapper\_test} simulates the wrapper and the
controller against the Micron models under Vivado's simulator. It
calibrates, then writes and reads back a burst of sixteen, a word in
another row, a word through an address with its top bits set, half a
word under its strobes, and a byte at an address that is not a
word's. It prints the latencies above. It is manual.
\code{//ddr3/sim:example\_test} is AMD's own proof of the controller as
configured, the native example design's traffic generator writing the
memory and reading it back through it; manual too.
\item \code{//cpu/vreteno/board/sim:}\allowbreak\code{board\_test} simulates the whole
board, this controller included, against two Micron models. It is
manual.
\end{itemize}

\subsection*{Used by}

\code{Board}, as \code{ddr3}, on the router's port for
\code{0x4000\_0000}, with no tracker in front.

\subsection*{Limits}

The Rust run of \code{Ddr3} is a model and not the controller. A
lowered netlist needs \code{//ddr3:hdl}, the wrapper, beside it, and
the controller's library, \code{//ddr3:ddr3\_mig\_axi}, which the build
generates under Vivado and which stays out of the tree, being AMD's.
The controller was generated without narrow bursts, so every beat is
the full thirty-two bits, with strobes for less; a byte's access is a
full-width beat at the byte's address under that byte's strobe, which
\code{wrapper\_test} checks. The path through the controller's own port
has run in simulation; its run on the board is the fourth step of
issue~\#1023.

% SPDX-License-Identifier: Apache-2.0
% covers: Entropy, RingOsc
\datasheet{Entropy}{The entropy source whole: eight ring oscillators as a foreign module, and the peripheral that watches, debiases and buffers their samples.}

\dsfacts{\code{lib/parts/src/trng.rs}}{\code{txhdl\_parts::trng::Entropy}}%
{\code{//docs:examples}, \code{//docs:vreteno}}

\subsection*{Function}

\code{Entropy} is a unit of units with two children. \code{ring} is
\code{RingOsc}, the physics: eight ring oscillators, each a loop of an
odd number of inverters through a gate that holds it still while
\code{en} is low, each sampled by a flip-flop on the clock. A ring
runs at a rate set by its own gates, wandering with temperature and
supply, and where the clock's edge catches it in its period is the
jitter that is the randomness. \code{trng} is \code{Trng}, which folds
the eight samples, watches, debiases and buffers them. The samples go
one way between the two and the run bit the other, so that a design
holds one unit and its netlist holds the module inside.

\code{RingOsc} is a foreign unit. Its netlist is an instance of the
module \code{ring\_osc} in \code{lib/parts/hdl/ring\_osc.v}, which
the lowering names and does not write, with the attributes that keep
Vivado from folding the loops or merging the rings, and
\code{ring\_osc.xdc} beside it, which turns the loop check into a
warning and keeps the timer off the loops. Under a simulator, where a
loop of zero-delay gates is an infinite event loop, the module's own
branch without \code{SYNTHESIS} defined replaces each ring with a
shift register with feedback; the unit's \code{run} in a Rust run is
the same model, seeded apart per ring as the module is, and the code
says so.

On the link side the unit is an AXI-Lite peripheral whose link is one
port, \code{bus}, a \code{LitePort}: \code{bus\_aw}, \code{bus\_ar}
and \code{bus\_w} in, \code{bus\_b} and \code{bus\_r} out. It has no
other port: the rings are inside.

\subsection*{Parameters}

None. The module takes \code{N}, the rings, and \code{L}, the
inverters in a ring, and the lowering passes \code{RINGS}, eight, and
\code{RING\_LENGTH}, seven, from \code{trng.rs}.

\subsection*{Register map}

The four words of \code{Trng}: \code{data}, \code{status},
\code{ctrl} and \code{raw}, at \code{0x0} to \code{0xc}.

\dstables{Entropy}

\subsection*{Behaviour}

\begin{itemize}
\item The rings run while \code{Trng}'s run bit is set and are held
still otherwise, their samples at zero.
\item Each ring's sample is a flip-flop that may go metastable, which
is the point of it; \code{Trng} registers the bit again before it
reads it.
\item The join runs the rings before the peripheral, so a sample the
rings drive in a step is what the peripheral folds in that step.
\end{itemize}

\subsection*{Verification}

\begin{itemize}
\item \code{//lib/parts:parts\_test} runs
\code{the\_netlist\_holds\_the\_rings\_as\_a\_module}: the Verilog
and the VHDL of \code{Entropy} instantiate \code{ring\_osc} with
\code{N} eight and do not write it.
\item \code{//cpu/vreteno:board\_test} runs
\code{the\_netlist\_holds\_the\_controller}, which also checks that
the board's netlist instantiates \code{ring\_osc} and does not write
it, and \code{the\_entropy\_source\_answers\_on\_the\_board}.
\item \code{//cpu/vreteno/board/sim:board\_test} simulates the board
with the module's simulation branch inside; manual.
\end{itemize}

\subsection*{Used by}

\code{Board}, as \code{entropy} on the sixth slot of the peripheral
page at \code{0x3500}, behind the page's \code{LiteBridge}, and Zephyr's
entropy driver \code{zephyr/drivers/entropy/entropy\_vreteno.c} through
it, which conditions the words with SHA-256 and refuses on either of
\code{Trng}'s health test faults (issues 780 and 918).
\code{ex\_trng} holds the two children apart, so that the samples are
a traced port of the lowered \code{Trng}.

\subsection*{Limits}

The Rust run and every simulation are a model of the rings, and a
model is evidence of the machinery and of nothing else. Whether the
rings are random is measured on the board and nowhere else, and
issue 458 records how far that has gone. A netlist that holds this
unit needs \code{//lib/parts:hdl}, the module, beside it, and
\code{//lib/parts:hdl\_xdc} among its constraints.

% SPDX-License-Identifier: Apache-2.0
% covers: Trng
\datasheet{Trng}{The entropy source's hardware: the rings' samples XORed to one bit, watched by two health tests, debiased and buffered as words behind AXI-Lite.}

\dsfacts{\code{lib/parts/src/trng.rs}}{\code{txhdl\_parts::trng::Trng}}%
{\code{//docs:examples}}

\subsection*{Function}

\code{Trng} takes eight samples a cycle on \code{raw}, one from each
ring oscillator of \code{RingOsc}, and makes words of entropy a host
reads. Each cycle it folds the eight to one bit by exclusive or, which
is less biased than any one of them, and does three things with the
stream. It watches it with the two continuous health tests of NIST SP
800-90B. The repetition count test, a count of identical bits in a
row, raises a sticky fault at \code{RCT\_CUTOFF}, forty. The adaptive
proportion test counts, in each window of \code{APT\_WINDOW}, 1024,
how many bits equal the window's first, and raises a sticky fault of
its own at \code{APT\_CUTOFF}, 793 (issue 918). Either fault stops the
buffer filling. It debiases
it: von Neumann's extractor takes the bits in pairs, \code{01} a zero,
\code{10} a one, \code{00} and \code{11} nothing, so a bias in the
source cancels at the cost of three bits in four. And it buffers it:
the bits that survive fill a word from the top, and whole words wait
in a buffer of \code{WORDS}, four, which a host takes one at a time.
A word that finds the buffer full is dropped.
The words are not handed out as they are: Zephyr's driver conditions
them with SHA-256, 32 words to each 32 bytes (issue 780), because the
extractor removes bias and not dependence.

Beside the words it keeps a capture for measuring the source (issue
817). A write of \code{cap} fills a buffer of \code{CAP\_WORDS}, 256
words, from the samples as they come, one word every 32 cycles
while the source runs, so the capture is 8192 samples in a row with
none missing. A host reads it out afterwards at any pace, a word at a
time through \code{capidx} and \code{capword}. \code{raw} alone cannot
do this: it holds the last 32 samples, and a program on the board
cannot read it again within 32 cycles.

It is an AXI-Lite peripheral whose link is one port, \code{bus}, a
\code{LitePort} of the five channels: \code{bus\_aw}, \code{bus\_ar}
and \code{bus\_w} in, \code{bus\_b} and \code{bus\_r} out. Its other
two ports face the rings: \code{raw} in, and \code{en} out, the run
bit, which holds the rings still while it is low.

\subsection*{Parameters}

None. \code{RINGS}, eight, \code{RCT\_CUTOFF}, forty,
\code{APT\_WINDOW}, 1024, \code{APT\_CUTOFF}, 793, and \code{WORDS},
four, are constants of \code{trng.rs}. The repetition count cutoff is
\code{1 + 20 / H} for a false alarm rate of one in a million at a
min-entropy \code{H} of half a bit per sample, rounded to forty. The
proportion cutoff is \code{1 + CRITBINOM(1024, 2\textsuperscript{-H},
1 - 2\textsuperscript{-20})} at the same \code{H}, section 4.4.2's
formula; the window is the size the standard sets for a binary source.

\subsection*{Register map}

Offsets from the peripheral's base. The port selects the word by bits
4 to 2 of the address.

\dsregs{Trng}

A read of \code{data} takes the oldest word, and reads zero when none
is waiting. \code{ctrl}'s \code{clear} is written only: a one clears
both faults and starts the proportion test's window again, and it
reads back zero. Writes to \code{data},
\code{status}, \code{raw} and \code{capword} are ignored. \code{raw}
is the samples before the extractor. \code{cap}'s \code{start} is
written only, and its \code{busy} and \code{done} are read only.

\dstables{Trng}

\subsection*{Behaviour}

\begin{itemize}
\item With the run bit clear nothing is sampled, \code{en} is low and
the buffer keeps what it holds, so a host stops the source and reads
out what is waiting.
\item Both tests run on every sample while the run bit is set, before
the extractor. A repetition count trip in the cycle of a clear stays
tripped; a proportion trip in that cycle is dropped, since the clear
starts a new window. While either fault is set no pair is taken, so no
word is made.
\item A read of \code{data} takes the oldest word when one is waiting
and answers zero otherwise; a word arriving in the cycle of a read
that empties the buffer is kept.
\item A read is answered on \code{bus\_r} in the cycle its address
beat is taken, with \code{OKAY}. A write needs its address beat, its
data beat and room on \code{bus\_b} in the same cycle, and is answered
\code{OKAY} in that cycle.
\item \code{raw} shifts in one sample a cycle while the run
bit is set, newest lowest, and is what a program on the board reads
in a loop to measure the source.
\item A capture fills only while the run bit is set: each sample goes
into the word being filled, and the word is written whole when it
holds 32, the same 32 \code{raw} reads then, so the words are samples
in a row with none missing and none twice. \code{busy} is set while it
fills and \code{done} when all 256 words are whole. The extractor and
the health tests run on regardless.
\end{itemize}

\subsection*{Verification}

\begin{itemize}
\item \code{//lib/parts:parts\_test} drives it over AXI-Lite from the
model rings and from stated streams: nothing comes before the run bit,
eight words come and differ, the buffer holds four and drains oldest
first, an alternating stream gives all zeros or all ones by the order
of its pairs, a stuck stream trips the fault after the cutoff and a
clear takes while a stream just under the cutoff does not; a stuck
stream trips both tests, a stream of ones four cycles in five with no run
past four trips the proportion test alone and stops the buffer, three
in four stays under its cutoff, and the model rings trip neither.
\item \code{//docs:trng\_apt\_sim\_trng\_apt\_tb\_test} and
\code{//docs:trng\_apt\_vsim\_test} simulate the netlist against the
run of \code{ex\_trng\_apt}, where a biased source trips the
proportion test and a clear clears it, so the test is checked in the
netlist through the trip.
\item \code{//docs:trng\_sim\_trng\_tb\_test} and
\code{//docs:trng\_vsim\_test} simulate the netlist against the run of
\code{ex\_trng} under nvc and Verilator, with the model rings' samples
as the run recorded them. The run takes a capture as well, reads it
out and finds it whole among the model's samples, in a row.
\item \code{//cpu/vreteno:board\_test} runs
\code{the\_entropy\_source\_answers\_on\_the\_board}: a program on the
core turns the source on, takes four words and says they differed,
then joins overlapping windows of \code{raw} into a run that must be
found whole among the model's samples (issue 805).
\end{itemize}

\subsection*{Used by}

\code{Entropy}, as \code{trng}, and through it \code{Board} at
\code{0x3500}, where Zephyr's driver
(\code{zephyr/drivers/entropy/entropy\_vreteno.c}) reads it and refuses
on either fault. \code{ex\_trng} and \code{ex\_trng\_apt}.

\subsection*{Limits}

Nothing here is evidence of randomness: every run is against a model
of the rings, and what the tests prove is the folding, the test, the
extractor and the buffer. The source's bias and entropy are measured
on the board, from \code{raw}, the capture and \code{data}, and
issues 458 and 780 record how far that has gone. Nothing here is
certified to SP 800-90B or FIPS 140: the two tests are the standard's,
at its cutoffs for half a bit a sample, and the conditioning is in
software, in the driver.

% SPDX-License-Identifier: Apache-2.0
% covers: Watch
\datasheet{Watch}{The three wires two cores drive, compared every cycle, and every disagreement gathered into one line.}

\dsfacts{\code{cpu/vreteno/src/pair.rs}}{\code{vreteno32::pair::Watch}}%
{\code{//docs:vreteno}}

\subsection*{Function}

A channel has a handshake, so two copies of it can be compared when a
word moves, which is what \code{Check} does. A wire has none: it holds
a value in every cycle and nothing says when to look. \code{Watch}
therefore compares the three wires a \code{Vreteno} core drives,
\code{halt}, \code{instr} and its \code{Writeback}, in every cycle, and
passes the first core's on.

It is also where the lines meet. A unit of units has no logic of its
own, so the three \code{Check} lines are brought here and gathered with
the wire comparison into the one line a design outside reads.

\subsection*{Parameters}

None. The ports are the core's own types, so the unit is written for
the core rather than for any pair of things.

\dstables{Watch}

\subsection*{Behaviour}

\begin{itemize}
\item Every cycle it compares the two halts, the two instruction words
and the two writebacks, and takes the three check lines as they are
this cycle.
\item \code{differed} is set in the first such cycle and is cleared
only by \code{rst}. The \code{differs} port is that register or this
cycle's disagreement, so the line is up in the cycle the cores stop
agreeing rather than the one after.
\item \code{halt}, \code{instr} and \code{wb} are the first core's,
driven combinationally, so a design around the pair sees one core and
no added cycle.
\end{itemize}

\subsection*{Verification}

\begin{itemize}
\item Through \code{Pair}, in \code{//cpu/vreteno:pair\_test}: eight random
programs with the two cores agreeing, and then one cycle of
\code{fault}, after which the line is up and stays up to the end of the
run.
\item It is never lowered on its own, so it has no netlist and no
co-simulation of its own.
\end{itemize}

\subsection*{Used by}

\code{Pair}, once.

\subsection*{Limits}

One bit of report for six comparisons, so nothing outside can tell
which of them fired. It compares what it is given and knows nothing of
the channels. Only \code{rst} clears it.

% SPDX-License-Identifier: Apache-2.0
% covers: Inject
\datasheet{Inject}{The second core's reset, held while `fault` is high: the self-test's one difference.}

\dsfacts{\code{cpu/vreteno/src/pair.rs}}{\code{vreteno32::pair::Inject}}%
{\code{//docs:vreteno}}

\subsection*{Function}

A comparator nobody has ever seen fire is a comparator nobody knows the
state of. \code{Inject} is how the pair is made to fire on purpose: it
gives the second core a reset of its own, the pair's reset or
\code{fault}, so that raising \code{fault} for one cycle holds that
core in reset a cycle longer than the first.

A core a cycle behind disagrees on everything it says afterwards, so
that one line runs through every part of the comparison path: the three
checks, the wire comparison and the gathering.

\subsection*{Parameters}

None.

\dstables{Inject}

\subsection*{Behaviour}

\begin{itemize}
\item \code{rst2} is \code{rst} or \code{fault}, combinational, with
nothing held between cycles. The second core is in reset for exactly
the cycles \code{fault} is high.
\end{itemize}

\subsection*{Verification}

\begin{itemize}
\item \code{a\_core\_held\_in\_reset\_a\_cycle\_longer\_is\_caught} and
\code{the\_line\_stays\_up\_after\_a\_disagreement}, in
\code{//cpu/vreteno:pair\_test}, raise \code{fault} for the twentieth cycle
of a six hundred cycle run and find the pair's line up, and still up at
the end.
\item It is never lowered on its own, so it has no netlist and no
co-simulation of its own.
\end{itemize}

\subsection*{Used by}

\code{Pair}, once.

\subsection*{Limits}

One shape of fault, and one core. It cannot corrupt a word, flip a bit
in a channel or fault the first core, so the self-test proves the
comparison path and not the comparison's coverage of every failure a
core can have. It holds no state, so it cannot stretch or debounce
\code{fault}.

% SPDX-License-Identifier: Apache-2.0
% covers: Pair
\datasheet{Pair}{Two Vreteno cores in step, with everything they say compared and one line that says they stopped agreeing.}

\dsfacts{\code{cpu/vreteno/src/pair.rs}}{\code{vreteno32::pair::Pair}}%
{\code{//docs:vreteno}}

\subsection*{Function}

\code{Pair<IW>} is two \code{Vreteno} cores running one program, with
one line that says they disagreed.

Every input the pair is given goes to both: the wires directly, since a
wire has no handshake, and each of the three channels through a
\code{Tee}, which hands a word to both cores in one cycle or to
neither. Every output is compared: the three channels through a
\code{Check} each, and the three wires through \code{Watch}.

What leaves the pair is the first core's, so a design puts a pair where
it would have put a core and wires nothing differently. What the pair
adds is one input and one output: \code{fault}, the self-test, and
\code{differs}.

\subsection*{Parameters}

\begin{center}\footnotesize
\begin{tabular}{@{}lp{0.7\textwidth}@{}}
\toprule
Parameter & Meaning \\
\midrule
\code{IW} & The width of the AXI identifier the cores use, in bits,
passed to both cores and to the channels between them. The tables use
2, which is the only width in the tree. \\
\bottomrule
\end{tabular}
\end{center}

\dstables{Pair}

\subsection*{Behaviour}

\begin{itemize}
\item The pair holds no register of its own. Every register in it
belongs to a core, a tee, a check or the watch.
\item The three tees put one cycle between a word reaching the pair and
a core being offered it. The delay is the same for both cores, which is
the point of it, and neither core can be offered a word the other has
not had.
\item A disagreement raises \code{differs} in the cycle it happens, on
a channel or on a wire alike, and the line stays up until \code{rst}.
\item Nothing else happens. Neither core is stopped, no channel is
held, no trap is raised, and the first core's word still goes out. The
pair says that there was a disagreement; what to do about it, between a
reset, a trap and a stop, belongs to the system around it.
\item A core that stops offering on a channel stalls that channel
rather than raising the line, since a check compares only words that
move. Such a core shows up on the wires, which are compared every
cycle.
\end{itemize}

\subsection*{Verification}

\begin{itemize}
\item \code{//cpu/vreteno:pair\_test} runs the pair on the lockstep test's
machine, with a tracker, a router, the data memory, the interrupt
controller and the serial port.
\code{two\_cores\_given\_the\_same\_run\_agree} runs eight random
programs and finds the line never up.
\code{a\_core\_held\_in\_reset\_a\_cycle\_longer\_is\_caught} raises
\code{fault} for one cycle and finds it up.
\code{the\_line\_stays\_up\_after\_a\_disagreement} reads it again at the
end of six hundred cycles and finds it still up.
\item The pair is lowered, and \code{//tools/datasheet} lowers
\code{Pair<2>} for this sheet's tables, but no waveform target
co-simulates its netlist and no target synthesises it. Its two parts from \code{//lib/parts}
are simulated on their own through \code{ex\_redundant}.
\end{itemize}

\subsection*{Used by}

Nothing in the tree but its own test. It is the shape a design takes
when it must not be quietly wrong, and no design here has asked for it
yet.

\subsection*{Limits}

One bit of report, with no record of which comparison fired or when.
The only fault the design can inject is holding the second core in
reset. The two cores are two instances of one Rust type, so a fault in
the design is in both and neither the comparison nor the self-test can
see it; what a pair catches is a fault that strikes one copy. The pair
has no address map and no registers of its own, so a program cannot
read the line; a design brings it out as a wire.

% SPDX-License-Identifier: Apache-2.0
% covers: Board
\datasheet{Board}{The Vreteno machine for the Alinx AX7A200B, as one lowered unit of units.}

\dsfacts{\code{cpu/vreteno/src/board.rs}}{\code{vreteno32::board::Board}}%
{\code{//docs:vreteno}}

\subsection*{Function}

\code{Board} is everything that computes on the board, written as a
unit of units so that \code{\#[lower]} makes its netlist as one module
holding the rest. It holds seven hosts, an arbiter that puts them onto
one link, a router of eight, and the peripherals behind its eight
ports. The hosts are the core's tracker; Vivado's JTAG-to-AXI master
on the board's top (issue 241), whose pins arrive under \code{jtag\_}
and are joined by an \code{AxiPins}, sharing one host through
\code{jarb} with the bridge of the RISC-V debug transport (issue 154);
the Ethernet port's two engines, each behind an \code{AxiHost} of its
own (issue 151); the video scanout's fetch (issue 151); and the SD card
host's two engines, sharing one host through \code{sdarb} (issue 912);
and Razboj's rasteriser behind \code{rhost} (issue 985).
The arbiter is an \code{Arbiter7<32, 32, 4, 2, 5, 0>}, so the
peripheral side of the bus allocates five bits of identifier, two for a
host's own and three for the port, which is room for eight hosts
(issue 1022). The debug transport, \code{dtm}, runs on the cable's
clock behind the top's \code{BSCANE2}, and \code{dreq} and
\code{dans} transfer its accesses across to the board's clock and back.
The data memory and the timer each sit behind an \code{AxiPer}. Ten
small peripherals share the page at \code{0x3000} behind the AXI-Lite
bridge \code{LiteBridge<10, ..>}, a sixteenth of the page each: the
serial port at
\code{0x3000}, the pulse width modulator at \code{0x3100}, whatever a
board hangs on the slot at \code{0x3200}, the remote peripheral at
\code{0x3300}, the Ethernet port's registers at \code{0x3400}, the
entropy source at \code{0x3500} (issue 458), the configuration
flash's SPI master at \code{0x3600} (issue 312), the Ethernet PHY's
management interface at \code{0x3700} (issue 864), the SD card host
at \code{0x3800} (issue 153), and Razboj's doorbell at \code{0x3900}
(issue 985). The DDR3 memory takes the router's port
with no tracker in front, and holds AMD's MIG controller with its own
AXI4 port as a foreign module, which makes the design's clock from the
board's. The interrupt
controller, \code{Plic3}, and the debug module, \code{Dm}, each sit
behind a \code{LiteBridge<1, ..>} of their own; the module's lines go to
the core's halt and resume requests and its register access, and take
the core's debug mode and the word read (issue 154).

The configuration flash, the chip that holds the bitstream, is reached
twice (issue 312): through that SPI master, \code{spi}, for commands,
and through \code{flashwin}, a \code{FlashWin<FLASH\_DIV, 5>} behind
its own \code{AxiPer}, where an ordinary read of the 16~MiB from
\code{0x2000\_0000} is answered with what the flash holds.
\code{cfgflash}, a \code{CfgFlash<1>}, stands between both and the
chip: it reaches the clock pin through \code{STARTUPE2}, waits for the
end of configuration, holds the window in reset until then, and
refuses write enable from the master.

\code{run.rs} wires the same parts for a simulation. \code{Board} is
that wiring as a unit, so a board top has to supply only the clocks,
the resets and the pins. The top for the AX7A200B is
\code{cpu/vreteno/board/vreteno\_board.v}.

\subsection*{Parameters}

\begin{center}\footnotesize
\begin{tabular}{@{}lp{0.7\textwidth}@{}}
\toprule
Parameter & Meaning \\
\midrule
\code{DIV} & The serial port's cycles per bit, as \code{Uart<DIV>}.
The tables use 868, the value \code{board\_netlist board} writes for
115200 baud at 100~MHz. \\
\bottomrule
\end{tabular}
\end{center}

\code{board\_netlist} lowers \code{Board<868>} for the board and
\code{Board<16>} for simulation. \code{//cpu/vreteno:board\_test}
runs \code{Board<4>}.

\subsection*{Address map}

\code{BoardRouter} is \code{Router<8, BoardMap, 32, 32, 4, 5>}, and
\code{BoardMap} beside it in \code{board.rs} states the map. A
peripheral gets an address when the address under the mask equals the
base, the first range that matches deciding.

\begin{center}\footnotesize
\begin{tabular}{@{}lllp{0.35\textwidth}@{}}
\toprule
Base & Mask & Child & Behind \\
\midrule
\code{0x0000\_0000} & \code{0xffff\_f000} & \code{rom}, \code{Rom<5>} &
\code{prom}, \code{AxiPer}: the boot memory, readable and not writable \\
\code{0x1000} & \code{0xffff\_f000} & \code{dmem}, \code{Dmem<5>} &
\code{pdmem}, \code{AxiPer} \\
\code{0x0200\_0000} & \code{0xffff\_0000} & \code{timer},
\code{Timer<5>} & \code{ptimer}, \code{AxiPer} \\
\code{0x3000} & \code{0xffff\_f000} & \code{uart}, \code{Uart<DIV>} &
\code{puart}, \code{LiteBridge<10, ..>} at \code{0x3000} \\
& \code{0xffff\_ff00} & \code{pwm}, \code{Pwm} &
the same bridge, at \code{0x3100} \\
& \code{0xffff\_ff00} & the board's, on its own clock &
the same bridge, at \code{0x3200} \\
& \code{0xffff\_ff00} & \code{remote},
\code{Remote<REMOTE\_WAIT>} & the same bridge, at \code{0x3300} \\
& \code{0xffff\_ff00} & \code{eth}, \code{EthSlots} &
the same bridge, at \code{0x3400} \\
& \code{0xffff\_ff00} & \code{entropy}, \code{Entropy} &
the same bridge, at \code{0x3500} \\
& \code{0xffff\_ff00} & \code{spi}, \code{Spi} &
the same bridge, at \code{0x3600} \\
& \code{0xffff\_ff00} & \code{mdio}, \code{Mdio} &
the same bridge, at \code{0x3700} \\
& \code{0xffff\_ff00} & \code{sd}, \code{Sd} &
the same bridge, at \code{0x3800} \\
& \code{0xffff\_ff00} & \code{doorbell}, \code{Doorbell} &
the same bridge, at \code{0x3900} \\
\code{0x4000\_0000} & \code{0xc000\_0000} & \code{ddr3},
\code{Ddr3Per} & nothing: the router's port goes to it \\
\code{0x0c00\_0000} & \code{0xfc00\_0000} & \code{plic},
\code{Plic3<0>} & \code{pplic}, \code{LiteBridge<1, ..>} at
\code{0x0c00\_0000} \\
\code{0x1000\_0000} & \code{0xffff\_0000} & \code{dmod}, \code{Dm} &
\code{pdm}, \code{LiteBridge<1, ..>} at \code{0x1000\_0000}; not nearer the
interrupt controller, whose 64~MiB window the router would match first \\
\code{0x2000\_0000} & \code{0xff00\_0000} & \code{flashwin},
\code{FlashWin<FLASH\_DIV, 5>} & \code{pflash}, \code{AxiPer}: the
configuration flash, read only \\
\bottomrule
\end{tabular}
\end{center}

The router answers a burst that matches no range with \code{DecErr}.
The boot memory is on the bus at zero as well (issue 268), readable and
not writable; the core fetches from its own copy below \code{0x1000}
and from the bus above it.

\dstables{Board}

\subsection*{Behaviour}

\begin{itemize}
\item The core's tracker is \code{AxiHost<32, 32, 4, 2, 4>}: 32-bit
addresses and words, four lanes, two-bit identifiers, four of them.
\item \code{rst} goes to the core, the timer, the serial port and the
interrupt controller. The DDR3 controller has its own reset,
\code{sys\_rst}, active high, and hands back \code{ui\_rst}, high
until the clock it makes is good.
\item The interrupt controller's source 1 is the serial port's
interrupt line, its source 2 is \code{irq}, and its source 3 is the
Ethernet port's arrival; all three ask while high.
Its line is the core's external interrupt. The timer's line goes to
the core's \code{tirq}.
\item Of the core's outputs, only \code{halt} is a port. \code{instr}
and the \code{Writeback} go nowhere; \code{dbg} and \code{dbg\_rdata}
go to the debug module, whose requests and register access come back
to the core.
\item The join runs the timer before the core, since the core reads
the timer's line in the same step, and the serial port before the
interrupt controller before the core, for the same reason. \code{Ddr3Per} runs its controller
before its pins part for the same reason.
\item \code{calib} and the memory's pins are \code{Ddr3Per}'s,
brought out as the unit's own ports.
\item The slot at \code{0x3200} does not reach a field. It leaves as
five channel ports, \code{vaw}, \code{var} and \code{vw} out and
\code{vb} and \code{vr} back, because what sits there runs on a clock
of its own and the crossing is the board top's business. A board with
nothing there ties them off, and a program that touched \code{0x3200}
on such a board would wait for ever.
\item The remote peripheral at \code{0x3300} is the other way about:
\code{remote} and \code{link} are fields, because both run on the
core's clock. What leaves the unit is the frame stream, \code{net\_tx}
out and \code{net\_rx} in, a byte and a last bit each, which is what
the MAC speaks, and which the link shares with the Ethernet port
through \code{share}. \code{REMOTE\_DEV} is 1 and \code{REMOTE\_WAIT} is
100 million cycles, one second at 100~MHz (issue 397), after which an
unanswered transaction is \code{SLVERR} on the bus.
\item \textbf{The Ethernet port moves frames itself.} \code{eth} is
the register map; behind it the fetch engine reads a frame out of its
transmit slot and \code{FrameOut} turns the words into bytes, and on
the way in \code{FrameLen} finds the frame's length again,
\code{FrameIn} packs its bytes into words and the store engine writes
them into a receive slot. The slots are in the DDR3 memory from
\code{0x4100\_0000}. Frames the driver has not acknowledged wait in
order, at most one a slot, and while both slots hold one \code{eth}
tells \code{fin}, which drops the next frame and counts it for
\code{rx\_errors} (issue 1313). \code{NoBeats} and \code{NoReads} hold the
direction of each engine's host that the engine never uses.
\item \textbf{The wire is shared by EtherType.} \code{share}, an
\code{EthShare<0x88b5>}, sends the remote peripheral's frames to its
link and every other frame to the Ethernet port, and merges the two
users' frames onto \code{net\_tx} a frame at a time.
\item The join runs the receive side, \code{share}, \code{flen} and
\code{fin}, before the register block, because the block reads the
slot and length they drive in the same step; and the store engine and
the sending side after it, because they read what it drives. With the
receive side after the block, the block read last step's slot, and
the board test caught it as a frame of the right length whose every
byte read back zero. \code{fin} reads the block's \code{rx\_full} a
step late in the Rust run for the same reason; that can only drop a
frame that starts in the cycle after an acknowledgement freed a slot,
which the netlist would store, and never lands one on a slot still
held.
\item \textbf{The scanout fetches its own lines.} \code{scan}, a
\code{ScanFetch<32, 16, 640>}, takes each line request that arrives on
\code{scan\_req} from the pixel clock and starts \code{vfetch}, a
\code{LineFetch} behind \code{vhost}, whose words leave on
\code{scan\_words} for the line pair. The crossings are the board
top's, as the video slot's are. \code{vnobeats} holds the write side
the fetch never uses.
\item \textbf{The SD card host moves blocks through memory.}
\code{sdstore} and \code{sdfetch}, each behind a host of its own with
one bit of identifier, share the arbiter's sixth port through
\code{sdarb}, an \code{Arbiter2<32, 32, 4, 1, 2, 0>}. One command is a
read or a write and never both, and the host starts no command until
the last one's engine has finished, so the two never offer at once
and the nest costs neither a turn; the host checks that their busy
lines are never high together. \code{sdnoreads} and \code{sdnobeats}
hold the directions they never use.
\item \textbf{Razboj draws into the frame the scanout shows.}
\code{raster}, a \code{Raster<32, 2, 10, 480, 0x4200\_0000,
0x4280\_0000, 0x3900>} behind \code{rhost}, takes the arbiter's seventh
port. It reads its display list at \code{0x4280\_0000} and writes
rows of 1024 pixels from \code{0x4200\_0000}, the scanout's stride. Its
count is \code{doorbell}'s register on the tenth slot, which drives the
rasteriser's \code{ring} while the count is not zero, so an idle
rasteriser reads nothing; the join runs the rasteriser before the
doorbell, which reads its \code{idle} line.
\item The join runs the SPI master and the window before
\code{cfgflash}, since the pins read the lines both drive in the same
step. The flash's data out, \code{fl\_miso}, goes to both directly. The
flash's select and data in leave as \code{fl\_cs\_n} and
\code{fl\_mosi}; \code{fl\_cclk} is the startup block's model of the
clock pin, which a board leaves open, and \code{fl\_refused} says the
pins refused write enable. The master's interrupt is not wired.
\item \code{board\_netlist} puts a program in the core's
\code{imem} and in the data memory's \code{lane0} to \code{lane3} of
the netlist, since nothing on the machine loads them at run time.
\end{itemize}

The top, \code{vreteno\_board.v}, does the following around the
lowered \code{board}:

\begin{itemize}
\item It hands the board's 200~MHz differential clock to the memory
controller on \code{sys\_clk}, which makes the design's 100~MHz,
\code{ui\_clk}, from it; the board has no clock generator of its own.
\item The controller's reset, \code{sys\_rst}, is a pulse at power-on
counted on the board's clock and nothing else. The core's reset
follows the button, \code{ui\_rst} being low and \code{calib}, through
two flip-flops on the design's clock.
\item It ties \code{irq} low.
\item It puts the flash's select and data on the configuration pins
FCS\_B, D00 and D01, and holds D02 and D03, the chip's write protect
and hold, high. The clock, CCLK, is reached through \code{STARTUPE2}
inside the netlist.
\item Its four LEDs, lit when driven low, show the modulator's first
channel or the core halted, the serial line driven at least once, the
memory calibrated, and a heartbeat from bit 25 of a counter, which
says the controller's clock is good.
\end{itemize}

\subsection*{Verification}

\begin{itemize}
\item \code{//cpu/vreteno:board\_test} runs \code{Board<4>} in
Rust. \code{the\_greeting\_runs\_on\_the\_board} must print
\code{hello from rust} and halt within 8000 cycles.
\code{the\_memory\_test\_runs\_on\_the\_board} must print
\code{ddr3 ok} and halt within 40000 cycles.
\code{input\_comes\_by\_interrupt\_one\_byte\_each} runs a program
that takes four typed bytes through the interrupt controller, one
interrupt each, and must print \code{got ping in 4 interrupts} and
halt within 20000 cycles.
\code{a\_frame\_received\_lands\_in\_memory\_and\_the\_core\_reads\_it\_back}
puts a frame of 21 bytes on \code{net\_rx}, and the core must print its
length and every byte from the slot it landed in.
\code{two\_frames\_written\_to\_memory\_leave\_the\_port\_as\_written\_in\_order}
has the core write two frames of 23 and 18 bytes and ask for both back
to back, and the bytes that leave on \code{net\_tx} must be the two
frames, whole and in order.
\code{the\_entropy\_source\_answers\_on\_the\_board} must print
\code{trng ok}, place every pair of raw windows and find the run it
joins whole among the model's samples, within 120000 cycles.
\code{the\_configuration\_flash\_answers\_on\_the\_board} gives the
board a flash model holding a bitstream's first bytes; the core must
print the chip's identity, a clear write enable latch and where the
sync word is, and the chip must have taken no write enable.
The same target also runs the boot memory, the remote peripheral, the
loader and the debug module through the board.
\code{the\_netlist\_holds\_the\_controller} checks that the Verilog
has the top, the core and the DDR3's pins part as modules,
instantiates \code{ddr3\_axi32} and \code{ring\_osc} without writing
either, and has
the \code{dq} pads.
\item \code{//cpu/vreteno/board/sim:board\_test} simulates the top, the
netlist of \code{//cpu/vreteno:board\_sim\_v}, the controller and two
Micron models under Vivado's simulator. It is manual.
\code{//docs:vreteno} reports that in that run the core says
\code{ddr3 ok} on the serial port and halts.
\item \code{//cpu/vreteno:vreteno\_board\_synth} and
\code{//cpu/vreteno:vreteno\_board\_pnr} synthesise, place and route
the design with the board's pins. Both are manual.
\code{//docs:vreteno} reports that every timing constraint is met, with
0.307~ns worst setup slack and 0.033~ns worst hold slack, and 12\,296
LUTs, 11\,209 flip-flops, five block RAM tiles and four DSP blocks, for
the board at eed5b7f, with the configuration flash and the entropy
source's capture buffer; it has not been measured since.
\item \code{//cpu/vreteno:vreteno\_board\_prog} and
\code{//cpu/vreteno:vreteno\_board\_flash} program the FPGA and the
QSPI flash. Both are manual.
\code{//cpu/vreteno:vreteno\_board\_fits\_test}, also manual, checks
that the routed bitstream ends before \code{0x00A0\_0000}, where
programs begin in the flash.
\end{itemize}

No waveform target simulates \code{Board}'s netlist against a trace.

\subsection*{Used by}

\code{//cpu/vreteno:board\_netlist}, whose genrules
\code{//cpu/vreteno:board\_v} and \code{//cpu/vreteno:board\_sim\_v}
write the netlist, and \code{vreteno\_board.v}, which instantiates it
as \code{board}. The test \code{cpu/vreteno/tests/board.rs}.

\subsection*{Limits}

\code{//docs:vreteno} reports the design with the DDR3 memory and MIG
run on the board on September 28, 2026, written into the part over
JTAG, and saying \code{ddr3 ok}; the run from the flash after a power
cycle is owed (issue 188). It reports the configuration flash read on
the board as well: the identity, write enable refused and the
bitstream through the window. The program is fixed in the netlist. The board top ties \code{irq} low, so on the
board the interrupt controller's second source never asks.

The arbiter's ports are the core on 0, \code{jarb} on 1, the Ethernet
port's fetch engine on 2 and its store engine on 3, the scanout's fetch
on 4, \code{sdarb} on 5 and Razboj's rasteriser on 6. It takes turns,
so with every host offering each wins one turn in seven. Three bits of
port number leave room for an eighth host, which raises the count and
widens nothing past the arbiter.

% PCIe.
% SPDX-License-Identifier: Apache-2.0
% covers: BarRegs
\datasheet{BarRegs}{Four 64-bit registers behind a peripheral tracker: what a host reaches in the PCIe endpoint's BAR1.}

\dsfacts{\code{pcie/src/bar.rs}}{\code{pcie::bar::BarRegs}}%
{\code{//docs:pcie}}

\subsection*{Function}

\code{BarRegs} is a peripheral client of an \code{AxiPer}, and its link
is one port, \code{bus}, a \code{PerPort} of the four channels: requests
come on \code{bus\_req} and write beats on \code{bus\_w}, and answers go out on
\code{bus\_ans} and \code{bus\_r}. It holds four words a host reads and writes,
and drives two LEDs from one of them. The input \code{rst} clears the
words that can be written and the count.

\subsection*{Parameters}

\begin{center}\footnotesize
\begin{tabular}{@{}lp{0.7\textwidth}@{}}
\toprule
Parameter & Meaning \\
\midrule
\code{I} & The width of the AXI identifier on \code{bus\_req}, \code{bus\_ans}
and \code{bus\_r}, in bits. The tables use 4, as \code{PcieBar} does. \\
\bottomrule
\end{tabular}
\end{center}

\subsection*{Register map}

Offsets from BAR1.

\begin{center}\footnotesize
\begin{tabular}{@{}llp{0.55\textwidth}@{}}
\toprule
Offset & Word & Access \\
\midrule
\code{0x00} & Identifier, \code{IDENT}, \code{0x5478\_4844\_4c00\_0001} &
Read only. \\
\code{0x08} & \code{scratch} & Read and write. \\
\code{0x10} & \code{lights}, bits 1 and 0 & Read and write; the other
bits read zero. \\
\code{0x18} & \code{writes} & Read only. \\
\bottomrule
\end{tabular}
\end{center}

It selects the word by bits 4 and 3 of the address and looks at no
other bits.

\dstables{BarRegs}

\subsection*{Behaviour}

\begin{itemize}
\item A read is taken when \code{bus\_r} has room and no write is waiting
for its beat. It is answered on \code{bus\_r} in the cycle it is taken,
with the word addressed, \code{OKAY} and \code{last} set.
\item A write is taken when no write is waiting. The word it names goes
into \code{psel} and its identifier into \code{pid}.
\item The held write's beat is taken when it has arrived and \code{bus\_ans}
has room. \code{writes} counts it. A write to \code{0x08} replaces
\code{scratch} with the beat's data, and one to \code{0x10} replaces
\code{lights} with its low two bits, whatever the strobe. A write to
\code{0x00} or \code{0x18} changes nothing else. \code{bus\_ans} sends
\code{OKAY}.
\item \code{leds} is \code{lights}, a register's output.
\item \code{rst} clears \code{pend}, \code{scratch}, \code{lights} and
\code{writes}.
\end{itemize}

\subsection*{Verification}

\code{a\_host\_reads\_and\_writes\_the\_bar}, in
\code{//pcie:bar\_test}, drives it inside \code{PcieBar} through the
endpoint's pins: the identifier, the scratch word read back, an LED,
a write to the identifier that changes nothing, and the count of writes.

\subsection*{Used by}

\code{PcieBar}, as its field \code{regs}.

\subsection*{Limits}

It serves bursts of one beat and ignores a write's strobe. One write is
held at a time, and no read is taken while it waits.

% SPDX-License-Identifier: Apache-2.0
% covers: PcieBar
\datasheet{PcieBar}{Everything behind the PCIe endpoint's BAR1, as one lowered unit of units.}

\dsfacts{\code{pcie/src/bar.rs}}{\code{pcie::bar::PcieBar}}%
{\code{//docs:pcie}}

\subsection*{Function}

\code{PcieBar} holds \code{AxiPins<32, 64, 8, 4>}, an
\code{AxiPer<32, 64, 8, 4>} and \code{BarRegs<4>}, joined by channels.
Its inputs, the struct \code{PcieBarIn}, are a reset and the 24 pins an
AXI4 host drives. Its outputs, \code{PcieBarOut}, are the eleven pins
the host reads and the LEDs' two bits. The board's top,
\code{pcie/board/pcie\_top.v}, wires the XDMA endpoint's master to it
pin for pin, the low 32 bits of the master's 64-bit addresses, and its
clock and reset to the endpoint's user clock and reset.

\subsection*{Parameters}

None. The widths are the endpoint's: 32-bit addresses on this side,
64-bit words, eight lanes, four-bit identifiers.

\subsection*{Address map}

The tracker sends every address to \code{BarRegs}, which decodes bits 4
and 3; see its sheet.

\dstables{PcieBar}

\subsection*{Behaviour}

\begin{itemize}
\item The pins' address phases and write beats go through
\code{AxiPins} onto the tracker's channels, and its responses and read
beats come back to the pins the same way.
\item The tracker hands each request to \code{BarRegs}.
\item \code{rst} goes to \code{BarRegs} only.
\end{itemize}

\subsection*{Verification}

\begin{itemize}
\item \code{a\_host\_reads\_and\_writes\_the\_bar}, in
\code{//pcie:bar\_test}, drives it through the pins with
\code{PinHost}.
\item \code{the\_netlist\_has\_the\_endpoint\_pins} checks that its
Verilog is one module, \code{pcie\_bar}, with the pins as ports.
\item \code{//pcie:endpoint\_synth} and \code{//pcie:endpoint\_pnr}
put it and the endpoint through Vivado with the board's pins. Both are
manual. \code{//docs:pcie} reports their results.
\end{itemize}

No waveform target simulates its netlist against a trace.

\subsection*{Used by}

\code{//pcie:bar\_netlist}, whose genrule \code{//pcie:bar\_v} writes
the netlist, and \code{pcie\_top.v}, which instantiates it as
\code{bar}.

\subsection*{Limits}

It has not run in a PC (issue 180). It is its own bus on the endpoint's
clock, and reaches nothing on Vreteno's.

% Razboj, the GPU.
% SPDX-License-Identifier: Apache-2.0
% covers: Fb
\datasheet{Fb}{Razboj's framebuffer: a memory of one word per pixel behind an AXI peripheral end.}

\dsfacts{\code{gpu/razboj/src/fb.rs}}{\code{razboj::fb::Fb}}%
{\code{//docs:razboj}, \code{//docs:axi}}

\subsection*{Function}

\code{Fb} is \code{N} words of memory, \code{px}, behind the
peripheral end of an AXI link. A pixel is one word. The rasteriser
writes pixels into it, and in Razboj's runs the same memory also
holds the display list and its count, which the rasteriser reads.

It is the peripheral client of an \code{AxiPer}, and its link is one
port, \code{bus}, a \code{PerPort} of the four channels: requests
come on \code{bus\_req} and write beats on \code{bus\_w}, and answers
go out on \code{bus\_ans} and \code{bus\_r}. It answers write
bursts and read bursts of any length. The rasteriser writes a run of a
row's pixels as one write burst (issue 987) and fetches a display list
entry as one read burst. A read's first beat is answered from the
memory in the cycle the request is taken and the rest follow a beat a
cycle, one burst at a time. A write is held while its beats arrive,
since AXI4 puts no identifier on the write data channel, and is
answered once, with its last beat. A beat writes only the bytes its
strobes name, so a beat with none writes nothing. It counts what it
served: \code{rbursts} read bursts taken, \code{rbeats} read beats
sent, \code{wbursts} write bursts taken and \code{wbeats} write beats
that wrote something, which a run reads to say how much of its time
the link was busy and how many pixels it wrote.

\subsection*{Parameters}

\begin{center}\footnotesize
\begin{tabular}{@{}lp{0.7\textwidth}@{}}
\toprule
Parameter & Meaning \\
\midrule
\code{A} & The width of the link's address, in bits. \\
\code{I} & The width of the AXI identifier, in bits. \\
\code{N} & Words of memory, a power of two. \\
\bottomrule
\end{tabular}
\end{center}

The tables use \code{Fb<16, 2, 1024>}, the framebuffer of
\code{//gpu/razboj:trace}.

\subsection*{Address map}

The word a burst names is its address shifted down by two, cut to
sixteen bits, and masked with \code{N}~$-$~1. The memory itself gives
no word a meaning. In \code{//gpu/razboj:trace} and the tests of
\code{gpu/razboj/src/sim.rs}, with \code{N} of 1024:

\begin{center}\footnotesize
\begin{tabular}{@{}lp{0.6\textwidth}@{}}
\toprule
Byte address & What \code{sim.rs} puts there \\
\midrule
\code{0x000} & The framebuffer, sixteen by sixteen pixels, a word
each. \\
\code{0x400} & The display list, eight words an instruction. \\
\code{0x600} & The count of instructions, written last. \\
\bottomrule
\end{tabular}
\end{center}

\dstables{Fb}

\subsection*{Behaviour}

\begin{itemize}
\item A read is taken when \code{bus\_r} has room and no write is waiting
for its beat. It is answered on \code{bus\_r} in the cycle it is taken,
with the word from \code{px}, \code{OKAY} and \code{last} set.
\item A write is taken when no write is held. Its first word and its
identifier go into \code{paddr} and \code{pid}, \code{pend} is set,
and \code{wbursts} counts it.
\item While a write is held, each of its beats is taken as it arrives,
and the last only when \code{bus\_ans} has room. A beat writes the
bytes its strobes name over the word at \code{paddr}, which then
moves to the next word; \code{wbeats} counts a beat with any strobe
set.
\item The last beat ends the write: \code{pend} clears and
\code{bus\_ans} sends \code{OKAY} with \code{pid}.
\item While a write is held, no request of either kind is taken.
\item An address beyond the memory wraps rather than ending the run.
\item \code{pixel} reads a word back for a run or a test.
\end{itemize}

\subsection*{Verification}

\begin{itemize}
\item \code{//gpu/razboj:razboj\_test} renders every scene on the rasteriser
through the link into \code{Fb}, reads the pixels back and compares
each with the model in \code{gpu/razboj/src/model.rs}. The tests are
\code{a\_scene\_of\_every\_kind\_agrees\_with\_the\_model},
\code{a\_clear\_says\_only\_its\_colour},
\code{a\_triangle\_is\_the\_same\_whichever\_way\_it\_is\_wound},
\code{a\_triangle\_hanging\_off}\allowbreak\code{\_the\_screen\_is\_clipped},
\code{a\_single\_pixel\_and\_an\_empty\_box} and
\code{pseudorandom\_scenes\_agree\_with\_the\_model}.
\item The build simulates \code{Fb<16, 2, 1024>}'s netlist, with the
display list and its count in \code{px}, against the trace of
\code{//gpu/razboj:trace} under nvc and Verilator:
\code{//docs:razboj\_sim\_fb\_fb\_tb\_test} and
\code{//docs:razboj\_vsim\_fb\_test}.
\item \code{//gpu/razboj:trace} and \code{//gpu/razboj:demo} assert that the
framebuffer equals the model's picture. The build runs
\code{//gpu/razboj:demo} to draw the picture in \code{//docs:razboj}.
\end{itemize}

\subsection*{Used by}

\code{razboj::sim::run} and \code{razboj::sim::run\_list}, and so Razboj's
tests, \code{//gpu/razboj:trace} and \code{//gpu/razboj:demo}. The system's
\code{//soc:demo} draws the core's display list a second time with
\code{run\_list}, on a plain link into \code{Fb}, and checks that the
picture is the same. On the lattice itself the memory is the
simulation-only \code{Ram}.

\subsection*{Limits}

It holds one write at a time and takes no read while it does, and a
read burst's beats are answered a cycle apart. A write burst's beats
are taken as they come; its length is not read, and its last beat ends
it. The
word index is sixteen bits wide, so the framebuffer reaches at most
65536 words.

% SPDX-License-Identifier: Apache-2.0
% covers: razboj::raster::Raster
\datasheet{Raster}{Razboj's rasteriser: it reads a display list over AXI and writes one pixel a cycle.}

\dsfacts{\code{gpu/razboj/src/raster.rs}}{\code{razboj::raster::Raster}}%
{\code{//docs:razboj}, \code{//docs:noc}}

\subsection*{Function}

\code{Raster} finds its work in memory and draws it there. While
\code{ring} is high it reads a count at \code{CTRL}, until the count
is not zero. It then reads the sixteen words of each instruction of
the display list at \code{DL}, as one read burst of sixteen beats, and
walks the instruction's box one pixel per cycle. It writes the
instruction's colour, its alpha in the top byte, at every pixel inside
the primitive: every pixel of a clear or a rectangle, and every pixel
of a triangle at which no edge function is negative. A run of a row's
pixels goes out as one write burst of up to sixteen beats, a beat a
pixel, and a beat whose pixel is outside the primitive has its strobes
off (issue 987). A shaded
triangle's colour is three planes the host worked out, each stepped
across the box as an edge is.

It draws list after list. When a list is drawn and every write has
been answered, it writes the count back to zero and reads it again.
On the board the count is \code{Doorbell}'s register, which drives
\code{ring} high while the count is not zero, so an idle rasteriser
puts nothing on the link (issue 985). Where nothing drives it,
\code{ring} is tied high and the count is read back to back.

It is an AXI host written as hardware. It issues bursts on
\code{issue}, sends a pixel's beat on \code{wbeat}, and hands
identifiers back on \code{release}. It receives \code{grant},
\code{done} and \code{rdata} from its tracker, and \code{ring}.
\code{idle} goes high when the list is drawn, its zero is answered and
no write is in flight.

\subsection*{Parameters}

\begin{center}\footnotesize
\begin{tabular}{@{}lp{0.7\textwidth}@{}}
\toprule
Parameter & Meaning \\
\midrule
\code{A} & The width of the link's address, in bits. \\
\code{I} & The width of the AXI identifier, in bits. \\
\code{LOGW} & The screen is \code{1 << LOGW} pixels wide. \\
\code{H} & The screen is \code{H} pixels high. \\
\code{BASE} & The byte address of the framebuffer's first word. \\
\code{DL} & The byte address of the display list's first word. \\
\code{CTRL} & The byte address of the count. \\
\bottomrule
\end{tabular}
\end{center}

The tables use \code{Raster<16, 2, 4, 16, 0, 0x400, 0x600>}, the
rasteriser of \code{//gpu/razboj:trace}: a sixteen by sixteen screen at
\code{0x0}, the list at \code{0x400} and the count at \code{0x600}. In
\code{soc} it is \code{Raster<32, 2, 7, 96, 0x8000, 0x2000, 0x4000>}.
On the flagship it is \code{Raster<32, 2, 10, 480, 0x4200\_0000,
0x4280\_0000, 0x3900>}: rows of 1024 pixels, the scanout's stride, of
which 640 are shown, and the count in \code{Doorbell}.

\subsection*{Address map}

\begin{center}\footnotesize
\begin{tabular}{@{}lp{0.6\textwidth}@{}}
\toprule
Address & What the rasteriser does there \\
\midrule
\code{CTRL} & Reads the count while \code{ring} is high, and writes
zero there when a list is drawn. \\
\code{DL + (i << 6) + (k << 2)} & Reads word \code{k}, 0 to 15, of
instruction \code{i}, as one burst. \\
\code{BASE + (((y << LOGW) + x) << 2)} & Writes pixel \code{x},
\code{y}. \\
\bottomrule
\end{tabular}
\end{center}

An instruction's words are in the format \code{gpu/razboj/src/dl.rs}
states.

\begin{center}\footnotesize
\begin{tabular}{@{}lp{0.42\textwidth}p{0.3\textwidth}@{}}
\toprule
Word & Low field & High field \\
\midrule
0 & bits 1 to 0: kind, 0 clear, 1 rectangle, 2 triangle, 3 shaded
triangle & bits 25 to 2: colour \\
1 & bits 9 to 0: \code{x0} & bits 25 to 16: \code{y0} \\
2 & bits 9 to 0: \code{x1} & bits 25 to 16: \code{y1} \\
3 & bits 15 to 0: \code{ax} & bits 31 to 16: \code{ay} \\
4 & bits 15 to 0: \code{bx} & bits 31 to 16: \code{by} \\
5 & bits 15 to 0: \code{cx} & bits 31 to 16: \code{cy} \\
6 to 14 & \multicolumn{2}{l}{a shaded triangle's red, green and blue
planes: value, step right, step down} \\
15 & bits 7 to 0: alpha & \\
\bottomrule
\end{tabular}
\end{center}

\dstables{Raster}

\subsection*{Behaviour}

The rasteriser is two processes. The first takes the link's answers
every cycle, and says whether the pixel under the walk is in the
primitive and whether the rasteriser is idle. The second is the front
end written as the sequence it is: poll the count until the list is
ready, then for each instruction fetch its sixteen words and walk its
box a pixel a turn, and when the list is drawn write the count back to
zero and poll again. The lowering
numbers the front end's waits into a state register the unit does not
declare, \code{at\_wait}, and keeps its loops' counts in \code{for0} to
\code{for3}; all are in the tables above.

\begin{itemize}
\item The front end waits for \code{ring}, then issues a read of the
count at \code{CTRL} once \code{issue} has room, and waits for its
beat. The count's low sixteen bits go into \code{left}. A count of
zero means wait for \code{ring} and read again.
\item For each of \code{left} instructions it waits a cycle, for the
instruction's index to read back, then reads its sixteen words as one
burst of sixteen beats, issued once \code{issue} has room. Each
word's fields are latched as its beat lands.
\item Once the last word has landed the walk is set up over three
cycles: the differences of the vertices, then the six products side
by side, then their sums. The products are computed there, and
nowhere else.
\item A clear's box is the whole screen, from its own type: all
\code{1 << LOGW} columns of \code{H} rows. A rectangle's and a
triangle's box is \code{x0}, \code{y0} to \code{x1}, \code{y1}, and
is not clipped to \code{H}. Vertices are sixteenths of a pixel in
sixteen bits of two's complement.
\item After a cycle for the box to read back, the walk goes over every
row of the box, a cycle at the start of each, and every column of the
row, a pixel a turn.
\item Each edge is kept as its value at this pixel, its value at the
start of this row and its two steps, in 32-bit two's complement. A
step to the right adds the column step; a new row reloads from the row
value plus the row step.
\item A pixel is sampled at its centre, and a vertex is in sixteenths
of a pixel. \code{hit} is high when the kind is not a triangle, or when
all three edge values are positive, or zero on a top or a left edge:
\code{tl0} to \code{tl2} say which edges are, from the signs of
their steps, set at the setup. So of two triangles sharing an edge,
exactly one draws a pixel whose centre lies on it.
\item A pixel with \code{hit} high and no burst under way starts one,
when both \code{issue} and \code{wbeat} have room, which adds one to
\code{issued}. The burst is size 2, \code{INCR}, and its length is the
pixels left in the box's row, the words left before the next 4\,KB
page, or sixteen, whichever is fewest, so no burst crosses a page.
\code{beats} holds the beats it still owes.
\item While a burst is under way, every pixel the walk steps over sends
the next beat when \code{wbeat} has room, whatever the pixel: the
alpha above the 24-bit colour, strobe \code{0xf} where \code{hit} is
high and \code{0x0} where it is not, and \code{last} on the burst's
last beat. A triangle's pixels in a row are one run, so a burst wastes
beats only past the run's end.
\item The walk steps when the pixel needs no write or its beat went
out, so backpressure holds the walk.
\item When the last instruction's box is walked and every write is
answered, the count at \code{CTRL} is written back to zero, and
\code{finished} is set as that write goes out. \code{inflight} is a wire, \code{issued} less \code{answered},
the writes whose response has not come back, and \code{idle} is high
when \code{finished} is set and nothing is in flight.
\item When a write response and a read beat arrive together, the
write's identifier goes back first and the read's waits a cycle.
Grants are taken and dropped.
\end{itemize}

\subsection*{Verification}

\begin{itemize}
\item \code{//gpu/razboj:razboj\_test} renders scenes on the rasteriser through
an \code{AxiHost}, an \code{AxiPer} and \code{Fb}, and compares every
pixel with the model of \code{gpu/razboj/src/model.rs}:
\code{a\_scene\_of\_every\_kind\_agrees\_with\_the\_model},
\code{a\_clear\_says\_only\_its\_colour},
\code{a\_triangle\_is\_the\_same\_whichever\_way\_it\_is\_wound},
\code{a\_triangle\_hanging\_off}\allowbreak\code{\_the\_screen\_is\_clipped},
\code{a\_single\_pixel\_and\_an\_empty\_box}, and
\code{pseudorandom\_scenes\_agree\_with\_the\_model}, twelve scenes of
up to four random rectangles and triangles.
\item \code{an\_instruction\_survives\_the\_format} in the same target
checks the display list's encoder against its decoder.
\item The build simulates the netlist of \code{Raster<16, 2, 4, 16, 0,
0x400, 0x600>} against the trace of \code{//gpu/razboj:trace} under nvc and
Verilator: \code{//docs:razboj\_sim\_raster\_raster\_tb\_test} and
\code{//docs:razboj\_vsim\_raster\_test}.
\item \code{//soc:demo\_test} runs
\code{a\_core\_draws\_a\_solid\_and\_a\_gpu\_fills\_it}. Vreteno writes
a display list on the lattice, the rasteriser draws it, and every pixel
must match the model's picture of that list.
\item It goes through Vivado out of context with the board's
parameters, \code{//gpu/razboj:razboj\_synth} and
\code{//gpu/razboj:razboj\_pnr}, and meets the board bus's 100~MHz
since its setup was split into three cycles (issue 1034); the size and
the slack of the last run are in \code{//docs:razboj}'s section on
Vivado.
\end{itemize}

\subsection*{Used by}

\code{razboj::sim::run} and \code{razboj::sim::run\_list}, and so Razboj's
tests, \code{//gpu/razboj:trace} and \code{//gpu/razboj:demo}. The system in
\code{soc}, where the rasteriser is the host at corner 1,0 of the
lattice.

\subsection*{Limits}

\code{//docs:razboj} lists them. It writes one pixel per cycle at most. It
has no depth buffer, so primitives are drawn in list order. It has no
interpolation, so a triangle is one colour, and no texture. It clips
no geometry, only boxes, and the assembler \code{Op::encode} does that
on the host. It has one read out at a time.

The code adds these. It has no reset port. A list holds at most 65535
entries, since \code{left} is sixteen bits. \code{BASE}, \code{DL}
and \code{CTRL} are constants of its type, so a design that draws
into two buffers moves the second one's rows past the first's, as the
flagship's icosahedron does (issue 986), rather than moving
\code{BASE}. A colour is 24 bits and its alpha 8.

% SPDX-License-Identifier: Apache-2.0
% covers: Doorbell
\datasheet{Doorbell}{Razboj's doorbell: the register a program writes a display list's count into, and the line that wakes the rasteriser.}

\dsfacts{\code{gpu/razboj/src/doorbell.rs}}{\code{razboj::doorbell::Doorbell}}%
{\code{//docs:vreteno}}

\subsection*{Function}

\code{Doorbell} holds the count \code{Raster} reads at its
\code{CTRL}: how many entries of the display list to draw, zero when
there is nothing to draw. A program writes the list into memory, then
the count; the rasteriser draws the list and writes the count back to
zero. The count is a register rather than a word of memory, so that
it holds zero from reset, and no word a loader left behind can start
a list nobody wrote (issue 985).

It drives \code{ring} high while the count is not zero, and the
rasteriser reads the count only then, so an idle rasteriser puts no
reads on the bus. A second word reads the rasteriser's \code{idle}
line.

It is an AXI-Lite peripheral whose link is one port, \code{bus}, a
\code{LitePort} of the five channels: \code{bus\_aw}, \code{bus\_ar}
and \code{bus\_w} in, \code{bus\_b} and \code{bus\_r} out. Its other
ports are \code{rst} and \code{idle} in, and \code{ring} out.

\subsection*{Parameters}

None.

\subsection*{Register map}

Offsets from the peripheral's base, \code{0x3900} on the flagship, the
peripheral page's tenth slot. The port selects the word by bit 2 of
the address.

\dsregs{Doorbell}

\dstables{Doorbell}

\subsection*{Behaviour}

\begin{itemize}
\item A read is answered on \code{bus\_r} in the cycle its address
beat is taken, with \code{OKAY}. A write needs its address beat, its
data beat and room on \code{bus\_b} in the same cycle, and is answered
at once, with \code{OKAY}.
\item A write of \code{count} sets it to the word's low sixteen bits.
A write of \code{status} is ignored.
\item \code{rst} high sets the count to zero, over a write in the same
cycle.
\item \code{ring} is high while the count is not zero.
\end{itemize}

A program waits for \code{count} to read zero, writes the list, then
writes \code{count}. A write while a list is being drawn starts
nothing: the rasteriser has read the count already, and writes zero
over it at the end.

\subsection*{Verification}

\begin{itemize}
\item \code{//gpu/razboj:razboj\_test} runs
\code{the\_count\_holds\_what\_is\_written\_and\_the\_status\_is\_the\_line}:
the count reads back what is written, \code{ring} follows it, the
status is the line, and a reset empties the count; and
\code{the\_doorbell\_lowers}.
\item \code{//cpu/vreteno:board\_test} runs
\code{razboj\_draws\_a\_list\_rung\_on\_the\_doorbell}: a program on the
board writes a list, rings, waits for the zero and for idle, and reads
the pixels back over the bus.
\end{itemize}

\subsection*{Used by}

\code{vreteno32::board::Board}, on the peripheral page's tenth slot,
beside \code{Raster<32, 2, 10, 480, 0x4200\_0000, 0x4280\_0000,
0x3900>}.

\subsection*{Limits}

The count is sixteen bits. A write's strobes are not looked at.


\end{document}
